
\documentclass[11pt]{report}

\usepackage[margin=1in]{geometry}
\usepackage[T1]{fontenc}
\usepackage[utf8]{inputenc}
\usepackage{lmodern}
\usepackage{microtype}
\usepackage{amsmath,amssymb,amsthm,mathtools}
\usepackage{array}
\usepackage{booktabs}
\usepackage{longtable}
\usepackage{enumitem}
\usepackage{xcolor}
\usepackage{hyperref}
\usepackage{listings}
\usepackage{stmaryrd}
\usepackage{tikz}
\usepackage{algorithm}
\usepackage{algpseudocode}

\hypersetup{colorlinks=true,linkcolor=blue!60!black,urlcolor=blue!60!black,citecolor=blue!60!black}

\lstdefinestyle{deppy}{
  basicstyle=\ttfamily\small,
  keywordstyle=\color{blue!60!black},
  commentstyle=\color{green!40!black},
  stringstyle=\color{red!50!black},
  showstringspaces=false,
  frame=single,
  breaklines=true,
  columns=fullflexible
}

\newcommand{\DepPy}{\textsc{DepPy}}
\newcommand{\TensorGuard}{\textsc{TensorGuard}}
\newcommand{\LiquidHaskell}{\textsc{LiquidHaskell}}
\newcommand{\Zthree}{\textsc{Z3}}
\newcommand{\Type}{\ensuremath{\mathsf{Type}}}
\newcommand{\Expr}{\ensuremath{\mathsf{Expr}}}
\newcommand{\Pred}{\ensuremath{\mathsf{Pred}}}
\newcommand{\Heap}{\ensuremath{\mathsf{Heap}}}
\newcommand{\Env}{\ensuremath{\mathsf{Env}}}
\newcommand{\Core}{\ensuremath{\mathsf{Core}}}
\newcommand{\Theory}{\ensuremath{\mathsf{Theory}}}
\newcommand{\Site}{\ensuremath{\mathsf{Site}}}
\newcommand{\Sem}{\ensuremath{\mathsf{Sem}}}
\newcommand{\Cover}{\ensuremath{\mathsf{Cover}}}
\newcommand{\Input}{\ensuremath{\mathsf{Input}}}
\newcommand{\Output}{\ensuremath{\mathsf{Output}}}
\newcommand{\ErrorSet}{\ensuremath{\mathsf{Err}}}
\newcommand{\Info}{\ensuremath{\mathsf{Info}}}
\newcommand{\cost}{\ensuremath{\mathsf{cost}}}
\newcommand{\gain}{\ensuremath{\mathsf{gain}}}
\newcommand{\trust}{\ensuremath{\mathsf{trust}}}
\newcommand{\support}{\ensuremath{\mathsf{support}}}
\DeclareRobustCommand{\py}[1]{\texttt{\detokenize{#1}}}
\newcommand{\sem}[1]{\llbracket #1 \rrbracket}
\newcommand{\liquid}[1]{\{\, #1 \,\}}
\newcommand{\Proved}{\ensuremath{\mathsf{Proved}}}
\newcommand{\Refuted}{\ensuremath{\mathsf{Refuted}}}
\newcommand{\Unknown}{\ensuremath{\mathsf{Unknown}}}
\newcommand{\Assumed}{\ensuremath{\mathsf{Assumed}}}
\newcommand{\Bounded}{\ensuremath{\mathsf{Bounded}}}

\theoremstyle{definition}
\newtheorem{definition}{Definition}[section]
\newtheorem{example}[definition]{Example}
\newtheorem{remark}[definition]{Remark}
\newtheorem{assumption}[definition]{Assumption}

\theoremstyle{plain}
\newtheorem{theorem}[definition]{Theorem}
\newtheorem{lemma}[definition]{Lemma}
\newtheorem{proposition}[definition]{Proposition}
\newtheorem{corollary}[definition]{Corollary}

\title{Sheaf-Descent Semantic Typing for Native Python\\
\large Maximal Automatic Inference of Input Requirements, Output Guarantees, Local Site Types, and Runtime-Error Explanations in \DepPy}
\author{Design and implementation monograph for the next \DepPy{} architecture}
\date{\today}

\begin{document}
\maketitle
\tableofcontents



\begin{abstract}
This monograph rewrites the \DepPy{} story around a different answer to the
question ``what is a type?'' A type is not a flat set of predicates chosen from a
candidate pool. A type is a globally glued semantic section obtained from local
evidence distributed across a cover of program sites: argument-boundary sites,
control-flow sites, SSA-value sites, call and wrapper sites, theory-trigger
sites, dynamic-trace sites, proof-module sites, heap/protocol sites, and
runtime-error sites. The practical motivation is the same as before but made
more ambitious and more precise. We want ordinary native Python and native
PyTorch code, written with minimal or zero additional annotation, to support the
automatic inference of the richest safe information available: input
requirements, output guarantees, intermediate local types, witness-carrying
facts, library-specific obligations, proof goals, and predicted runtime errors.
We also want the same system to admit an explicit proof-first surface that feels
closer to F* or theorem-carrying programming, yet still elaborates back into the
same semantic core and the same explicit trusted boundary.

The sheaf-descent perspective makes the slogan ``minimal annotation, maximum
information'' mathematically concrete. Minimal annotation is not merely fewer
clauses; it is the smallest amount of boundary information needed to eliminate
global ambiguity. Maximum information is not merely a high score on a frontier
of atoms; it is the richest globally glueable section compatible with local
solvability, proof transport, and library semantics. This change unlocks a
unified account of four things that are usually separated: the synthesis of safe
input refinements, the synthesis of informative output refinements, the typing
of local intermediate values, and the anticipation of runtime errors that remain
reachable under current assumptions. Backward propagation through the site cover
computes the weakest safe input neighborhoods that keep error-sensitive local
sites nonempty. Forward descent computes the strongest output semantics and the
richest local contracts justified by those safe neighborhoods together with the
available evidence. Residual ambiguity is represented not as an unstructured
list of failed obligations but as obstruction data that points to the exact site
or overlap where more information would be most valuable.

The implementation thesis of this document is intentionally concrete. It maps
the mathematics to a staged \DepPy{} architecture grounded in the current code
base: AST lowering, CFG and SSA construction, value-lineage recovery,
library-trigger normalization, theory-pack decomposition, interprocedural
summary transport, frontier extraction, repair guidance, and renderer-based user
surfaces. The existing modules such as \py{src/deppy/pipeline_monograph.py}, \py{src/deppy/surface/python_ssa.py}, \py{src/deppy/interprocedural/frontier_propagation.py}, and the split PyTorch
packs become the seed of a more principled implementation in which theory packs
create sites and compatibility maps rather than merely proposing atoms. The
monograph is therefore both a semantic proposal and a theory of implementation:
a blueprint for a system that can verify that a Torch module avoids
library-specific errors such as invalid reshape, out-of-bounds indexing,
broadcast mismatch, sort/index misuse, or device inconsistencies, while also
supporting proof-oriented capabilities reminiscent of F* but expressed in native
Python.
\end{abstract}


\section{Problem statement, target behavior, and acceptance criteria}

The problem is to build a dependent environment for native Python that infers as
much semantically useful information as possible without forcing users to leave
ordinary Python syntax. A successful system should analyze a module and return a
high-value semantic object describing at least six classes of information.
First, it should synthesize input requirements strong enough to avoid reachable
runtime errors and weak enough to remain useful. Second, it should synthesize
output guarantees strong enough to describe the actual semantic value of the
module: sortedness, permutation preservation, output shape, non-emptiness,
protocol availability, witness export, ownership, or algorithm-specific
correctness coordinates. Third, it should infer local-site types for
intermediate values rather than treating function boundaries as the only place
where meaning lives. Fourth, it should classify which runtime errors remain
reachable and under what input neighborhoods. Fifth, it should distinguish
facts that are solver-checked, proof-checked, dynamically observed, heuristically
suggested, or still residual. Sixth, it should do all of this under a library
model rich enough for real PyTorch code and a proof surface rich enough for
F*-like theorem-carrying development.

The user experience target is deliberately bimodal. A regular Python user should
be able to write a small Torch helper using plain control flow, local bindings,
wrappers, and a few natural assertions. Running the analyzer should surface a
module summary containing inferred preconditions, inferred postconditions, local
intermediate semantics, and concrete explanations of why certain runtime errors
have been eliminated or remain possible. A proof-oriented user should be able to
add theorem modules, transport lemmas, or witness-carrying helper functions in
native Python syntax, with optional decorators or helper combinators, and have
those artifacts tighten the same underlying semantic object rather than enter a
parallel tool. The trust boundary must remain explicit: every imported proof or
suggested refinement is either checked, reduced to a decidable fragment,
classified as bounded automation, or retained as a residual obligation.

This monograph therefore uses a stronger acceptance criterion than ``the tool
finds some useful refinements.'' The correct target is a contract transformer
that maps sparse boundary evidence to maximal safe semantic structure. For a
module \py{f}, the system should compute a pair $(\rho_{in}, \rho_{out})$ plus a
family of local sections and a residual runtime-error map. The input refinement
$\rho_{in}$ should be minimal among safe refinements in the selected lattice of
admissible input neighborhoods. The output refinement $\rho_{out}$ should be
maximal among globally glued postconditions compatible with $\rho_{in}$ and the
local evidence. The local-site family should explain why those boundary facts
hold, and the runtime-error map should explain which error-sensitive sites are
empty or nonempty under the synthesized assumptions. The result should be
renderable as native-Python contracts, diagnostics, proof goals, benchmark
telemetry, and semantic frontiers.

A final constraint is migration realism. The current \DepPy{} code base already
contains a serious prototype: semantic signatures, theory packs, split PyTorch
reasoning, CFG and SSA lowering, frontier enumeration, interprocedural
propagation, repair suggestions, and showcase examples. The rewrite must not
pretend that implementation begins from zero. Instead, it should explain how to
reorganize existing modules into a sheaf-descent architecture so that current
capabilities become special cases of a more coherent semantic core.


\section{Design stance: native Python first, rich semantics second, silence never}

The design stance has three parts. The first is syntactic humility. Users should
write native Python, not a pseudo-Python proof language that happens to borrow
Python keywords. This means normal functions, classes, local variables,
comprehensions, exceptions, higher-order helpers, and PyTorch operator calls
must remain first-class. Optional decorators such as \py{@deppy.requires},
\py{@deppy.ensures}, \py{@deppy.theorem}, or \py{@deppy.decreases} are allowed,
but the architecture must treat them as seeds, not as a precondition for
precision. The second is semantic ambition. Internally the system should support
full dependent structure: dependent functions, witness-carrying results,
indexed families, equality transports, module-level theorems, and rich library
contracts. The third is epistemic honesty. Any fact that cannot be justified by
the selected local logic, proof kernel, or marked empirical evidence must remain
visible as a residual. The system should never silently collapse unknowns into
weak defaults if doing so would mislead users about safety.

This stance immediately rules out two common failure modes. The first is the
``annotation cliff,'' where rich typing is available only once the user has
already rewritten the program in a theorem-prover style. The second is the
``heuristic mirage,'' where a tool produces attractive contracts but hides which
parts come from actual proof, which from a bounded SMT slice, and which from a
best-effort guess or observed execution. A sheaf-descent system can avoid both
because it records provenance at the site level. Local sections carry not only
semantic content but also support and trust metadata. When the global section is
glued, the renderer can explain which regions were settled by theory packs,
which by proofs, which by dynamic traces, which by imported call summaries, and
which remain ambiguous.

The native-Python stance also shapes the implementation vocabulary. Instead of
pretending that users manipulate global proof terms directly, the system should
harvest meaning from ordinary control-flow regions, SSA values, local assertions,
library operators, and wrapper structure. This is why the site notion is so
important. A local variable carrying the result of \py{torch.sort(xs).values}, a
branch guard such as \py{if xs.numel() > 0}, a proof-carrying helper returning a
witness pair, and a call to a protocol method are all legitimate places where
semantics becomes locally clear. A global contract is then something the system
glues from those local clarities, not something the user must narrate from
scratch.


\section{Running examples and the target shape of answers}

A useful monograph should pin the theory to a few representative programs. The
first running example is a Torch ranking helper that sorts scores, extracts the
top element, and then feeds the result through a lightweight wrapper chain. A
current prototype can already infer sortedness and permutation-like facts for
certain direct patterns. The rewritten system should do more. It should infer a
safe input requirement stating the tensor is nonempty if the function indexes
its top element. It should infer output guarantees stating that the returned
value comes from the original tensor, that the associated index lives in the
admissible source range, and that the intermediate sorted tensor satisfies a
local order section. It should explain why the index-out-of-range error site is
empty under the synthesized precondition and why the sortedness site remains
nonempty under the imported sort semantics.

The second running example is a proof-oriented algorithm module written in
native Python but using theorem decorators and ghost helpers. Think of a stable
partition helper or a graph relaxation routine. The user writes a few explicit
lemmas: perhaps one transport lemma about how a partition step preserves a
permutation witness, or one invariant lemma about a loop maintaining a monotone
potential. The sheaf-descent system should interpret these as transport maps
between proof-local sites and ordinary program sites. The user should not need
to restate the full correctness theorem on every wrapper. The system should
instead import the proof-carrying local sections, propagate them through the
cover, and tighten both input requirements and output guarantees. If some step
still exceeds the decidable fragment, the residual should be attached to the
specific overlap or transport site that failed, not dumped as an opaque theorem
hole.

The third running example is a library-verification use case. Suppose the goal
is simply to certify that a Torch module does not hit library-specific errors on
a selected input family. The user may care less about theorem-level postconditions
and more about excluding invalid reshape, mismatched indexing, empty reductions,
device mismatches, or unsupported dtype/operator combinations. The rewritten
architecture should treat these errors as first-class sites whose viability can
be checked locally and pulled back to the input boundary. The output of analysis
should then include a compact synthesized requirement such as ``rank at least
one, first dimension positive, target view cardinality equal to source
cardinality, and selector domain included in source extent.'' This is a highly
practical deliverable: a theorem-free user still gets a strong precondition that
prevents library-specific runtime failures.

These three examples fix the standard for the rest of the document. A good
implementation theory must explain how local semantics is extracted, how it is
represented, how it is checked, how it is glued, how it is pulled back to safe
inputs, how it is pushed forward to rich outputs, and how it is rendered back to
native Python users.


\section{Semantic kernel: site categories, presheaves, boundaries, and errors}

\begin{definition}[Observation site category]
Given a Python module, let $\mathcal{S}$ be a category whose objects are
observation sites and whose morphisms are restriction or reindexing maps.
Objects include module sites, function sites, CFG-block sites, SSA-value sites,
call sites, return sites, argument-boundary sites, output-boundary sites,
dynamic-trace sites, proof-module sites, heap/protocol sites, and
error-sensitive sites induced by partial operators or exceptional transitions.
\end{definition}

\begin{definition}[Semantic presheaf]
A semantic presheaf is a functor $\Sem : \mathcal{S}^{op} \to \mathsf{Poset}$
assigning to each site a partially ordered set of local semantic sections. Each
section contains a carrier type, a local refinement family, provenance metadata,
a proof-boundary classification, and an optional witness or transport payload.
Restriction maps preserve soundness and drop only those coordinates that are not
visible in the smaller site.
\end{definition}

The decisive shift from earlier formulations is that candidate atoms are no
longer primary. A flat atom is only a view of a richer local section. For a
Torch sort site, the local section may contain not only \py{sorted(result)} but
also a permutation witness, index-domain constraints, projected result types,
site provenance, and transport hooks to downstream indexing or view sites. For a
heap/protocol site, the local section may contain ownership, field consistency,
protocol availability, mutation frame information, and exception viability. When
the system renders a contract surface or a frontier, it projects from this
richer object rather than mistaking the projection for the semantic core.

Two boundaries are special. The input boundary $\partial_{in}$ contains sites for
formal parameters, receiver objects, ambient protocol assumptions, ghost inputs,
module configuration, and any dynamic assumptions the user or test trace is
willing to mark as admissible. The output boundary $\partial_{out}$ contains
sites for return values, raised exceptions, mutated heap regions, exported
witnesses, and theorem-level guarantees. Runtime errors are represented by a
family $\ErrorSet$ of error-sensitive sites. Each such site has a local notion of
viability; it is nonempty if the operation can proceed without entering a
library-specific or general runtime error state under the current boundary
assumptions.

\begin{definition}[Bidirectional refinement problem]
For a module or function $f$, the bidirectional refinement problem is to compute
an input section $\rho_{in}$ on $\partial_{in}$, an output section
$\rho_{out}$ on $\partial_{out}$, and a compatible global section $\sigma$ such
that: (1) every error-sensitive site reachable from $\rho_{in}$ is nonempty;
(2) $\rho_{out}$ is maximal among output sections extendable from $\rho_{in}$;
(3) $\rho_{in}$ is minimal among safe input sections in the chosen admissible
lattice; and (4) any remaining ambiguity is explained by explicit obstruction
data.
\end{definition}

This definition intentionally merges typing, verification, and runtime error
anticipation. A dependent environment should not force users to ask separate
tools for precondition synthesis, postcondition synthesis, intermediate typing,
and error prediction. All four are shadows cast by one semantic object viewed
from different boundaries.


\section{The information order and the optimization objective}

Although the sheaf-descent rewrite moves away from ``choose atoms'' as the
primary optimization variable, optimization remains central. There are usually
many compatible global sections, many possible input refinements, and many ways
to summarize local evidence. We therefore need an explicit information order and
a cost model.

For input refinements, the order is reverse permissiveness: a stronger
precondition is lower in the lattice because it excludes more behaviors. For
output refinements, the order is semantic informativeness: a richer postcondition
is higher because it settles more coordinates, exports more witnesses, or fixes
more relations. For local sections, the order combines ambiguity reduction,
proof status, and runtime-safety value. A useful implementation does not search
all these orders blindly. It uses site-local solvers, compatibility constraints,
and ranking heuristics to identify a finite frontier of globally meaningful
options.

A convenient objective for implementation is a lexicographic or weighted tuple:
first minimize residual runtime-error reachability; then minimize input
assumption cost; then maximize output information; then maximize local semantic
coverage; then minimize proof debt; then minimize explanation complexity. This
ordering matches the intended user experience. A system that infers a gorgeous
postcondition while leaving an index error reachable under ordinary inputs is not
acceptable. A system that eliminates the error only by inventing absurdly strong
preconditions is also not acceptable. A system that leaves huge residuals where
cheap local theory reasoning would suffice is wasting automation budget.

The existing \DepPy{} frontier machinery is therefore not discarded; it is
reinterpreted. Frontiers should now range over compact summaries of glued
sections rather than subsets of atoms. A frontier point may represent one safe
input refinement together with one maximal output summary and a small list of
key local witnesses or residual obstructions. This is a richer object than the
current packed signature, but it remains finite and renderable because the heavy
lifting stays local and the global glueing data is compressed before rendering.


\section{From native Python source to site covers}

A precise implementation theory must say exactly how the site category is built
from real Python. The correct pipeline begins with existing frontend and surface
infrastructure rather than bypassing it. The current code already exposes AST,
CFG, SSA, heap-sensitive lowering, and interprocedural summaries. The rewrite
should formalize those stages as cover synthesis.

First, parse the module and record ordinary syntax sites: function definitions,
class bodies, comprehensions, attribute accesses, literal values, call
expressions, assertions, exception edges, and return nodes. Second, lower to a
CFG so that control-dependent sites become explicit. Branch conditions, loop
headers, exceptional successors, and join points each matter because they define
where local facts are introduced, merged, or invalidated. Third, construct SSA
and heap-SSA. The current module \py{src/deppy/surface/python_ssa.py}
already computes versioned names, phi nodes, final returns, and value lineage.
The sheaf-descent rewrite should reinterpret these as indices for SSA-value
sites and overlap maps. A phi node is not merely a bookkeeping device; it is an
overlap site at which sections from predecessor blocks must agree or expose an
obstruction. A lineage chain is not merely a convenience for wrapper detection;
it is a witness that two textual values live in the same semantic neighborhood.

Fourth, normalize theory-triggering expressions. The shared normalization helper
\py{src/deppy/library_theories/pytorch_expr_normalization.py} already
shows the direction: theory packs should see normalized return or lineage views,
not shallow syntax. Fifth, generate library-specific sites. A call to
\py{torch.sort} should create order, permutation, index, and alias sites; a
\py{view} or \py{reshape} call should create shape-compatibility and output-shape
sites; an indexing operator should create selector-domain and source-target
relation sites; a protocol method access should create method-availability and
frame sites. Sixth, add proof sites for any theorem modules, lemma decorators,
proof combinators, or checked witness exports. Finally, add trace sites if the
user provides representative executions or recorded observations.

The output of this phase should be a first-class cover artifact. Concretely it
should be a data structure containing site identifiers, site kinds, local carrier
schemas, incoming and outgoing restriction maps, provenance handles back to AST
and SSA positions, and a dependency index for error-sensitive operations. The
monograph recommends introducing an explicit cover package under something like
\py{src/deppy/core/site_cover.py} or \py{src/deppy/observe/site_graph.py}. This
artifact would become the new substrate over which theory packs, local solvers,
descent, and rendering operate.


\section{Local stalk theories and decidable kernels}

The sheaf formulation does not grant license to build one gigantic global logic.
Quite the opposite: it makes local decidability more important. Every site kind
must declare a local logic or ``stalk theory'' appropriate to the neighborhood.
The following families are central.

Arithmetic and sequence sites support linear arithmetic, interval reasoning,
length equalities, non-emptiness, and simple finite-domain facts. Heap/protocol
sites support field availability, ownership frames, alias exclusions, optional
nullability, and effect summaries. Tensor-shape sites support rank, dimension,
shape tuples, reshape/view cardinality, broadcast legality, and reduction
preconditions. Tensor-order and permutation sites support sortedness, monotone
projection, permutation witnesses, argsort relations, and sorted-value/index
coupling. Indexing sites support selector-domain inclusion and source-target
value provenance. Proof sites support theorem statements, transport obligations,
residual lemmas, and imported witness objects. Trace sites support empirical
contraction of otherwise ambiguous neighborhoods without pretending to establish
a proof.

A precise implementation rule is that every theory pack must declare three
things. First, how it creates sites and what local carriers those sites use.
Second, what formulas or witness objects its local solver can discharge.
Third, how facts restrict across overlaps. This is the sheaf equivalent of the
current theory-pack schema in \py{src/deppy/core/theory_pack.py}. The
schema should therefore be extended from atom-generation callbacks to a richer
interface containing site constructors, local logic descriptors, viability
predicates, forward refinement extractors, backward precondition pullbacks, and
proof-boundary classifiers. The implementation should still allow pack
composition and dependencies, but the unit of composition becomes site families
and transports rather than isolated atoms.

Decidability remains explicit. A tensor-shape site may live in a solver-friendly
semialgebraic fragment. A permutation witness site may only support finite or
symbolic template reasoning. A proof transport site may require a checked theorem
module or stay residual. The renderer must preserve these distinctions. The
system should not claim that a theorem-level permutation property is as trusted
as a simple index-bound fact unless both truly pass the same kernel.


\section{Bidirectional synthesis: safe input requirements, rich local sections, and maximal output guarantees}

Bidirectional synthesis is the central algorithmic novelty. The backward pass
starts from error-sensitive sites and from user-prioritized output desires. Each
partial operator exports a viability predicate over the local incoming
neighborhood. For example, an indexing site says that a selector domain must lie
inside the source extent; a reshape site says that source and target cardinality
must match; a reduction site may require a nonempty dimension unless an identity
value is supplied; a protocol dispatch site says that the receiver neighborhood
must implement the requested method or field. These viability predicates are
transported backward along restriction maps toward the input boundary. The meet
of all pulled-back viability requirements is the synthesized safe input
refinement $\rho_{in}$, provided the meet is consistent. If the meet is
inconsistent, the system reports the obstruction class rather than silently
weakening the claim.

The forward pass begins from $\rho_{in}$ together with local syntax, theory-pack,
trace, and proof evidence. It solves each site locally, propagates solved
sections across overlaps, and glues the richest compatible family. Local sites
for intermediate values therefore receive precise semantics even when the user
never annotated them. The same pass yields the maximal output section
$\rho_{out}$ extendable from the safe input boundary. In practice this means the
analyzer can return both a synthesized \py{requires} surface and a synthesized
\py{ensures} surface, along with a catalogue of local intermediate facts and a
runtime-error report. The output side is not second-class. A user wants to know
not only what must hold for safety, but also what is gained in return.

The implementation should run these passes to a fixed point because they inform
one another. A stronger synthesized output section may reveal that some input
constraint is redundant because the gluing data already forces a stronger local
relation upstream. Conversely, a sharper input requirement may unlock stronger
local solves and therefore stronger outputs. The fixed-point engine should stop
when both the backward safe-input lattice element and the forward output frontier
stabilize.

A useful design principle is that every synthesized contract should be grounded
in a minimal explanation set. When the system proposes an input requirement, it
should attach the local error-sensitive sites that forced it. When it proposes an
output guarantee, it should attach the sites and overlaps that glued to justify
it. This gives users both trust and leverage: they can decide whether to accept
the synthesized contract, weaken it manually, or add one proof or assertion that
changes the fixed point.


\section{Runtime errors as first-class semantic sites}

Traditional type systems talk primarily about accepted and rejected programs.
For native Python, especially library-rich Python, this is not enough. We need a
structured account of runtime errors that can be anticipated, localized,
classified, and linked to synthesized preconditions. The sheaf-descent rewrite
therefore models runtime errors as first-class semantic sites rather than as an
afterthought.

Every error-sensitive operator introduces one or more error sites. Standard
Python introduces sites for \py{IndexError}, \py{KeyError}, \py{AttributeError},
\py{TypeError}, \py{ValueError}, \py{ZeroDivisionError}, assertion failure,
iteration exhaustion in contexts that expect a value, and protocol violations.
PyTorch introduces sites for invalid reshape or view, broadcast mismatch,
out-of-range gather or index-select, invalid reduction over an empty dimension,
device mismatch, unsupported dtype/operator combinations, invalid matrix shape
alignment, and library-specific semantic mismatches such as confusing sorted
values with sorted indices. Heapful object code introduces sites for missing
field initialization, stale alias use after mutation, frame escape, or
constructor invariant failure. Proof-oriented code introduces sites for missing
transport lemmas, unverified witness exports, or decreases obligations that fail
to justify recursion.

Each error site must provide four methods to the implementation. First, a
constructor that identifies where the site comes from and which operands it sees.
Second, a local viability predicate describing when the site is safe. Third, a
local explanation generator that can render the site to the user in ordinary
language. Fourth, a pullback rule that transforms the site's viability into a
constraint on upstream neighborhoods. The correct module boundary for this logic
is a mix of theory packs and a centralized error-site registry. The monograph
recommends adding a package such as \py{src/deppy/kernel/error_sites.py} or \py{src/deppy/observe/error_catalog.py} so that
runtime-error treatment is uniform across standard Python and library-specific
packs.

This approach has a subtle but important effect on the overall story. A module
analysis is no longer judged only by the postconditions it discovers. It is also
judged by how completely it explains the absence or presence of runtime errors.
A user verifying a Torch module for library-specific safety may care primarily
about the error-site frontier, not about full functional correctness. The same
sheaf-descent architecture supports both use cases.


\section{Interprocedural transport, wrappers, and local transparency}

Interprocedural behavior is where many prototypes lose semantic information. The
current \DepPy{} work already contains an important lesson: direct-return wrapper
propagation is too weak, but SSA lineage and alias-aware wrapper recognition make
interprocedural transport much more robust. The sheaf-descent rewrite turns that
lesson into a general rule.

A function summary should not be a single exported atom set. It should be a
boundary section plus an interior miniature cover containing especially valuable
local sites and transport maps. When a caller invokes the function, the system
instantiates the callee cover with an argument reindexing map, imports the
boundary and selected interior sections, and glues them into the caller cover at
the call-result and wrapper sites. The existing logic in \py{src/deppy/interprocedural/frontier_propagation.py} already performs a weaker version
of this idea for direct-return and alias-forwarding wrappers. The next step is to
lift it from atom import to section import.

Wrapper transparency is especially important for native-Python ergonomics. Users
constantly write tiny helpers, local wrappers, decorator-generated functions,
output extractors, and alias-preserving reorganizations. If the semantics does
not survive such structure, the system will feel fragile and non-Pythonic. A
site-based formulation provides a precise criterion for transparency: a wrapper
is semantically transparent to the extent that its induced reindexing map
preserves or refines the relevant local sections without introducing new
obstructions. This criterion is stronger and cleaner than a bag of textual
pattern recognizers, even though pattern recognizers may remain the first
implementation strategy.

The implementation plan should therefore extend interprocedural summaries in
three directions. First, summaries must retain local-site provenance so that
imported sections can explain themselves. Second, summaries must support
bidirectional transport: callee error sites should contribute to caller input
requirements, and caller input assumptions may tighten callee outputs. Third,
summary import must be lineage-aware so that wrappers built out of local aliases,
projections, and shape-preserving operators remain semantically visible.


\section{Proof-oriented capabilities in native Python: an F*-like surface without leaving Python}

The user explicitly asked for F*-like capabilities but in native Python and with
minimal extra annotation. The sheaf-descent design makes this plausible because
proof artifacts become just another family of local sites and transport maps.
The proof surface should expose optional tools, not a mandatory second language.
At minimum the surface should support theorem decorators, explicit ghost or proof
parameters, witness-returning helpers, local \py{assert} statements lifted into
proof seeds, optional \py{@deppy.requires}/\py{@deppy.ensures} decorators,
optional \py{@deppy.decreases} or \py{@deppy.variant} markers for recursive and
loop reasoning, and helper combinators for equality transport or witness
projection.

A key design choice is that proof artifacts should not merely add finished facts.
They should define transport maps. A theorem about a sorting helper should not
just inject \py{sorted(result)} as a leaf proposition. It should also describe
how that fact behaves under wrappers, views, projections, and index consumers.
A loop invariant should not just be checked at a function boundary. It should be
attached to loop-header and loop-body sites so that the global section can glue
across iterations. A witness-returning helper should create both a value site and
a proof site connected by a dependency relation. This is exactly how F*-like
proof terms acquire practical value for ordinary code: they change the geometry
of future inference rather than remaining isolated certificates.

The current proof modules in \py{src/deppy/surface/proof_dsl.py}, \py{src/deppy/surface/proof_modules.py}, and related elaboration files are an ideal seed.
The rewrite should recast them as frontends to a proof-site constructor and a
transport compiler. Residual proof obligations should remain visible when the
kernel or bounded solver fragment cannot discharge them. The renderer should be
able to show both proof-first users and ordinary users what semantic leverage a
given proof artifact provides. For example: ``this lemma resolves the only
remaining obstruction between the sort-index site and the downstream selector
site, which strengthens the output contract from range-safety to permutation
preservation.'' This kind of explanation is more compelling than a raw theorem
list.


\section{Torch as the flagship library theory}

PyTorch remains the flagship because it combines rich local semantics, real user
value, and the need to reason about library-specific runtime errors. A sheaf-
descent implementation should treat Torch packs as factories for semantic covers.
A call to \py{torch.sort} should create order, permutation, and index sites plus
compatibility maps to any downstream consumers. A \py{view} or \py{reshape} call
should create shape-compatibility, cardinality, alias, and output-shape sites. A
subscript or \py{gather} call should create selector-domain and source-target
value-relation sites. A device move or mixed-device operator should create
compatibility sites whose viability rules generate either safe input constraints
or concrete diagnostics.

The current split packs in \py{src/deppy/library_theories/pytorch_sort.py}, \py{src/deppy/library_theories/pytorch_shapes.py}, and \py{src/deppy/library_theories/pytorch_indexing.py} already point in the right direction.
The rewrite should keep the decomposition but strengthen the interface. Each pack
should export: site constructors, local solver hooks, backward viability
functions, forward output summarizers, residual explanation templates, and
transport declarations. This allows Torch support to grow monotonically.
Adding \py{topk}, \py{gather}, \py{take_along_dim}, \py{argmax}, or broadcasting
packs should create more local semantic neighborhoods without forcing a global
rearchitecture.

The practical target is strong and concrete. For a Torch module, the analyzer
should be able to synthesize a useful input contract explaining how to avoid
library-specific runtime errors, infer local semantics strong enough to catch
values-versus-indices confusion or illegal shape transformations, and emit an
output contract describing properties such as sortedness, range safety,
output-shape, or selector provenance. This is already a compelling verification
story even for users who do not care about full theorem-level proofs.


\section{General Python, protocols, heaps, classes, and exceptions}

Although Torch is the showcase, the architecture must be broader. Native Python
programs are full of protocols, attribute access, constructors, mutable objects,
control-dependent initialization, and exception-heavy APIs. A sheaf-descent
system should therefore treat object and protocol reasoning as first-class local
semantics. A constructor body defines sites for field initialization and
invariant establishment. Attribute reads and writes define heap/protocol sites.
Method dispatch defines protocol-availability sites. Branches or guards that
check \py{is None}, \py{hasattr}, membership, or length define dynamic-local
refinements. Exception edges define error-sensitive sites and possibly
exceptional output sections.

This model yields a practical way to support Pythonic object-oriented code
without pretending that Python's heap is simple. A function may be safe only if
a receiver object satisfies a frame condition or owns a certain initialized
field. Instead of encoding the entire heap globally, the system records local
heap sections where they matter and glues them through explicit overlaps and
frame maps. This keeps the logic modular and helps explanation: when an
attribute error remains reachable, the renderer can point to the exact protocol
or initialization site that failed to glue.

The existing example corpus already contains files such as \py{examples/heap_field_mismatch.py}, \py{examples/protocol_mismatch.py}, and \py{examples/constructor_usage_bug.py}. The rewritten monograph should treat them not as toy extras but as proof that the same site-based architecture used for Torch can handle ordinary Python safety, protocol conformance, and constructor invariants.


\section{Trust boundary, solver boundary, and explicit residuals}

The sheaf-descent rewrite should be more ambitious semantically but no less
strict about trust. Every local solver and every imported proof artifact must
advertise its status. The recommended status lattice is still: trusted automatic,
bounded automatic, dynamically observed, proof-checked, assumed, and residual.
However, the classifier now lives on local sections rather than on standalone
atoms. This is important because one local section may mix carriers settled by
parsing, arithmetic facts settled by \Zthree{}, a witness settled by a theorem
module, and an unresolved transport dependency.

The implementation should centralize this classification in the kernel layer.
The current proof-boundary logic can be refactored into something like a
section-boundary auditor. Whenever a site is solved, the auditor records which
coordinates are trusted and which remain residual. During gluing, the global
section accumulates this information instead of flattening it. Renderers should
therefore be able to say ``the output sortedness claim is solver-backed, the
permutation witness is imported from a checked theorem module, and the alias-
preservation step through a custom wrapper remains residual.'' This level of
clarity is essential if the system is to support both theorem-oriented users and
ordinary developers.

A sheaf semantics also clarifies what not to do. The system must not infer a
safe input requirement from purely empirical evidence and then present it as a
hard theorem. It may offer that requirement as an observed or suggested boundary
section, but the provenance must remain explicit. Similarly, a proof authored by
an LLM should count only after elaboration and kernel checking. The fact that a
rich local section exists at all does not mean every coordinate of that section
has the same trust status.


\section{Precise implementation blueprint and migration from the current code base}

The fastest way to make this rewrite useful is to build it as a disciplined
reorganization of the current \DepPy{} code base. The present prototype already
contains almost every conceptual ingredient, but they are organized around
signatures, atoms, and frontiers rather than site covers and descended sections.
A precise migration should therefore proceed in layers.

Layer one introduces an explicit cover artifact. The frontend and surface layers
already parse Python, lower CFGs, and build SSA plus lineage. Those outputs
should feed a new cover builder that creates site nodes and overlap edges. Layer
two enriches theory packs so that each pack builds site families instead of only
returning atoms. The existing sort, shape, and indexing packs can be migrated
first because they already expose semantically meaningful decomposition points.
Layer three adds local solver interfaces and error-site viability interfaces.
Layer four replaces or wraps current frontier extraction so that it operates on
section summaries. Layer five upgrades renderers to display synthesized input
requirements, output guarantees, local-site facts, and residual obstructions.
Layer six strengthens proof-surface integration by compiling theorem modules into
transport maps and proof sites.

The migration must be incremental. The old \py{DepPyPipeline} and the current
\py{MonographPipeline} should stay alive while the new site-cover pipeline is
introduced in parallel. Every step should come with tests and with corpus
examples that exercise exactly one new capability. This is especially important
for performance and explanation quality. A site cover can explode if constructed
naively. The implementation therefore needs locality heuristics, summary sites,
and pack-specific sparsification rules from the beginning.


\section{Package-by-package implementation blueprint}

The rest of this document is deliberately implementation-heavy. The goal is not merely to describe a beautiful semantic object, but to pin that object to concrete packages, data structures, and migration steps. The following subsections map the sheaf-descent architecture package by package.
\subsection{Package \py{src/deppy/assist}}

\paragraph{Current and continuing responsibility.}
The package \py{src/deppy/assist} should continue to own interactive explanation, prompt-shaping, and recommendation helpers. The
rewrite does not erase this role; it makes the role more explicit by
connecting it to site construction, local section solving, gluing, or
rendering.

\paragraph{Role under sheaf-descent semantics.}
Under the rewrite, the package should Convert obstruction and frontier objects into concise user guidance that points to exact sites, overlaps, and candidate annotations rather than generic advice. This means its public
interfaces must be rephrased in terms of site covers, boundary sections,
local sections, transports, error viability, or audited frontiers rather
than opaque atom sets.

\paragraph{Required additions.}
Add site-aware explanation templates, support for contrastive explanations between two frontiers, and active-query helpers that ask for the smallest boundary seed resolving a chosen obstruction.

\paragraph{Acceptance criteria.}
A package migration is complete when tests can exercise its behavior without
appealing to the full pipeline. In other words, the package should expose
artifacts that are meaningful on their own: a cover builder, a pack-local
site constructor, a local solver result, a transport map, a frontier summary,
or a renderer output whose provenance is site-aware.

\subsection{Package \py{src/deppy/contracts}}

\paragraph{Current and continuing responsibility.}
The package \py{src/deppy/contracts} should continue to own surface syntax for explicit user contracts. The
rewrite does not erase this role; it makes the role more explicit by
connecting it to site construction, local section solving, gluing, or
rendering.

\paragraph{Role under sheaf-descent semantics.}
Under the rewrite, the package should Continue hosting requires, ensures, witness, and theorem-facing contract objects, but redefine them as boundary seeds and local section constraints instead of global atom bags. This means its public
interfaces must be rephrased in terms of site covers, boundary sections,
local sections, transports, error viability, or audited frontiers rather
than opaque atom sets.

\paragraph{Required additions.}
Add explicit input-boundary and output-boundary seed classes, transport-seed descriptors, and provenance fields that record whether a seed came from source text, proof code, or synthesized repair.

\paragraph{Acceptance criteria.}
A package migration is complete when tests can exercise its behavior without
appealing to the full pipeline. In other words, the package should expose
artifacts that are meaningful on their own: a cover builder, a pack-local
site constructor, a local solver result, a transport map, a frontier summary,
or a renderer output whose provenance is site-aware.

\subsection{Package \py{src/deppy/core}}

\paragraph{Current and continuing responsibility.}
The package \py{src/deppy/core} should continue to own semantic datatypes and the stable trusted substrate. The
rewrite does not erase this role; it makes the role more explicit by
connecting it to site construction, local section solving, gluing, or
rendering.

\paragraph{Role under sheaf-descent semantics.}
Under the rewrite, the package should Own the site, section, boundary, obstruction, and frontier summary data structures that every other package manipulates. This means its public
interfaces must be rephrased in terms of site covers, boundary sections,
local sections, transports, error viability, or audited frontiers rather
than opaque atom sets.

\paragraph{Required additions.}
Create modules such as \py{site_cover}, \py{local_section}, \py{obstruction_basis}, \py{boundary_section}, and \py{frontier_summary} while keeping compatibility shims for current semantic signatures.

\paragraph{Acceptance criteria.}
A package migration is complete when tests can exercise its behavior without
appealing to the full pipeline. In other words, the package should expose
artifacts that are meaningful on their own: a cover builder, a pack-local
site constructor, a local solver result, a transport map, a frontier summary,
or a renderer output whose provenance is site-aware.

\subsection{Package \py{src/deppy/frontend}}

\paragraph{Current and continuing responsibility.}
The package \py{src/deppy/frontend} should continue to own module loading and source normalization. The
rewrite does not erase this role; it makes the role more explicit by
connecting it to site construction, local section solving, gluing, or
rendering.

\paragraph{Role under sheaf-descent semantics.}
Under the rewrite, the package should Remain responsible for turning Python modules, proof files, and example corpora into normalized function and module models. This means its public
interfaces must be rephrased in terms of site covers, boundary sections,
local sections, transports, error viability, or audited frontiers rather
than opaque atom sets.

\paragraph{Required additions.}
Expose richer hooks so the cover builder sees decorators, theorem markers, docstring hints, and inline assertions as possible seed sources.

\paragraph{Acceptance criteria.}
A package migration is complete when tests can exercise its behavior without
appealing to the full pipeline. In other words, the package should expose
artifacts that are meaningful on their own: a cover builder, a pack-local
site constructor, a local solver result, a transport map, a frontier summary,
or a renderer output whose provenance is site-aware.

\subsection{Package \py{src/deppy/generate}}

\paragraph{Current and continuing responsibility.}
The package \py{src/deppy/generate} should continue to own candidate and suggestion generation. The
rewrite does not erase this role; it makes the role more explicit by
connecting it to site construction, local section solving, gluing, or
rendering.

\paragraph{Role under sheaf-descent semantics.}
Under the rewrite, the package should Shift from raw candidate-atom generation to candidate-site and candidate-transport generation. This means its public
interfaces must be rephrased in terms of site covers, boundary sections,
local sections, transports, error viability, or audited frontiers rather
than opaque atom sets.

\paragraph{Required additions.}
Preserve current candidate-bank ideas but make them produce boundary-seed options, local witness templates, and transport hypotheses ranked by obstruction impact.

\paragraph{Acceptance criteria.}
A package migration is complete when tests can exercise its behavior without
appealing to the full pipeline. In other words, the package should expose
artifacts that are meaningful on their own: a cover builder, a pack-local
site constructor, a local solver result, a transport map, a frontier summary,
or a renderer output whose provenance is site-aware.

\subsection{Package \py{src/deppy/harvest}}

\paragraph{Current and continuing responsibility.}
The package \py{src/deppy/harvest} should continue to own extracting guards and latent semantic hints from source. The
rewrite does not erase this role; it makes the role more explicit by
connecting it to site construction, local section solving, gluing, or
rendering.

\paragraph{Role under sheaf-descent semantics.}
Under the rewrite, the package should Harvest input-side guards, output-side assertions, loop invariants, protocol checks, and assertion sites from ordinary Python as first-class local evidence. This means its public
interfaces must be rephrased in terms of site covers, boundary sections,
local sections, transports, error viability, or audited frontiers rather
than opaque atom sets.

\paragraph{Required additions.}
Add structured harvesting of guard scopes, traceability to SSA values, and canonicalization so harvested evidence becomes a section constructor rather than an untyped string.

\paragraph{Acceptance criteria.}
A package migration is complete when tests can exercise its behavior without
appealing to the full pipeline. In other words, the package should expose
artifacts that are meaningful on their own: a cover builder, a pack-local
site constructor, a local solver result, a transport map, a frontier summary,
or a renderer output whose provenance is site-aware.

\subsection{Package \py{src/deppy/interprocedural}}

\paragraph{Current and continuing responsibility.}
The package \py{src/deppy/interprocedural} should continue to own summary propagation across call graphs. The
rewrite does not erase this role; it makes the role more explicit by
connecting it to site construction, local section solving, gluing, or
rendering.

\paragraph{Role under sheaf-descent semantics.}
Under the rewrite, the package should Promote existing call summaries and frontier propagation into section transport across caller and callee covers. This means its public
interfaces must be rephrased in terms of site covers, boundary sections,
local sections, transports, error viability, or audited frontiers rather
than opaque atom sets.

\paragraph{Required additions.}
Extend summary objects so they carry boundary sections, selected interior sites, local provenance, and backward error pullbacks.

\paragraph{Acceptance criteria.}
A package migration is complete when tests can exercise its behavior without
appealing to the full pipeline. In other words, the package should expose
artifacts that are meaningful on their own: a cover builder, a pack-local
site constructor, a local solver result, a transport map, a frontier summary,
or a renderer output whose provenance is site-aware.

\subsection{Package \py{src/deppy/kernel}}

\paragraph{Current and continuing responsibility.}
The package \py{src/deppy/kernel} should continue to own trusted classification and checking. The
rewrite does not erase this role; it makes the role more explicit by
connecting it to site construction, local section solving, gluing, or
rendering.

\paragraph{Role under sheaf-descent semantics.}
Under the rewrite, the package should Centralize section-boundary classification, local proof checking, solver bridges, and residual marking. This means its public
interfaces must be rephrased in terms of site covers, boundary sections,
local sections, transports, error viability, or audited frontiers rather
than opaque atom sets.

\paragraph{Required additions.}
Introduce a section auditor that classifies each coordinate of a local section, stores proof obligations, and exposes audited gluing results to renderers.

\paragraph{Acceptance criteria.}
A package migration is complete when tests can exercise its behavior without
appealing to the full pipeline. In other words, the package should expose
artifacts that are meaningful on their own: a cover builder, a pack-local
site constructor, a local solver result, a transport map, a frontier summary,
or a renderer output whose provenance is site-aware.

\subsection{Package \py{src/deppy/library_theories}}

\paragraph{Current and continuing responsibility.}
The package \py{src/deppy/library_theories} should continue to own pack-specific site generation and local solvers. The
rewrite does not erase this role; it makes the role more explicit by
connecting it to site construction, local section solving, gluing, or
rendering.

\paragraph{Role under sheaf-descent semantics.}
Under the rewrite, the package should Remain the home of PyTorch and future library reasoning, but elevate theory packs from atom factories to cover factories. This means its public
interfaces must be rephrased in terms of site covers, boundary sections,
local sections, transports, error viability, or audited frontiers rather
than opaque atom sets.

\paragraph{Required additions.}
Each pack should export site constructors, viability predicates, forward summarizers, local solver adapters, and transport declarations.

\paragraph{Acceptance criteria.}
A package migration is complete when tests can exercise its behavior without
appealing to the full pipeline. In other words, the package should expose
artifacts that are meaningful on their own: a cover builder, a pack-local
site constructor, a local solver result, a transport map, a frontier summary,
or a renderer output whose provenance is site-aware.

\subsection{Package \py{src/deppy/native_python}}

\paragraph{Current and continuing responsibility.}
The package \py{src/deppy/native_python} should continue to own Pythonic user surface and diagnostics. The
rewrite does not erase this role; it makes the role more explicit by
connecting it to site construction, local section solving, gluing, or
rendering.

\paragraph{Role under sheaf-descent semantics.}
Under the rewrite, the package should Own lightweight decorator APIs, diagnostics integration, and user-facing translations of boundary sections into native-Python syntax. This means its public
interfaces must be rephrased in terms of site covers, boundary sections,
local sections, transports, error viability, or audited frontiers rather
than opaque atom sets.

\paragraph{Required additions.}
Add utilities that render synthesized requires and ensures clauses, optional theorem stubs, and precise runtime-error narratives into forms familiar to Python developers.

\paragraph{Acceptance criteria.}
A package migration is complete when tests can exercise its behavior without
appealing to the full pipeline. In other words, the package should expose
artifacts that are meaningful on their own: a cover builder, a pack-local
site constructor, a local solver result, a transport map, a frontier summary,
or a renderer output whose provenance is site-aware.

\subsection{Package \py{src/deppy/normalize}}

\paragraph{Current and continuing responsibility.}
The package \py{src/deppy/normalize} should continue to own canonical expression and contract normalization. The
rewrite does not erase this role; it makes the role more explicit by
connecting it to site construction, local section solving, gluing, or
rendering.

\paragraph{Role under sheaf-descent semantics.}
Under the rewrite, the package should Provide canonicalized expressions, lineage-aware names, theorem skeleton normalization, and error-site normalization. This means its public
interfaces must be rephrased in terms of site covers, boundary sections,
local sections, transports, error viability, or audited frontiers rather
than opaque atom sets.

\paragraph{Required additions.}
The success criterion is that theory packs, error sites, and proof transport all consume one normalization vocabulary instead of ad hoc AST walkers.

\paragraph{Acceptance criteria.}
A package migration is complete when tests can exercise its behavior without
appealing to the full pipeline. In other words, the package should expose
artifacts that are meaningful on their own: a cover builder, a pack-local
site constructor, a local solver result, a transport map, a frontier summary,
or a renderer output whose provenance is site-aware.

\subsection{Package \py{src/deppy/observe}}

\paragraph{Current and continuing responsibility.}
The package \py{src/deppy/observe} should continue to own dynamic trace ingestion and empirical contraction. The
rewrite does not erase this role; it makes the role more explicit by
connecting it to site construction, local section solving, gluing, or
rendering.

\paragraph{Role under sheaf-descent semantics.}
Under the rewrite, the package should Treat traces, logged tensor shapes, observed branch outcomes, and replayed assertions as trace-site constructors and local section contractors. This means its public
interfaces must be rephrased in terms of site covers, boundary sections,
local sections, transports, error viability, or audited frontiers rather
than opaque atom sets.

\paragraph{Required additions.}
Add explicit provenance tags so empirical contractions are useful without being confused with proof-backed facts.

\paragraph{Acceptance criteria.}
A package migration is complete when tests can exercise its behavior without
appealing to the full pipeline. In other words, the package should expose
artifacts that are meaningful on their own: a cover builder, a pack-local
site constructor, a local solver result, a transport map, a frontier summary,
or a renderer output whose provenance is site-aware.

\subsection{Package \py{src/deppy/optimize}}

\paragraph{Current and continuing responsibility.}
The package \py{src/deppy/optimize} should continue to own frontiers, repair, and cost-aware search. The
rewrite does not erase this role; it makes the role more explicit by
connecting it to site construction, local section solving, gluing, or
rendering.

\paragraph{Role under sheaf-descent semantics.}
Under the rewrite, the package should Reframe optimization over summaries of descended sections and obstruction-resolving actions. This means its public
interfaces must be rephrased in terms of site covers, boundary sections,
local sections, transports, error viability, or audited frontiers rather
than opaque atom sets.

\paragraph{Required additions.}
Existing MaxSMT or enumeration machinery should be retained where useful, but the optimized object becomes a frontier point over input sections, output sections, and residual debt.

\paragraph{Acceptance criteria.}
A package migration is complete when tests can exercise its behavior without
appealing to the full pipeline. In other words, the package should expose
artifacts that are meaningful on their own: a cover builder, a pack-local
site constructor, a local solver result, a transport map, a frontier summary,
or a renderer output whose provenance is site-aware.

\subsection{Package \py{src/deppy/predicates}}

\paragraph{Current and continuing responsibility.}
The package \py{src/deppy/predicates} should continue to own reusable logical vocabularies. The
rewrite does not erase this role; it makes the role more explicit by
connecting it to site construction, local section solving, gluing, or
rendering.

\paragraph{Role under sheaf-descent semantics.}
Under the rewrite, the package should Keep reusable predicates, but attach them to site families and local logic descriptors rather than storing them as context-free strings. This means its public
interfaces must be rephrased in terms of site covers, boundary sections,
local sections, transports, error viability, or audited frontiers rather
than opaque atom sets.

\paragraph{Required additions.}
This package should become the vocabulary layer from which packs and renderers derive printable contracts.

\paragraph{Acceptance criteria.}
A package migration is complete when tests can exercise its behavior without
appealing to the full pipeline. In other words, the package should expose
artifacts that are meaningful on their own: a cover builder, a pack-local
site constructor, a local solver result, a transport map, a frontier summary,
or a renderer output whose provenance is site-aware.

\subsection{Package \py{src/deppy/proof}}

\paragraph{Current and continuing responsibility.}
The package \py{src/deppy/proof} should continue to own theorem-oriented APIs and proof transport compilation. The
rewrite does not erase this role; it makes the role more explicit by
connecting it to site construction, local section solving, gluing, or
rendering.

\paragraph{Role under sheaf-descent semantics.}
Under the rewrite, the package should Compile theorem modules, witness combinators, and explicit proof annotations into proof sites and transport maps. This means its public
interfaces must be rephrased in terms of site covers, boundary sections,
local sections, transports, error viability, or audited frontiers rather
than opaque atom sets.

\paragraph{Required additions.}
Implement proof elaboration as site construction plus kernel obligations instead of as direct fact insertion.

\paragraph{Acceptance criteria.}
A package migration is complete when tests can exercise its behavior without
appealing to the full pipeline. In other words, the package should expose
artifacts that are meaningful on their own: a cover builder, a pack-local
site constructor, a local solver result, a transport map, a frontier summary,
or a renderer output whose provenance is site-aware.

\subsection{Package \py{src/deppy/render}}

\paragraph{Current and continuing responsibility.}
The package \py{src/deppy/render} should continue to own contract, diagnostic, and proof-goal views. The
rewrite does not erase this role; it makes the role more explicit by
connecting it to site construction, local section solving, gluing, or
rendering.

\paragraph{Role under sheaf-descent semantics.}
Under the rewrite, the package should Render synthesized input requirements, output guarantees, local-site sections, error-site reports, and obstruction explanations in multiple user-facing forms. This means its public
interfaces must be rephrased in terms of site covers, boundary sections,
local sections, transports, error viability, or audited frontiers rather
than opaque atom sets.

\paragraph{Required additions.}
Add sheaf-aware views such as overlap conflicts, site-level provenance cards, and frontier comparisons.

\paragraph{Acceptance criteria.}
A package migration is complete when tests can exercise its behavior without
appealing to the full pipeline. In other words, the package should expose
artifacts that are meaningful on their own: a cover builder, a pack-local
site constructor, a local solver result, a transport map, a frontier summary,
or a renderer output whose provenance is site-aware.

\subsection{Package \py{src/deppy/solver}}

\paragraph{Current and continuing responsibility.}
The package \py{src/deppy/solver} should continue to own bridges to SMT and other local engines. The
rewrite does not erase this role; it makes the role more explicit by
connecting it to site construction, local section solving, gluing, or
rendering.

\paragraph{Role under sheaf-descent semantics.}
Under the rewrite, the package should Stay local: solve arithmetic, shape, and finite witness subproblems attached to specific site families. This means its public
interfaces must be rephrased in terms of site covers, boundary sections,
local sections, transports, error viability, or audited frontiers rather
than opaque atom sets.

\paragraph{Required additions.}
Expose per-site kernels and explicit failure reasons so gluing can distinguish solver failure from genuine semantic ambiguity.

\paragraph{Acceptance criteria.}
A package migration is complete when tests can exercise its behavior without
appealing to the full pipeline. In other words, the package should expose
artifacts that are meaningful on their own: a cover builder, a pack-local
site constructor, a local solver result, a transport map, a frontier summary,
or a renderer output whose provenance is site-aware.

\subsection{Package \py{src/deppy/static_analysis}}

\paragraph{Current and continuing responsibility.}
The package \py{src/deppy/static_analysis} should continue to own supporting classical analyses reused by cover synthesis. The
rewrite does not erase this role; it makes the role more explicit by
connecting it to site construction, local section solving, gluing, or
rendering.

\paragraph{Role under sheaf-descent semantics.}
Under the rewrite, the package should Retain conventional analyses that feed the cover, such as call graph construction, simple data flow, and control-dependence summaries. This means its public
interfaces must be rephrased in terms of site covers, boundary sections,
local sections, transports, error viability, or audited frontiers rather
than opaque atom sets.

\paragraph{Required additions.}
The key addition is stable identifiers that let site construction and renderers refer back to traditional analysis results.

\paragraph{Acceptance criteria.}
A package migration is complete when tests can exercise its behavior without
appealing to the full pipeline. In other words, the package should expose
artifacts that are meaningful on their own: a cover builder, a pack-local
site constructor, a local solver result, a transport map, a frontier summary,
or a renderer output whose provenance is site-aware.

\subsection{Package \py{src/deppy/surface}}

\paragraph{Current and continuing responsibility.}
The package \py{src/deppy/surface} should continue to own ordinary Python and proof-surface lowering. The
rewrite does not erase this role; it makes the role more explicit by
connecting it to site construction, local section solving, gluing, or
rendering.

\paragraph{Role under sheaf-descent semantics.}
Under the rewrite, the package should Remain the main lowering entry point from syntax into CFG, SSA, heap-SSA, proof syntax, and now site covers. This means its public
interfaces must be rephrased in terms of site covers, boundary sections,
local sections, transports, error viability, or audited frontiers rather
than opaque atom sets.

\paragraph{Required additions.}
This package should expose a cover builder API built on top of current AST, CFG, and SSA datatypes so tests can target it directly.

\paragraph{Acceptance criteria.}
A package migration is complete when tests can exercise its behavior without
appealing to the full pipeline. In other words, the package should expose
artifacts that are meaningful on their own: a cover builder, a pack-local
site constructor, a local solver result, a transport map, a frontier summary,
or a renderer output whose provenance is site-aware.

\subsection{Package \py{src/deppy/types}}

\paragraph{Current and continuing responsibility.}
The package \py{src/deppy/types} should continue to own base carrier and type-language definitions. The
rewrite does not erase this role; it makes the role more explicit by
connecting it to site construction, local section solving, gluing, or
rendering.

\paragraph{Role under sheaf-descent semantics.}
Under the rewrite, the package should Hold the base carrier language, indexed families, witness-bearing result forms, and nominal/structural type descriptors used by local sections. This means its public
interfaces must be rephrased in terms of site covers, boundary sections,
local sections, transports, error viability, or audited frontiers rather
than opaque atom sets.

\paragraph{Required additions.}
The rewrite should distinguish carrier types from local section refinements, allowing the same carrier to participate in many site-specific sections.

\paragraph{Acceptance criteria.}
A package migration is complete when tests can exercise its behavior without
appealing to the full pipeline. In other words, the package should expose
artifacts that are meaningful on their own: a cover builder, a pack-local
site constructor, a local solver result, a transport map, a frontier summary,
or a renderer output whose provenance is site-aware.


\section{Site-family catalogue}

The site-family catalogue below is implementation guidance. Each family should have an explicit constructor, a local carrier schema, and clear rules for backward and forward propagation.
\subsection{Argument-boundary sites}

\paragraph{Program fragment represented.}
Argument-boundary sites represent formal parameters, receiver objects, ambient module state, and optional ghost inputs. They matter because the cover must preserve the
places where meaning becomes locally obvious or locally dangerous.

\paragraph{Primary semantic purpose.}
Their main role is input admissibility and safe precondition synthesis. In the sheaf-descent view a site family is
justified only if it reduces annotation demand, improves runtime-safety
prediction, or enables stronger output synthesis.

\paragraph{Concrete representation.}
The natural implementation home is surface, contracts, and harvest layers. The local payload should contain
at least parameter carriers, harvested guards, protocol assumptions, and call-context projections together with provenance back to source and an audited
trust classification.

\paragraph{Backward and forward behavior.}
Every family must specify how viability requirements pull back toward the
input boundary and how solved local sections push forward toward outputs.
Families that do neither are not mature enough for inclusion in the core
architecture.

\subsection{Return and output-boundary sites}

\paragraph{Program fragment represented.}
Return and output-boundary sites represent returns, raised exceptions, mutated heap regions, and exported witnesses. They matter because the cover must preserve the
places where meaning becomes locally obvious or locally dangerous.

\paragraph{Primary semantic purpose.}
Their main role is maximal output refinement synthesis. In the sheaf-descent view a site family is
justified only if it reduces annotation demand, improves runtime-safety
prediction, or enables stronger output synthesis.

\paragraph{Concrete representation.}
The natural implementation home is core and render layers. The local payload should contain
at least boundary sections, postcondition summaries, and exceptional outputs together with provenance back to source and an audited
trust classification.

\paragraph{Backward and forward behavior.}
Every family must specify how viability requirements pull back toward the
input boundary and how solved local sections push forward toward outputs.
Families that do neither are not mature enough for inclusion in the core
architecture.

\subsection{SSA-value sites}

\paragraph{Program fragment represented.}
SSA-value sites represent versioned locals, tuple projections, alias chains, and phi merges. They matter because the cover must preserve the
places where meaning becomes locally obvious or locally dangerous.

\paragraph{Primary semantic purpose.}
Their main role is local intermediate typing and wrapper transparency. In the sheaf-descent view a site family is
justified only if it reduces annotation demand, improves runtime-safety
prediction, or enables stronger output synthesis.

\paragraph{Concrete representation.}
The natural implementation home is \py{surface/python_ssa.py} and interprocedural transport. The local payload should contain
at least carrier types, lineage identifiers, local refinements, and overlap maps together with provenance back to source and an audited
trust classification.

\paragraph{Backward and forward behavior.}
Every family must specify how viability requirements pull back toward the
input boundary and how solved local sections push forward toward outputs.
Families that do neither are not mature enough for inclusion in the core
architecture.

\subsection{Branch-guard sites}

\paragraph{Program fragment represented.}
Branch-guard sites represent control-dependent refinements introduced by if-statements, assertions, and match-like structures. They matter because the cover must preserve the
places where meaning becomes locally obvious or locally dangerous.

\paragraph{Primary semantic purpose.}
Their main role is guard harvesting and local narrowing. In the sheaf-descent view a site family is
justified only if it reduces annotation demand, improves runtime-safety
prediction, or enables stronger output synthesis.

\paragraph{Concrete representation.}
The natural implementation home is harvest plus \py{static_analysis}. The local payload should contain
at least predicate implications, branch polarity, and merge-side obstruction hints together with provenance back to source and an audited
trust classification.

\paragraph{Backward and forward behavior.}
Every family must specify how viability requirements pull back toward the
input boundary and how solved local sections push forward toward outputs.
Families that do neither are not mature enough for inclusion in the core
architecture.

\subsection{Call-result sites}

\paragraph{Program fragment represented.}
Call-result sites represent semantic neighborhoods created at function and method calls. They matter because the cover must preserve the
places where meaning becomes locally obvious or locally dangerous.

\paragraph{Primary semantic purpose.}
Their main role is summary import and error propagation. In the sheaf-descent view a site family is
justified only if it reduces annotation demand, improves runtime-safety
prediction, or enables stronger output synthesis.

\paragraph{Concrete representation.}
The natural implementation home is interprocedural summaries and library packs. The local payload should contain
at least callee boundary import, actual-to-formal reindexing, and local result sections together with provenance back to source and an audited
trust classification.

\paragraph{Backward and forward behavior.}
Every family must specify how viability requirements pull back toward the
input boundary and how solved local sections push forward toward outputs.
Families that do neither are not mature enough for inclusion in the core
architecture.

\subsection{Tensor-shape sites}

\paragraph{Program fragment represented.}
Tensor-shape sites represent rank, shape tuple, broadcasting, reshape/view legality, and reduction preconditions. They matter because the cover must preserve the
places where meaning becomes locally obvious or locally dangerous.

\paragraph{Primary semantic purpose.}
Their main role is Torch-specific safety and rich tensor outputs. In the sheaf-descent view a site family is
justified only if it reduces annotation demand, improves runtime-safety
prediction, or enables stronger output synthesis.

\paragraph{Concrete representation.}
The natural implementation home is \py{library_theories/pytorch_shapes.py} and solver bridges. The local payload should contain
at least shape equations, cardinality facts, rank bounds, and view legality witnesses together with provenance back to source and an audited
trust classification.

\paragraph{Backward and forward behavior.}
Every family must specify how viability requirements pull back toward the
input boundary and how solved local sections push forward toward outputs.
Families that do neither are not mature enough for inclusion in the core
architecture.

\subsection{Tensor-order sites}

\paragraph{Program fragment represented.}
Tensor-order sites represent sortedness, argsort, permutation, monotone slices, and value-index coupling. They matter because the cover must preserve the
places where meaning becomes locally obvious or locally dangerous.

\paragraph{Primary semantic purpose.}
Their main role is semantics around sort-like operators. In the sheaf-descent view a site family is
justified only if it reduces annotation demand, improves runtime-safety
prediction, or enables stronger output synthesis.

\paragraph{Concrete representation.}
The natural implementation home is \py{library_theories/pytorch_sort.py}. The local payload should contain
at least order predicates, permutation witnesses, and transport hooks to selectors together with provenance back to source and an audited
trust classification.

\paragraph{Backward and forward behavior.}
Every family must specify how viability requirements pull back toward the
input boundary and how solved local sections push forward toward outputs.
Families that do neither are not mature enough for inclusion in the core
architecture.

\subsection{Tensor-indexing sites}

\paragraph{Program fragment represented.}
Tensor-indexing sites represent selectors, slices, gather/index-select domains, and source-target relations. They matter because the cover must preserve the
places where meaning becomes locally obvious or locally dangerous.

\paragraph{Primary semantic purpose.}
Their main role is index safety and provenance-preserving outputs. In the sheaf-descent view a site family is
justified only if it reduces annotation demand, improves runtime-safety
prediction, or enables stronger output synthesis.

\paragraph{Concrete representation.}
The natural implementation home is \py{library_theories/pytorch_indexing.py}. The local payload should contain
at least selector-domain sections, source extent constraints, and target-origin facts together with provenance back to source and an audited
trust classification.

\paragraph{Backward and forward behavior.}
Every family must specify how viability requirements pull back toward the
input boundary and how solved local sections push forward toward outputs.
Families that do neither are not mature enough for inclusion in the core
architecture.

\subsection{Heap/protocol sites}

\paragraph{Program fragment represented.}
Heap/protocol sites represent attribute existence, field initialization, alias frames, protocol method availability, and ownership. They matter because the cover must preserve the
places where meaning becomes locally obvious or locally dangerous.

\paragraph{Primary semantic purpose.}
Their main role is ordinary Python object safety. In the sheaf-descent view a site family is
justified only if it reduces annotation demand, improves runtime-safety
prediction, or enables stronger output synthesis.

\paragraph{Concrete representation.}
The natural implementation home is surface, predicates, and kernel support. The local payload should contain
at least field maps, protocol obligations, frame conditions, and constructor invariants together with provenance back to source and an audited
trust classification.

\paragraph{Backward and forward behavior.}
Every family must specify how viability requirements pull back toward the
input boundary and how solved local sections push forward toward outputs.
Families that do neither are not mature enough for inclusion in the core
architecture.

\subsection{Proof sites}

\paragraph{Program fragment represented.}
Proof sites represent theorem declarations, lemma bodies, witness exports, decreases obligations, and transport seeds. They matter because the cover must preserve the
places where meaning becomes locally obvious or locally dangerous.

\paragraph{Primary semantic purpose.}
Their main role is F*-like capabilities in native Python. In the sheaf-descent view a site family is
justified only if it reduces annotation demand, improves runtime-safety
prediction, or enables stronger output synthesis.

\paragraph{Concrete representation.}
The natural implementation home is proof and kernel packages. The local payload should contain
at least proof obligations, witness carriers, and checked transport maps together with provenance back to source and an audited
trust classification.

\paragraph{Backward and forward behavior.}
Every family must specify how viability requirements pull back toward the
input boundary and how solved local sections push forward toward outputs.
Families that do neither are not mature enough for inclusion in the core
architecture.

\subsection{Trace sites}

\paragraph{Program fragment represented.}
Trace sites represent observed shapes, branch outcomes, assertion passes, and sampled value relations. They matter because the cover must preserve the
places where meaning becomes locally obvious or locally dangerous.

\paragraph{Primary semantic purpose.}
Their main role is empirical contraction of ambiguity classes. In the sheaf-descent view a site family is
justified only if it reduces annotation demand, improves runtime-safety
prediction, or enables stronger output synthesis.

\paragraph{Concrete representation.}
The natural implementation home is observe package. The local payload should contain
at least observation sets, provenance labels, and non-theorem trust markers together with provenance back to source and an audited
trust classification.

\paragraph{Backward and forward behavior.}
Every family must specify how viability requirements pull back toward the
input boundary and how solved local sections push forward toward outputs.
Families that do neither are not mature enough for inclusion in the core
architecture.

\subsection{Error sites}

\paragraph{Program fragment represented.}
Error sites represent partial operations or exceptional transitions that may fail. They matter because the cover must preserve the
places where meaning becomes locally obvious or locally dangerous.

\paragraph{Primary semantic purpose.}
Their main role is runtime-error prediction and safe-input synthesis. In the sheaf-descent view a site family is
justified only if it reduces annotation demand, improves runtime-safety
prediction, or enables stronger output synthesis.

\paragraph{Concrete representation.}
The natural implementation home is kernel error catalog and packs. The local payload should contain
at least viability predicates, explanation templates, and upstream pullback rules together with provenance back to source and an audited
trust classification.

\paragraph{Backward and forward behavior.}
Every family must specify how viability requirements pull back toward the
input boundary and how solved local sections push forward toward outputs.
Families that do neither are not mature enough for inclusion in the core
architecture.

\subsection{Loop-header and invariant sites}

\paragraph{Program fragment represented.}
Loop-header and invariant sites represent repeated states at loop heads, candidate invariants, and ranking data. They matter because the cover must preserve the
places where meaning becomes locally obvious or locally dangerous.

\paragraph{Primary semantic purpose.}
Their main role is local-to-global reasoning for iterative algorithms. In the sheaf-descent view a site family is
justified only if it reduces annotation demand, improves runtime-safety
prediction, or enables stronger output synthesis.

\paragraph{Concrete representation.}
The natural implementation home is surface, proof, and optimize. The local payload should contain
at least candidate invariant sections, decreases facts, and merge compatibility constraints together with provenance back to source and an audited
trust classification.

\paragraph{Backward and forward behavior.}
Every family must specify how viability requirements pull back toward the
input boundary and how solved local sections push forward toward outputs.
Families that do neither are not mature enough for inclusion in the core
architecture.

\subsection{Module-summary sites}

\paragraph{Program fragment represented.}
Module-summary sites represent module-wide exports, public API contracts, and cached summaries. They matter because the cover must preserve the
places where meaning becomes locally obvious or locally dangerous.

\paragraph{Primary semantic purpose.}
Their main role is cross-module reuse and distribution. In the sheaf-descent view a site family is
justified only if it reduces annotation demand, improves runtime-safety
prediction, or enables stronger output synthesis.

\paragraph{Concrete representation.}
The natural implementation home is pipeline and render packages. The local payload should contain
at least public boundary sections, exposed residuals, and reusable transport assets together with provenance back to source and an audited
trust classification.

\paragraph{Backward and forward behavior.}
Every family must specify how viability requirements pull back toward the
input boundary and how solved local sections push forward toward outputs.
Families that do neither are not mature enough for inclusion in the core
architecture.


\section{Theory-family catalogue}

Theories remain open-world and additive, but the unit of addition is now a site-aware local calculus rather than an atom bank.
\subsection{Linear arithmetic and intervals}

\paragraph{Logical vocabulary.}
This theory family should cover non-negativity, bounds, lengths, simple affine equalities. The vocabulary must be printable to
users, normalizable for local solving, and explicit about provenance.

\paragraph{Program evidence source.}
It is triggered primarily by argument bounds, list lengths, loop counters, and tensor extents. The trigger mechanism should use the
shared normalization and lineage vocabulary so that aliasing and wrappers do
not erase semantics.

\paragraph{Local solving discipline.}
The recommended local engine is local SMT over arithmetic fragments. Even when heuristic ranking or
proof suggestions help, the trust boundary should be expressed through a
precise local kernel judgment.

\paragraph{Why it belongs in the core story.}
The semantic payoff is IndexError, ZeroDivisionError, and shape-side precondition pullback. A theory family that cannot influence safe
input synthesis, maximal output synthesis, local typing quality, or residual
explanation quality should be treated as optional rather than foundational.

\subsection{Boolean guard theory}

\paragraph{Logical vocabulary.}
This theory family should cover guard polarity, assertion scope, and branch-conditioned refinement. The vocabulary must be printable to
users, normalizable for local solving, and explicit about provenance.

\paragraph{Program evidence source.}
It is triggered primarily by if-statements, asserts, early returns, and exception guards. The trigger mechanism should use the
shared normalization and lineage vocabulary so that aliasing and wrappers do
not erase semantics.

\paragraph{Local solving discipline.}
The recommended local engine is normalization plus implication checks. Even when heuristic ranking or
proof suggestions help, the trust boundary should be expressed through a
precise local kernel judgment.

\paragraph{Why it belongs in the core story.}
The semantic payoff is control-dependent safety and lightweight proof seed extraction. A theory family that cannot influence safe
input synthesis, maximal output synthesis, local typing quality, or residual
explanation quality should be treated as optional rather than foundational.

\subsection{Nullability and optionality}

\paragraph{Logical vocabulary.}
This theory family should cover None checks, existence guards, and optional protocol availability. The vocabulary must be printable to
users, normalizable for local solving, and explicit about provenance.

\paragraph{Program evidence source.}
It is triggered primarily by ordinary Python branches and constructor code. The trigger mechanism should use the
shared normalization and lineage vocabulary so that aliasing and wrappers do
not erase semantics.

\paragraph{Local solving discipline.}
The recommended local engine is simple symbolic narrowing. Even when heuristic ranking or
proof suggestions help, the trust boundary should be expressed through a
precise local kernel judgment.

\paragraph{Why it belongs in the core story.}
The semantic payoff is AttributeError prevention and stronger post-branch outputs. A theory family that cannot influence safe
input synthesis, maximal output synthesis, local typing quality, or residual
explanation quality should be treated as optional rather than foundational.

\subsection{Heap and protocol theory}

\paragraph{Logical vocabulary.}
This theory family should cover field initialization, protocol method existence, frame conditions, and alias exclusions. The vocabulary must be printable to
users, normalizable for local solving, and explicit about provenance.

\paragraph{Program evidence source.}
It is triggered primarily by classes, protocols, mutable objects, and resource wrappers. The trigger mechanism should use the
shared normalization and lineage vocabulary so that aliasing and wrappers do
not erase semantics.

\paragraph{Local solving discipline.}
The recommended local engine is bounded symbolic checks plus explicit residuals. Even when heuristic ranking or
proof suggestions help, the trust boundary should be expressed through a
precise local kernel judgment.

\paragraph{Why it belongs in the core story.}
The semantic payoff is constructor safety, attribute access, and method-dispatch correctness. A theory family that cannot influence safe
input synthesis, maximal output synthesis, local typing quality, or residual
explanation quality should be treated as optional rather than foundational.

\subsection{Tensor shape theory}

\paragraph{Logical vocabulary.}
This theory family should cover rank, dimension, shape tuples, cardinality, broadcast legality, and reduction support. The vocabulary must be printable to
users, normalizable for local solving, and explicit about provenance.

\paragraph{Program evidence source.}
It is triggered primarily by Torch tensor operators and shape-sensitive wrappers. The trigger mechanism should use the
shared normalization and lineage vocabulary so that aliasing and wrappers do
not erase semantics.

\paragraph{Local solving discipline.}
The recommended local engine is solver-backed local shape kernels. Even when heuristic ranking or
proof suggestions help, the trust boundary should be expressed through a
precise local kernel judgment.

\paragraph{Why it belongs in the core story.}
The semantic payoff is invalid reshape, broadcast mismatch, and exact output-shape synthesis. A theory family that cannot influence safe
input synthesis, maximal output synthesis, local typing quality, or residual
explanation quality should be treated as optional rather than foundational.

\subsection{Tensor order theory}

\paragraph{Logical vocabulary.}
This theory family should cover sortedness, monotonicity, permutation witnesses, argsort relations. The vocabulary must be printable to
users, normalizable for local solving, and explicit about provenance.

\paragraph{Program evidence source.}
It is triggered primarily by sort-like operators and order-preserving wrappers. The trigger mechanism should use the
shared normalization and lineage vocabulary so that aliasing and wrappers do
not erase semantics.

\paragraph{Local solving discipline.}
The recommended local engine is finite witness templates plus solver-backed subchecks. Even when heuristic ranking or
proof suggestions help, the trust boundary should be expressed through a
precise local kernel judgment.

\paragraph{Why it belongs in the core story.}
The semantic payoff is values-vs-indices confusion and richer sorted outputs. A theory family that cannot influence safe
input synthesis, maximal output synthesis, local typing quality, or residual
explanation quality should be treated as optional rather than foundational.

\subsection{Tensor indexing theory}

\paragraph{Logical vocabulary.}
This theory family should cover selector domains, slice extents, gather/index-select legality, and value provenance. The vocabulary must be printable to
users, normalizable for local solving, and explicit about provenance.

\paragraph{Program evidence source.}
It is triggered primarily by subscripts, slices, and advanced indexing. The trigger mechanism should use the
shared normalization and lineage vocabulary so that aliasing and wrappers do
not erase semantics.

\paragraph{Local solving discipline.}
The recommended local engine is arithmetic plus pack-specific rules. Even when heuristic ranking or
proof suggestions help, the trust boundary should be expressed through a
precise local kernel judgment.

\paragraph{Why it belongs in the core story.}
The semantic payoff is out-of-range selector detection and source-target output relations. A theory family that cannot influence safe
input synthesis, maximal output synthesis, local typing quality, or residual
explanation quality should be treated as optional rather than foundational.

\subsection{Device and dtype compatibility theory}

\paragraph{Logical vocabulary.}
This theory family should cover device equality, admissible implicit casts, and operator support matrices. The vocabulary must be printable to
users, normalizable for local solving, and explicit about provenance.

\paragraph{Program evidence source.}
It is triggered primarily by mixed-device or mixed-dtype Torch programs. The trigger mechanism should use the
shared normalization and lineage vocabulary so that aliasing and wrappers do
not erase semantics.

\paragraph{Local solving discipline.}
The recommended local engine is table-driven local checks. Even when heuristic ranking or
proof suggestions help, the trust boundary should be expressed through a
precise local kernel judgment.

\paragraph{Why it belongs in the core story.}
The semantic payoff is device mismatch and unsupported operator combinations. A theory family that cannot influence safe
input synthesis, maximal output synthesis, local typing quality, or residual
explanation quality should be treated as optional rather than foundational.

\subsection{Sequence and collection theory}

\paragraph{Logical vocabulary.}
This theory family should cover non-emptiness, membership, stable length relations, and head/tail legality. The vocabulary must be printable to
users, normalizable for local solving, and explicit about provenance.

\paragraph{Program evidence source.}
It is triggered primarily by lists, tuples, deques, and iterator adapters. The trigger mechanism should use the
shared normalization and lineage vocabulary so that aliasing and wrappers do
not erase semantics.

\paragraph{Local solving discipline.}
The recommended local engine is simple algebraic checks. Even when heuristic ranking or
proof suggestions help, the trust boundary should be expressed through a
precise local kernel judgment.

\paragraph{Why it belongs in the core story.}
The semantic payoff is head-of-empty and unpacking safety. A theory family that cannot influence safe
input synthesis, maximal output synthesis, local typing quality, or residual
explanation quality should be treated as optional rather than foundational.

\subsection{Exception theory}

\paragraph{Logical vocabulary.}
This theory family should cover exceptional output channels, throw sites, and caught-vs-uncaught distinctions. The vocabulary must be printable to
users, normalizable for local solving, and explicit about provenance.

\paragraph{Program evidence source.}
It is triggered primarily by try/except blocks, library raises, and proof obligations that may fail. The trigger mechanism should use the
shared normalization and lineage vocabulary so that aliasing and wrappers do
not erase semantics.

\paragraph{Local solving discipline.}
The recommended local engine is control-flow plus error-site hooks. Even when heuristic ranking or
proof suggestions help, the trust boundary should be expressed through a
precise local kernel judgment.

\paragraph{Why it belongs in the core story.}
The semantic payoff is precise reporting of exceptional contracts and remaining failures. A theory family that cannot influence safe
input synthesis, maximal output synthesis, local typing quality, or residual
explanation quality should be treated as optional rather than foundational.

\subsection{Witness and equality transport theory}

\paragraph{Logical vocabulary.}
This theory family should cover dependent pairs, equality witnesses, and explicit transports. The vocabulary must be printable to
users, normalizable for local solving, and explicit about provenance.

\paragraph{Program evidence source.}
It is triggered primarily by proof-oriented code and witness-returning helpers. The trigger mechanism should use the
shared normalization and lineage vocabulary so that aliasing and wrappers do
not erase semantics.

\paragraph{Local solving discipline.}
The recommended local engine is proof-kernel backed local checking. Even when heuristic ranking or
proof suggestions help, the trust boundary should be expressed through a
precise local kernel judgment.

\paragraph{Why it belongs in the core story.}
The semantic payoff is proof reuse and output enrichment. A theory family that cannot influence safe
input synthesis, maximal output synthesis, local typing quality, or residual
explanation quality should be treated as optional rather than foundational.

\subsection{Loop invariant and decreases theory}

\paragraph{Logical vocabulary.}
This theory family should cover candidate invariants, ranking relations, and preservation constraints. The vocabulary must be printable to
users, normalizable for local solving, and explicit about provenance.

\paragraph{Program evidence source.}
It is triggered primarily by loop-heavy algorithms and recursive helpers. The trigger mechanism should use the
shared normalization and lineage vocabulary so that aliasing and wrappers do
not erase semantics.

\paragraph{Local solving discipline.}
The recommended local engine is template generation plus bounded checks. Even when heuristic ranking or
proof suggestions help, the trust boundary should be expressed through a
precise local kernel judgment.

\paragraph{Why it belongs in the core story.}
The semantic payoff is F*-like capability without abandoning native Python. A theory family that cannot influence safe
input synthesis, maximal output synthesis, local typing quality, or residual
explanation quality should be treated as optional rather than foundational.


\section{Algorithmic realization}

The algorithms below should exist as explicit pipeline stages, each with standalone tests and telemetry.
\subsection{Cover synthesis}

\begin{algorithm}[H]
\caption{Sheaf-descent stage: Cover synthesis}
\begin{algorithmic}[1]
\State Parse the module and lower every function to CFG plus SSA plus heap-sensitive summaries.
\State Create site nodes for arguments, returns, SSA values, branch guards, calls, library triggers, proof artifacts, traces, and error-sensitive operations.
\State Insert overlap edges whenever two sites share lineage, control provenance, or pack-declared transport structure.
\State Audit the resulting cover for sparsity so that low-value sites can be summarized before local solving begins.
\end{algorithmic}
\end{algorithm}


\paragraph{Implementation note.}
The stage \emph{Cover synthesis} should emit a durable artifact rather than merely
mutating hidden pipeline state. This is essential for debugging, benchmarking,
and regression testing against the existing example corpus.

\subsection{Local solving}

\begin{algorithm}[H]
\caption{Sheaf-descent stage: Local solving}
\begin{algorithmic}[1]
\State Dispatch each site to its declared local solver or checker.
\State Store both solved coordinates and unsolved residual coordinates with trust classification.
\State Normalize witness objects so that downstream transports can consume them without pack-specific decoding.
\State Emit viability information for every error-sensitive or partially defined local site.
\end{algorithmic}
\end{algorithm}


\paragraph{Implementation note.}
The stage \emph{Local solving} should emit a durable artifact rather than merely
mutating hidden pipeline state. This is essential for debugging, benchmarking,
and regression testing against the existing example corpus.

\subsection{Backward safe-input synthesis}

\begin{algorithm}[H]
\caption{Sheaf-descent stage: Backward safe-input synthesis}
\begin{algorithmic}[1]
\State Initialize a requirement set from every nontrivial error-sensitive site and every user-prioritized output guarantee that depends on an input-side support condition.
\State Pull those requirements backward through overlap and restriction maps toward the input boundary.
\State Meet the pulled-back requirements in the input lattice and detect inconsistency as an obstruction.
\State Render the surviving meet as synthesized requires-style contracts plus a provenance report.
\end{algorithmic}
\end{algorithm}


\paragraph{Implementation note.}
The stage \emph{Backward safe-input synthesis} should emit a durable artifact rather than merely
mutating hidden pipeline state. This is essential for debugging, benchmarking,
and regression testing against the existing example corpus.

\subsection{Forward output synthesis}

\begin{algorithm}[H]
\caption{Sheaf-descent stage: Forward output synthesis}
\begin{algorithmic}[1]
\State Solve local sites under the synthesized safe input section.
\State Propagate solved sections forward through overlaps and wrapper transparencies.
\State Compute maximal output summaries consistent with the global glueing constraints.
\State Render postconditions, witness exports, local-site facts, and residuals separately by trust class.
\end{algorithmic}
\end{algorithm}


\paragraph{Implementation note.}
The stage \emph{Forward output synthesis} should emit a durable artifact rather than merely
mutating hidden pipeline state. This is essential for debugging, benchmarking,
and regression testing against the existing example corpus.

\subsection{Obstruction extraction}

\begin{algorithm}[H]
\caption{Sheaf-descent stage: Obstruction extraction}
\begin{algorithmic}[1]
\State When gluing fails or remains ambiguous, collect the precise overlaps whose restrictions disagree or whose transports remain unproven.
\State Compress these disagreements into an obstruction basis that is small enough to explain to users.
\State Rank candidate seeds, proofs, or theory activations by how many obstruction coordinates they resolve.
\State Expose the ranked repairs through the assist and render layers.
\end{algorithmic}
\end{algorithm}


\paragraph{Implementation note.}
The stage \emph{Obstruction extraction} should emit a durable artifact rather than merely
mutating hidden pipeline state. This is essential for debugging, benchmarking,
and regression testing against the existing example corpus.

\subsection{Interprocedural section transport}

\begin{algorithm}[H]
\caption{Sheaf-descent stage: Interprocedural section transport}
\begin{algorithmic}[1]
\State Import callee boundary and selected interior sections at call sites using actual-to-formal reindexing.
\State Push back callee error viability to caller inputs when the caller has not already enforced stronger conditions.
\State Track alias-forwarding wrappers via SSA lineage so that helper structure does not erase semantics.
\State Cache imported summaries at module-summary sites for reuse and benchmarking.
\end{algorithmic}
\end{algorithm}


\paragraph{Implementation note.}
The stage \emph{Interprocedural section transport} should emit a durable artifact rather than merely
mutating hidden pipeline state. This is essential for debugging, benchmarking,
and regression testing against the existing example corpus.


\section{Native-Python contract and proof surface}

A practical system needs a surface story that ordinary Python developers can
recognize. The best design is sparse and layered: zero-annotation mode should be
fully supported; decorator-based seeds should be optional; and theorem-style
helpers should be available for users who want them.

\begin{lstlisting}[style=deppy,language=Python]
import deppy
import torch

@deppy.requires(lambda scores: scores.dim() == 1)
@deppy.ensures(lambda result, scores: deppy.sorted_tensor(result.values))
def rank_scores(scores: torch.Tensor):
    values, idx = torch.sort(scores)
    return values, idx
\end{lstlisting}

The rewrite, however, should make clear that such decorators are only one way to
seed the cover. The same module should be analyzable without them, perhaps using
only local assertions or representative traces. For proof-oriented users, the
surface may look like the following.

\begin{lstlisting}[style=deppy,language=Python]
import deppy

@deppy.theorem
def sort_projection_transport(xs):
    values, idx = torch.sort(xs)
    deppy.assume(deppy.sorted_tensor(values))
    return deppy.transport(
        source=deppy.sorted_tensor(values),
        target=deppy.argsort_relation(idx, xs),
    )
\end{lstlisting}

The theorem syntax remains native Python, but the elaborator should convert it
into proof sites, local obligations, and transport objects. Users who never want
such features can ignore them entirely. Users who do want them should receive
leverage from surprisingly small annotations because the proof artifacts change
future descent rather than merely satisfying one immediate goal.


\section{Example corpus and benchmark expectations}

The existing example corpus is a strategic asset. The rewrite should use it as a standing regression suite for semantic ambition, not just for parser correctness.
\subsection{\py{examples/protocol_mismatch.py}}

\paragraph{Primary focus.}
This example stresses structural protocol availability and attribute safety. It is valuable because it forces the cover,
theory packs, and renderer to cooperate on a task that ordinary Python users
actually care about.

\paragraph{Desired success condition.}
The analyzer should synthesize receiver-side requirements or report the exact protocol site that remains unglued. The example should produce synthesized boundary contracts,
local-site sections, and residual explanations that are stable enough to be
snapshot-tested.

\subsection{\py{examples/heap_field_mismatch.py}}

\paragraph{Primary focus.}
This example stresses field initialization and mutation-sensitive heap sections. It is valuable because it forces the cover,
theory packs, and renderer to cooperate on a task that ordinary Python users
actually care about.

\paragraph{Desired success condition.}
The output should explain which constructor or branch failed to establish a required heap section. The example should produce synthesized boundary contracts,
local-site sections, and residual explanations that are stable enough to be
snapshot-tested.

\subsection{\py{examples/safe_guarded_head.py}}

\paragraph{Primary focus.}
This example stresses guard harvesting for head-of-sequence safety. It is valuable because it forces the cover,
theory packs, and renderer to cooperate on a task that ordinary Python users
actually care about.

\paragraph{Desired success condition.}
The backward pass should recover the minimal non-emptiness requirement when the guard is absent and discharge the error site when the guard is present. The example should produce synthesized boundary contracts,
local-site sections, and residual explanations that are stable enough to be
snapshot-tested.

\subsection{\py{examples/pytorch_shape_bug_from_tensorguard.py}}

\paragraph{Primary focus.}
This example stresses invalid shape transformations in tensor code. It is valuable because it forces the cover,
theory packs, and renderer to cooperate on a task that ordinary Python users
actually care about.

\paragraph{Desired success condition.}
A good run synthesizes cardinality or rank requirements and a precise shape-oriented error explanation. The example should produce synthesized boundary contracts,
local-site sections, and residual explanations that are stable enough to be
snapshot-tested.

\subsection{\py{examples/pytorch_sort_bad_return.py}}

\paragraph{Primary focus.}
This example stresses values-versus-indices confusion after sort. It is valuable because it forces the cover,
theory packs, and renderer to cooperate on a task that ordinary Python users
actually care about.

\paragraph{Desired success condition.}
The local order and selector sites should separate the semantic roles of values and index outputs. The example should produce synthesized boundary contracts,
local-site sections, and residual explanations that are stable enough to be
snapshot-tested.

\subsection{\py{examples/pytorch_semantic_showcase.py}}

\paragraph{Primary focus.}
This example stresses interprocedural sort, indexing, and shape composition. It is valuable because it forces the cover,
theory packs, and renderer to cooperate on a task that ordinary Python users
actually care about.

\paragraph{Desired success condition.}
This file is the showcase for imported transport across realistic wrapper chains. The example should produce synthesized boundary contracts,
local-site sections, and residual explanations that are stable enough to be
snapshot-tested.

\subsection{\py{examples/pytorch_long_capability_audit.py}}

\paragraph{Primary focus.}
This example stresses long composed pipelines with nested shape and index operations. It is valuable because it forces the cover,
theory packs, and renderer to cooperate on a task that ordinary Python users
actually care about.

\paragraph{Desired success condition.}
The benchmark should show how far local semantics survives under aliasing and composition. The example should produce synthesized boundary contracts,
local-site sections, and residual explanations that are stable enough to be
snapshot-tested.

\subsection{\py{examples/dependent_sigma_witness.py}}

\paragraph{Primary focus.}
This example stresses witness-carrying results and sigma-like exports. It is valuable because it forces the cover,
theory packs, and renderer to cooperate on a task that ordinary Python users
actually care about.

\paragraph{Desired success condition.}
The proof and output-boundary sites should surface explicit witness structure rather than flattening it to text. The example should produce synthesized boundary contracts,
local-site sections, and residual explanations that are stable enough to be
snapshot-tested.

\subsection{\py{examples/constructor_usage_bug.py}}

\paragraph{Primary focus.}
This example stresses constructor invariants and object misuse. It is valuable because it forces the cover,
theory packs, and renderer to cooperate on a task that ordinary Python users
actually care about.

\paragraph{Desired success condition.}
A sheaf-descent implementation should pinpoint the constructor or usage site that leaves the object state unsafe. The example should produce synthesized boundary contracts,
local-site sections, and residual explanations that are stable enough to be
snapshot-tested.

\subsection{\py{examples/heap_protocol_branch_bug.py}}

\paragraph{Primary focus.}
This example stresses branch-sensitive protocol and heap reasoning. It is valuable because it forces the cover,
theory packs, and renderer to cooperate on a task that ordinary Python users
actually care about.

\paragraph{Desired success condition.}
Guard sites and heap/protocol sites should cooperate rather than existing in separate analyses. The example should produce synthesized boundary contracts,
local-site sections, and residual explanations that are stable enough to be
snapshot-tested.

\subsection{\py{examples/mixed_mode_sort_equivalence.py}}

\paragraph{Primary focus.}
This example stresses proof-oriented and ordinary-Python semantics meeting in one pipeline. It is valuable because it forces the cover,
theory packs, and renderer to cooperate on a task that ordinary Python users
actually care about.

\paragraph{Desired success condition.}
The example should demonstrate that theorem-level transport improves unannotated client typing. The example should produce synthesized boundary contracts,
local-site sections, and residual explanations that are stable enough to be
snapshot-tested.

\subsection{\py{examples/theorem_carrying_benchmark.proof}}

\paragraph{Primary focus.}
This example stresses proof-first transport reuse. It is valuable because it forces the cover,
theory packs, and renderer to cooperate on a task that ordinary Python users
actually care about.

\paragraph{Desired success condition.}
This benchmark should quantify how much ordinary-user annotation rate drops once proof sites are available. The example should produce synthesized boundary contracts,
local-site sections, and residual explanations that are stable enough to be
snapshot-tested.


\section{Runtime-error catalogue and safety synthesis}

Error handling is central because safe-input synthesis is one of the main user goals.
\subsection{IndexError on standard sequences}

\paragraph{Local cause.}
The site becomes problematic when selector-domain exceeds sequence extent In the site catalogue this should be
represented explicitly rather than buried in an exception string.

\paragraph{Backward viability strategy.}
The correct response is to arithmetically pull back range bounds to the input boundary or to upstream guards. This is the mechanism by which runtime
error prediction becomes contract synthesis.

\paragraph{Expected rendered outcome.}
The final user-facing artifact should include input non-emptiness or index-range synthesis plus local provenance explanations. This is how the same
semantic machinery serves both verification-heavy and diagnostics-heavy use
cases.

\subsection{KeyError on mappings}

\paragraph{Local cause.}
The site becomes problematic when key-membership cannot be established at a mapping access site In the site catalogue this should be
represented explicitly rather than buried in an exception string.

\paragraph{Backward viability strategy.}
The correct response is to harvest membership guards and propagate key-domain requirements. This is the mechanism by which runtime
error prediction becomes contract synthesis.

\paragraph{Expected rendered outcome.}
The final user-facing artifact should include synthesized input membership requirements or explicit residual key assumptions. This is how the same
semantic machinery serves both verification-heavy and diagnostics-heavy use
cases.

\subsection{AttributeError}

\paragraph{Local cause.}
The site becomes problematic when receiver protocol or initialization site lacks the requested member In the site catalogue this should be
represented explicitly rather than buried in an exception string.

\paragraph{Backward viability strategy.}
The correct response is to pull back method-availability or field-initialization requirements. This is the mechanism by which runtime
error prediction becomes contract synthesis.

\paragraph{Expected rendered outcome.}
The final user-facing artifact should include protocol-aware diagnostics and constructor-side repairs. This is how the same
semantic machinery serves both verification-heavy and diagnostics-heavy use
cases.

\subsection{TypeError from wrong callable or wrong operand}

\paragraph{Local cause.}
The site becomes problematic when carrier mismatch at a call or operator site In the site catalogue this should be
represented explicitly rather than buried in an exception string.

\paragraph{Backward viability strategy.}
The correct response is to strengthen input carrier assumptions or propagate proof-backed coercion transports. This is the mechanism by which runtime
error prediction becomes contract synthesis.

\paragraph{Expected rendered outcome.}
The final user-facing artifact should include clear distinction between type-carrier failure and refinement failure. This is how the same
semantic machinery serves both verification-heavy and diagnostics-heavy use
cases.

\subsection{ValueError from invalid reshape or view}

\paragraph{Local cause.}
The site becomes problematic when shape cardinality or contiguity assumptions fail locally In the site catalogue this should be
represented explicitly rather than buried in an exception string.

\paragraph{Backward viability strategy.}
The correct response is to derive cardinality, rank, or layout requirements on tensor inputs. This is the mechanism by which runtime
error prediction becomes contract synthesis.

\paragraph{Expected rendered outcome.}
The final user-facing artifact should include library-specific preconditions and exact offending site reports. This is how the same
semantic machinery serves both verification-heavy and diagnostics-heavy use
cases.

\subsection{RuntimeError from broadcast mismatch}

\paragraph{Local cause.}
The site becomes problematic when tensor extents are not broadcast-compatible In the site catalogue this should be
represented explicitly rather than buried in an exception string.

\paragraph{Backward viability strategy.}
The correct response is to solve local broadcast equations and pull them back to argument shapes. This is the mechanism by which runtime
error prediction becomes contract synthesis.

\paragraph{Expected rendered outcome.}
The final user-facing artifact should include strong shape requirements and output-shape guarantees when safe. This is how the same
semantic machinery serves both verification-heavy and diagnostics-heavy use
cases.

\subsection{Out-of-range gather/index-select}

\paragraph{Local cause.}
The site becomes problematic when advanced selector domains exceed tensor extents In the site catalogue this should be
represented explicitly rather than buried in an exception string.

\paragraph{Backward viability strategy.}
The correct response is to combine indexing theory with lineage-aware source extents. This is the mechanism by which runtime
error prediction becomes contract synthesis.

\paragraph{Expected rendered outcome.}
The final user-facing artifact should include selector-range requirements and source-target provenance. This is how the same
semantic machinery serves both verification-heavy and diagnostics-heavy use
cases.

\subsection{Reduction on an empty dimension}

\paragraph{Local cause.}
The site becomes problematic when an operator assumes nonempty input along a reduced axis In the site catalogue this should be
represented explicitly rather than buried in an exception string.

\paragraph{Backward viability strategy.}
The correct response is to infer non-emptiness on the relevant dimension or track identity-value exceptions. This is the mechanism by which runtime
error prediction becomes contract synthesis.

\paragraph{Expected rendered outcome.}
The final user-facing artifact should include safer tensor analytics and better top-k style checks. This is how the same
semantic machinery serves both verification-heavy and diagnostics-heavy use
cases.

\subsection{Device mismatch}

\paragraph{Local cause.}
The site becomes problematic when operators expect tensors on compatible devices but receive mixed inputs In the site catalogue this should be
represented explicitly rather than buried in an exception string.

\paragraph{Backward viability strategy.}
The correct response is to pull back device equality or admissible transfer assumptions. This is the mechanism by which runtime
error prediction becomes contract synthesis.

\paragraph{Expected rendered outcome.}
The final user-facing artifact should include device contracts that ordinary users can actually read. This is how the same
semantic machinery serves both verification-heavy and diagnostics-heavy use
cases.

\subsection{Unsupported dtype/operator combination}

\paragraph{Local cause.}
The site becomes problematic when local operator support matrix excludes the current dtype family In the site catalogue this should be
represented explicitly rather than buried in an exception string.

\paragraph{Backward viability strategy.}
The correct response is to derive dtype compatibility or explicit cast requirements. This is the mechanism by which runtime
error prediction becomes contract synthesis.

\paragraph{Expected rendered outcome.}
The final user-facing artifact should include library-aware contracts instead of opaque runtime surprises. This is how the same
semantic machinery serves both verification-heavy and diagnostics-heavy use
cases.

\subsection{ZeroDivisionError}

\paragraph{Local cause.}
The site becomes problematic when denominator site admits zero In the site catalogue this should be
represented explicitly rather than buried in an exception string.

\paragraph{Backward viability strategy.}
The correct response is to backward viability over arithmetic sites. This is the mechanism by which runtime
error prediction becomes contract synthesis.

\paragraph{Expected rendered outcome.}
The final user-facing artifact should include small but essential safety contracts for numeric Python. This is how the same
semantic machinery serves both verification-heavy and diagnostics-heavy use
cases.

\subsection{Assertion failure}

\paragraph{Local cause.}
The site becomes problematic when user assertion site remains nontrivial after local solving In the site catalogue this should be
represented explicitly rather than buried in an exception string.

\paragraph{Backward viability strategy.}
The correct response is to treat assertions as either proven local sections or residual error sites. This is the mechanism by which runtime
error prediction becomes contract synthesis.

\paragraph{Expected rendered outcome.}
The final user-facing artifact should include precise proof-goal rendering rather than generic assertion text. This is how the same
semantic machinery serves both verification-heavy and diagnostics-heavy use
cases.

\subsection{Constructor invariant failure}

\paragraph{Local cause.}
The site becomes problematic when object fields do not glue to a safe initialized state In the site catalogue this should be
represented explicitly rather than buried in an exception string.

\paragraph{Backward viability strategy.}
The correct response is to pull back field-initialization and branch constraints to constructor entry. This is the mechanism by which runtime
error prediction becomes contract synthesis.

\paragraph{Expected rendered outcome.}
The final user-facing artifact should include ordinary OO safety with explicit field provenance. This is how the same
semantic machinery serves both verification-heavy and diagnostics-heavy use
cases.

\subsection{Protocol misuse in method dispatch}

\paragraph{Local cause.}
The site becomes problematic when receiver lacks a required method family or semantic contract In the site catalogue this should be
represented explicitly rather than buried in an exception string.

\paragraph{Backward viability strategy.}
The correct response is to synthesize protocol requirements or residual obligations. This is the mechanism by which runtime
error prediction becomes contract synthesis.

\paragraph{Expected rendered outcome.}
The final user-facing artifact should include structural typing that remains native to Python. This is how the same
semantic machinery serves both verification-heavy and diagnostics-heavy use
cases.

\subsection{Proof transport failure}

\paragraph{Local cause.}
The site becomes problematic when a theorem site exists but the transport relation remains unverified In the site catalogue this should be
represented explicitly rather than buried in an exception string.

\paragraph{Backward viability strategy.}
The correct response is to emit a residual proof site rather than pretending a guarantee exists. This is the mechanism by which runtime
error prediction becomes contract synthesis.

\paragraph{Expected rendered outcome.}
The final user-facing artifact should include clear separation of proof debt from ordinary runtime errors. This is how the same
semantic machinery serves both verification-heavy and diagnostics-heavy use
cases.

\subsection{Decreases or invariant failure}

\paragraph{Local cause.}
The site becomes problematic when recursive or loop reasoning cannot justify progress or preservation In the site catalogue this should be
represented explicitly rather than buried in an exception string.

\paragraph{Backward viability strategy.}
The correct response is to generate proof-side obligations and optionally strengthen input assumptions or invariants. This is the mechanism by which runtime
error prediction becomes contract synthesis.

\paragraph{Expected rendered outcome.}
The final user-facing artifact should include F*-like feedback within native Python syntax. This is how the same
semantic machinery serves both verification-heavy and diagnostics-heavy use
cases.


\section{Cross-component interaction profiles}

The following interaction notes are intentionally concrete. They describe how particular site families and theory families should cooperate in an actual implementation.
\subsection{Interaction: Argument-boundary sites with Linear arithmetic and intervals}

\paragraph{Why the interaction matters.}
Argument-boundary sites are one of the places where linear arithmetic and intervals can become
genuinely useful to ordinary users. If the interaction is modeled poorly,
either the system loses semantic information or it asks for annotations that
should have been avoidable.

\paragraph{Concrete implementation rule.}
The cover builder should create enough overlap structure that the local
vocabulary of linear arithmetic and intervals can be restricted onto the relevant
coordinates of argument-boundary sites. The local solver should emit not only
solved facts but also viability and transport metadata. The backward pass
should then decide whether the interaction contributes to safe input
synthesis, while the forward pass decides whether it contributes to output
enrichment.

\paragraph{Typical failure mode if omitted.}
Without this interaction the system is likely either to over-approximate and
lose useful guarantees, or to under-approximate and miss a reachable runtime
error. In both cases the user experience degrades because local evidence
fails to become global semantics.

\paragraph{Suggested test shape.}
Add one corpus example in which a fact from linear arithmetic and intervals is only
visible because of the presence of argument-boundary sites. The test should
assert both a synthesized contract effect and a provenance explanation.

\subsection{Interaction: Argument-boundary sites with Boolean guard theory}

\paragraph{Why the interaction matters.}
Argument-boundary sites are one of the places where boolean guard theory can become
genuinely useful to ordinary users. If the interaction is modeled poorly,
either the system loses semantic information or it asks for annotations that
should have been avoidable.

\paragraph{Concrete implementation rule.}
The cover builder should create enough overlap structure that the local
vocabulary of boolean guard theory can be restricted onto the relevant
coordinates of argument-boundary sites. The local solver should emit not only
solved facts but also viability and transport metadata. The backward pass
should then decide whether the interaction contributes to safe input
synthesis, while the forward pass decides whether it contributes to output
enrichment.

\paragraph{Typical failure mode if omitted.}
Without this interaction the system is likely either to over-approximate and
lose useful guarantees, or to under-approximate and miss a reachable runtime
error. In both cases the user experience degrades because local evidence
fails to become global semantics.

\paragraph{Suggested test shape.}
Add one corpus example in which a fact from boolean guard theory is only
visible because of the presence of argument-boundary sites. The test should
assert both a synthesized contract effect and a provenance explanation.

\subsection{Interaction: Argument-boundary sites with Nullability and optionality}

\paragraph{Why the interaction matters.}
Argument-boundary sites are one of the places where nullability and optionality can become
genuinely useful to ordinary users. If the interaction is modeled poorly,
either the system loses semantic information or it asks for annotations that
should have been avoidable.

\paragraph{Concrete implementation rule.}
The cover builder should create enough overlap structure that the local
vocabulary of nullability and optionality can be restricted onto the relevant
coordinates of argument-boundary sites. The local solver should emit not only
solved facts but also viability and transport metadata. The backward pass
should then decide whether the interaction contributes to safe input
synthesis, while the forward pass decides whether it contributes to output
enrichment.

\paragraph{Typical failure mode if omitted.}
Without this interaction the system is likely either to over-approximate and
lose useful guarantees, or to under-approximate and miss a reachable runtime
error. In both cases the user experience degrades because local evidence
fails to become global semantics.

\paragraph{Suggested test shape.}
Add one corpus example in which a fact from nullability and optionality is only
visible because of the presence of argument-boundary sites. The test should
assert both a synthesized contract effect and a provenance explanation.

\subsection{Interaction: Argument-boundary sites with Heap and protocol theory}

\paragraph{Why the interaction matters.}
Argument-boundary sites are one of the places where heap and protocol theory can become
genuinely useful to ordinary users. If the interaction is modeled poorly,
either the system loses semantic information or it asks for annotations that
should have been avoidable.

\paragraph{Concrete implementation rule.}
The cover builder should create enough overlap structure that the local
vocabulary of heap and protocol theory can be restricted onto the relevant
coordinates of argument-boundary sites. The local solver should emit not only
solved facts but also viability and transport metadata. The backward pass
should then decide whether the interaction contributes to safe input
synthesis, while the forward pass decides whether it contributes to output
enrichment.

\paragraph{Typical failure mode if omitted.}
Without this interaction the system is likely either to over-approximate and
lose useful guarantees, or to under-approximate and miss a reachable runtime
error. In both cases the user experience degrades because local evidence
fails to become global semantics.

\paragraph{Suggested test shape.}
Add one corpus example in which a fact from heap and protocol theory is only
visible because of the presence of argument-boundary sites. The test should
assert both a synthesized contract effect and a provenance explanation.

\subsection{Interaction: Argument-boundary sites with Tensor shape theory}

\paragraph{Why the interaction matters.}
Argument-boundary sites are one of the places where tensor shape theory can become
genuinely useful to ordinary users. If the interaction is modeled poorly,
either the system loses semantic information or it asks for annotations that
should have been avoidable.

\paragraph{Concrete implementation rule.}
The cover builder should create enough overlap structure that the local
vocabulary of tensor shape theory can be restricted onto the relevant
coordinates of argument-boundary sites. The local solver should emit not only
solved facts but also viability and transport metadata. The backward pass
should then decide whether the interaction contributes to safe input
synthesis, while the forward pass decides whether it contributes to output
enrichment.

\paragraph{Typical failure mode if omitted.}
Without this interaction the system is likely either to over-approximate and
lose useful guarantees, or to under-approximate and miss a reachable runtime
error. In both cases the user experience degrades because local evidence
fails to become global semantics.

\paragraph{Suggested test shape.}
Add one corpus example in which a fact from tensor shape theory is only
visible because of the presence of argument-boundary sites. The test should
assert both a synthesized contract effect and a provenance explanation.

\subsection{Interaction: Argument-boundary sites with Tensor order theory}

\paragraph{Why the interaction matters.}
Argument-boundary sites are one of the places where tensor order theory can become
genuinely useful to ordinary users. If the interaction is modeled poorly,
either the system loses semantic information or it asks for annotations that
should have been avoidable.

\paragraph{Concrete implementation rule.}
The cover builder should create enough overlap structure that the local
vocabulary of tensor order theory can be restricted onto the relevant
coordinates of argument-boundary sites. The local solver should emit not only
solved facts but also viability and transport metadata. The backward pass
should then decide whether the interaction contributes to safe input
synthesis, while the forward pass decides whether it contributes to output
enrichment.

\paragraph{Typical failure mode if omitted.}
Without this interaction the system is likely either to over-approximate and
lose useful guarantees, or to under-approximate and miss a reachable runtime
error. In both cases the user experience degrades because local evidence
fails to become global semantics.

\paragraph{Suggested test shape.}
Add one corpus example in which a fact from tensor order theory is only
visible because of the presence of argument-boundary sites. The test should
assert both a synthesized contract effect and a provenance explanation.

\subsection{Interaction: Argument-boundary sites with Tensor indexing theory}

\paragraph{Why the interaction matters.}
Argument-boundary sites are one of the places where tensor indexing theory can become
genuinely useful to ordinary users. If the interaction is modeled poorly,
either the system loses semantic information or it asks for annotations that
should have been avoidable.

\paragraph{Concrete implementation rule.}
The cover builder should create enough overlap structure that the local
vocabulary of tensor indexing theory can be restricted onto the relevant
coordinates of argument-boundary sites. The local solver should emit not only
solved facts but also viability and transport metadata. The backward pass
should then decide whether the interaction contributes to safe input
synthesis, while the forward pass decides whether it contributes to output
enrichment.

\paragraph{Typical failure mode if omitted.}
Without this interaction the system is likely either to over-approximate and
lose useful guarantees, or to under-approximate and miss a reachable runtime
error. In both cases the user experience degrades because local evidence
fails to become global semantics.

\paragraph{Suggested test shape.}
Add one corpus example in which a fact from tensor indexing theory is only
visible because of the presence of argument-boundary sites. The test should
assert both a synthesized contract effect and a provenance explanation.

\subsection{Interaction: Argument-boundary sites with Device and dtype compatibility theory}

\paragraph{Why the interaction matters.}
Argument-boundary sites are one of the places where device and dtype compatibility theory can become
genuinely useful to ordinary users. If the interaction is modeled poorly,
either the system loses semantic information or it asks for annotations that
should have been avoidable.

\paragraph{Concrete implementation rule.}
The cover builder should create enough overlap structure that the local
vocabulary of device and dtype compatibility theory can be restricted onto the relevant
coordinates of argument-boundary sites. The local solver should emit not only
solved facts but also viability and transport metadata. The backward pass
should then decide whether the interaction contributes to safe input
synthesis, while the forward pass decides whether it contributes to output
enrichment.

\paragraph{Typical failure mode if omitted.}
Without this interaction the system is likely either to over-approximate and
lose useful guarantees, or to under-approximate and miss a reachable runtime
error. In both cases the user experience degrades because local evidence
fails to become global semantics.

\paragraph{Suggested test shape.}
Add one corpus example in which a fact from device and dtype compatibility theory is only
visible because of the presence of argument-boundary sites. The test should
assert both a synthesized contract effect and a provenance explanation.

\subsection{Interaction: Argument-boundary sites with Sequence and collection theory}

\paragraph{Why the interaction matters.}
Argument-boundary sites are one of the places where sequence and collection theory can become
genuinely useful to ordinary users. If the interaction is modeled poorly,
either the system loses semantic information or it asks for annotations that
should have been avoidable.

\paragraph{Concrete implementation rule.}
The cover builder should create enough overlap structure that the local
vocabulary of sequence and collection theory can be restricted onto the relevant
coordinates of argument-boundary sites. The local solver should emit not only
solved facts but also viability and transport metadata. The backward pass
should then decide whether the interaction contributes to safe input
synthesis, while the forward pass decides whether it contributes to output
enrichment.

\paragraph{Typical failure mode if omitted.}
Without this interaction the system is likely either to over-approximate and
lose useful guarantees, or to under-approximate and miss a reachable runtime
error. In both cases the user experience degrades because local evidence
fails to become global semantics.

\paragraph{Suggested test shape.}
Add one corpus example in which a fact from sequence and collection theory is only
visible because of the presence of argument-boundary sites. The test should
assert both a synthesized contract effect and a provenance explanation.

\subsection{Interaction: Argument-boundary sites with Exception theory}

\paragraph{Why the interaction matters.}
Argument-boundary sites are one of the places where exception theory can become
genuinely useful to ordinary users. If the interaction is modeled poorly,
either the system loses semantic information or it asks for annotations that
should have been avoidable.

\paragraph{Concrete implementation rule.}
The cover builder should create enough overlap structure that the local
vocabulary of exception theory can be restricted onto the relevant
coordinates of argument-boundary sites. The local solver should emit not only
solved facts but also viability and transport metadata. The backward pass
should then decide whether the interaction contributes to safe input
synthesis, while the forward pass decides whether it contributes to output
enrichment.

\paragraph{Typical failure mode if omitted.}
Without this interaction the system is likely either to over-approximate and
lose useful guarantees, or to under-approximate and miss a reachable runtime
error. In both cases the user experience degrades because local evidence
fails to become global semantics.

\paragraph{Suggested test shape.}
Add one corpus example in which a fact from exception theory is only
visible because of the presence of argument-boundary sites. The test should
assert both a synthesized contract effect and a provenance explanation.

\subsection{Interaction: Argument-boundary sites with Witness and equality transport theory}

\paragraph{Why the interaction matters.}
Argument-boundary sites are one of the places where witness and equality transport theory can become
genuinely useful to ordinary users. If the interaction is modeled poorly,
either the system loses semantic information or it asks for annotations that
should have been avoidable.

\paragraph{Concrete implementation rule.}
The cover builder should create enough overlap structure that the local
vocabulary of witness and equality transport theory can be restricted onto the relevant
coordinates of argument-boundary sites. The local solver should emit not only
solved facts but also viability and transport metadata. The backward pass
should then decide whether the interaction contributes to safe input
synthesis, while the forward pass decides whether it contributes to output
enrichment.

\paragraph{Typical failure mode if omitted.}
Without this interaction the system is likely either to over-approximate and
lose useful guarantees, or to under-approximate and miss a reachable runtime
error. In both cases the user experience degrades because local evidence
fails to become global semantics.

\paragraph{Suggested test shape.}
Add one corpus example in which a fact from witness and equality transport theory is only
visible because of the presence of argument-boundary sites. The test should
assert both a synthesized contract effect and a provenance explanation.

\subsection{Interaction: Argument-boundary sites with Loop invariant and decreases theory}

\paragraph{Why the interaction matters.}
Argument-boundary sites are one of the places where loop invariant and decreases theory can become
genuinely useful to ordinary users. If the interaction is modeled poorly,
either the system loses semantic information or it asks for annotations that
should have been avoidable.

\paragraph{Concrete implementation rule.}
The cover builder should create enough overlap structure that the local
vocabulary of loop invariant and decreases theory can be restricted onto the relevant
coordinates of argument-boundary sites. The local solver should emit not only
solved facts but also viability and transport metadata. The backward pass
should then decide whether the interaction contributes to safe input
synthesis, while the forward pass decides whether it contributes to output
enrichment.

\paragraph{Typical failure mode if omitted.}
Without this interaction the system is likely either to over-approximate and
lose useful guarantees, or to under-approximate and miss a reachable runtime
error. In both cases the user experience degrades because local evidence
fails to become global semantics.

\paragraph{Suggested test shape.}
Add one corpus example in which a fact from loop invariant and decreases theory is only
visible because of the presence of argument-boundary sites. The test should
assert both a synthesized contract effect and a provenance explanation.

\subsection{Interaction: Return and output-boundary sites with Linear arithmetic and intervals}

\paragraph{Why the interaction matters.}
Return and output-boundary sites are one of the places where linear arithmetic and intervals can become
genuinely useful to ordinary users. If the interaction is modeled poorly,
either the system loses semantic information or it asks for annotations that
should have been avoidable.

\paragraph{Concrete implementation rule.}
The cover builder should create enough overlap structure that the local
vocabulary of linear arithmetic and intervals can be restricted onto the relevant
coordinates of return and output-boundary sites. The local solver should emit not only
solved facts but also viability and transport metadata. The backward pass
should then decide whether the interaction contributes to safe input
synthesis, while the forward pass decides whether it contributes to output
enrichment.

\paragraph{Typical failure mode if omitted.}
Without this interaction the system is likely either to over-approximate and
lose useful guarantees, or to under-approximate and miss a reachable runtime
error. In both cases the user experience degrades because local evidence
fails to become global semantics.

\paragraph{Suggested test shape.}
Add one corpus example in which a fact from linear arithmetic and intervals is only
visible because of the presence of return and output-boundary sites. The test should
assert both a synthesized contract effect and a provenance explanation.

\subsection{Interaction: Return and output-boundary sites with Boolean guard theory}

\paragraph{Why the interaction matters.}
Return and output-boundary sites are one of the places where boolean guard theory can become
genuinely useful to ordinary users. If the interaction is modeled poorly,
either the system loses semantic information or it asks for annotations that
should have been avoidable.

\paragraph{Concrete implementation rule.}
The cover builder should create enough overlap structure that the local
vocabulary of boolean guard theory can be restricted onto the relevant
coordinates of return and output-boundary sites. The local solver should emit not only
solved facts but also viability and transport metadata. The backward pass
should then decide whether the interaction contributes to safe input
synthesis, while the forward pass decides whether it contributes to output
enrichment.

\paragraph{Typical failure mode if omitted.}
Without this interaction the system is likely either to over-approximate and
lose useful guarantees, or to under-approximate and miss a reachable runtime
error. In both cases the user experience degrades because local evidence
fails to become global semantics.

\paragraph{Suggested test shape.}
Add one corpus example in which a fact from boolean guard theory is only
visible because of the presence of return and output-boundary sites. The test should
assert both a synthesized contract effect and a provenance explanation.

\subsection{Interaction: Return and output-boundary sites with Nullability and optionality}

\paragraph{Why the interaction matters.}
Return and output-boundary sites are one of the places where nullability and optionality can become
genuinely useful to ordinary users. If the interaction is modeled poorly,
either the system loses semantic information or it asks for annotations that
should have been avoidable.

\paragraph{Concrete implementation rule.}
The cover builder should create enough overlap structure that the local
vocabulary of nullability and optionality can be restricted onto the relevant
coordinates of return and output-boundary sites. The local solver should emit not only
solved facts but also viability and transport metadata. The backward pass
should then decide whether the interaction contributes to safe input
synthesis, while the forward pass decides whether it contributes to output
enrichment.

\paragraph{Typical failure mode if omitted.}
Without this interaction the system is likely either to over-approximate and
lose useful guarantees, or to under-approximate and miss a reachable runtime
error. In both cases the user experience degrades because local evidence
fails to become global semantics.

\paragraph{Suggested test shape.}
Add one corpus example in which a fact from nullability and optionality is only
visible because of the presence of return and output-boundary sites. The test should
assert both a synthesized contract effect and a provenance explanation.

\subsection{Interaction: Return and output-boundary sites with Heap and protocol theory}

\paragraph{Why the interaction matters.}
Return and output-boundary sites are one of the places where heap and protocol theory can become
genuinely useful to ordinary users. If the interaction is modeled poorly,
either the system loses semantic information or it asks for annotations that
should have been avoidable.

\paragraph{Concrete implementation rule.}
The cover builder should create enough overlap structure that the local
vocabulary of heap and protocol theory can be restricted onto the relevant
coordinates of return and output-boundary sites. The local solver should emit not only
solved facts but also viability and transport metadata. The backward pass
should then decide whether the interaction contributes to safe input
synthesis, while the forward pass decides whether it contributes to output
enrichment.

\paragraph{Typical failure mode if omitted.}
Without this interaction the system is likely either to over-approximate and
lose useful guarantees, or to under-approximate and miss a reachable runtime
error. In both cases the user experience degrades because local evidence
fails to become global semantics.

\paragraph{Suggested test shape.}
Add one corpus example in which a fact from heap and protocol theory is only
visible because of the presence of return and output-boundary sites. The test should
assert both a synthesized contract effect and a provenance explanation.

\subsection{Interaction: Return and output-boundary sites with Tensor shape theory}

\paragraph{Why the interaction matters.}
Return and output-boundary sites are one of the places where tensor shape theory can become
genuinely useful to ordinary users. If the interaction is modeled poorly,
either the system loses semantic information or it asks for annotations that
should have been avoidable.

\paragraph{Concrete implementation rule.}
The cover builder should create enough overlap structure that the local
vocabulary of tensor shape theory can be restricted onto the relevant
coordinates of return and output-boundary sites. The local solver should emit not only
solved facts but also viability and transport metadata. The backward pass
should then decide whether the interaction contributes to safe input
synthesis, while the forward pass decides whether it contributes to output
enrichment.

\paragraph{Typical failure mode if omitted.}
Without this interaction the system is likely either to over-approximate and
lose useful guarantees, or to under-approximate and miss a reachable runtime
error. In both cases the user experience degrades because local evidence
fails to become global semantics.

\paragraph{Suggested test shape.}
Add one corpus example in which a fact from tensor shape theory is only
visible because of the presence of return and output-boundary sites. The test should
assert both a synthesized contract effect and a provenance explanation.

\subsection{Interaction: Return and output-boundary sites with Tensor order theory}

\paragraph{Why the interaction matters.}
Return and output-boundary sites are one of the places where tensor order theory can become
genuinely useful to ordinary users. If the interaction is modeled poorly,
either the system loses semantic information or it asks for annotations that
should have been avoidable.

\paragraph{Concrete implementation rule.}
The cover builder should create enough overlap structure that the local
vocabulary of tensor order theory can be restricted onto the relevant
coordinates of return and output-boundary sites. The local solver should emit not only
solved facts but also viability and transport metadata. The backward pass
should then decide whether the interaction contributes to safe input
synthesis, while the forward pass decides whether it contributes to output
enrichment.

\paragraph{Typical failure mode if omitted.}
Without this interaction the system is likely either to over-approximate and
lose useful guarantees, or to under-approximate and miss a reachable runtime
error. In both cases the user experience degrades because local evidence
fails to become global semantics.

\paragraph{Suggested test shape.}
Add one corpus example in which a fact from tensor order theory is only
visible because of the presence of return and output-boundary sites. The test should
assert both a synthesized contract effect and a provenance explanation.

\subsection{Interaction: Return and output-boundary sites with Tensor indexing theory}

\paragraph{Why the interaction matters.}
Return and output-boundary sites are one of the places where tensor indexing theory can become
genuinely useful to ordinary users. If the interaction is modeled poorly,
either the system loses semantic information or it asks for annotations that
should have been avoidable.

\paragraph{Concrete implementation rule.}
The cover builder should create enough overlap structure that the local
vocabulary of tensor indexing theory can be restricted onto the relevant
coordinates of return and output-boundary sites. The local solver should emit not only
solved facts but also viability and transport metadata. The backward pass
should then decide whether the interaction contributes to safe input
synthesis, while the forward pass decides whether it contributes to output
enrichment.

\paragraph{Typical failure mode if omitted.}
Without this interaction the system is likely either to over-approximate and
lose useful guarantees, or to under-approximate and miss a reachable runtime
error. In both cases the user experience degrades because local evidence
fails to become global semantics.

\paragraph{Suggested test shape.}
Add one corpus example in which a fact from tensor indexing theory is only
visible because of the presence of return and output-boundary sites. The test should
assert both a synthesized contract effect and a provenance explanation.

\subsection{Interaction: Return and output-boundary sites with Device and dtype compatibility theory}

\paragraph{Why the interaction matters.}
Return and output-boundary sites are one of the places where device and dtype compatibility theory can become
genuinely useful to ordinary users. If the interaction is modeled poorly,
either the system loses semantic information or it asks for annotations that
should have been avoidable.

\paragraph{Concrete implementation rule.}
The cover builder should create enough overlap structure that the local
vocabulary of device and dtype compatibility theory can be restricted onto the relevant
coordinates of return and output-boundary sites. The local solver should emit not only
solved facts but also viability and transport metadata. The backward pass
should then decide whether the interaction contributes to safe input
synthesis, while the forward pass decides whether it contributes to output
enrichment.

\paragraph{Typical failure mode if omitted.}
Without this interaction the system is likely either to over-approximate and
lose useful guarantees, or to under-approximate and miss a reachable runtime
error. In both cases the user experience degrades because local evidence
fails to become global semantics.

\paragraph{Suggested test shape.}
Add one corpus example in which a fact from device and dtype compatibility theory is only
visible because of the presence of return and output-boundary sites. The test should
assert both a synthesized contract effect and a provenance explanation.

\subsection{Interaction: Return and output-boundary sites with Sequence and collection theory}

\paragraph{Why the interaction matters.}
Return and output-boundary sites are one of the places where sequence and collection theory can become
genuinely useful to ordinary users. If the interaction is modeled poorly,
either the system loses semantic information or it asks for annotations that
should have been avoidable.

\paragraph{Concrete implementation rule.}
The cover builder should create enough overlap structure that the local
vocabulary of sequence and collection theory can be restricted onto the relevant
coordinates of return and output-boundary sites. The local solver should emit not only
solved facts but also viability and transport metadata. The backward pass
should then decide whether the interaction contributes to safe input
synthesis, while the forward pass decides whether it contributes to output
enrichment.

\paragraph{Typical failure mode if omitted.}
Without this interaction the system is likely either to over-approximate and
lose useful guarantees, or to under-approximate and miss a reachable runtime
error. In both cases the user experience degrades because local evidence
fails to become global semantics.

\paragraph{Suggested test shape.}
Add one corpus example in which a fact from sequence and collection theory is only
visible because of the presence of return and output-boundary sites. The test should
assert both a synthesized contract effect and a provenance explanation.

\subsection{Interaction: Return and output-boundary sites with Exception theory}

\paragraph{Why the interaction matters.}
Return and output-boundary sites are one of the places where exception theory can become
genuinely useful to ordinary users. If the interaction is modeled poorly,
either the system loses semantic information or it asks for annotations that
should have been avoidable.

\paragraph{Concrete implementation rule.}
The cover builder should create enough overlap structure that the local
vocabulary of exception theory can be restricted onto the relevant
coordinates of return and output-boundary sites. The local solver should emit not only
solved facts but also viability and transport metadata. The backward pass
should then decide whether the interaction contributes to safe input
synthesis, while the forward pass decides whether it contributes to output
enrichment.

\paragraph{Typical failure mode if omitted.}
Without this interaction the system is likely either to over-approximate and
lose useful guarantees, or to under-approximate and miss a reachable runtime
error. In both cases the user experience degrades because local evidence
fails to become global semantics.

\paragraph{Suggested test shape.}
Add one corpus example in which a fact from exception theory is only
visible because of the presence of return and output-boundary sites. The test should
assert both a synthesized contract effect and a provenance explanation.

\subsection{Interaction: Return and output-boundary sites with Witness and equality transport theory}

\paragraph{Why the interaction matters.}
Return and output-boundary sites are one of the places where witness and equality transport theory can become
genuinely useful to ordinary users. If the interaction is modeled poorly,
either the system loses semantic information or it asks for annotations that
should have been avoidable.

\paragraph{Concrete implementation rule.}
The cover builder should create enough overlap structure that the local
vocabulary of witness and equality transport theory can be restricted onto the relevant
coordinates of return and output-boundary sites. The local solver should emit not only
solved facts but also viability and transport metadata. The backward pass
should then decide whether the interaction contributes to safe input
synthesis, while the forward pass decides whether it contributes to output
enrichment.

\paragraph{Typical failure mode if omitted.}
Without this interaction the system is likely either to over-approximate and
lose useful guarantees, or to under-approximate and miss a reachable runtime
error. In both cases the user experience degrades because local evidence
fails to become global semantics.

\paragraph{Suggested test shape.}
Add one corpus example in which a fact from witness and equality transport theory is only
visible because of the presence of return and output-boundary sites. The test should
assert both a synthesized contract effect and a provenance explanation.

\subsection{Interaction: Return and output-boundary sites with Loop invariant and decreases theory}

\paragraph{Why the interaction matters.}
Return and output-boundary sites are one of the places where loop invariant and decreases theory can become
genuinely useful to ordinary users. If the interaction is modeled poorly,
either the system loses semantic information or it asks for annotations that
should have been avoidable.

\paragraph{Concrete implementation rule.}
The cover builder should create enough overlap structure that the local
vocabulary of loop invariant and decreases theory can be restricted onto the relevant
coordinates of return and output-boundary sites. The local solver should emit not only
solved facts but also viability and transport metadata. The backward pass
should then decide whether the interaction contributes to safe input
synthesis, while the forward pass decides whether it contributes to output
enrichment.

\paragraph{Typical failure mode if omitted.}
Without this interaction the system is likely either to over-approximate and
lose useful guarantees, or to under-approximate and miss a reachable runtime
error. In both cases the user experience degrades because local evidence
fails to become global semantics.

\paragraph{Suggested test shape.}
Add one corpus example in which a fact from loop invariant and decreases theory is only
visible because of the presence of return and output-boundary sites. The test should
assert both a synthesized contract effect and a provenance explanation.

\subsection{Interaction: SSA-value sites with Linear arithmetic and intervals}

\paragraph{Why the interaction matters.}
SSA-value sites are one of the places where linear arithmetic and intervals can become
genuinely useful to ordinary users. If the interaction is modeled poorly,
either the system loses semantic information or it asks for annotations that
should have been avoidable.

\paragraph{Concrete implementation rule.}
The cover builder should create enough overlap structure that the local
vocabulary of linear arithmetic and intervals can be restricted onto the relevant
coordinates of ssa-value sites. The local solver should emit not only
solved facts but also viability and transport metadata. The backward pass
should then decide whether the interaction contributes to safe input
synthesis, while the forward pass decides whether it contributes to output
enrichment.

\paragraph{Typical failure mode if omitted.}
Without this interaction the system is likely either to over-approximate and
lose useful guarantees, or to under-approximate and miss a reachable runtime
error. In both cases the user experience degrades because local evidence
fails to become global semantics.

\paragraph{Suggested test shape.}
Add one corpus example in which a fact from linear arithmetic and intervals is only
visible because of the presence of ssa-value sites. The test should
assert both a synthesized contract effect and a provenance explanation.

\subsection{Interaction: SSA-value sites with Boolean guard theory}

\paragraph{Why the interaction matters.}
SSA-value sites are one of the places where boolean guard theory can become
genuinely useful to ordinary users. If the interaction is modeled poorly,
either the system loses semantic information or it asks for annotations that
should have been avoidable.

\paragraph{Concrete implementation rule.}
The cover builder should create enough overlap structure that the local
vocabulary of boolean guard theory can be restricted onto the relevant
coordinates of ssa-value sites. The local solver should emit not only
solved facts but also viability and transport metadata. The backward pass
should then decide whether the interaction contributes to safe input
synthesis, while the forward pass decides whether it contributes to output
enrichment.

\paragraph{Typical failure mode if omitted.}
Without this interaction the system is likely either to over-approximate and
lose useful guarantees, or to under-approximate and miss a reachable runtime
error. In both cases the user experience degrades because local evidence
fails to become global semantics.

\paragraph{Suggested test shape.}
Add one corpus example in which a fact from boolean guard theory is only
visible because of the presence of ssa-value sites. The test should
assert both a synthesized contract effect and a provenance explanation.

\subsection{Interaction: SSA-value sites with Nullability and optionality}

\paragraph{Why the interaction matters.}
SSA-value sites are one of the places where nullability and optionality can become
genuinely useful to ordinary users. If the interaction is modeled poorly,
either the system loses semantic information or it asks for annotations that
should have been avoidable.

\paragraph{Concrete implementation rule.}
The cover builder should create enough overlap structure that the local
vocabulary of nullability and optionality can be restricted onto the relevant
coordinates of ssa-value sites. The local solver should emit not only
solved facts but also viability and transport metadata. The backward pass
should then decide whether the interaction contributes to safe input
synthesis, while the forward pass decides whether it contributes to output
enrichment.

\paragraph{Typical failure mode if omitted.}
Without this interaction the system is likely either to over-approximate and
lose useful guarantees, or to under-approximate and miss a reachable runtime
error. In both cases the user experience degrades because local evidence
fails to become global semantics.

\paragraph{Suggested test shape.}
Add one corpus example in which a fact from nullability and optionality is only
visible because of the presence of ssa-value sites. The test should
assert both a synthesized contract effect and a provenance explanation.

\subsection{Interaction: SSA-value sites with Heap and protocol theory}

\paragraph{Why the interaction matters.}
SSA-value sites are one of the places where heap and protocol theory can become
genuinely useful to ordinary users. If the interaction is modeled poorly,
either the system loses semantic information or it asks for annotations that
should have been avoidable.

\paragraph{Concrete implementation rule.}
The cover builder should create enough overlap structure that the local
vocabulary of heap and protocol theory can be restricted onto the relevant
coordinates of ssa-value sites. The local solver should emit not only
solved facts but also viability and transport metadata. The backward pass
should then decide whether the interaction contributes to safe input
synthesis, while the forward pass decides whether it contributes to output
enrichment.

\paragraph{Typical failure mode if omitted.}
Without this interaction the system is likely either to over-approximate and
lose useful guarantees, or to under-approximate and miss a reachable runtime
error. In both cases the user experience degrades because local evidence
fails to become global semantics.

\paragraph{Suggested test shape.}
Add one corpus example in which a fact from heap and protocol theory is only
visible because of the presence of ssa-value sites. The test should
assert both a synthesized contract effect and a provenance explanation.

\subsection{Interaction: SSA-value sites with Tensor shape theory}

\paragraph{Why the interaction matters.}
SSA-value sites are one of the places where tensor shape theory can become
genuinely useful to ordinary users. If the interaction is modeled poorly,
either the system loses semantic information or it asks for annotations that
should have been avoidable.

\paragraph{Concrete implementation rule.}
The cover builder should create enough overlap structure that the local
vocabulary of tensor shape theory can be restricted onto the relevant
coordinates of ssa-value sites. The local solver should emit not only
solved facts but also viability and transport metadata. The backward pass
should then decide whether the interaction contributes to safe input
synthesis, while the forward pass decides whether it contributes to output
enrichment.

\paragraph{Typical failure mode if omitted.}
Without this interaction the system is likely either to over-approximate and
lose useful guarantees, or to under-approximate and miss a reachable runtime
error. In both cases the user experience degrades because local evidence
fails to become global semantics.

\paragraph{Suggested test shape.}
Add one corpus example in which a fact from tensor shape theory is only
visible because of the presence of ssa-value sites. The test should
assert both a synthesized contract effect and a provenance explanation.

\subsection{Interaction: SSA-value sites with Tensor order theory}

\paragraph{Why the interaction matters.}
SSA-value sites are one of the places where tensor order theory can become
genuinely useful to ordinary users. If the interaction is modeled poorly,
either the system loses semantic information or it asks for annotations that
should have been avoidable.

\paragraph{Concrete implementation rule.}
The cover builder should create enough overlap structure that the local
vocabulary of tensor order theory can be restricted onto the relevant
coordinates of ssa-value sites. The local solver should emit not only
solved facts but also viability and transport metadata. The backward pass
should then decide whether the interaction contributes to safe input
synthesis, while the forward pass decides whether it contributes to output
enrichment.

\paragraph{Typical failure mode if omitted.}
Without this interaction the system is likely either to over-approximate and
lose useful guarantees, or to under-approximate and miss a reachable runtime
error. In both cases the user experience degrades because local evidence
fails to become global semantics.

\paragraph{Suggested test shape.}
Add one corpus example in which a fact from tensor order theory is only
visible because of the presence of ssa-value sites. The test should
assert both a synthesized contract effect and a provenance explanation.

\subsection{Interaction: SSA-value sites with Tensor indexing theory}

\paragraph{Why the interaction matters.}
SSA-value sites are one of the places where tensor indexing theory can become
genuinely useful to ordinary users. If the interaction is modeled poorly,
either the system loses semantic information or it asks for annotations that
should have been avoidable.

\paragraph{Concrete implementation rule.}
The cover builder should create enough overlap structure that the local
vocabulary of tensor indexing theory can be restricted onto the relevant
coordinates of ssa-value sites. The local solver should emit not only
solved facts but also viability and transport metadata. The backward pass
should then decide whether the interaction contributes to safe input
synthesis, while the forward pass decides whether it contributes to output
enrichment.

\paragraph{Typical failure mode if omitted.}
Without this interaction the system is likely either to over-approximate and
lose useful guarantees, or to under-approximate and miss a reachable runtime
error. In both cases the user experience degrades because local evidence
fails to become global semantics.

\paragraph{Suggested test shape.}
Add one corpus example in which a fact from tensor indexing theory is only
visible because of the presence of ssa-value sites. The test should
assert both a synthesized contract effect and a provenance explanation.

\subsection{Interaction: SSA-value sites with Device and dtype compatibility theory}

\paragraph{Why the interaction matters.}
SSA-value sites are one of the places where device and dtype compatibility theory can become
genuinely useful to ordinary users. If the interaction is modeled poorly,
either the system loses semantic information or it asks for annotations that
should have been avoidable.

\paragraph{Concrete implementation rule.}
The cover builder should create enough overlap structure that the local
vocabulary of device and dtype compatibility theory can be restricted onto the relevant
coordinates of ssa-value sites. The local solver should emit not only
solved facts but also viability and transport metadata. The backward pass
should then decide whether the interaction contributes to safe input
synthesis, while the forward pass decides whether it contributes to output
enrichment.

\paragraph{Typical failure mode if omitted.}
Without this interaction the system is likely either to over-approximate and
lose useful guarantees, or to under-approximate and miss a reachable runtime
error. In both cases the user experience degrades because local evidence
fails to become global semantics.

\paragraph{Suggested test shape.}
Add one corpus example in which a fact from device and dtype compatibility theory is only
visible because of the presence of ssa-value sites. The test should
assert both a synthesized contract effect and a provenance explanation.

\subsection{Interaction: SSA-value sites with Sequence and collection theory}

\paragraph{Why the interaction matters.}
SSA-value sites are one of the places where sequence and collection theory can become
genuinely useful to ordinary users. If the interaction is modeled poorly,
either the system loses semantic information or it asks for annotations that
should have been avoidable.

\paragraph{Concrete implementation rule.}
The cover builder should create enough overlap structure that the local
vocabulary of sequence and collection theory can be restricted onto the relevant
coordinates of ssa-value sites. The local solver should emit not only
solved facts but also viability and transport metadata. The backward pass
should then decide whether the interaction contributes to safe input
synthesis, while the forward pass decides whether it contributes to output
enrichment.

\paragraph{Typical failure mode if omitted.}
Without this interaction the system is likely either to over-approximate and
lose useful guarantees, or to under-approximate and miss a reachable runtime
error. In both cases the user experience degrades because local evidence
fails to become global semantics.

\paragraph{Suggested test shape.}
Add one corpus example in which a fact from sequence and collection theory is only
visible because of the presence of ssa-value sites. The test should
assert both a synthesized contract effect and a provenance explanation.

\subsection{Interaction: SSA-value sites with Exception theory}

\paragraph{Why the interaction matters.}
SSA-value sites are one of the places where exception theory can become
genuinely useful to ordinary users. If the interaction is modeled poorly,
either the system loses semantic information or it asks for annotations that
should have been avoidable.

\paragraph{Concrete implementation rule.}
The cover builder should create enough overlap structure that the local
vocabulary of exception theory can be restricted onto the relevant
coordinates of ssa-value sites. The local solver should emit not only
solved facts but also viability and transport metadata. The backward pass
should then decide whether the interaction contributes to safe input
synthesis, while the forward pass decides whether it contributes to output
enrichment.

\paragraph{Typical failure mode if omitted.}
Without this interaction the system is likely either to over-approximate and
lose useful guarantees, or to under-approximate and miss a reachable runtime
error. In both cases the user experience degrades because local evidence
fails to become global semantics.

\paragraph{Suggested test shape.}
Add one corpus example in which a fact from exception theory is only
visible because of the presence of ssa-value sites. The test should
assert both a synthesized contract effect and a provenance explanation.

\subsection{Interaction: SSA-value sites with Witness and equality transport theory}

\paragraph{Why the interaction matters.}
SSA-value sites are one of the places where witness and equality transport theory can become
genuinely useful to ordinary users. If the interaction is modeled poorly,
either the system loses semantic information or it asks for annotations that
should have been avoidable.

\paragraph{Concrete implementation rule.}
The cover builder should create enough overlap structure that the local
vocabulary of witness and equality transport theory can be restricted onto the relevant
coordinates of ssa-value sites. The local solver should emit not only
solved facts but also viability and transport metadata. The backward pass
should then decide whether the interaction contributes to safe input
synthesis, while the forward pass decides whether it contributes to output
enrichment.

\paragraph{Typical failure mode if omitted.}
Without this interaction the system is likely either to over-approximate and
lose useful guarantees, or to under-approximate and miss a reachable runtime
error. In both cases the user experience degrades because local evidence
fails to become global semantics.

\paragraph{Suggested test shape.}
Add one corpus example in which a fact from witness and equality transport theory is only
visible because of the presence of ssa-value sites. The test should
assert both a synthesized contract effect and a provenance explanation.

\subsection{Interaction: SSA-value sites with Loop invariant and decreases theory}

\paragraph{Why the interaction matters.}
SSA-value sites are one of the places where loop invariant and decreases theory can become
genuinely useful to ordinary users. If the interaction is modeled poorly,
either the system loses semantic information or it asks for annotations that
should have been avoidable.

\paragraph{Concrete implementation rule.}
The cover builder should create enough overlap structure that the local
vocabulary of loop invariant and decreases theory can be restricted onto the relevant
coordinates of ssa-value sites. The local solver should emit not only
solved facts but also viability and transport metadata. The backward pass
should then decide whether the interaction contributes to safe input
synthesis, while the forward pass decides whether it contributes to output
enrichment.

\paragraph{Typical failure mode if omitted.}
Without this interaction the system is likely either to over-approximate and
lose useful guarantees, or to under-approximate and miss a reachable runtime
error. In both cases the user experience degrades because local evidence
fails to become global semantics.

\paragraph{Suggested test shape.}
Add one corpus example in which a fact from loop invariant and decreases theory is only
visible because of the presence of ssa-value sites. The test should
assert both a synthesized contract effect and a provenance explanation.

\subsection{Interaction: Branch-guard sites with Linear arithmetic and intervals}

\paragraph{Why the interaction matters.}
Branch-guard sites are one of the places where linear arithmetic and intervals can become
genuinely useful to ordinary users. If the interaction is modeled poorly,
either the system loses semantic information or it asks for annotations that
should have been avoidable.

\paragraph{Concrete implementation rule.}
The cover builder should create enough overlap structure that the local
vocabulary of linear arithmetic and intervals can be restricted onto the relevant
coordinates of branch-guard sites. The local solver should emit not only
solved facts but also viability and transport metadata. The backward pass
should then decide whether the interaction contributes to safe input
synthesis, while the forward pass decides whether it contributes to output
enrichment.

\paragraph{Typical failure mode if omitted.}
Without this interaction the system is likely either to over-approximate and
lose useful guarantees, or to under-approximate and miss a reachable runtime
error. In both cases the user experience degrades because local evidence
fails to become global semantics.

\paragraph{Suggested test shape.}
Add one corpus example in which a fact from linear arithmetic and intervals is only
visible because of the presence of branch-guard sites. The test should
assert both a synthesized contract effect and a provenance explanation.

\subsection{Interaction: Branch-guard sites with Boolean guard theory}

\paragraph{Why the interaction matters.}
Branch-guard sites are one of the places where boolean guard theory can become
genuinely useful to ordinary users. If the interaction is modeled poorly,
either the system loses semantic information or it asks for annotations that
should have been avoidable.

\paragraph{Concrete implementation rule.}
The cover builder should create enough overlap structure that the local
vocabulary of boolean guard theory can be restricted onto the relevant
coordinates of branch-guard sites. The local solver should emit not only
solved facts but also viability and transport metadata. The backward pass
should then decide whether the interaction contributes to safe input
synthesis, while the forward pass decides whether it contributes to output
enrichment.

\paragraph{Typical failure mode if omitted.}
Without this interaction the system is likely either to over-approximate and
lose useful guarantees, or to under-approximate and miss a reachable runtime
error. In both cases the user experience degrades because local evidence
fails to become global semantics.

\paragraph{Suggested test shape.}
Add one corpus example in which a fact from boolean guard theory is only
visible because of the presence of branch-guard sites. The test should
assert both a synthesized contract effect and a provenance explanation.

\subsection{Interaction: Branch-guard sites with Nullability and optionality}

\paragraph{Why the interaction matters.}
Branch-guard sites are one of the places where nullability and optionality can become
genuinely useful to ordinary users. If the interaction is modeled poorly,
either the system loses semantic information or it asks for annotations that
should have been avoidable.

\paragraph{Concrete implementation rule.}
The cover builder should create enough overlap structure that the local
vocabulary of nullability and optionality can be restricted onto the relevant
coordinates of branch-guard sites. The local solver should emit not only
solved facts but also viability and transport metadata. The backward pass
should then decide whether the interaction contributes to safe input
synthesis, while the forward pass decides whether it contributes to output
enrichment.

\paragraph{Typical failure mode if omitted.}
Without this interaction the system is likely either to over-approximate and
lose useful guarantees, or to under-approximate and miss a reachable runtime
error. In both cases the user experience degrades because local evidence
fails to become global semantics.

\paragraph{Suggested test shape.}
Add one corpus example in which a fact from nullability and optionality is only
visible because of the presence of branch-guard sites. The test should
assert both a synthesized contract effect and a provenance explanation.

\subsection{Interaction: Branch-guard sites with Heap and protocol theory}

\paragraph{Why the interaction matters.}
Branch-guard sites are one of the places where heap and protocol theory can become
genuinely useful to ordinary users. If the interaction is modeled poorly,
either the system loses semantic information or it asks for annotations that
should have been avoidable.

\paragraph{Concrete implementation rule.}
The cover builder should create enough overlap structure that the local
vocabulary of heap and protocol theory can be restricted onto the relevant
coordinates of branch-guard sites. The local solver should emit not only
solved facts but also viability and transport metadata. The backward pass
should then decide whether the interaction contributes to safe input
synthesis, while the forward pass decides whether it contributes to output
enrichment.

\paragraph{Typical failure mode if omitted.}
Without this interaction the system is likely either to over-approximate and
lose useful guarantees, or to under-approximate and miss a reachable runtime
error. In both cases the user experience degrades because local evidence
fails to become global semantics.

\paragraph{Suggested test shape.}
Add one corpus example in which a fact from heap and protocol theory is only
visible because of the presence of branch-guard sites. The test should
assert both a synthesized contract effect and a provenance explanation.

\subsection{Interaction: Branch-guard sites with Tensor shape theory}

\paragraph{Why the interaction matters.}
Branch-guard sites are one of the places where tensor shape theory can become
genuinely useful to ordinary users. If the interaction is modeled poorly,
either the system loses semantic information or it asks for annotations that
should have been avoidable.

\paragraph{Concrete implementation rule.}
The cover builder should create enough overlap structure that the local
vocabulary of tensor shape theory can be restricted onto the relevant
coordinates of branch-guard sites. The local solver should emit not only
solved facts but also viability and transport metadata. The backward pass
should then decide whether the interaction contributes to safe input
synthesis, while the forward pass decides whether it contributes to output
enrichment.

\paragraph{Typical failure mode if omitted.}
Without this interaction the system is likely either to over-approximate and
lose useful guarantees, or to under-approximate and miss a reachable runtime
error. In both cases the user experience degrades because local evidence
fails to become global semantics.

\paragraph{Suggested test shape.}
Add one corpus example in which a fact from tensor shape theory is only
visible because of the presence of branch-guard sites. The test should
assert both a synthesized contract effect and a provenance explanation.

\subsection{Interaction: Branch-guard sites with Tensor order theory}

\paragraph{Why the interaction matters.}
Branch-guard sites are one of the places where tensor order theory can become
genuinely useful to ordinary users. If the interaction is modeled poorly,
either the system loses semantic information or it asks for annotations that
should have been avoidable.

\paragraph{Concrete implementation rule.}
The cover builder should create enough overlap structure that the local
vocabulary of tensor order theory can be restricted onto the relevant
coordinates of branch-guard sites. The local solver should emit not only
solved facts but also viability and transport metadata. The backward pass
should then decide whether the interaction contributes to safe input
synthesis, while the forward pass decides whether it contributes to output
enrichment.

\paragraph{Typical failure mode if omitted.}
Without this interaction the system is likely either to over-approximate and
lose useful guarantees, or to under-approximate and miss a reachable runtime
error. In both cases the user experience degrades because local evidence
fails to become global semantics.

\paragraph{Suggested test shape.}
Add one corpus example in which a fact from tensor order theory is only
visible because of the presence of branch-guard sites. The test should
assert both a synthesized contract effect and a provenance explanation.

\subsection{Interaction: Branch-guard sites with Tensor indexing theory}

\paragraph{Why the interaction matters.}
Branch-guard sites are one of the places where tensor indexing theory can become
genuinely useful to ordinary users. If the interaction is modeled poorly,
either the system loses semantic information or it asks for annotations that
should have been avoidable.

\paragraph{Concrete implementation rule.}
The cover builder should create enough overlap structure that the local
vocabulary of tensor indexing theory can be restricted onto the relevant
coordinates of branch-guard sites. The local solver should emit not only
solved facts but also viability and transport metadata. The backward pass
should then decide whether the interaction contributes to safe input
synthesis, while the forward pass decides whether it contributes to output
enrichment.

\paragraph{Typical failure mode if omitted.}
Without this interaction the system is likely either to over-approximate and
lose useful guarantees, or to under-approximate and miss a reachable runtime
error. In both cases the user experience degrades because local evidence
fails to become global semantics.

\paragraph{Suggested test shape.}
Add one corpus example in which a fact from tensor indexing theory is only
visible because of the presence of branch-guard sites. The test should
assert both a synthesized contract effect and a provenance explanation.

\subsection{Interaction: Branch-guard sites with Device and dtype compatibility theory}

\paragraph{Why the interaction matters.}
Branch-guard sites are one of the places where device and dtype compatibility theory can become
genuinely useful to ordinary users. If the interaction is modeled poorly,
either the system loses semantic information or it asks for annotations that
should have been avoidable.

\paragraph{Concrete implementation rule.}
The cover builder should create enough overlap structure that the local
vocabulary of device and dtype compatibility theory can be restricted onto the relevant
coordinates of branch-guard sites. The local solver should emit not only
solved facts but also viability and transport metadata. The backward pass
should then decide whether the interaction contributes to safe input
synthesis, while the forward pass decides whether it contributes to output
enrichment.

\paragraph{Typical failure mode if omitted.}
Without this interaction the system is likely either to over-approximate and
lose useful guarantees, or to under-approximate and miss a reachable runtime
error. In both cases the user experience degrades because local evidence
fails to become global semantics.

\paragraph{Suggested test shape.}
Add one corpus example in which a fact from device and dtype compatibility theory is only
visible because of the presence of branch-guard sites. The test should
assert both a synthesized contract effect and a provenance explanation.

\subsection{Interaction: Branch-guard sites with Sequence and collection theory}

\paragraph{Why the interaction matters.}
Branch-guard sites are one of the places where sequence and collection theory can become
genuinely useful to ordinary users. If the interaction is modeled poorly,
either the system loses semantic information or it asks for annotations that
should have been avoidable.

\paragraph{Concrete implementation rule.}
The cover builder should create enough overlap structure that the local
vocabulary of sequence and collection theory can be restricted onto the relevant
coordinates of branch-guard sites. The local solver should emit not only
solved facts but also viability and transport metadata. The backward pass
should then decide whether the interaction contributes to safe input
synthesis, while the forward pass decides whether it contributes to output
enrichment.

\paragraph{Typical failure mode if omitted.}
Without this interaction the system is likely either to over-approximate and
lose useful guarantees, or to under-approximate and miss a reachable runtime
error. In both cases the user experience degrades because local evidence
fails to become global semantics.

\paragraph{Suggested test shape.}
Add one corpus example in which a fact from sequence and collection theory is only
visible because of the presence of branch-guard sites. The test should
assert both a synthesized contract effect and a provenance explanation.

\subsection{Interaction: Branch-guard sites with Exception theory}

\paragraph{Why the interaction matters.}
Branch-guard sites are one of the places where exception theory can become
genuinely useful to ordinary users. If the interaction is modeled poorly,
either the system loses semantic information or it asks for annotations that
should have been avoidable.

\paragraph{Concrete implementation rule.}
The cover builder should create enough overlap structure that the local
vocabulary of exception theory can be restricted onto the relevant
coordinates of branch-guard sites. The local solver should emit not only
solved facts but also viability and transport metadata. The backward pass
should then decide whether the interaction contributes to safe input
synthesis, while the forward pass decides whether it contributes to output
enrichment.

\paragraph{Typical failure mode if omitted.}
Without this interaction the system is likely either to over-approximate and
lose useful guarantees, or to under-approximate and miss a reachable runtime
error. In both cases the user experience degrades because local evidence
fails to become global semantics.

\paragraph{Suggested test shape.}
Add one corpus example in which a fact from exception theory is only
visible because of the presence of branch-guard sites. The test should
assert both a synthesized contract effect and a provenance explanation.

\subsection{Interaction: Branch-guard sites with Witness and equality transport theory}

\paragraph{Why the interaction matters.}
Branch-guard sites are one of the places where witness and equality transport theory can become
genuinely useful to ordinary users. If the interaction is modeled poorly,
either the system loses semantic information or it asks for annotations that
should have been avoidable.

\paragraph{Concrete implementation rule.}
The cover builder should create enough overlap structure that the local
vocabulary of witness and equality transport theory can be restricted onto the relevant
coordinates of branch-guard sites. The local solver should emit not only
solved facts but also viability and transport metadata. The backward pass
should then decide whether the interaction contributes to safe input
synthesis, while the forward pass decides whether it contributes to output
enrichment.

\paragraph{Typical failure mode if omitted.}
Without this interaction the system is likely either to over-approximate and
lose useful guarantees, or to under-approximate and miss a reachable runtime
error. In both cases the user experience degrades because local evidence
fails to become global semantics.

\paragraph{Suggested test shape.}
Add one corpus example in which a fact from witness and equality transport theory is only
visible because of the presence of branch-guard sites. The test should
assert both a synthesized contract effect and a provenance explanation.

\subsection{Interaction: Branch-guard sites with Loop invariant and decreases theory}

\paragraph{Why the interaction matters.}
Branch-guard sites are one of the places where loop invariant and decreases theory can become
genuinely useful to ordinary users. If the interaction is modeled poorly,
either the system loses semantic information or it asks for annotations that
should have been avoidable.

\paragraph{Concrete implementation rule.}
The cover builder should create enough overlap structure that the local
vocabulary of loop invariant and decreases theory can be restricted onto the relevant
coordinates of branch-guard sites. The local solver should emit not only
solved facts but also viability and transport metadata. The backward pass
should then decide whether the interaction contributes to safe input
synthesis, while the forward pass decides whether it contributes to output
enrichment.

\paragraph{Typical failure mode if omitted.}
Without this interaction the system is likely either to over-approximate and
lose useful guarantees, or to under-approximate and miss a reachable runtime
error. In both cases the user experience degrades because local evidence
fails to become global semantics.

\paragraph{Suggested test shape.}
Add one corpus example in which a fact from loop invariant and decreases theory is only
visible because of the presence of branch-guard sites. The test should
assert both a synthesized contract effect and a provenance explanation.

\subsection{Interaction: Call-result sites with Linear arithmetic and intervals}

\paragraph{Why the interaction matters.}
Call-result sites are one of the places where linear arithmetic and intervals can become
genuinely useful to ordinary users. If the interaction is modeled poorly,
either the system loses semantic information or it asks for annotations that
should have been avoidable.

\paragraph{Concrete implementation rule.}
The cover builder should create enough overlap structure that the local
vocabulary of linear arithmetic and intervals can be restricted onto the relevant
coordinates of call-result sites. The local solver should emit not only
solved facts but also viability and transport metadata. The backward pass
should then decide whether the interaction contributes to safe input
synthesis, while the forward pass decides whether it contributes to output
enrichment.

\paragraph{Typical failure mode if omitted.}
Without this interaction the system is likely either to over-approximate and
lose useful guarantees, or to under-approximate and miss a reachable runtime
error. In both cases the user experience degrades because local evidence
fails to become global semantics.

\paragraph{Suggested test shape.}
Add one corpus example in which a fact from linear arithmetic and intervals is only
visible because of the presence of call-result sites. The test should
assert both a synthesized contract effect and a provenance explanation.

\subsection{Interaction: Call-result sites with Boolean guard theory}

\paragraph{Why the interaction matters.}
Call-result sites are one of the places where boolean guard theory can become
genuinely useful to ordinary users. If the interaction is modeled poorly,
either the system loses semantic information or it asks for annotations that
should have been avoidable.

\paragraph{Concrete implementation rule.}
The cover builder should create enough overlap structure that the local
vocabulary of boolean guard theory can be restricted onto the relevant
coordinates of call-result sites. The local solver should emit not only
solved facts but also viability and transport metadata. The backward pass
should then decide whether the interaction contributes to safe input
synthesis, while the forward pass decides whether it contributes to output
enrichment.

\paragraph{Typical failure mode if omitted.}
Without this interaction the system is likely either to over-approximate and
lose useful guarantees, or to under-approximate and miss a reachable runtime
error. In both cases the user experience degrades because local evidence
fails to become global semantics.

\paragraph{Suggested test shape.}
Add one corpus example in which a fact from boolean guard theory is only
visible because of the presence of call-result sites. The test should
assert both a synthesized contract effect and a provenance explanation.

\subsection{Interaction: Call-result sites with Nullability and optionality}

\paragraph{Why the interaction matters.}
Call-result sites are one of the places where nullability and optionality can become
genuinely useful to ordinary users. If the interaction is modeled poorly,
either the system loses semantic information or it asks for annotations that
should have been avoidable.

\paragraph{Concrete implementation rule.}
The cover builder should create enough overlap structure that the local
vocabulary of nullability and optionality can be restricted onto the relevant
coordinates of call-result sites. The local solver should emit not only
solved facts but also viability and transport metadata. The backward pass
should then decide whether the interaction contributes to safe input
synthesis, while the forward pass decides whether it contributes to output
enrichment.

\paragraph{Typical failure mode if omitted.}
Without this interaction the system is likely either to over-approximate and
lose useful guarantees, or to under-approximate and miss a reachable runtime
error. In both cases the user experience degrades because local evidence
fails to become global semantics.

\paragraph{Suggested test shape.}
Add one corpus example in which a fact from nullability and optionality is only
visible because of the presence of call-result sites. The test should
assert both a synthesized contract effect and a provenance explanation.

\subsection{Interaction: Call-result sites with Heap and protocol theory}

\paragraph{Why the interaction matters.}
Call-result sites are one of the places where heap and protocol theory can become
genuinely useful to ordinary users. If the interaction is modeled poorly,
either the system loses semantic information or it asks for annotations that
should have been avoidable.

\paragraph{Concrete implementation rule.}
The cover builder should create enough overlap structure that the local
vocabulary of heap and protocol theory can be restricted onto the relevant
coordinates of call-result sites. The local solver should emit not only
solved facts but also viability and transport metadata. The backward pass
should then decide whether the interaction contributes to safe input
synthesis, while the forward pass decides whether it contributes to output
enrichment.

\paragraph{Typical failure mode if omitted.}
Without this interaction the system is likely either to over-approximate and
lose useful guarantees, or to under-approximate and miss a reachable runtime
error. In both cases the user experience degrades because local evidence
fails to become global semantics.

\paragraph{Suggested test shape.}
Add one corpus example in which a fact from heap and protocol theory is only
visible because of the presence of call-result sites. The test should
assert both a synthesized contract effect and a provenance explanation.

\subsection{Interaction: Call-result sites with Tensor shape theory}

\paragraph{Why the interaction matters.}
Call-result sites are one of the places where tensor shape theory can become
genuinely useful to ordinary users. If the interaction is modeled poorly,
either the system loses semantic information or it asks for annotations that
should have been avoidable.

\paragraph{Concrete implementation rule.}
The cover builder should create enough overlap structure that the local
vocabulary of tensor shape theory can be restricted onto the relevant
coordinates of call-result sites. The local solver should emit not only
solved facts but also viability and transport metadata. The backward pass
should then decide whether the interaction contributes to safe input
synthesis, while the forward pass decides whether it contributes to output
enrichment.

\paragraph{Typical failure mode if omitted.}
Without this interaction the system is likely either to over-approximate and
lose useful guarantees, or to under-approximate and miss a reachable runtime
error. In both cases the user experience degrades because local evidence
fails to become global semantics.

\paragraph{Suggested test shape.}
Add one corpus example in which a fact from tensor shape theory is only
visible because of the presence of call-result sites. The test should
assert both a synthesized contract effect and a provenance explanation.

\subsection{Interaction: Call-result sites with Tensor order theory}

\paragraph{Why the interaction matters.}
Call-result sites are one of the places where tensor order theory can become
genuinely useful to ordinary users. If the interaction is modeled poorly,
either the system loses semantic information or it asks for annotations that
should have been avoidable.

\paragraph{Concrete implementation rule.}
The cover builder should create enough overlap structure that the local
vocabulary of tensor order theory can be restricted onto the relevant
coordinates of call-result sites. The local solver should emit not only
solved facts but also viability and transport metadata. The backward pass
should then decide whether the interaction contributes to safe input
synthesis, while the forward pass decides whether it contributes to output
enrichment.

\paragraph{Typical failure mode if omitted.}
Without this interaction the system is likely either to over-approximate and
lose useful guarantees, or to under-approximate and miss a reachable runtime
error. In both cases the user experience degrades because local evidence
fails to become global semantics.

\paragraph{Suggested test shape.}
Add one corpus example in which a fact from tensor order theory is only
visible because of the presence of call-result sites. The test should
assert both a synthesized contract effect and a provenance explanation.

\subsection{Interaction: Call-result sites with Tensor indexing theory}

\paragraph{Why the interaction matters.}
Call-result sites are one of the places where tensor indexing theory can become
genuinely useful to ordinary users. If the interaction is modeled poorly,
either the system loses semantic information or it asks for annotations that
should have been avoidable.

\paragraph{Concrete implementation rule.}
The cover builder should create enough overlap structure that the local
vocabulary of tensor indexing theory can be restricted onto the relevant
coordinates of call-result sites. The local solver should emit not only
solved facts but also viability and transport metadata. The backward pass
should then decide whether the interaction contributes to safe input
synthesis, while the forward pass decides whether it contributes to output
enrichment.

\paragraph{Typical failure mode if omitted.}
Without this interaction the system is likely either to over-approximate and
lose useful guarantees, or to under-approximate and miss a reachable runtime
error. In both cases the user experience degrades because local evidence
fails to become global semantics.

\paragraph{Suggested test shape.}
Add one corpus example in which a fact from tensor indexing theory is only
visible because of the presence of call-result sites. The test should
assert both a synthesized contract effect and a provenance explanation.

\subsection{Interaction: Call-result sites with Device and dtype compatibility theory}

\paragraph{Why the interaction matters.}
Call-result sites are one of the places where device and dtype compatibility theory can become
genuinely useful to ordinary users. If the interaction is modeled poorly,
either the system loses semantic information or it asks for annotations that
should have been avoidable.

\paragraph{Concrete implementation rule.}
The cover builder should create enough overlap structure that the local
vocabulary of device and dtype compatibility theory can be restricted onto the relevant
coordinates of call-result sites. The local solver should emit not only
solved facts but also viability and transport metadata. The backward pass
should then decide whether the interaction contributes to safe input
synthesis, while the forward pass decides whether it contributes to output
enrichment.

\paragraph{Typical failure mode if omitted.}
Without this interaction the system is likely either to over-approximate and
lose useful guarantees, or to under-approximate and miss a reachable runtime
error. In both cases the user experience degrades because local evidence
fails to become global semantics.

\paragraph{Suggested test shape.}
Add one corpus example in which a fact from device and dtype compatibility theory is only
visible because of the presence of call-result sites. The test should
assert both a synthesized contract effect and a provenance explanation.

\subsection{Interaction: Call-result sites with Sequence and collection theory}

\paragraph{Why the interaction matters.}
Call-result sites are one of the places where sequence and collection theory can become
genuinely useful to ordinary users. If the interaction is modeled poorly,
either the system loses semantic information or it asks for annotations that
should have been avoidable.

\paragraph{Concrete implementation rule.}
The cover builder should create enough overlap structure that the local
vocabulary of sequence and collection theory can be restricted onto the relevant
coordinates of call-result sites. The local solver should emit not only
solved facts but also viability and transport metadata. The backward pass
should then decide whether the interaction contributes to safe input
synthesis, while the forward pass decides whether it contributes to output
enrichment.

\paragraph{Typical failure mode if omitted.}
Without this interaction the system is likely either to over-approximate and
lose useful guarantees, or to under-approximate and miss a reachable runtime
error. In both cases the user experience degrades because local evidence
fails to become global semantics.

\paragraph{Suggested test shape.}
Add one corpus example in which a fact from sequence and collection theory is only
visible because of the presence of call-result sites. The test should
assert both a synthesized contract effect and a provenance explanation.

\subsection{Interaction: Call-result sites with Exception theory}

\paragraph{Why the interaction matters.}
Call-result sites are one of the places where exception theory can become
genuinely useful to ordinary users. If the interaction is modeled poorly,
either the system loses semantic information or it asks for annotations that
should have been avoidable.

\paragraph{Concrete implementation rule.}
The cover builder should create enough overlap structure that the local
vocabulary of exception theory can be restricted onto the relevant
coordinates of call-result sites. The local solver should emit not only
solved facts but also viability and transport metadata. The backward pass
should then decide whether the interaction contributes to safe input
synthesis, while the forward pass decides whether it contributes to output
enrichment.

\paragraph{Typical failure mode if omitted.}
Without this interaction the system is likely either to over-approximate and
lose useful guarantees, or to under-approximate and miss a reachable runtime
error. In both cases the user experience degrades because local evidence
fails to become global semantics.

\paragraph{Suggested test shape.}
Add one corpus example in which a fact from exception theory is only
visible because of the presence of call-result sites. The test should
assert both a synthesized contract effect and a provenance explanation.

\subsection{Interaction: Call-result sites with Witness and equality transport theory}

\paragraph{Why the interaction matters.}
Call-result sites are one of the places where witness and equality transport theory can become
genuinely useful to ordinary users. If the interaction is modeled poorly,
either the system loses semantic information or it asks for annotations that
should have been avoidable.

\paragraph{Concrete implementation rule.}
The cover builder should create enough overlap structure that the local
vocabulary of witness and equality transport theory can be restricted onto the relevant
coordinates of call-result sites. The local solver should emit not only
solved facts but also viability and transport metadata. The backward pass
should then decide whether the interaction contributes to safe input
synthesis, while the forward pass decides whether it contributes to output
enrichment.

\paragraph{Typical failure mode if omitted.}
Without this interaction the system is likely either to over-approximate and
lose useful guarantees, or to under-approximate and miss a reachable runtime
error. In both cases the user experience degrades because local evidence
fails to become global semantics.

\paragraph{Suggested test shape.}
Add one corpus example in which a fact from witness and equality transport theory is only
visible because of the presence of call-result sites. The test should
assert both a synthesized contract effect and a provenance explanation.

\subsection{Interaction: Call-result sites with Loop invariant and decreases theory}

\paragraph{Why the interaction matters.}
Call-result sites are one of the places where loop invariant and decreases theory can become
genuinely useful to ordinary users. If the interaction is modeled poorly,
either the system loses semantic information or it asks for annotations that
should have been avoidable.

\paragraph{Concrete implementation rule.}
The cover builder should create enough overlap structure that the local
vocabulary of loop invariant and decreases theory can be restricted onto the relevant
coordinates of call-result sites. The local solver should emit not only
solved facts but also viability and transport metadata. The backward pass
should then decide whether the interaction contributes to safe input
synthesis, while the forward pass decides whether it contributes to output
enrichment.

\paragraph{Typical failure mode if omitted.}
Without this interaction the system is likely either to over-approximate and
lose useful guarantees, or to under-approximate and miss a reachable runtime
error. In both cases the user experience degrades because local evidence
fails to become global semantics.

\paragraph{Suggested test shape.}
Add one corpus example in which a fact from loop invariant and decreases theory is only
visible because of the presence of call-result sites. The test should
assert both a synthesized contract effect and a provenance explanation.

\subsection{Interaction: Tensor-shape sites with Linear arithmetic and intervals}

\paragraph{Why the interaction matters.}
Tensor-shape sites are one of the places where linear arithmetic and intervals can become
genuinely useful to ordinary users. If the interaction is modeled poorly,
either the system loses semantic information or it asks for annotations that
should have been avoidable.

\paragraph{Concrete implementation rule.}
The cover builder should create enough overlap structure that the local
vocabulary of linear arithmetic and intervals can be restricted onto the relevant
coordinates of tensor-shape sites. The local solver should emit not only
solved facts but also viability and transport metadata. The backward pass
should then decide whether the interaction contributes to safe input
synthesis, while the forward pass decides whether it contributes to output
enrichment.

\paragraph{Typical failure mode if omitted.}
Without this interaction the system is likely either to over-approximate and
lose useful guarantees, or to under-approximate and miss a reachable runtime
error. In both cases the user experience degrades because local evidence
fails to become global semantics.

\paragraph{Suggested test shape.}
Add one corpus example in which a fact from linear arithmetic and intervals is only
visible because of the presence of tensor-shape sites. The test should
assert both a synthesized contract effect and a provenance explanation.

\subsection{Interaction: Tensor-shape sites with Boolean guard theory}

\paragraph{Why the interaction matters.}
Tensor-shape sites are one of the places where boolean guard theory can become
genuinely useful to ordinary users. If the interaction is modeled poorly,
either the system loses semantic information or it asks for annotations that
should have been avoidable.

\paragraph{Concrete implementation rule.}
The cover builder should create enough overlap structure that the local
vocabulary of boolean guard theory can be restricted onto the relevant
coordinates of tensor-shape sites. The local solver should emit not only
solved facts but also viability and transport metadata. The backward pass
should then decide whether the interaction contributes to safe input
synthesis, while the forward pass decides whether it contributes to output
enrichment.

\paragraph{Typical failure mode if omitted.}
Without this interaction the system is likely either to over-approximate and
lose useful guarantees, or to under-approximate and miss a reachable runtime
error. In both cases the user experience degrades because local evidence
fails to become global semantics.

\paragraph{Suggested test shape.}
Add one corpus example in which a fact from boolean guard theory is only
visible because of the presence of tensor-shape sites. The test should
assert both a synthesized contract effect and a provenance explanation.

\subsection{Interaction: Tensor-shape sites with Nullability and optionality}

\paragraph{Why the interaction matters.}
Tensor-shape sites are one of the places where nullability and optionality can become
genuinely useful to ordinary users. If the interaction is modeled poorly,
either the system loses semantic information or it asks for annotations that
should have been avoidable.

\paragraph{Concrete implementation rule.}
The cover builder should create enough overlap structure that the local
vocabulary of nullability and optionality can be restricted onto the relevant
coordinates of tensor-shape sites. The local solver should emit not only
solved facts but also viability and transport metadata. The backward pass
should then decide whether the interaction contributes to safe input
synthesis, while the forward pass decides whether it contributes to output
enrichment.

\paragraph{Typical failure mode if omitted.}
Without this interaction the system is likely either to over-approximate and
lose useful guarantees, or to under-approximate and miss a reachable runtime
error. In both cases the user experience degrades because local evidence
fails to become global semantics.

\paragraph{Suggested test shape.}
Add one corpus example in which a fact from nullability and optionality is only
visible because of the presence of tensor-shape sites. The test should
assert both a synthesized contract effect and a provenance explanation.

\subsection{Interaction: Tensor-shape sites with Heap and protocol theory}

\paragraph{Why the interaction matters.}
Tensor-shape sites are one of the places where heap and protocol theory can become
genuinely useful to ordinary users. If the interaction is modeled poorly,
either the system loses semantic information or it asks for annotations that
should have been avoidable.

\paragraph{Concrete implementation rule.}
The cover builder should create enough overlap structure that the local
vocabulary of heap and protocol theory can be restricted onto the relevant
coordinates of tensor-shape sites. The local solver should emit not only
solved facts but also viability and transport metadata. The backward pass
should then decide whether the interaction contributes to safe input
synthesis, while the forward pass decides whether it contributes to output
enrichment.

\paragraph{Typical failure mode if omitted.}
Without this interaction the system is likely either to over-approximate and
lose useful guarantees, or to under-approximate and miss a reachable runtime
error. In both cases the user experience degrades because local evidence
fails to become global semantics.

\paragraph{Suggested test shape.}
Add one corpus example in which a fact from heap and protocol theory is only
visible because of the presence of tensor-shape sites. The test should
assert both a synthesized contract effect and a provenance explanation.

\subsection{Interaction: Tensor-shape sites with Tensor shape theory}

\paragraph{Why the interaction matters.}
Tensor-shape sites are one of the places where tensor shape theory can become
genuinely useful to ordinary users. If the interaction is modeled poorly,
either the system loses semantic information or it asks for annotations that
should have been avoidable.

\paragraph{Concrete implementation rule.}
The cover builder should create enough overlap structure that the local
vocabulary of tensor shape theory can be restricted onto the relevant
coordinates of tensor-shape sites. The local solver should emit not only
solved facts but also viability and transport metadata. The backward pass
should then decide whether the interaction contributes to safe input
synthesis, while the forward pass decides whether it contributes to output
enrichment.

\paragraph{Typical failure mode if omitted.}
Without this interaction the system is likely either to over-approximate and
lose useful guarantees, or to under-approximate and miss a reachable runtime
error. In both cases the user experience degrades because local evidence
fails to become global semantics.

\paragraph{Suggested test shape.}
Add one corpus example in which a fact from tensor shape theory is only
visible because of the presence of tensor-shape sites. The test should
assert both a synthesized contract effect and a provenance explanation.

\subsection{Interaction: Tensor-shape sites with Tensor order theory}

\paragraph{Why the interaction matters.}
Tensor-shape sites are one of the places where tensor order theory can become
genuinely useful to ordinary users. If the interaction is modeled poorly,
either the system loses semantic information or it asks for annotations that
should have been avoidable.

\paragraph{Concrete implementation rule.}
The cover builder should create enough overlap structure that the local
vocabulary of tensor order theory can be restricted onto the relevant
coordinates of tensor-shape sites. The local solver should emit not only
solved facts but also viability and transport metadata. The backward pass
should then decide whether the interaction contributes to safe input
synthesis, while the forward pass decides whether it contributes to output
enrichment.

\paragraph{Typical failure mode if omitted.}
Without this interaction the system is likely either to over-approximate and
lose useful guarantees, or to under-approximate and miss a reachable runtime
error. In both cases the user experience degrades because local evidence
fails to become global semantics.

\paragraph{Suggested test shape.}
Add one corpus example in which a fact from tensor order theory is only
visible because of the presence of tensor-shape sites. The test should
assert both a synthesized contract effect and a provenance explanation.

\subsection{Interaction: Tensor-shape sites with Tensor indexing theory}

\paragraph{Why the interaction matters.}
Tensor-shape sites are one of the places where tensor indexing theory can become
genuinely useful to ordinary users. If the interaction is modeled poorly,
either the system loses semantic information or it asks for annotations that
should have been avoidable.

\paragraph{Concrete implementation rule.}
The cover builder should create enough overlap structure that the local
vocabulary of tensor indexing theory can be restricted onto the relevant
coordinates of tensor-shape sites. The local solver should emit not only
solved facts but also viability and transport metadata. The backward pass
should then decide whether the interaction contributes to safe input
synthesis, while the forward pass decides whether it contributes to output
enrichment.

\paragraph{Typical failure mode if omitted.}
Without this interaction the system is likely either to over-approximate and
lose useful guarantees, or to under-approximate and miss a reachable runtime
error. In both cases the user experience degrades because local evidence
fails to become global semantics.

\paragraph{Suggested test shape.}
Add one corpus example in which a fact from tensor indexing theory is only
visible because of the presence of tensor-shape sites. The test should
assert both a synthesized contract effect and a provenance explanation.

\subsection{Interaction: Tensor-shape sites with Device and dtype compatibility theory}

\paragraph{Why the interaction matters.}
Tensor-shape sites are one of the places where device and dtype compatibility theory can become
genuinely useful to ordinary users. If the interaction is modeled poorly,
either the system loses semantic information or it asks for annotations that
should have been avoidable.

\paragraph{Concrete implementation rule.}
The cover builder should create enough overlap structure that the local
vocabulary of device and dtype compatibility theory can be restricted onto the relevant
coordinates of tensor-shape sites. The local solver should emit not only
solved facts but also viability and transport metadata. The backward pass
should then decide whether the interaction contributes to safe input
synthesis, while the forward pass decides whether it contributes to output
enrichment.

\paragraph{Typical failure mode if omitted.}
Without this interaction the system is likely either to over-approximate and
lose useful guarantees, or to under-approximate and miss a reachable runtime
error. In both cases the user experience degrades because local evidence
fails to become global semantics.

\paragraph{Suggested test shape.}
Add one corpus example in which a fact from device and dtype compatibility theory is only
visible because of the presence of tensor-shape sites. The test should
assert both a synthesized contract effect and a provenance explanation.

\subsection{Interaction: Tensor-shape sites with Sequence and collection theory}

\paragraph{Why the interaction matters.}
Tensor-shape sites are one of the places where sequence and collection theory can become
genuinely useful to ordinary users. If the interaction is modeled poorly,
either the system loses semantic information or it asks for annotations that
should have been avoidable.

\paragraph{Concrete implementation rule.}
The cover builder should create enough overlap structure that the local
vocabulary of sequence and collection theory can be restricted onto the relevant
coordinates of tensor-shape sites. The local solver should emit not only
solved facts but also viability and transport metadata. The backward pass
should then decide whether the interaction contributes to safe input
synthesis, while the forward pass decides whether it contributes to output
enrichment.

\paragraph{Typical failure mode if omitted.}
Without this interaction the system is likely either to over-approximate and
lose useful guarantees, or to under-approximate and miss a reachable runtime
error. In both cases the user experience degrades because local evidence
fails to become global semantics.

\paragraph{Suggested test shape.}
Add one corpus example in which a fact from sequence and collection theory is only
visible because of the presence of tensor-shape sites. The test should
assert both a synthesized contract effect and a provenance explanation.

\subsection{Interaction: Tensor-shape sites with Exception theory}

\paragraph{Why the interaction matters.}
Tensor-shape sites are one of the places where exception theory can become
genuinely useful to ordinary users. If the interaction is modeled poorly,
either the system loses semantic information or it asks for annotations that
should have been avoidable.

\paragraph{Concrete implementation rule.}
The cover builder should create enough overlap structure that the local
vocabulary of exception theory can be restricted onto the relevant
coordinates of tensor-shape sites. The local solver should emit not only
solved facts but also viability and transport metadata. The backward pass
should then decide whether the interaction contributes to safe input
synthesis, while the forward pass decides whether it contributes to output
enrichment.

\paragraph{Typical failure mode if omitted.}
Without this interaction the system is likely either to over-approximate and
lose useful guarantees, or to under-approximate and miss a reachable runtime
error. In both cases the user experience degrades because local evidence
fails to become global semantics.

\paragraph{Suggested test shape.}
Add one corpus example in which a fact from exception theory is only
visible because of the presence of tensor-shape sites. The test should
assert both a synthesized contract effect and a provenance explanation.


\section{Evaluation, benchmarks, and success criteria}

A monograph of this scope should end with measurable criteria. The rewrite
succeeds if it improves four things simultaneously. First, annotation economy:
ordinary users should obtain strong contracts and strong error explanations with
little or no added text. Second, semantic richness: the analyzer should recover
input refinements, output guarantees, local intermediate sections, and theorem-
level transport effects that the current atom-centric architecture would either
miss or flatten. Third, honesty: the system must clearly separate solver-backed,
proof-backed, trace-backed, and residual coordinates. Fourth, migration
plausibility: new artifacts such as site covers, local sections, and obstruction
bases must fit naturally around the current packages and example corpus.

Benchmarking should therefore record more than pass/fail. Useful metrics include:
number of constructed sites after sparsification; overlap density; number of
error-sensitive sites proved empty under synthesized preconditions; size of the
minimal safe input section; information score of the maximal output section;
number of local intermediate values carrying nontrivial sections; count of
residual obstruction coordinates before and after user seeds; leverage of proof
artifacts on unannotated client modules; stability of synthesized contracts under
small wrapper refactorings; and runtime overhead of cover construction, local
solving, and gluing. For Torch specifically, the benchmark suite should measure
how often the system prevents library-specific errors without manual annotation
and how often it recovers semantically rich output summaries such as sortedness,
selector provenance, or exact shape.

The final success criterion is conceptual: users should come away believing that
native Python can support serious dependent reasoning without ceasing to be
native Python. If the sheaf-descent rewrite works, the explanation will be
simple. Rich semantics becomes possible not because the tool guessed a lucky set
of predicates, but because the program, proofs, library calls, and traces jointly
determine a global semantic object whose ambiguity can be measured, reduced, and
rendered honestly.


\section{Staged migration plan}

The migration plan should be staged so that every phase lands with stable tests
and without invalidating the current user-facing pipeline.

\begin{enumerate}[label=\textbf{Phase \arabic*:}]
  \item Introduce explicit site-cover and local-section datatypes in the core,
  with compatibility shims from current semantic signatures.
  \item Build cover synthesis from the existing AST, CFG, SSA, and lineage
  infrastructure, and snapshot the resulting cover on selected examples.
  \item Upgrade split PyTorch packs to export site constructors and viability
  hooks while preserving current example behavior.
  \item Add backward safe-input synthesis and forward output synthesis on a small
  family of site kinds: sort, shape, and indexing.
  \item Generalize runtime-error sites across standard Python and heap/protocol
  reasoning.
  \item Compile theorem modules into transport maps and proof sites.
  \item Replace or wrap current frontier extraction with section-summary
  frontiers and obstruction-aware repairs.
  \item Migrate renderers so every analysis can be viewed as synthesized Python
  contracts, diagnostics, proof goals, or telemetry.
  \item Expand the corpus and benchmark suite until the new architecture is the
  default path for both Torch-heavy and ordinary-Python examples.
\end{enumerate}

Each phase should preserve one invariant: if the new architecture makes a claim,
that claim must be at least as explicit about trust status as the current one.
This invariant matters more than temporary performance regressions, because once
trust becomes blurry the rest of the design loses its point.


\section{Extended implementation profiles}

This appendix exists to make the implementation story extremely explicit. It records package-theory pairings that matter for a realistic build.


\subsection{Package \py{src/deppy/assist} with Linear arithmetic and intervals}

\paragraph{Integration objective.}
The package \py{src/deppy/assist} must interact cleanly with the
theory family linear arithmetic and intervals. The reason is practical: the
package owns part of the data path through which non-negativity, bounds, lengths, simple affine equalities becomes
visible, audited, and eventually rendered.

\paragraph{Required interface.}
At minimum the integration should expose normalized inputs, site-aware
outputs, provenance, and failure reasons. The package should neither hide
theory-relevant structure from the cover builder nor flatten theory
results so aggressively that renderers lose meaning.

\paragraph{Typical evidence source.}
In practice the evidence will come from argument bounds, list lengths, loop counters, and tensor extents. The integration should
preserve enough of that source structure that local solvers using local SMT over arithmetic fragments
can operate without pack-specific hacks.

\paragraph{Expected user-facing payoff.}
If implemented well, the payoff is IndexError, ZeroDivisionError, and shape-side precondition pullback. This should appear not only
in benchmark metrics but also in concrete rendered contracts,
diagnostics, or proof goals.


\subsection{Package \py{src/deppy/assist} with Boolean guard theory}

\paragraph{Integration objective.}
The package \py{src/deppy/assist} must interact cleanly with the
theory family boolean guard theory. The reason is practical: the
package owns part of the data path through which guard polarity, assertion scope, and branch-conditioned refinement becomes
visible, audited, and eventually rendered.

\paragraph{Required interface.}
At minimum the integration should expose normalized inputs, site-aware
outputs, provenance, and failure reasons. The package should neither hide
theory-relevant structure from the cover builder nor flatten theory
results so aggressively that renderers lose meaning.

\paragraph{Typical evidence source.}
In practice the evidence will come from if-statements, asserts, early returns, and exception guards. The integration should
preserve enough of that source structure that local solvers using normalization plus implication checks
can operate without pack-specific hacks.

\paragraph{Expected user-facing payoff.}
If implemented well, the payoff is control-dependent safety and lightweight proof seed extraction. This should appear not only
in benchmark metrics but also in concrete rendered contracts,
diagnostics, or proof goals.


\subsection{Package \py{src/deppy/assist} with Nullability and optionality}

\paragraph{Integration objective.}
The package \py{src/deppy/assist} must interact cleanly with the
theory family nullability and optionality. The reason is practical: the
package owns part of the data path through which none checks, existence guards, and optional protocol availability becomes
visible, audited, and eventually rendered.

\paragraph{Required interface.}
At minimum the integration should expose normalized inputs, site-aware
outputs, provenance, and failure reasons. The package should neither hide
theory-relevant structure from the cover builder nor flatten theory
results so aggressively that renderers lose meaning.

\paragraph{Typical evidence source.}
In practice the evidence will come from ordinary Python branches and constructor code. The integration should
preserve enough of that source structure that local solvers using simple symbolic narrowing
can operate without pack-specific hacks.

\paragraph{Expected user-facing payoff.}
If implemented well, the payoff is AttributeError prevention and stronger post-branch outputs. This should appear not only
in benchmark metrics but also in concrete rendered contracts,
diagnostics, or proof goals.


\subsection{Package \py{src/deppy/assist} with Heap and protocol theory}

\paragraph{Integration objective.}
The package \py{src/deppy/assist} must interact cleanly with the
theory family heap and protocol theory. The reason is practical: the
package owns part of the data path through which field initialization, protocol method existence, frame conditions, and alias exclusions becomes
visible, audited, and eventually rendered.

\paragraph{Required interface.}
At minimum the integration should expose normalized inputs, site-aware
outputs, provenance, and failure reasons. The package should neither hide
theory-relevant structure from the cover builder nor flatten theory
results so aggressively that renderers lose meaning.

\paragraph{Typical evidence source.}
In practice the evidence will come from classes, protocols, mutable objects, and resource wrappers. The integration should
preserve enough of that source structure that local solvers using bounded symbolic checks plus explicit residuals
can operate without pack-specific hacks.

\paragraph{Expected user-facing payoff.}
If implemented well, the payoff is constructor safety, attribute access, and method-dispatch correctness. This should appear not only
in benchmark metrics but also in concrete rendered contracts,
diagnostics, or proof goals.


\subsection{Package \py{src/deppy/assist} with Tensor shape theory}

\paragraph{Integration objective.}
The package \py{src/deppy/assist} must interact cleanly with the
theory family tensor shape theory. The reason is practical: the
package owns part of the data path through which rank, dimension, shape tuples, cardinality, broadcast legality, and reduction support becomes
visible, audited, and eventually rendered.

\paragraph{Required interface.}
At minimum the integration should expose normalized inputs, site-aware
outputs, provenance, and failure reasons. The package should neither hide
theory-relevant structure from the cover builder nor flatten theory
results so aggressively that renderers lose meaning.

\paragraph{Typical evidence source.}
In practice the evidence will come from Torch tensor operators and shape-sensitive wrappers. The integration should
preserve enough of that source structure that local solvers using solver-backed local shape kernels
can operate without pack-specific hacks.

\paragraph{Expected user-facing payoff.}
If implemented well, the payoff is invalid reshape, broadcast mismatch, and exact output-shape synthesis. This should appear not only
in benchmark metrics but also in concrete rendered contracts,
diagnostics, or proof goals.


\subsection{Package \py{src/deppy/assist} with Tensor order theory}

\paragraph{Integration objective.}
The package \py{src/deppy/assist} must interact cleanly with the
theory family tensor order theory. The reason is practical: the
package owns part of the data path through which sortedness, monotonicity, permutation witnesses, argsort relations becomes
visible, audited, and eventually rendered.

\paragraph{Required interface.}
At minimum the integration should expose normalized inputs, site-aware
outputs, provenance, and failure reasons. The package should neither hide
theory-relevant structure from the cover builder nor flatten theory
results so aggressively that renderers lose meaning.

\paragraph{Typical evidence source.}
In practice the evidence will come from sort-like operators and order-preserving wrappers. The integration should
preserve enough of that source structure that local solvers using finite witness templates plus solver-backed subchecks
can operate without pack-specific hacks.

\paragraph{Expected user-facing payoff.}
If implemented well, the payoff is values-vs-indices confusion and richer sorted outputs. This should appear not only
in benchmark metrics but also in concrete rendered contracts,
diagnostics, or proof goals.


\subsection{Package \py{src/deppy/assist} with Tensor indexing theory}

\paragraph{Integration objective.}
The package \py{src/deppy/assist} must interact cleanly with the
theory family tensor indexing theory. The reason is practical: the
package owns part of the data path through which selector domains, slice extents, gather/index-select legality, and value provenance becomes
visible, audited, and eventually rendered.

\paragraph{Required interface.}
At minimum the integration should expose normalized inputs, site-aware
outputs, provenance, and failure reasons. The package should neither hide
theory-relevant structure from the cover builder nor flatten theory
results so aggressively that renderers lose meaning.

\paragraph{Typical evidence source.}
In practice the evidence will come from subscripts, slices, and advanced indexing. The integration should
preserve enough of that source structure that local solvers using arithmetic plus pack-specific rules
can operate without pack-specific hacks.

\paragraph{Expected user-facing payoff.}
If implemented well, the payoff is out-of-range selector detection and source-target output relations. This should appear not only
in benchmark metrics but also in concrete rendered contracts,
diagnostics, or proof goals.


\subsection{Package \py{src/deppy/assist} with Device and dtype compatibility theory}

\paragraph{Integration objective.}
The package \py{src/deppy/assist} must interact cleanly with the
theory family device and dtype compatibility theory. The reason is practical: the
package owns part of the data path through which device equality, admissible implicit casts, and operator support matrices becomes
visible, audited, and eventually rendered.

\paragraph{Required interface.}
At minimum the integration should expose normalized inputs, site-aware
outputs, provenance, and failure reasons. The package should neither hide
theory-relevant structure from the cover builder nor flatten theory
results so aggressively that renderers lose meaning.

\paragraph{Typical evidence source.}
In practice the evidence will come from mixed-device or mixed-dtype Torch programs. The integration should
preserve enough of that source structure that local solvers using table-driven local checks
can operate without pack-specific hacks.

\paragraph{Expected user-facing payoff.}
If implemented well, the payoff is device mismatch and unsupported operator combinations. This should appear not only
in benchmark metrics but also in concrete rendered contracts,
diagnostics, or proof goals.


\subsection{Package \py{src/deppy/assist} with Sequence and collection theory}

\paragraph{Integration objective.}
The package \py{src/deppy/assist} must interact cleanly with the
theory family sequence and collection theory. The reason is practical: the
package owns part of the data path through which non-emptiness, membership, stable length relations, and head/tail legality becomes
visible, audited, and eventually rendered.

\paragraph{Required interface.}
At minimum the integration should expose normalized inputs, site-aware
outputs, provenance, and failure reasons. The package should neither hide
theory-relevant structure from the cover builder nor flatten theory
results so aggressively that renderers lose meaning.

\paragraph{Typical evidence source.}
In practice the evidence will come from lists, tuples, deques, and iterator adapters. The integration should
preserve enough of that source structure that local solvers using simple algebraic checks
can operate without pack-specific hacks.

\paragraph{Expected user-facing payoff.}
If implemented well, the payoff is head-of-empty and unpacking safety. This should appear not only
in benchmark metrics but also in concrete rendered contracts,
diagnostics, or proof goals.


\subsection{Package \py{src/deppy/assist} with Exception theory}

\paragraph{Integration objective.}
The package \py{src/deppy/assist} must interact cleanly with the
theory family exception theory. The reason is practical: the
package owns part of the data path through which exceptional output channels, throw sites, and caught-vs-uncaught distinctions becomes
visible, audited, and eventually rendered.

\paragraph{Required interface.}
At minimum the integration should expose normalized inputs, site-aware
outputs, provenance, and failure reasons. The package should neither hide
theory-relevant structure from the cover builder nor flatten theory
results so aggressively that renderers lose meaning.

\paragraph{Typical evidence source.}
In practice the evidence will come from try/except blocks, library raises, and proof obligations that may fail. The integration should
preserve enough of that source structure that local solvers using control-flow plus error-site hooks
can operate without pack-specific hacks.

\paragraph{Expected user-facing payoff.}
If implemented well, the payoff is precise reporting of exceptional contracts and remaining failures. This should appear not only
in benchmark metrics but also in concrete rendered contracts,
diagnostics, or proof goals.


\subsection{Package \py{src/deppy/assist} with Witness and equality transport theory}

\paragraph{Integration objective.}
The package \py{src/deppy/assist} must interact cleanly with the
theory family witness and equality transport theory. The reason is practical: the
package owns part of the data path through which dependent pairs, equality witnesses, and explicit transports becomes
visible, audited, and eventually rendered.

\paragraph{Required interface.}
At minimum the integration should expose normalized inputs, site-aware
outputs, provenance, and failure reasons. The package should neither hide
theory-relevant structure from the cover builder nor flatten theory
results so aggressively that renderers lose meaning.

\paragraph{Typical evidence source.}
In practice the evidence will come from proof-oriented code and witness-returning helpers. The integration should
preserve enough of that source structure that local solvers using proof-kernel backed local checking
can operate without pack-specific hacks.

\paragraph{Expected user-facing payoff.}
If implemented well, the payoff is proof reuse and output enrichment. This should appear not only
in benchmark metrics but also in concrete rendered contracts,
diagnostics, or proof goals.


\subsection{Package \py{src/deppy/assist} with Loop invariant and decreases theory}

\paragraph{Integration objective.}
The package \py{src/deppy/assist} must interact cleanly with the
theory family loop invariant and decreases theory. The reason is practical: the
package owns part of the data path through which candidate invariants, ranking relations, and preservation constraints becomes
visible, audited, and eventually rendered.

\paragraph{Required interface.}
At minimum the integration should expose normalized inputs, site-aware
outputs, provenance, and failure reasons. The package should neither hide
theory-relevant structure from the cover builder nor flatten theory
results so aggressively that renderers lose meaning.

\paragraph{Typical evidence source.}
In practice the evidence will come from loop-heavy algorithms and recursive helpers. The integration should
preserve enough of that source structure that local solvers using template generation plus bounded checks
can operate without pack-specific hacks.

\paragraph{Expected user-facing payoff.}
If implemented well, the payoff is F*-like capability without abandoning native Python. This should appear not only
in benchmark metrics but also in concrete rendered contracts,
diagnostics, or proof goals.


\subsection{Package \py{src/deppy/contracts} with Linear arithmetic and intervals}

\paragraph{Integration objective.}
The package \py{src/deppy/contracts} must interact cleanly with the
theory family linear arithmetic and intervals. The reason is practical: the
package owns part of the data path through which non-negativity, bounds, lengths, simple affine equalities becomes
visible, audited, and eventually rendered.

\paragraph{Required interface.}
At minimum the integration should expose normalized inputs, site-aware
outputs, provenance, and failure reasons. The package should neither hide
theory-relevant structure from the cover builder nor flatten theory
results so aggressively that renderers lose meaning.

\paragraph{Typical evidence source.}
In practice the evidence will come from argument bounds, list lengths, loop counters, and tensor extents. The integration should
preserve enough of that source structure that local solvers using local SMT over arithmetic fragments
can operate without pack-specific hacks.

\paragraph{Expected user-facing payoff.}
If implemented well, the payoff is IndexError, ZeroDivisionError, and shape-side precondition pullback. This should appear not only
in benchmark metrics but also in concrete rendered contracts,
diagnostics, or proof goals.


\subsection{Package \py{src/deppy/contracts} with Boolean guard theory}

\paragraph{Integration objective.}
The package \py{src/deppy/contracts} must interact cleanly with the
theory family boolean guard theory. The reason is practical: the
package owns part of the data path through which guard polarity, assertion scope, and branch-conditioned refinement becomes
visible, audited, and eventually rendered.

\paragraph{Required interface.}
At minimum the integration should expose normalized inputs, site-aware
outputs, provenance, and failure reasons. The package should neither hide
theory-relevant structure from the cover builder nor flatten theory
results so aggressively that renderers lose meaning.

\paragraph{Typical evidence source.}
In practice the evidence will come from if-statements, asserts, early returns, and exception guards. The integration should
preserve enough of that source structure that local solvers using normalization plus implication checks
can operate without pack-specific hacks.

\paragraph{Expected user-facing payoff.}
If implemented well, the payoff is control-dependent safety and lightweight proof seed extraction. This should appear not only
in benchmark metrics but also in concrete rendered contracts,
diagnostics, or proof goals.


\subsection{Package \py{src/deppy/contracts} with Nullability and optionality}

\paragraph{Integration objective.}
The package \py{src/deppy/contracts} must interact cleanly with the
theory family nullability and optionality. The reason is practical: the
package owns part of the data path through which none checks, existence guards, and optional protocol availability becomes
visible, audited, and eventually rendered.

\paragraph{Required interface.}
At minimum the integration should expose normalized inputs, site-aware
outputs, provenance, and failure reasons. The package should neither hide
theory-relevant structure from the cover builder nor flatten theory
results so aggressively that renderers lose meaning.

\paragraph{Typical evidence source.}
In practice the evidence will come from ordinary Python branches and constructor code. The integration should
preserve enough of that source structure that local solvers using simple symbolic narrowing
can operate without pack-specific hacks.

\paragraph{Expected user-facing payoff.}
If implemented well, the payoff is AttributeError prevention and stronger post-branch outputs. This should appear not only
in benchmark metrics but also in concrete rendered contracts,
diagnostics, or proof goals.


\subsection{Package \py{src/deppy/contracts} with Heap and protocol theory}

\paragraph{Integration objective.}
The package \py{src/deppy/contracts} must interact cleanly with the
theory family heap and protocol theory. The reason is practical: the
package owns part of the data path through which field initialization, protocol method existence, frame conditions, and alias exclusions becomes
visible, audited, and eventually rendered.

\paragraph{Required interface.}
At minimum the integration should expose normalized inputs, site-aware
outputs, provenance, and failure reasons. The package should neither hide
theory-relevant structure from the cover builder nor flatten theory
results so aggressively that renderers lose meaning.

\paragraph{Typical evidence source.}
In practice the evidence will come from classes, protocols, mutable objects, and resource wrappers. The integration should
preserve enough of that source structure that local solvers using bounded symbolic checks plus explicit residuals
can operate without pack-specific hacks.

\paragraph{Expected user-facing payoff.}
If implemented well, the payoff is constructor safety, attribute access, and method-dispatch correctness. This should appear not only
in benchmark metrics but also in concrete rendered contracts,
diagnostics, or proof goals.


\subsection{Package \py{src/deppy/contracts} with Tensor shape theory}

\paragraph{Integration objective.}
The package \py{src/deppy/contracts} must interact cleanly with the
theory family tensor shape theory. The reason is practical: the
package owns part of the data path through which rank, dimension, shape tuples, cardinality, broadcast legality, and reduction support becomes
visible, audited, and eventually rendered.

\paragraph{Required interface.}
At minimum the integration should expose normalized inputs, site-aware
outputs, provenance, and failure reasons. The package should neither hide
theory-relevant structure from the cover builder nor flatten theory
results so aggressively that renderers lose meaning.

\paragraph{Typical evidence source.}
In practice the evidence will come from Torch tensor operators and shape-sensitive wrappers. The integration should
preserve enough of that source structure that local solvers using solver-backed local shape kernels
can operate without pack-specific hacks.

\paragraph{Expected user-facing payoff.}
If implemented well, the payoff is invalid reshape, broadcast mismatch, and exact output-shape synthesis. This should appear not only
in benchmark metrics but also in concrete rendered contracts,
diagnostics, or proof goals.


\subsection{Package \py{src/deppy/contracts} with Tensor order theory}

\paragraph{Integration objective.}
The package \py{src/deppy/contracts} must interact cleanly with the
theory family tensor order theory. The reason is practical: the
package owns part of the data path through which sortedness, monotonicity, permutation witnesses, argsort relations becomes
visible, audited, and eventually rendered.

\paragraph{Required interface.}
At minimum the integration should expose normalized inputs, site-aware
outputs, provenance, and failure reasons. The package should neither hide
theory-relevant structure from the cover builder nor flatten theory
results so aggressively that renderers lose meaning.

\paragraph{Typical evidence source.}
In practice the evidence will come from sort-like operators and order-preserving wrappers. The integration should
preserve enough of that source structure that local solvers using finite witness templates plus solver-backed subchecks
can operate without pack-specific hacks.

\paragraph{Expected user-facing payoff.}
If implemented well, the payoff is values-vs-indices confusion and richer sorted outputs. This should appear not only
in benchmark metrics but also in concrete rendered contracts,
diagnostics, or proof goals.


\subsection{Package \py{src/deppy/contracts} with Tensor indexing theory}

\paragraph{Integration objective.}
The package \py{src/deppy/contracts} must interact cleanly with the
theory family tensor indexing theory. The reason is practical: the
package owns part of the data path through which selector domains, slice extents, gather/index-select legality, and value provenance becomes
visible, audited, and eventually rendered.

\paragraph{Required interface.}
At minimum the integration should expose normalized inputs, site-aware
outputs, provenance, and failure reasons. The package should neither hide
theory-relevant structure from the cover builder nor flatten theory
results so aggressively that renderers lose meaning.

\paragraph{Typical evidence source.}
In practice the evidence will come from subscripts, slices, and advanced indexing. The integration should
preserve enough of that source structure that local solvers using arithmetic plus pack-specific rules
can operate without pack-specific hacks.

\paragraph{Expected user-facing payoff.}
If implemented well, the payoff is out-of-range selector detection and source-target output relations. This should appear not only
in benchmark metrics but also in concrete rendered contracts,
diagnostics, or proof goals.


\subsection{Package \py{src/deppy/contracts} with Device and dtype compatibility theory}

\paragraph{Integration objective.}
The package \py{src/deppy/contracts} must interact cleanly with the
theory family device and dtype compatibility theory. The reason is practical: the
package owns part of the data path through which device equality, admissible implicit casts, and operator support matrices becomes
visible, audited, and eventually rendered.

\paragraph{Required interface.}
At minimum the integration should expose normalized inputs, site-aware
outputs, provenance, and failure reasons. The package should neither hide
theory-relevant structure from the cover builder nor flatten theory
results so aggressively that renderers lose meaning.

\paragraph{Typical evidence source.}
In practice the evidence will come from mixed-device or mixed-dtype Torch programs. The integration should
preserve enough of that source structure that local solvers using table-driven local checks
can operate without pack-specific hacks.

\paragraph{Expected user-facing payoff.}
If implemented well, the payoff is device mismatch and unsupported operator combinations. This should appear not only
in benchmark metrics but also in concrete rendered contracts,
diagnostics, or proof goals.


\subsection{Package \py{src/deppy/contracts} with Sequence and collection theory}

\paragraph{Integration objective.}
The package \py{src/deppy/contracts} must interact cleanly with the
theory family sequence and collection theory. The reason is practical: the
package owns part of the data path through which non-emptiness, membership, stable length relations, and head/tail legality becomes
visible, audited, and eventually rendered.

\paragraph{Required interface.}
At minimum the integration should expose normalized inputs, site-aware
outputs, provenance, and failure reasons. The package should neither hide
theory-relevant structure from the cover builder nor flatten theory
results so aggressively that renderers lose meaning.

\paragraph{Typical evidence source.}
In practice the evidence will come from lists, tuples, deques, and iterator adapters. The integration should
preserve enough of that source structure that local solvers using simple algebraic checks
can operate without pack-specific hacks.

\paragraph{Expected user-facing payoff.}
If implemented well, the payoff is head-of-empty and unpacking safety. This should appear not only
in benchmark metrics but also in concrete rendered contracts,
diagnostics, or proof goals.


\subsection{Package \py{src/deppy/contracts} with Exception theory}

\paragraph{Integration objective.}
The package \py{src/deppy/contracts} must interact cleanly with the
theory family exception theory. The reason is practical: the
package owns part of the data path through which exceptional output channels, throw sites, and caught-vs-uncaught distinctions becomes
visible, audited, and eventually rendered.

\paragraph{Required interface.}
At minimum the integration should expose normalized inputs, site-aware
outputs, provenance, and failure reasons. The package should neither hide
theory-relevant structure from the cover builder nor flatten theory
results so aggressively that renderers lose meaning.

\paragraph{Typical evidence source.}
In practice the evidence will come from try/except blocks, library raises, and proof obligations that may fail. The integration should
preserve enough of that source structure that local solvers using control-flow plus error-site hooks
can operate without pack-specific hacks.

\paragraph{Expected user-facing payoff.}
If implemented well, the payoff is precise reporting of exceptional contracts and remaining failures. This should appear not only
in benchmark metrics but also in concrete rendered contracts,
diagnostics, or proof goals.


\subsection{Package \py{src/deppy/contracts} with Witness and equality transport theory}

\paragraph{Integration objective.}
The package \py{src/deppy/contracts} must interact cleanly with the
theory family witness and equality transport theory. The reason is practical: the
package owns part of the data path through which dependent pairs, equality witnesses, and explicit transports becomes
visible, audited, and eventually rendered.

\paragraph{Required interface.}
At minimum the integration should expose normalized inputs, site-aware
outputs, provenance, and failure reasons. The package should neither hide
theory-relevant structure from the cover builder nor flatten theory
results so aggressively that renderers lose meaning.

\paragraph{Typical evidence source.}
In practice the evidence will come from proof-oriented code and witness-returning helpers. The integration should
preserve enough of that source structure that local solvers using proof-kernel backed local checking
can operate without pack-specific hacks.

\paragraph{Expected user-facing payoff.}
If implemented well, the payoff is proof reuse and output enrichment. This should appear not only
in benchmark metrics but also in concrete rendered contracts,
diagnostics, or proof goals.


\subsection{Package \py{src/deppy/contracts} with Loop invariant and decreases theory}

\paragraph{Integration objective.}
The package \py{src/deppy/contracts} must interact cleanly with the
theory family loop invariant and decreases theory. The reason is practical: the
package owns part of the data path through which candidate invariants, ranking relations, and preservation constraints becomes
visible, audited, and eventually rendered.

\paragraph{Required interface.}
At minimum the integration should expose normalized inputs, site-aware
outputs, provenance, and failure reasons. The package should neither hide
theory-relevant structure from the cover builder nor flatten theory
results so aggressively that renderers lose meaning.

\paragraph{Typical evidence source.}
In practice the evidence will come from loop-heavy algorithms and recursive helpers. The integration should
preserve enough of that source structure that local solvers using template generation plus bounded checks
can operate without pack-specific hacks.

\paragraph{Expected user-facing payoff.}
If implemented well, the payoff is F*-like capability without abandoning native Python. This should appear not only
in benchmark metrics but also in concrete rendered contracts,
diagnostics, or proof goals.


\subsection{Package \py{src/deppy/core} with Linear arithmetic and intervals}

\paragraph{Integration objective.}
The package \py{src/deppy/core} must interact cleanly with the
theory family linear arithmetic and intervals. The reason is practical: the
package owns part of the data path through which non-negativity, bounds, lengths, simple affine equalities becomes
visible, audited, and eventually rendered.

\paragraph{Required interface.}
At minimum the integration should expose normalized inputs, site-aware
outputs, provenance, and failure reasons. The package should neither hide
theory-relevant structure from the cover builder nor flatten theory
results so aggressively that renderers lose meaning.

\paragraph{Typical evidence source.}
In practice the evidence will come from argument bounds, list lengths, loop counters, and tensor extents. The integration should
preserve enough of that source structure that local solvers using local SMT over arithmetic fragments
can operate without pack-specific hacks.

\paragraph{Expected user-facing payoff.}
If implemented well, the payoff is IndexError, ZeroDivisionError, and shape-side precondition pullback. This should appear not only
in benchmark metrics but also in concrete rendered contracts,
diagnostics, or proof goals.


\subsection{Package \py{src/deppy/core} with Boolean guard theory}

\paragraph{Integration objective.}
The package \py{src/deppy/core} must interact cleanly with the
theory family boolean guard theory. The reason is practical: the
package owns part of the data path through which guard polarity, assertion scope, and branch-conditioned refinement becomes
visible, audited, and eventually rendered.

\paragraph{Required interface.}
At minimum the integration should expose normalized inputs, site-aware
outputs, provenance, and failure reasons. The package should neither hide
theory-relevant structure from the cover builder nor flatten theory
results so aggressively that renderers lose meaning.

\paragraph{Typical evidence source.}
In practice the evidence will come from if-statements, asserts, early returns, and exception guards. The integration should
preserve enough of that source structure that local solvers using normalization plus implication checks
can operate without pack-specific hacks.

\paragraph{Expected user-facing payoff.}
If implemented well, the payoff is control-dependent safety and lightweight proof seed extraction. This should appear not only
in benchmark metrics but also in concrete rendered contracts,
diagnostics, or proof goals.


\subsection{Package \py{src/deppy/core} with Nullability and optionality}

\paragraph{Integration objective.}
The package \py{src/deppy/core} must interact cleanly with the
theory family nullability and optionality. The reason is practical: the
package owns part of the data path through which none checks, existence guards, and optional protocol availability becomes
visible, audited, and eventually rendered.

\paragraph{Required interface.}
At minimum the integration should expose normalized inputs, site-aware
outputs, provenance, and failure reasons. The package should neither hide
theory-relevant structure from the cover builder nor flatten theory
results so aggressively that renderers lose meaning.

\paragraph{Typical evidence source.}
In practice the evidence will come from ordinary Python branches and constructor code. The integration should
preserve enough of that source structure that local solvers using simple symbolic narrowing
can operate without pack-specific hacks.

\paragraph{Expected user-facing payoff.}
If implemented well, the payoff is AttributeError prevention and stronger post-branch outputs. This should appear not only
in benchmark metrics but also in concrete rendered contracts,
diagnostics, or proof goals.


\subsection{Package \py{src/deppy/core} with Heap and protocol theory}

\paragraph{Integration objective.}
The package \py{src/deppy/core} must interact cleanly with the
theory family heap and protocol theory. The reason is practical: the
package owns part of the data path through which field initialization, protocol method existence, frame conditions, and alias exclusions becomes
visible, audited, and eventually rendered.

\paragraph{Required interface.}
At minimum the integration should expose normalized inputs, site-aware
outputs, provenance, and failure reasons. The package should neither hide
theory-relevant structure from the cover builder nor flatten theory
results so aggressively that renderers lose meaning.

\paragraph{Typical evidence source.}
In practice the evidence will come from classes, protocols, mutable objects, and resource wrappers. The integration should
preserve enough of that source structure that local solvers using bounded symbolic checks plus explicit residuals
can operate without pack-specific hacks.

\paragraph{Expected user-facing payoff.}
If implemented well, the payoff is constructor safety, attribute access, and method-dispatch correctness. This should appear not only
in benchmark metrics but also in concrete rendered contracts,
diagnostics, or proof goals.


\subsection{Package \py{src/deppy/core} with Tensor shape theory}

\paragraph{Integration objective.}
The package \py{src/deppy/core} must interact cleanly with the
theory family tensor shape theory. The reason is practical: the
package owns part of the data path through which rank, dimension, shape tuples, cardinality, broadcast legality, and reduction support becomes
visible, audited, and eventually rendered.

\paragraph{Required interface.}
At minimum the integration should expose normalized inputs, site-aware
outputs, provenance, and failure reasons. The package should neither hide
theory-relevant structure from the cover builder nor flatten theory
results so aggressively that renderers lose meaning.

\paragraph{Typical evidence source.}
In practice the evidence will come from Torch tensor operators and shape-sensitive wrappers. The integration should
preserve enough of that source structure that local solvers using solver-backed local shape kernels
can operate without pack-specific hacks.

\paragraph{Expected user-facing payoff.}
If implemented well, the payoff is invalid reshape, broadcast mismatch, and exact output-shape synthesis. This should appear not only
in benchmark metrics but also in concrete rendered contracts,
diagnostics, or proof goals.


\subsection{Package \py{src/deppy/core} with Tensor order theory}

\paragraph{Integration objective.}
The package \py{src/deppy/core} must interact cleanly with the
theory family tensor order theory. The reason is practical: the
package owns part of the data path through which sortedness, monotonicity, permutation witnesses, argsort relations becomes
visible, audited, and eventually rendered.

\paragraph{Required interface.}
At minimum the integration should expose normalized inputs, site-aware
outputs, provenance, and failure reasons. The package should neither hide
theory-relevant structure from the cover builder nor flatten theory
results so aggressively that renderers lose meaning.

\paragraph{Typical evidence source.}
In practice the evidence will come from sort-like operators and order-preserving wrappers. The integration should
preserve enough of that source structure that local solvers using finite witness templates plus solver-backed subchecks
can operate without pack-specific hacks.

\paragraph{Expected user-facing payoff.}
If implemented well, the payoff is values-vs-indices confusion and richer sorted outputs. This should appear not only
in benchmark metrics but also in concrete rendered contracts,
diagnostics, or proof goals.


\subsection{Package \py{src/deppy/core} with Tensor indexing theory}

\paragraph{Integration objective.}
The package \py{src/deppy/core} must interact cleanly with the
theory family tensor indexing theory. The reason is practical: the
package owns part of the data path through which selector domains, slice extents, gather/index-select legality, and value provenance becomes
visible, audited, and eventually rendered.

\paragraph{Required interface.}
At minimum the integration should expose normalized inputs, site-aware
outputs, provenance, and failure reasons. The package should neither hide
theory-relevant structure from the cover builder nor flatten theory
results so aggressively that renderers lose meaning.

\paragraph{Typical evidence source.}
In practice the evidence will come from subscripts, slices, and advanced indexing. The integration should
preserve enough of that source structure that local solvers using arithmetic plus pack-specific rules
can operate without pack-specific hacks.

\paragraph{Expected user-facing payoff.}
If implemented well, the payoff is out-of-range selector detection and source-target output relations. This should appear not only
in benchmark metrics but also in concrete rendered contracts,
diagnostics, or proof goals.


\subsection{Package \py{src/deppy/core} with Device and dtype compatibility theory}

\paragraph{Integration objective.}
The package \py{src/deppy/core} must interact cleanly with the
theory family device and dtype compatibility theory. The reason is practical: the
package owns part of the data path through which device equality, admissible implicit casts, and operator support matrices becomes
visible, audited, and eventually rendered.

\paragraph{Required interface.}
At minimum the integration should expose normalized inputs, site-aware
outputs, provenance, and failure reasons. The package should neither hide
theory-relevant structure from the cover builder nor flatten theory
results so aggressively that renderers lose meaning.

\paragraph{Typical evidence source.}
In practice the evidence will come from mixed-device or mixed-dtype Torch programs. The integration should
preserve enough of that source structure that local solvers using table-driven local checks
can operate without pack-specific hacks.

\paragraph{Expected user-facing payoff.}
If implemented well, the payoff is device mismatch and unsupported operator combinations. This should appear not only
in benchmark metrics but also in concrete rendered contracts,
diagnostics, or proof goals.


\subsection{Package \py{src/deppy/core} with Sequence and collection theory}

\paragraph{Integration objective.}
The package \py{src/deppy/core} must interact cleanly with the
theory family sequence and collection theory. The reason is practical: the
package owns part of the data path through which non-emptiness, membership, stable length relations, and head/tail legality becomes
visible, audited, and eventually rendered.

\paragraph{Required interface.}
At minimum the integration should expose normalized inputs, site-aware
outputs, provenance, and failure reasons. The package should neither hide
theory-relevant structure from the cover builder nor flatten theory
results so aggressively that renderers lose meaning.

\paragraph{Typical evidence source.}
In practice the evidence will come from lists, tuples, deques, and iterator adapters. The integration should
preserve enough of that source structure that local solvers using simple algebraic checks
can operate without pack-specific hacks.

\paragraph{Expected user-facing payoff.}
If implemented well, the payoff is head-of-empty and unpacking safety. This should appear not only
in benchmark metrics but also in concrete rendered contracts,
diagnostics, or proof goals.


\subsection{Package \py{src/deppy/core} with Exception theory}

\paragraph{Integration objective.}
The package \py{src/deppy/core} must interact cleanly with the
theory family exception theory. The reason is practical: the
package owns part of the data path through which exceptional output channels, throw sites, and caught-vs-uncaught distinctions becomes
visible, audited, and eventually rendered.

\paragraph{Required interface.}
At minimum the integration should expose normalized inputs, site-aware
outputs, provenance, and failure reasons. The package should neither hide
theory-relevant structure from the cover builder nor flatten theory
results so aggressively that renderers lose meaning.

\paragraph{Typical evidence source.}
In practice the evidence will come from try/except blocks, library raises, and proof obligations that may fail. The integration should
preserve enough of that source structure that local solvers using control-flow plus error-site hooks
can operate without pack-specific hacks.

\paragraph{Expected user-facing payoff.}
If implemented well, the payoff is precise reporting of exceptional contracts and remaining failures. This should appear not only
in benchmark metrics but also in concrete rendered contracts,
diagnostics, or proof goals.


\subsection{Package \py{src/deppy/core} with Witness and equality transport theory}

\paragraph{Integration objective.}
The package \py{src/deppy/core} must interact cleanly with the
theory family witness and equality transport theory. The reason is practical: the
package owns part of the data path through which dependent pairs, equality witnesses, and explicit transports becomes
visible, audited, and eventually rendered.

\paragraph{Required interface.}
At minimum the integration should expose normalized inputs, site-aware
outputs, provenance, and failure reasons. The package should neither hide
theory-relevant structure from the cover builder nor flatten theory
results so aggressively that renderers lose meaning.

\paragraph{Typical evidence source.}
In practice the evidence will come from proof-oriented code and witness-returning helpers. The integration should
preserve enough of that source structure that local solvers using proof-kernel backed local checking
can operate without pack-specific hacks.

\paragraph{Expected user-facing payoff.}
If implemented well, the payoff is proof reuse and output enrichment. This should appear not only
in benchmark metrics but also in concrete rendered contracts,
diagnostics, or proof goals.


\subsection{Package \py{src/deppy/core} with Loop invariant and decreases theory}

\paragraph{Integration objective.}
The package \py{src/deppy/core} must interact cleanly with the
theory family loop invariant and decreases theory. The reason is practical: the
package owns part of the data path through which candidate invariants, ranking relations, and preservation constraints becomes
visible, audited, and eventually rendered.

\paragraph{Required interface.}
At minimum the integration should expose normalized inputs, site-aware
outputs, provenance, and failure reasons. The package should neither hide
theory-relevant structure from the cover builder nor flatten theory
results so aggressively that renderers lose meaning.

\paragraph{Typical evidence source.}
In practice the evidence will come from loop-heavy algorithms and recursive helpers. The integration should
preserve enough of that source structure that local solvers using template generation plus bounded checks
can operate without pack-specific hacks.

\paragraph{Expected user-facing payoff.}
If implemented well, the payoff is F*-like capability without abandoning native Python. This should appear not only
in benchmark metrics but also in concrete rendered contracts,
diagnostics, or proof goals.


\subsection{Package \py{src/deppy/frontend} with Linear arithmetic and intervals}

\paragraph{Integration objective.}
The package \py{src/deppy/frontend} must interact cleanly with the
theory family linear arithmetic and intervals. The reason is practical: the
package owns part of the data path through which non-negativity, bounds, lengths, simple affine equalities becomes
visible, audited, and eventually rendered.

\paragraph{Required interface.}
At minimum the integration should expose normalized inputs, site-aware
outputs, provenance, and failure reasons. The package should neither hide
theory-relevant structure from the cover builder nor flatten theory
results so aggressively that renderers lose meaning.

\paragraph{Typical evidence source.}
In practice the evidence will come from argument bounds, list lengths, loop counters, and tensor extents. The integration should
preserve enough of that source structure that local solvers using local SMT over arithmetic fragments
can operate without pack-specific hacks.

\paragraph{Expected user-facing payoff.}
If implemented well, the payoff is IndexError, ZeroDivisionError, and shape-side precondition pullback. This should appear not only
in benchmark metrics but also in concrete rendered contracts,
diagnostics, or proof goals.


\subsection{Package \py{src/deppy/frontend} with Boolean guard theory}

\paragraph{Integration objective.}
The package \py{src/deppy/frontend} must interact cleanly with the
theory family boolean guard theory. The reason is practical: the
package owns part of the data path through which guard polarity, assertion scope, and branch-conditioned refinement becomes
visible, audited, and eventually rendered.

\paragraph{Required interface.}
At minimum the integration should expose normalized inputs, site-aware
outputs, provenance, and failure reasons. The package should neither hide
theory-relevant structure from the cover builder nor flatten theory
results so aggressively that renderers lose meaning.

\paragraph{Typical evidence source.}
In practice the evidence will come from if-statements, asserts, early returns, and exception guards. The integration should
preserve enough of that source structure that local solvers using normalization plus implication checks
can operate without pack-specific hacks.

\paragraph{Expected user-facing payoff.}
If implemented well, the payoff is control-dependent safety and lightweight proof seed extraction. This should appear not only
in benchmark metrics but also in concrete rendered contracts,
diagnostics, or proof goals.


\subsection{Package \py{src/deppy/frontend} with Nullability and optionality}

\paragraph{Integration objective.}
The package \py{src/deppy/frontend} must interact cleanly with the
theory family nullability and optionality. The reason is practical: the
package owns part of the data path through which none checks, existence guards, and optional protocol availability becomes
visible, audited, and eventually rendered.

\paragraph{Required interface.}
At minimum the integration should expose normalized inputs, site-aware
outputs, provenance, and failure reasons. The package should neither hide
theory-relevant structure from the cover builder nor flatten theory
results so aggressively that renderers lose meaning.

\paragraph{Typical evidence source.}
In practice the evidence will come from ordinary Python branches and constructor code. The integration should
preserve enough of that source structure that local solvers using simple symbolic narrowing
can operate without pack-specific hacks.

\paragraph{Expected user-facing payoff.}
If implemented well, the payoff is AttributeError prevention and stronger post-branch outputs. This should appear not only
in benchmark metrics but also in concrete rendered contracts,
diagnostics, or proof goals.


\subsection{Package \py{src/deppy/frontend} with Heap and protocol theory}

\paragraph{Integration objective.}
The package \py{src/deppy/frontend} must interact cleanly with the
theory family heap and protocol theory. The reason is practical: the
package owns part of the data path through which field initialization, protocol method existence, frame conditions, and alias exclusions becomes
visible, audited, and eventually rendered.

\paragraph{Required interface.}
At minimum the integration should expose normalized inputs, site-aware
outputs, provenance, and failure reasons. The package should neither hide
theory-relevant structure from the cover builder nor flatten theory
results so aggressively that renderers lose meaning.

\paragraph{Typical evidence source.}
In practice the evidence will come from classes, protocols, mutable objects, and resource wrappers. The integration should
preserve enough of that source structure that local solvers using bounded symbolic checks plus explicit residuals
can operate without pack-specific hacks.

\paragraph{Expected user-facing payoff.}
If implemented well, the payoff is constructor safety, attribute access, and method-dispatch correctness. This should appear not only
in benchmark metrics but also in concrete rendered contracts,
diagnostics, or proof goals.


\subsection{Package \py{src/deppy/frontend} with Tensor shape theory}

\paragraph{Integration objective.}
The package \py{src/deppy/frontend} must interact cleanly with the
theory family tensor shape theory. The reason is practical: the
package owns part of the data path through which rank, dimension, shape tuples, cardinality, broadcast legality, and reduction support becomes
visible, audited, and eventually rendered.

\paragraph{Required interface.}
At minimum the integration should expose normalized inputs, site-aware
outputs, provenance, and failure reasons. The package should neither hide
theory-relevant structure from the cover builder nor flatten theory
results so aggressively that renderers lose meaning.

\paragraph{Typical evidence source.}
In practice the evidence will come from Torch tensor operators and shape-sensitive wrappers. The integration should
preserve enough of that source structure that local solvers using solver-backed local shape kernels
can operate without pack-specific hacks.

\paragraph{Expected user-facing payoff.}
If implemented well, the payoff is invalid reshape, broadcast mismatch, and exact output-shape synthesis. This should appear not only
in benchmark metrics but also in concrete rendered contracts,
diagnostics, or proof goals.


\subsection{Package \py{src/deppy/frontend} with Tensor order theory}

\paragraph{Integration objective.}
The package \py{src/deppy/frontend} must interact cleanly with the
theory family tensor order theory. The reason is practical: the
package owns part of the data path through which sortedness, monotonicity, permutation witnesses, argsort relations becomes
visible, audited, and eventually rendered.

\paragraph{Required interface.}
At minimum the integration should expose normalized inputs, site-aware
outputs, provenance, and failure reasons. The package should neither hide
theory-relevant structure from the cover builder nor flatten theory
results so aggressively that renderers lose meaning.

\paragraph{Typical evidence source.}
In practice the evidence will come from sort-like operators and order-preserving wrappers. The integration should
preserve enough of that source structure that local solvers using finite witness templates plus solver-backed subchecks
can operate without pack-specific hacks.

\paragraph{Expected user-facing payoff.}
If implemented well, the payoff is values-vs-indices confusion and richer sorted outputs. This should appear not only
in benchmark metrics but also in concrete rendered contracts,
diagnostics, or proof goals.


\subsection{Package \py{src/deppy/frontend} with Tensor indexing theory}

\paragraph{Integration objective.}
The package \py{src/deppy/frontend} must interact cleanly with the
theory family tensor indexing theory. The reason is practical: the
package owns part of the data path through which selector domains, slice extents, gather/index-select legality, and value provenance becomes
visible, audited, and eventually rendered.

\paragraph{Required interface.}
At minimum the integration should expose normalized inputs, site-aware
outputs, provenance, and failure reasons. The package should neither hide
theory-relevant structure from the cover builder nor flatten theory
results so aggressively that renderers lose meaning.

\paragraph{Typical evidence source.}
In practice the evidence will come from subscripts, slices, and advanced indexing. The integration should
preserve enough of that source structure that local solvers using arithmetic plus pack-specific rules
can operate without pack-specific hacks.

\paragraph{Expected user-facing payoff.}
If implemented well, the payoff is out-of-range selector detection and source-target output relations. This should appear not only
in benchmark metrics but also in concrete rendered contracts,
diagnostics, or proof goals.


\subsection{Package \py{src/deppy/frontend} with Device and dtype compatibility theory}

\paragraph{Integration objective.}
The package \py{src/deppy/frontend} must interact cleanly with the
theory family device and dtype compatibility theory. The reason is practical: the
package owns part of the data path through which device equality, admissible implicit casts, and operator support matrices becomes
visible, audited, and eventually rendered.

\paragraph{Required interface.}
At minimum the integration should expose normalized inputs, site-aware
outputs, provenance, and failure reasons. The package should neither hide
theory-relevant structure from the cover builder nor flatten theory
results so aggressively that renderers lose meaning.

\paragraph{Typical evidence source.}
In practice the evidence will come from mixed-device or mixed-dtype Torch programs. The integration should
preserve enough of that source structure that local solvers using table-driven local checks
can operate without pack-specific hacks.

\paragraph{Expected user-facing payoff.}
If implemented well, the payoff is device mismatch and unsupported operator combinations. This should appear not only
in benchmark metrics but also in concrete rendered contracts,
diagnostics, or proof goals.


\subsection{Package \py{src/deppy/frontend} with Sequence and collection theory}

\paragraph{Integration objective.}
The package \py{src/deppy/frontend} must interact cleanly with the
theory family sequence and collection theory. The reason is practical: the
package owns part of the data path through which non-emptiness, membership, stable length relations, and head/tail legality becomes
visible, audited, and eventually rendered.

\paragraph{Required interface.}
At minimum the integration should expose normalized inputs, site-aware
outputs, provenance, and failure reasons. The package should neither hide
theory-relevant structure from the cover builder nor flatten theory
results so aggressively that renderers lose meaning.

\paragraph{Typical evidence source.}
In practice the evidence will come from lists, tuples, deques, and iterator adapters. The integration should
preserve enough of that source structure that local solvers using simple algebraic checks
can operate without pack-specific hacks.

\paragraph{Expected user-facing payoff.}
If implemented well, the payoff is head-of-empty and unpacking safety. This should appear not only
in benchmark metrics but also in concrete rendered contracts,
diagnostics, or proof goals.


\subsection{Package \py{src/deppy/frontend} with Exception theory}

\paragraph{Integration objective.}
The package \py{src/deppy/frontend} must interact cleanly with the
theory family exception theory. The reason is practical: the
package owns part of the data path through which exceptional output channels, throw sites, and caught-vs-uncaught distinctions becomes
visible, audited, and eventually rendered.

\paragraph{Required interface.}
At minimum the integration should expose normalized inputs, site-aware
outputs, provenance, and failure reasons. The package should neither hide
theory-relevant structure from the cover builder nor flatten theory
results so aggressively that renderers lose meaning.

\paragraph{Typical evidence source.}
In practice the evidence will come from try/except blocks, library raises, and proof obligations that may fail. The integration should
preserve enough of that source structure that local solvers using control-flow plus error-site hooks
can operate without pack-specific hacks.

\paragraph{Expected user-facing payoff.}
If implemented well, the payoff is precise reporting of exceptional contracts and remaining failures. This should appear not only
in benchmark metrics but also in concrete rendered contracts,
diagnostics, or proof goals.


\subsection{Package \py{src/deppy/frontend} with Witness and equality transport theory}

\paragraph{Integration objective.}
The package \py{src/deppy/frontend} must interact cleanly with the
theory family witness and equality transport theory. The reason is practical: the
package owns part of the data path through which dependent pairs, equality witnesses, and explicit transports becomes
visible, audited, and eventually rendered.

\paragraph{Required interface.}
At minimum the integration should expose normalized inputs, site-aware
outputs, provenance, and failure reasons. The package should neither hide
theory-relevant structure from the cover builder nor flatten theory
results so aggressively that renderers lose meaning.

\paragraph{Typical evidence source.}
In practice the evidence will come from proof-oriented code and witness-returning helpers. The integration should
preserve enough of that source structure that local solvers using proof-kernel backed local checking
can operate without pack-specific hacks.

\paragraph{Expected user-facing payoff.}
If implemented well, the payoff is proof reuse and output enrichment. This should appear not only
in benchmark metrics but also in concrete rendered contracts,
diagnostics, or proof goals.


\subsection{Package \py{src/deppy/frontend} with Loop invariant and decreases theory}

\paragraph{Integration objective.}
The package \py{src/deppy/frontend} must interact cleanly with the
theory family loop invariant and decreases theory. The reason is practical: the
package owns part of the data path through which candidate invariants, ranking relations, and preservation constraints becomes
visible, audited, and eventually rendered.

\paragraph{Required interface.}
At minimum the integration should expose normalized inputs, site-aware
outputs, provenance, and failure reasons. The package should neither hide
theory-relevant structure from the cover builder nor flatten theory
results so aggressively that renderers lose meaning.

\paragraph{Typical evidence source.}
In practice the evidence will come from loop-heavy algorithms and recursive helpers. The integration should
preserve enough of that source structure that local solvers using template generation plus bounded checks
can operate without pack-specific hacks.

\paragraph{Expected user-facing payoff.}
If implemented well, the payoff is F*-like capability without abandoning native Python. This should appear not only
in benchmark metrics but also in concrete rendered contracts,
diagnostics, or proof goals.


\subsection{Package \py{src/deppy/generate} with Linear arithmetic and intervals}

\paragraph{Integration objective.}
The package \py{src/deppy/generate} must interact cleanly with the
theory family linear arithmetic and intervals. The reason is practical: the
package owns part of the data path through which non-negativity, bounds, lengths, simple affine equalities becomes
visible, audited, and eventually rendered.

\paragraph{Required interface.}
At minimum the integration should expose normalized inputs, site-aware
outputs, provenance, and failure reasons. The package should neither hide
theory-relevant structure from the cover builder nor flatten theory
results so aggressively that renderers lose meaning.

\paragraph{Typical evidence source.}
In practice the evidence will come from argument bounds, list lengths, loop counters, and tensor extents. The integration should
preserve enough of that source structure that local solvers using local SMT over arithmetic fragments
can operate without pack-specific hacks.

\paragraph{Expected user-facing payoff.}
If implemented well, the payoff is IndexError, ZeroDivisionError, and shape-side precondition pullback. This should appear not only
in benchmark metrics but also in concrete rendered contracts,
diagnostics, or proof goals.


\subsection{Package \py{src/deppy/generate} with Boolean guard theory}

\paragraph{Integration objective.}
The package \py{src/deppy/generate} must interact cleanly with the
theory family boolean guard theory. The reason is practical: the
package owns part of the data path through which guard polarity, assertion scope, and branch-conditioned refinement becomes
visible, audited, and eventually rendered.

\paragraph{Required interface.}
At minimum the integration should expose normalized inputs, site-aware
outputs, provenance, and failure reasons. The package should neither hide
theory-relevant structure from the cover builder nor flatten theory
results so aggressively that renderers lose meaning.

\paragraph{Typical evidence source.}
In practice the evidence will come from if-statements, asserts, early returns, and exception guards. The integration should
preserve enough of that source structure that local solvers using normalization plus implication checks
can operate without pack-specific hacks.

\paragraph{Expected user-facing payoff.}
If implemented well, the payoff is control-dependent safety and lightweight proof seed extraction. This should appear not only
in benchmark metrics but also in concrete rendered contracts,
diagnostics, or proof goals.


\subsection{Package \py{src/deppy/generate} with Nullability and optionality}

\paragraph{Integration objective.}
The package \py{src/deppy/generate} must interact cleanly with the
theory family nullability and optionality. The reason is practical: the
package owns part of the data path through which none checks, existence guards, and optional protocol availability becomes
visible, audited, and eventually rendered.

\paragraph{Required interface.}
At minimum the integration should expose normalized inputs, site-aware
outputs, provenance, and failure reasons. The package should neither hide
theory-relevant structure from the cover builder nor flatten theory
results so aggressively that renderers lose meaning.

\paragraph{Typical evidence source.}
In practice the evidence will come from ordinary Python branches and constructor code. The integration should
preserve enough of that source structure that local solvers using simple symbolic narrowing
can operate without pack-specific hacks.

\paragraph{Expected user-facing payoff.}
If implemented well, the payoff is AttributeError prevention and stronger post-branch outputs. This should appear not only
in benchmark metrics but also in concrete rendered contracts,
diagnostics, or proof goals.


\subsection{Package \py{src/deppy/generate} with Heap and protocol theory}

\paragraph{Integration objective.}
The package \py{src/deppy/generate} must interact cleanly with the
theory family heap and protocol theory. The reason is practical: the
package owns part of the data path through which field initialization, protocol method existence, frame conditions, and alias exclusions becomes
visible, audited, and eventually rendered.

\paragraph{Required interface.}
At minimum the integration should expose normalized inputs, site-aware
outputs, provenance, and failure reasons. The package should neither hide
theory-relevant structure from the cover builder nor flatten theory
results so aggressively that renderers lose meaning.

\paragraph{Typical evidence source.}
In practice the evidence will come from classes, protocols, mutable objects, and resource wrappers. The integration should
preserve enough of that source structure that local solvers using bounded symbolic checks plus explicit residuals
can operate without pack-specific hacks.

\paragraph{Expected user-facing payoff.}
If implemented well, the payoff is constructor safety, attribute access, and method-dispatch correctness. This should appear not only
in benchmark metrics but also in concrete rendered contracts,
diagnostics, or proof goals.


\subsection{Package \py{src/deppy/generate} with Tensor shape theory}

\paragraph{Integration objective.}
The package \py{src/deppy/generate} must interact cleanly with the
theory family tensor shape theory. The reason is practical: the
package owns part of the data path through which rank, dimension, shape tuples, cardinality, broadcast legality, and reduction support becomes
visible, audited, and eventually rendered.

\paragraph{Required interface.}
At minimum the integration should expose normalized inputs, site-aware
outputs, provenance, and failure reasons. The package should neither hide
theory-relevant structure from the cover builder nor flatten theory
results so aggressively that renderers lose meaning.

\paragraph{Typical evidence source.}
In practice the evidence will come from Torch tensor operators and shape-sensitive wrappers. The integration should
preserve enough of that source structure that local solvers using solver-backed local shape kernels
can operate without pack-specific hacks.

\paragraph{Expected user-facing payoff.}
If implemented well, the payoff is invalid reshape, broadcast mismatch, and exact output-shape synthesis. This should appear not only
in benchmark metrics but also in concrete rendered contracts,
diagnostics, or proof goals.


\subsection{Package \py{src/deppy/generate} with Tensor order theory}

\paragraph{Integration objective.}
The package \py{src/deppy/generate} must interact cleanly with the
theory family tensor order theory. The reason is practical: the
package owns part of the data path through which sortedness, monotonicity, permutation witnesses, argsort relations becomes
visible, audited, and eventually rendered.

\paragraph{Required interface.}
At minimum the integration should expose normalized inputs, site-aware
outputs, provenance, and failure reasons. The package should neither hide
theory-relevant structure from the cover builder nor flatten theory
results so aggressively that renderers lose meaning.

\paragraph{Typical evidence source.}
In practice the evidence will come from sort-like operators and order-preserving wrappers. The integration should
preserve enough of that source structure that local solvers using finite witness templates plus solver-backed subchecks
can operate without pack-specific hacks.

\paragraph{Expected user-facing payoff.}
If implemented well, the payoff is values-vs-indices confusion and richer sorted outputs. This should appear not only
in benchmark metrics but also in concrete rendered contracts,
diagnostics, or proof goals.


\subsection{Package \py{src/deppy/generate} with Tensor indexing theory}

\paragraph{Integration objective.}
The package \py{src/deppy/generate} must interact cleanly with the
theory family tensor indexing theory. The reason is practical: the
package owns part of the data path through which selector domains, slice extents, gather/index-select legality, and value provenance becomes
visible, audited, and eventually rendered.

\paragraph{Required interface.}
At minimum the integration should expose normalized inputs, site-aware
outputs, provenance, and failure reasons. The package should neither hide
theory-relevant structure from the cover builder nor flatten theory
results so aggressively that renderers lose meaning.

\paragraph{Typical evidence source.}
In practice the evidence will come from subscripts, slices, and advanced indexing. The integration should
preserve enough of that source structure that local solvers using arithmetic plus pack-specific rules
can operate without pack-specific hacks.

\paragraph{Expected user-facing payoff.}
If implemented well, the payoff is out-of-range selector detection and source-target output relations. This should appear not only
in benchmark metrics but also in concrete rendered contracts,
diagnostics, or proof goals.


\subsection{Package \py{src/deppy/generate} with Device and dtype compatibility theory}

\paragraph{Integration objective.}
The package \py{src/deppy/generate} must interact cleanly with the
theory family device and dtype compatibility theory. The reason is practical: the
package owns part of the data path through which device equality, admissible implicit casts, and operator support matrices becomes
visible, audited, and eventually rendered.

\paragraph{Required interface.}
At minimum the integration should expose normalized inputs, site-aware
outputs, provenance, and failure reasons. The package should neither hide
theory-relevant structure from the cover builder nor flatten theory
results so aggressively that renderers lose meaning.

\paragraph{Typical evidence source.}
In practice the evidence will come from mixed-device or mixed-dtype Torch programs. The integration should
preserve enough of that source structure that local solvers using table-driven local checks
can operate without pack-specific hacks.

\paragraph{Expected user-facing payoff.}
If implemented well, the payoff is device mismatch and unsupported operator combinations. This should appear not only
in benchmark metrics but also in concrete rendered contracts,
diagnostics, or proof goals.


\subsection{Package \py{src/deppy/generate} with Sequence and collection theory}

\paragraph{Integration objective.}
The package \py{src/deppy/generate} must interact cleanly with the
theory family sequence and collection theory. The reason is practical: the
package owns part of the data path through which non-emptiness, membership, stable length relations, and head/tail legality becomes
visible, audited, and eventually rendered.

\paragraph{Required interface.}
At minimum the integration should expose normalized inputs, site-aware
outputs, provenance, and failure reasons. The package should neither hide
theory-relevant structure from the cover builder nor flatten theory
results so aggressively that renderers lose meaning.

\paragraph{Typical evidence source.}
In practice the evidence will come from lists, tuples, deques, and iterator adapters. The integration should
preserve enough of that source structure that local solvers using simple algebraic checks
can operate without pack-specific hacks.

\paragraph{Expected user-facing payoff.}
If implemented well, the payoff is head-of-empty and unpacking safety. This should appear not only
in benchmark metrics but also in concrete rendered contracts,
diagnostics, or proof goals.


\subsection{Package \py{src/deppy/generate} with Exception theory}

\paragraph{Integration objective.}
The package \py{src/deppy/generate} must interact cleanly with the
theory family exception theory. The reason is practical: the
package owns part of the data path through which exceptional output channels, throw sites, and caught-vs-uncaught distinctions becomes
visible, audited, and eventually rendered.

\paragraph{Required interface.}
At minimum the integration should expose normalized inputs, site-aware
outputs, provenance, and failure reasons. The package should neither hide
theory-relevant structure from the cover builder nor flatten theory
results so aggressively that renderers lose meaning.

\paragraph{Typical evidence source.}
In practice the evidence will come from try/except blocks, library raises, and proof obligations that may fail. The integration should
preserve enough of that source structure that local solvers using control-flow plus error-site hooks
can operate without pack-specific hacks.

\paragraph{Expected user-facing payoff.}
If implemented well, the payoff is precise reporting of exceptional contracts and remaining failures. This should appear not only
in benchmark metrics but also in concrete rendered contracts,
diagnostics, or proof goals.


\subsection{Package \py{src/deppy/generate} with Witness and equality transport theory}

\paragraph{Integration objective.}
The package \py{src/deppy/generate} must interact cleanly with the
theory family witness and equality transport theory. The reason is practical: the
package owns part of the data path through which dependent pairs, equality witnesses, and explicit transports becomes
visible, audited, and eventually rendered.

\paragraph{Required interface.}
At minimum the integration should expose normalized inputs, site-aware
outputs, provenance, and failure reasons. The package should neither hide
theory-relevant structure from the cover builder nor flatten theory
results so aggressively that renderers lose meaning.

\paragraph{Typical evidence source.}
In practice the evidence will come from proof-oriented code and witness-returning helpers. The integration should
preserve enough of that source structure that local solvers using proof-kernel backed local checking
can operate without pack-specific hacks.

\paragraph{Expected user-facing payoff.}
If implemented well, the payoff is proof reuse and output enrichment. This should appear not only
in benchmark metrics but also in concrete rendered contracts,
diagnostics, or proof goals.


\subsection{Package \py{src/deppy/generate} with Loop invariant and decreases theory}

\paragraph{Integration objective.}
The package \py{src/deppy/generate} must interact cleanly with the
theory family loop invariant and decreases theory. The reason is practical: the
package owns part of the data path through which candidate invariants, ranking relations, and preservation constraints becomes
visible, audited, and eventually rendered.

\paragraph{Required interface.}
At minimum the integration should expose normalized inputs, site-aware
outputs, provenance, and failure reasons. The package should neither hide
theory-relevant structure from the cover builder nor flatten theory
results so aggressively that renderers lose meaning.

\paragraph{Typical evidence source.}
In practice the evidence will come from loop-heavy algorithms and recursive helpers. The integration should
preserve enough of that source structure that local solvers using template generation plus bounded checks
can operate without pack-specific hacks.

\paragraph{Expected user-facing payoff.}
If implemented well, the payoff is F*-like capability without abandoning native Python. This should appear not only
in benchmark metrics but also in concrete rendered contracts,
diagnostics, or proof goals.

\section{Deep dive: list semantics, bounds safety, and length-dependent contracts}

Python lists are the most common mutable sequence and one of the richest sources
of runtime errors in everyday code. A sheaf-descent system must model lists not
as opaque containers but as semantic neighborhoods carrying length, element-type,
ordering, and mutability information at every site where a list value is live.

\subsection{Length as a first-class refinement coordinate}

The single most valuable fact about a list is its length. An enormous fraction of
Python runtime errors reduce to a length violation: indexing past the end,
unpacking the wrong number of elements, calling \py{pop()} on an empty list,
slicing with out-of-range bounds, or passing a list of unexpected size to a
library function that silently assumes a particular shape. A sheaf-descent system
should therefore treat length as a first-class coordinate on every list-valued
site.

Concretely, every SSA value whose carrier type includes \py{list} should have a
local section containing at least a symbolic length variable. When a list is
constructed via a literal, the length is exactly known. When a list is
constructed via comprehension, the length may depend on the length of the source
iterable and the filtering predicate. When a list is received as a function
argument, the length is initially unconstrained but may be refined by guards,
assertions, or downstream usage patterns.

\paragraph{Guard harvesting for length.}
The harvest layer should recognize the following patterns as length evidence:
\begin{itemize}[nosep]
  \item \py{if len(xs) > 0}, \py{if xs}, \py{if not xs} --- non-emptiness or
        emptiness guards that refine the length coordinate to $> 0$ or $= 0$.
  \item \py{if len(xs) == n}, \py{if len(xs) >= k} --- exact or bounded length
        guards.
  \item \py{assert len(xs) > 0} --- assertion-based refinement.
  \item \py{for x in xs: ...} --- the loop body site sees a non-exhausted
        iterator, which implies the original list was nonempty if the body
        executes at all.
  \item \py{xs = list(range(n))} or \py{xs = [0] * n} --- construction sites
        where the length is directly computable.
\end{itemize}

\paragraph{Backward synthesis of length requirements.}
When a downstream site indexes a list at position $i$, the error site for
\py{IndexError} has viability predicate $0 \le i < \mathrm{len}(xs)$. The
backward pass should pull this requirement to the input boundary. If $i$ is a
constant, the requirement becomes $\mathrm{len}(xs) > i$. If $i$ is symbolic
but bounded by an upstream guard or computation, the requirement becomes the
appropriate symbolic inequality.

The same logic applies to slice bounds, unpacking sites, and reduction calls.
For example, \py{xs[0]} requires $\mathrm{len}(xs) \ge 1$;
\py{a, b, c = xs} requires $\mathrm{len}(xs) = 3$; \py{min(xs)} requires
$\mathrm{len}(xs) \ge 1$ unless an explicit default is provided.

\paragraph{Forward synthesis of length guarantees.}
After a list comprehension \py{ys = [f(x) for x in xs]}, the output site should
carry $\mathrm{len}(ys) = \mathrm{len}(xs)$. After a filter comprehension
\py{ys = [x for x in xs if p(x)]}, the output site carries
$\mathrm{len}(ys) \le \mathrm{len}(xs)$. After \py{xs.append(v)}, the
post-mutation site carries $\mathrm{len}(xs') = \mathrm{len}(xs) + 1$. These
forward facts become part of the output section and flow into downstream sites.

\subsection{Element-type refinement and homogeneity}

Python lists are heterogeneous in principle, but in practice most lists are
homogeneous. The system should track element-type information when it is
statically visible or dynamically observed. A list literal \py{[1, 2, 3]} has
element type \py{int}. A typed comprehension \py{[int(x) for x in xs]} has
element type \py{int}. A list received as a parameter with annotation
\py{list[str]} has element type \py{str}. Element-type information flows into
downstream operations: \py{sum(xs)} requires numeric elements; \py{"".join(xs)}
requires string elements; \py{sorted(xs)} requires comparable elements.

\subsection{Mutation, aliasing, and frame conditions for lists}

Lists are mutable, so the sheaf-descent system must track mutation sites.
An \py{xs.append(v)} call creates a post-mutation site where the length
has increased by one and the last element is \py{v}. An \py{xs.pop()} call
creates a post-mutation site where the length has decreased by one, with an
error site requiring $\mathrm{len}(xs) \ge 1$. An \py{xs[i] = v} call creates
a mutation site where the element at position $i$ has changed, with an error
site requiring $0 \le i < \mathrm{len}(xs)$.

Aliasing matters because two variables may refer to the same list object. If
\py{ys = xs} and then \py{ys.append(v)}, the mutation is visible through
\py{xs}. The heap-SSA layer should track this aliasing so that the sheaf-descent
system can propagate mutation effects through aliases. The frame condition for a
function that receives a list should state which list arguments may be mutated
and which are treated as read-only.

\subsection{Sorting, searching, and order-sensitive list operations}

The list method \py{xs.sort()} is an in-place mutation that establishes a
sortedness invariant on the post-mutation site. The built-in \py{sorted(xs)}
returns a new sorted list. In both cases, the output site should carry a
\py{sorted} predicate. Downstream operations like \py{bisect.bisect_left(xs, v)}
require a sorted input; the backward pass should synthesize a sortedness
requirement or verify that the upstream site already carries the invariant.

\subsection{List-specific error sites}

The following error sites should be registered for list operations:
\begin{itemize}[nosep]
  \item \py{IndexError}: access at index $i$ when $i \notin [0, \mathrm{len})$
        or $i \notin [-\mathrm{len}, -1]$ for negative indexing.
  \item \py{ValueError}: \py{xs.remove(v)} when $v \notin xs$.
  \item \py{IndexError}: \py{xs.pop(i)} when $i$ is out of range.
  \item Unpacking mismatch: \py{a, b = xs} when $\mathrm{len}(xs) \neq 2$.
  \item Empty-sequence reduction: \py{min(xs)}, \py{max(xs)} when
        $\mathrm{len}(xs) = 0$ and no default is provided.
\end{itemize}

Each error site should have an explicit viability predicate, a backward pullback
rule, and an explanation template so that the renderer can say exactly why an
\py{IndexError} is possible or has been excluded.

\subsection{Interaction with the broader sheaf architecture}

List sites participate in the same cover as all other program sites. A function
that receives a list, guards its length, indexes it, and returns an element
should produce a global section where the input boundary carries a length
requirement, the guard site narrows the local section, the indexing site is
verified safe, and the output boundary carries an element-type guarantee. The
same backward and forward passes that handle tensor shapes should handle list
lengths, because both are instances of the same viability and extension
machinery. The only difference is the local theory: tensor shapes use
semialgebraic arithmetic over multi-dimensional extents, while list lengths use
simpler linear arithmetic over a single dimension.

\section{Deep dive: dict semantics, key safety, and structural contracts}

Python dictionaries are the backbone of configuration, serialization, keyword
dispatch, and ad hoc record types. They are also a major source of
\py{KeyError}, \py{TypeError}, and silent logic bugs when expected keys are
missing or carry unexpected types. A sheaf-descent system should model
dictionaries with enough structure to synthesize key-membership requirements,
value-type guarantees, and structural contracts.

\subsection{Key domain as a refinement coordinate}

The most valuable fact about a dictionary is its key domain: the set of keys
that are definitely present. A dict literal \py{\{"a": 1, "b": 2\}} has exact
key domain $\{\text{"a"}, \text{"b"}\}$. A dict built by successive
\py{d[k] = v} assignments accumulates keys. A dict received as a parameter may
have an unknown key domain, but guards like \py{if "x" in d} or
\py{d.get("x", default)} refine the local section.

The system should distinguish three levels of key-domain knowledge:
\begin{enumerate}[nosep]
  \item \emph{Exact domain}: all keys are known, as in a literal or a fully
        traced construction.
  \item \emph{Lower-bounded domain}: certain keys are known present due to
        guards, construction patterns, or TypedDict annotations, but additional
        keys may exist.
  \item \emph{Unknown domain}: no key-membership information is available.
\end{enumerate}

\paragraph{Guard harvesting for key membership.}
The harvest layer should recognize:
\begin{itemize}[nosep]
  \item \py{if k in d} --- establishes key presence in the true branch.
  \item \py{if k not in d} --- establishes key absence in the true branch.
  \item \py{d.get(k, default)} --- safe access that avoids \py{KeyError} but
        introduces an optionality section on the result.
  \item \py{d.setdefault(k, v)} --- establishes key presence after the call.
  \item \py{try: v = d[k] except KeyError: ...} --- the try body site sees
        key presence; the except body sees key absence.
\end{itemize}

\subsection{Value-type refinement per key}

For dictionaries used as records or configuration objects, the system should
track per-key value types. A \py{TypedDict} annotation provides this directly.
Without annotations, the system can infer per-key types from construction
patterns: \py{d = \{"name": "Alice", "age": 30\}} gives
$d[\text{"name"}] : \py{str}$ and $d[\text{"age"}] : \py{int}$.

When a function receives a dict and accesses specific keys with type-sensitive
operations, the backward pass should synthesize per-key type requirements. For
example, if the function calls \py{int(d["count"])}, the backward pass should
require that \py{"count"} is present and that its value supports conversion to
\py{int}.

\subsection{Dict-specific error sites}

The following error sites should be registered:
\begin{itemize}[nosep]
  \item \py{KeyError}: \py{d[k]} when $k$ is not in the key domain and no
        guard or \py{get} alternative is used.
  \item \py{TypeError}: using an unhashable key type.
  \item \py{AttributeError}: calling dict-specific methods on a non-dict
        value that was expected to be a dict.
  \item Structural mismatch: passing a dict that lacks required keys to a
        function expecting a particular record shape.
\end{itemize}

\paragraph{Backward synthesis of key requirements.}
When the function accesses \py{d["x"]}, \py{d["y"]}, and \py{d["z"]}, the
backward pass should synthesize the input requirement
$\{\text{"x"}, \text{"y"}, \text{"z"}\} \subseteq \mathrm{keys}(d)$.
If any access is guarded, only the unguarded accesses contribute to the
mandatory requirement.

\paragraph{Forward synthesis of structural guarantees.}
When a function constructs and returns a dict, the output section should
describe the key domain and per-key value types. This is especially valuable
for JSON-like pipelines and configuration builders where downstream consumers
depend on specific keys being present.

\subsection{TypedDict and structural dict contracts}

Python's \py{TypedDict} provides a static annotation for dict structure. The
sheaf-descent system should treat \py{TypedDict} annotations as boundary seeds
that pin the key domain and per-key types. But the system should also be able
to infer TypedDict-like structural contracts from unannotated code. If a
function always constructs a dict with keys \py{"status"}, \py{"data"}, and
\py{"error"}, the output section should reflect that structure even without an
explicit \py{TypedDict} annotation.

\subsection{Dict comprehensions and transformations}

Dict comprehensions like \py{\{k: f(v) for k, v in items\}} should propagate
key-domain information from the source. If the source has a known key domain,
the result has the same domain with transformed values. Filtering
comprehensions like \py{\{k: v for k, v in d.items() if p(k)\}} produce a
subset domain.

The \py{dict.update()} method and the \py{|} merge operator should be modeled
as domain-union operations. After \py{d1.update(d2)}, the post-mutation site
has $\mathrm{keys}(d1') \supseteq \mathrm{keys}(d1) \cup \mathrm{keys}(d2)$
and value types from $d2$ override those from $d1$ for shared keys.

\subsection{Interaction with JSON, serialization, and API boundaries}

Many Python programs receive dicts from JSON parsing, API responses, or
configuration files. These are precisely the places where key-membership
errors are most common and most painful. The sheaf-descent system should
treat deserialization boundaries as input sites with rich structural
requirements. If downstream code accesses specific keys, the backward pass
should synthesize the structural contract that the input JSON must satisfy.
This is a high-value practical deliverable that requires no theorem-level
sophistication---just careful key-domain tracking and backward propagation.

\section{Deep dive: heap semantics, object state, and mutation tracking}

Python is a heap-heavy language. Objects carry mutable fields, classes define
initialization protocols, methods read and write instance state, and aliasing is
pervasive. A serious dependent environment must model enough heap structure to
catch real bugs---missing field initialization, stale alias reads, frame
violations---without attempting a full separation-logic encoding that would be
impractical for a Python tool.

\subsection{Object state as a local section}

Every object-valued SSA site should carry a local section describing the
object's known state: which fields are definitely initialized, what types those
fields carry, whether the object satisfies a particular protocol, and what
mutation frame applies. This is the heap analog of carrying a length coordinate
on a list site or a shape coordinate on a tensor site.

\paragraph{Field initialization tracking.}
The constructor body (\py{__init__}) is the primary site where fields are
established. Each \py{self.field = value} statement creates a post-assignment
site where the field is marked initialized with a particular type. At the end
of the constructor, the object should carry a section listing all initialized
fields. If a field is conditionally initialized (inside an \py{if} branch), the
post-constructor section should reflect that the field is only \emph{possibly}
initialized, creating an obstruction that the backward pass can convert into an
input requirement on the branch condition.

\paragraph{Field read safety.}
An \py{AttributeError} error site should be created at every field read
\py{self.field} or \py{obj.field}. The viability predicate requires that the
field is initialized in the current heap section. The backward pass should pull
this requirement to the constructor or to the call site that created the object.

\paragraph{Example: conditional initialization bug.}
Consider a class whose constructor initializes \py{self.data} only when a flag
is true. Any method that reads \py{self.data} unconditionally has a reachable
\py{AttributeError}. The sheaf-descent system should detect this by noting that
the field-initialization site is branch-dependent while the field-read site is
unconditional, producing an obstruction at the overlap between constructor
output and method input.

\subsection{Mutation frames and aliasing}

When a method mutates an object's field, the post-mutation site should update
the field's type and value information. When two variables alias the same object
(\py{y = x}), mutations through one alias must be visible at the other. The
heap-SSA layer should track aliasing so that the cover builder can insert
appropriate overlap edges between sites that share aliased heap locations.

A practical simplification is to track \emph{frame conditions} rather than full
alias analysis. A frame condition for a function states which heap locations may
be modified and which are preserved. If a function's frame condition says it
does not modify \py{obj.data}, then callers can assume that \py{obj.data}
retains its pre-call section after the call returns.

\subsection{Ownership and resource semantics}

Some Python objects represent resources with ownership semantics: file handles,
database connections, locks, GPU memory allocations. The sheaf-descent system
should support ownership sections that track whether a resource is open, closed,
acquired, or released. An error site should be created at every use of a
resource that requires it to be in a particular state. For example, calling
\py{f.read()} after \py{f.close()} should produce an error site with viability
predicate ``file is open.''

Context managers (\py{with} statements) provide natural ownership boundaries.
The \py{__enter__} site establishes the resource as acquired; the
\py{__exit__} site marks it as released. Between them, the resource is
available. After the \py{with} block, the resource should not be used. The
system should synthesize these ownership transitions automatically for standard
context-manager patterns.

\subsection{Class hierarchies and method resolution}

When a method is called on an object, the system must resolve which
implementation will execute. For concrete classes this is straightforward. For
abstract base classes or protocol-typed parameters, the resolution depends on
the runtime type. The sheaf-descent system should model method resolution as a
dispatch site where the viability predicate requires that the receiver's type
supports the called method. This integrates naturally with protocol reasoning
(covered in the next section) because protocol conformance is exactly the
statement that the dispatch site is viable.

\subsection{Dataclasses, named tuples, and structured constructors}

Python's \py{dataclasses}, \py{typing.NamedTuple}, and \py{attrs} classes
provide structured constructors with known fields. These should be treated as
special cases where the constructor section is fully determined by the class
definition. Every field is initialized, every field has a declared type, and
the post-construction section is exact. This makes dataclass-heavy code
especially amenable to automatic analysis.

\subsection{Heap-related error sites}

\begin{itemize}[nosep]
  \item \py{AttributeError}: field not initialized or not present on the
        receiver's type.
  \item Stale alias: reading a field through an alias after it was mutated
        through another path, leading to unexpected values.
  \item Resource use after release: calling methods on a closed file, released
        lock, or deallocated buffer.
  \item Frame violation: a function modifies heap state that its caller assumed
        was preserved.
  \item Constructor invariant failure: an object is used before its constructor
        has established the required field set.
  \item Slot mismatch: accessing a field not declared in \py{__slots__}.
\end{itemize}

Each of these should have a viability predicate, a backward pullback rule to
the constructor or caller boundary, and an explanation template suitable for
rendering to ordinary Python developers.

\section{Deep dive: protocols, structural typing, and duck typing}

Python's type system is fundamentally structural in practice: if an object has
the right methods and attributes, it works, regardless of its nominal class
hierarchy. The \py{typing.Protocol} mechanism formalizes this, but most Python
code relies on duck typing without explicit protocol annotations. A
sheaf-descent system must handle both annotated and unannotated structural
typing with enough precision to catch protocol mismatches and synthesize
structural requirements.

\subsection{Protocol conformance as a local section}

When a function calls \py{obj.read()}, \py{obj.write(data)}, or
\py{len(obj)}, it implicitly requires that the receiver conforms to a protocol:
the object must have a \py{read} method, a \py{write} method, or a
\py{__len__} method. Each such method access creates a protocol-dispatch site
with a viability predicate requiring the method's existence and correct
signature.

The sheaf-descent system should infer these protocol requirements automatically
from method calls, attribute accesses, and special-method invocations
(\py{__iter__}, \py{__getitem__}, \py{__contains__}, \py{__enter__},
\py{__exit__}, etc.). The backward pass should collect all required methods
and attributes into a synthesized protocol section on the input boundary.

\paragraph{Example: synthesized protocol requirement.}
Given a function:
\begin{lstlisting}[style=deppy,language=Python]
def process(stream):
    header = stream.readline()
    data = stream.read()
    stream.close()
    return header, data
\end{lstlisting}
The system should synthesize the input requirement that \py{stream} must have
methods \py{readline() -> str}, \py{read() -> str}, and \py{close() -> None}.
This is equivalent to conformance with a protocol containing those three methods.

\subsection{Explicit Protocol annotations as boundary seeds}

When a user provides \py{typing.Protocol} annotations, those should be treated
as boundary seeds that pin the protocol section. The system should verify that
the function body's usage is consistent with the declared protocol and that
callers provide objects conforming to it. If the function uses methods beyond
those declared in the protocol, the system should report an obstruction: the
declared protocol is weaker than what the body requires.

\subsection{Protocol subtyping and refinement}

Protocol subtyping is structural: a protocol $P$ is a subtype of $Q$ if $P$
has all the methods of $Q$ with compatible signatures. The sheaf-descent system
should compute protocol subtyping at overlap sites where a value flows from a
site with one protocol assumption to a site with another. If the downstream
protocol is not a subtype of the upstream protocol, the overlap produces an
obstruction.

Protocol refinement goes further: a refined protocol may not only require
methods but also attach semantic contracts to them. For example, an
\py{OrderedContainer} protocol might require not just \py{__getitem__} but also
that the elements are accessible in sorted order. This is where protocol
reasoning connects to the broader dependent typing story.

\subsection{Duck typing without annotations}

Most Python code uses duck typing without explicit \py{Protocol} declarations.
The sheaf-descent system should handle this by inferring implicit protocols from
usage patterns. If a function calls \py{obj.fit(X, y)} and
\py{obj.predict(X)}, the system infers a ``sklearn-like estimator'' protocol
even without an explicit annotation. This inferred protocol becomes a local
section on the input boundary.

The quality of this inference depends on how completely the system harvests
method calls, attribute accesses, and special-method invocations from the
function body. The harvest layer should be thorough, and the resulting protocol
section should be renderable as a concrete \py{Protocol} class definition that
the user could paste into their code if they wanted to make the requirement
explicit.

\subsection{Protocol-related error sites}

\begin{itemize}[nosep]
  \item \py{AttributeError}: the receiver does not have the required method
        or attribute.
  \item \py{TypeError}: the method exists but has an incompatible signature
        (wrong number of arguments, wrong argument types).
  \item Protocol mismatch: a caller passes an object that satisfies some but
        not all of the required protocol methods.
  \item Semantic protocol violation: the method exists and has the right
        signature but does not satisfy a semantic contract (e.g., a
        \py{__lt__} that does not define a total order).
\end{itemize}

\subsection{Iterable, iterator, and sequence protocols}

The iteration protocol (\py{__iter__}, \py{__next__}) deserves special
treatment because it is used in \py{for} loops, comprehensions, unpacking,
and many library functions. The system should model:
\begin{itemize}[nosep]
  \item Whether an object is iterable (has \py{__iter__}).
  \item Whether iteration is finite (known or bounded length) or potentially
        infinite.
  \item The element type yielded by iteration.
  \item Whether the object supports random access (\py{__getitem__} with
        integer keys), which upgrades it from \py{Iterable} to \py{Sequence}.
  \item Whether the object supports \py{__len__}, which enables length-based
        reasoning.
\end{itemize}

These protocol facts should flow into downstream sites. A \py{for x in xs}
loop body site should see that \py{x} has the element type of \py{xs}. A
\py{list(xs)} call should produce a list with the same element type and, if
the length is known, the same length.

\subsection{Callable protocols and higher-order functions}

Functions are first-class objects in Python, and higher-order patterns are
common: callbacks, decorators, \py{map}, \py{filter}, \py{functools.reduce}.
The system should model callable protocols:
\begin{itemize}[nosep]
  \item The number and types of parameters.
  \item The return type.
  \item Side effects (mutation frame, I/O, exceptions).
  \item Semantic contracts on the callable (e.g., a comparator must define a
        total order; a predicate must return a boolean).
\end{itemize}

When a higher-order function receives a callable argument, the backward pass
should synthesize a callable-protocol requirement. When the callable is applied
inside the function, the forward pass should propagate the callable's return
type and semantic properties into the result section.

\section{Deep dive: decorators, metaprogramming, and dynamic class construction}

Python's decorator mechanism, metaclasses, and dynamic class construction
patterns are among the most challenging features for any static analysis tool.
A sheaf-descent system cannot pretend they do not exist; it must model them
well enough to handle common patterns and degrade gracefully on exotic ones.

\subsection{Decorators as semantic transformers}

A decorator \py{@deco} applied to a function \py{f} replaces \py{f} with
\py{deco(f)}. The system must model how this transformation affects the
function's semantic section. Several important cases arise:

\paragraph{Identity-preserving decorators.}
Some decorators (\py{@staticmethod}, \py{@classmethod}, \py{@property},
\py{@functools.wraps}, many logging or timing decorators) preserve the
essential signature and semantics of the wrapped function. The system should
recognize these as transparent wrappers and propagate the inner function's
semantic section through the decorator site.

\paragraph{Signature-changing decorators.}
Some decorators change the function's signature: adding parameters, removing
parameters, changing return types, or wrapping the return value. For example,
\py{@click.command()} transforms a function into a CLI command object;
\py{@torch.no_grad()} wraps a function to disable gradient computation; retry
decorators may change the exception profile. The system should model these
through decorator-specific theory packs that describe how the output section
relates to the input section.

\paragraph{Validation decorators.}
Decorators like \py{@pydantic.validate_call}, \py{@beartype.beartype}, or
custom input-validation decorators add runtime checks. These should be treated
as guard-insertion sites: the decorated function's body site sees inputs that
have passed the validation checks, strengthening the local section.

\paragraph{DepPy's own decorators.}
The \py{@deppy.requires}, \py{@deppy.ensures}, \py{@deppy.theorem}, and
\py{@deppy.decreases} decorators are boundary seeds in the sheaf-descent
architecture. They should be parsed by the frontend, converted into boundary
section constraints, and removed from the runtime execution path. The system
should verify that the function body is consistent with the declared contracts
and that callers satisfy the declared requirements.

\subsection{Metaclasses and dynamic class construction}

Python supports metaclasses (\py{type} subclasses), \py{__init_subclass__},
\py{__set_name__}, and dynamic class creation via \py{type("Name", bases,
namespace)}. These features are used by frameworks like Django, SQLAlchemy,
and dataclasses. A full treatment is beyond the scope of a first
implementation, but the system should handle the common cases:

\begin{itemize}[nosep]
  \item \py{@dataclass}: treat as a structured constructor with known fields
        and types.
  \item \py{NamedTuple}: treat as an immutable record with known fields.
  \item \py{Enum}: treat as a finite value domain with known members.
  \item Django/SQLAlchemy models: treat as objects with fields declared in the
        class body, analogous to dataclasses but with ORM-specific semantics.
\end{itemize}

For unrecognized metaclass patterns, the system should emit a residual
explaining that the class's field structure cannot be determined statically.
This is honest and useful: the user learns exactly where the analysis reaches
its limits.

\subsection{Dynamic attribute assignment and \py{__getattr__}}

Python allows dynamic attribute assignment (\py{obj.new_field = value}) and
custom attribute resolution (\py{__getattr__}, \py{__getattribute__}). These
make heap reasoning harder because the set of fields is not statically
bounded. The system should handle three tiers:

\begin{enumerate}[nosep]
  \item Classes with \py{__slots__}: the field set is exactly known.
  \item Classes without \py{__slots__} but with a standard \py{__init__}:
        the field set is determined by analyzing the constructor.
  \item Classes with \py{__getattr__} or metaclass magic: the field set is
        unknown and field accesses should generate residual protocol
        requirements rather than concrete field-initialization checks.
\end{enumerate}

\subsection{Monkey-patching and module-level mutation}

Python allows monkey-patching: modifying classes or modules at runtime. The
system should not attempt to track arbitrary monkey-patches, but it should
recognize common patterns:
\begin{itemize}[nosep]
  \item Module-level configuration: \py{config.DEBUG = True} at import time.
  \item Test fixtures that replace methods or attributes for testing.
  \item Plugin systems that register handlers dynamically.
\end{itemize}

For these patterns, the system should model the mutation effect at the
module-initialization site and propagate the resulting state to downstream
sites. If a monkey-patch invalidates a previously established section, the
system should report the obstruction.

\subsection{Decorator stacking and composition}

Multiple decorators can be stacked:
\begin{lstlisting}[style=deppy,language=Python]
@deppy.ensures(lambda r: r >= 0)
@functools.lru_cache(maxsize=128)
@deppy.requires(lambda n: n >= 0)
def fibonacci(n: int) -> int:
    ...
\end{lstlisting}
The system should process decorators from innermost to outermost. The
\py{@deppy.requires} creates an input boundary seed. The \py{@lru_cache}
wraps the function but preserves its semantic section (caching does not change
the function's logical contract). The \py{@deppy.ensures} creates an output
boundary seed. The resulting section should carry both the input requirement
and the output guarantee.

\section{Deep dive: the optional annotation system and annotation economy}

A central thesis of this monograph is that annotations should be optional seeds,
not mandatory prerequisites. The sheaf-descent architecture gives this thesis a
precise meaning: an annotation is a user-supplied local section or boundary
constraint that eliminates ambiguity, strengthens the global section, or
resolves an obstruction. The system should work at every annotation level from
zero to fully explicit, with smooth and predictable improvements as the user
adds more.

\subsection{Zero-annotation mode}

In zero-annotation mode, the system relies entirely on syntax, control flow,
library semantics, and optional dynamic traces. The cover is built from AST
and SSA structure. Theory packs create sites from recognized library patterns.
Guard harvesting extracts implicit refinements from branches and assertions.
The backward and forward passes run to compute synthesized input requirements
and output guarantees.

The quality of zero-annotation analysis depends on three factors:
\begin{enumerate}[nosep]
  \item How completely the harvest layer extracts guards, assertions, and
        construction patterns.
  \item How rich the library theory packs are for the libraries actually used.
  \item How deeply interprocedural transport preserves semantics across call
        boundaries.
\end{enumerate}

For simple Torch modules, zero-annotation analysis should already catch shape
errors, index errors, and values-versus-indices confusion. For complex
algorithm code, zero-annotation analysis may leave significant residuals, but
it should still identify which sites carry ambiguity and what kind of seed
would resolve it.

\subsection{Lightweight annotation: type hints}

Python's standard type hints (\py{def f(x: int) -> str}) provide carrier-type
information. These are already seeds: they pin the carrier coordinate of the
input and output boundary sections. The system should use them to:
\begin{itemize}[nosep]
  \item Narrow element types for generic containers: \py{list[int]},
        \py{dict[str, float]}, \py{Optional[Tensor]}.
  \item Establish protocol conformance: \py{Iterable[T]}, \py{Callable[[int],
        str]}, \py{Sequence[float]}.
  \item Enable richer downstream inference: knowing that \py{x: int} allows
        arithmetic-site solvers to fire; knowing that \py{xs: list[int]}
        enables length and element-type reasoning.
\end{itemize}

Standard type hints are cheap and familiar. The system should make the most of
them by treating them as lightweight but valuable boundary seeds rather than
ignoring them because they are not full dependent types.

\subsection{Medium annotation: deppy decorators}

The \py{@deppy.requires} and \py{@deppy.ensures} decorators allow users to
state refinement predicates:
\begin{lstlisting}[style=deppy,language=Python]
@deppy.requires(lambda xs: len(xs) > 0)
@deppy.ensures(lambda result, xs: result in xs)
def first_element(xs: list[int]) -> int:
    return xs[0]
\end{lstlisting}
These decorators are boundary seeds that resolve specific obstructions. The
\py{requires} clause pins the input boundary section so that the
\py{IndexError} site at \py{xs[0]} becomes empty. The \py{ensures} clause
pins the output boundary section so that downstream consumers know the result
is an element of the original list.

The system should not require both. If only \py{requires} is provided, the
system should synthesize the strongest output section justified by the
requirement. If only \py{ensures} is provided, the system should synthesize
the weakest input section needed to support the guarantee. If neither is
provided, both are synthesized. If both are provided, the system should verify
consistency and report any obstruction between them.

\subsection{Heavy annotation: theorem decorators and proof modules}

For proof-oriented users, the system supports theorem decorators, ghost
parameters, witness-returning functions, and explicit transport lemmas. These
are heavy annotations in the sense that they carry high semantic payload, but
they should still be expressed in native Python syntax:
\begin{lstlisting}[style=deppy,language=Python]
@deppy.theorem
def sort_preserves_permutation(xs: list[int]):
    result = sorted(xs)
    deppy.witness(deppy.permutation(result, xs))
    return result
\end{lstlisting}

The theorem decorator tells the system to treat this function as a proof
artifact rather than a runtime computation. The \py{deppy.witness} call
exports a witness section. The proof elaborator converts this into a proof
site with a transport map that downstream consumers can import.

\subsection{Annotation economy: measuring the value of each annotation}

The sheaf-descent architecture enables a precise measurement of annotation
value. When a user adds an annotation, the system can compute:
\begin{itemize}[nosep]
  \item How many obstruction coordinates the annotation resolves.
  \item How many error sites become provably safe.
  \item How many output-section coordinates become stronger.
  \item How many downstream functions benefit through interprocedural transport.
\end{itemize}

This information should be exposed to the user through the renderer. A
diagnostic like ``adding \py{@deppy.requires(lambda xs: len(xs) > 0)} would
resolve 3 error sites and strengthen 2 output contracts'' is far more useful
than a generic ``add a precondition.'' This is annotation economy: helping
users invest their annotation effort where it provides the most semantic
return.

\subsection{LLM-authored annotations}

Large language models can generate annotations, contracts, and even theorem
modules. The system should treat LLM-generated annotations the same as
user-written ones: they are boundary seeds that must be verified. If an LLM
suggests \py{@deppy.ensures(lambda r: sorted(r))}, the system should check
whether the function body actually supports this claim. If verification
succeeds, the annotation is accepted. If it fails, the system should report
the specific obstruction that the LLM's suggestion did not address.

This creates a natural feedback loop: the LLM suggests, the system verifies,
and unresolved obstructions guide the next suggestion. The key architectural
point is that LLM outputs never bypass the trust boundary. They are treated as
unverified seeds until the kernel or solver confirms them.

\subsection{Gradual adoption path}

The annotation system should support gradual adoption:
\begin{enumerate}[nosep]
  \item Start with zero annotations; the system infers what it can.
  \item Add standard type hints to improve carrier-type reasoning.
  \item Add \py{@deppy.requires}/\py{@deppy.ensures} where the system
        reports high-value obstructions.
  \item Add theorem modules for algorithmic correctness where full dependent
        guarantees are desired.
\end{enumerate}

At each step, the system should report what has improved and what remains. The
user should never feel that they must annotate everything or nothing. The
sheaf-descent architecture makes this possible because every annotation is a
local operation on the cover, not a global mode switch.

\section{Deep dive: PyTorch shape theory and symbolic shape inference}

Shape errors are the most common runtime failure mode in tensor programs. A
reshape with incompatible cardinality, a matmul with misaligned inner
dimensions, a broadcast between incompatible extents, or a reduction over
a nonexistent axis---these all produce opaque \py{RuntimeError} messages that
could have been caught statically. The sheaf-descent system should provide a
comprehensive shape theory that can infer output shapes as symbolic functions
of input shapes and synthesize the weakest input shape requirements that
prevent all shape-related runtime errors.

\subsection{Shape representation}

A tensor's shape is a tuple of non-negative integers. The system should
represent shapes symbolically so that relationships between input and output
shapes can be expressed as equations and inequalities.

\begin{definition}[Symbolic shape]
A symbolic shape is a tuple $(d_1, d_2, \ldots, d_r)$ where each $d_i$ is
either a concrete non-negative integer, a symbolic variable, or a symbolic
expression over variables. The rank $r$ may itself be symbolic or concrete.
\end{definition}

For example, if a function receives a tensor \py{x} with unknown shape, the
system represents its shape as $(d_1, d_2, \ldots, d_r)$ with symbolic
variables. After \py{y = x.view(d_1, -1)}, the output shape is
$(d_1, \prod_{i=2}^{r} d_i)$, which can be expressed symbolically as
$(d_1, d_1^{-1} \cdot \prod_i d_i)$.

\subsection{Shape rules for core operations}

Each tensor operation has a shape rule that relates input shapes to output
shapes. The system should encode these rules as site-local viability predicates
and output-shape functions.

\paragraph{Element-wise operations.}
For element-wise operations (\py{+}, \py{-}, \py{*}, \py{/}, \py{torch.relu},
etc.), the output shape equals the broadcast-compatible shape of the inputs.
The viability predicate requires that the inputs are broadcast-compatible: for
each dimension, either the extents match, or one of them is 1. The output
extent in each dimension is the maximum of the input extents.

\paragraph{Matrix multiplication.}
For \py{torch.matmul(A, B)} or \py{A @ B}:
\begin{itemize}[nosep]
  \item If $A$ has shape $(\ldots, m, k)$ and $B$ has shape $(\ldots, k, n)$,
        the output has shape $(\ldots, m, n)$.
  \item The viability predicate requires that the inner dimensions match:
        $A.\mathrm{shape}[-1] = B.\mathrm{shape}[-2]$.
  \item Batch dimensions must be broadcast-compatible.
\end{itemize}

\paragraph{Reshape and view.}
For \py{x.view(s_1, s_2, \ldots, s_k)} or \py{x.reshape(s_1, \ldots, s_k)}:
\begin{itemize}[nosep]
  \item The viability predicate requires $\prod_i s_i = \prod_j d_j$ where
        $d_j$ are the input dimensions (with at most one $s_i = -1$,
        representing the inferred dimension).
  \item For \py{view} specifically, contiguity of the input tensor is also
        required; if the input is not contiguous, a \py{RuntimeError} occurs.
  \item The output shape is $(s_1, s_2, \ldots, s_k)$ with the $-1$ dimension
        resolved to $\prod_j d_j / \prod_{i \neq -1} s_i$.
\end{itemize}

\paragraph{Transpose and permute.}
For \py{x.transpose(dim0, dim1)}, the output shape swaps dimensions
\py{dim0} and \py{dim1}. The viability predicate requires that both
dimension indices are valid: $0 \le \py{dim0}, \py{dim1} < r$. For
\py{x.permute(dims)}, the output shape is the input shape permuted
according to \py{dims}, with viability requiring that \py{dims} is a
valid permutation of $\{0, \ldots, r-1\}$.

\paragraph{Reduction operations.}
For \py{x.sum(dim=d)}, \py{x.mean(dim=d)}, etc.:
\begin{itemize}[nosep]
  \item If \py{keepdim=False}, the output shape is the input shape with
        dimension $d$ removed.
  \item If \py{keepdim=True}, the output shape has dimension $d$ set to 1.
  \item The viability predicate requires $0 \le d < r$.
  \item For operations that require nonempty input along the reduced dimension
        (e.g., \py{argmax}, \py{argmin}), the viability predicate additionally
        requires $d_{\py{dim}} > 0$.
\end{itemize}

\paragraph{Concatenation and stacking.}
For \py{torch.cat(tensors, dim=d)}:
\begin{itemize}[nosep]
  \item All tensors must have the same number of dimensions.
  \item All dimensions except $d$ must match across all tensors.
  \item The output dimension $d$ is the sum of the input dimensions at $d$.
\end{itemize}
For \py{torch.stack(tensors, dim=d)}, a new dimension is inserted at
position $d$ with extent equal to the number of tensors.

\paragraph{Indexing and gathering.}
For \py{x[index]}, \py{torch.gather(x, dim, index)}, and
\py{torch.index_select(x, dim, index)}:
\begin{itemize}[nosep]
  \item The viability predicate requires that all index values lie within
        $[0, x.\mathrm{shape}[\py{dim}])$.
  \item The output shape depends on the indexing mode: basic indexing may
        reduce dimensions; advanced indexing may reshape according to the
        index tensor's shape; \py{gather} preserves the index tensor's shape.
\end{itemize}

\paragraph{Convolution.}
For \py{torch.nn.Conv2d(in_channels, out_channels, kernel_size, ...)}
applied to input of shape $(N, C_{in}, H, W)$:
\begin{itemize}[nosep]
  \item Viability requires $C_{in}$ matches the layer's \py{in_channels}.
  \item Output shape: $(N, C_{out}, H_{out}, W_{out})$ where
        $H_{out} = \lfloor (H + 2p - k) / s \rfloor + 1$ with padding $p$,
        kernel size $k$, and stride $s$.
  \item Viability requires $H + 2p - k \ge 0$ and similarly for $W$.
\end{itemize}

\subsection{Symbolic shape propagation through a module}

A PyTorch module typically chains several operations. The shape theory should
propagate symbolic shapes through the entire chain, accumulating equations and
viability predicates at each step. For example:
\begin{lstlisting}[style=deppy,language=Python]
def forward(self, x):          # x: (batch, channels, height, width)
    x = self.conv1(x)          # (batch, 32, h1, w1)
    x = x.view(x.size(0), -1) # (batch, 32 * h1 * w1)
    x = self.fc(x)             # (batch, num_classes)
    return x
\end{lstlisting}

The shape theory should:
\begin{enumerate}[nosep]
  \item Start with symbolic input shape $(b, c, h, w)$.
  \item After \py{conv1}: output $(b, 32, h_1, w_1)$ where $h_1, w_1$ are
        functions of $h, w$ and the convolution parameters.
  \item After \py{view}: output $(b, 32 \cdot h_1 \cdot w_1)$. The viability
        predicate is automatically satisfied because the product is preserved.
  \item After \py{fc}: output $(b, n_c)$ with viability requiring
        $32 \cdot h_1 \cdot w_1 = \py{fc.in_features}$.
\end{enumerate}

The backward pass should pull the viability requirement at \py{fc} back to the
input: the input dimensions $(h, w)$ must satisfy
$32 \cdot h_1(h) \cdot w_1(w) = \py{fc.in_features}$. This is the synthesized
input shape requirement for the module.

\subsection{Output shape as a refinement type}

The output section should express the output shape as a symbolic function of
the input shape. For the module above, the output refinement type would be
something like:
$$\py{result} : \py{Tensor}[\text{shape} = (b, n_c)]$$
where $b$ is the batch dimension from the input. This is a genuine dependent
type: the output shape depends on the input shape.

For more complex modules, the output shape may be a piecewise function (if
there are shape-dependent branches) or may contain symbolic expressions
involving input dimensions, layer parameters, and arithmetic operations.

\subsection{Shape error sites and diagnostics}

Each shape-violating operation should create an error site with:
\begin{itemize}[nosep]
  \item The exact viability predicate that was violated.
  \item The symbolic shape values at the point of failure.
  \item A backward trace showing which input shape constraint would prevent
        the error.
  \item A suggested fix in terms of input shape requirements, layer parameter
        changes, or inserted reshape operations.
\end{itemize}

This is the flagship deliverable for Torch users: run the analyzer on a module
and get a precise report of all potential shape errors, together with the
exact input shape constraints that would prevent them.

\subsection{Broadcasting rules and their subtleties}

Broadcasting is one of the most error-prone aspects of tensor programming
because it is implicit and can mask dimension mismatches. The shape theory
should model broadcasting precisely:
\begin{enumerate}[nosep]
  \item Right-align the shapes of the operands.
  \item For each dimension, check that the extents are either equal or one
        of them is 1.
  \item If neither condition holds, report a broadcast mismatch error site.
  \item The output extent in each dimension is the maximum of the operand
        extents.
\end{enumerate}

The backward pass should synthesize input shape requirements that guarantee
broadcast compatibility. If a function adds tensors of shapes $(a, b)$ and
$(c, d)$, the requirement is $a = c \lor a = 1 \lor c = 1$ and
$b = d \lor b = 1 \lor d = 1$.

\subsection{Device and dtype as secondary shape coordinates}

While shape is the primary concern, device and dtype are closely related
coordinates that affect runtime safety:
\begin{itemize}[nosep]
  \item Operations on tensors from different devices produce a
        \py{RuntimeError}.
  \item Some operations are not supported for certain dtypes (e.g., integer
        tensors do not support \py{torch.sin}).
  \item Mixed-precision operations may produce unexpected results.
\end{itemize}

The shape theory should carry device and dtype as secondary coordinates on
tensor sites. The viability predicate for any binary operation should include
device equality, and the output section should specify the result dtype
according to PyTorch's promotion rules.

\section{Deep dive: library-specific theories beyond PyTorch}

The sheaf-descent architecture is not limited to PyTorch. Any library whose
operators have structured semantics---shape rules, protocol requirements,
ordering assumptions, domain constraints---can be modeled through theory packs
that create local sites, viability predicates, and transport maps. This section
sketches theories for several important Python libraries.

\subsection{NumPy shape and broadcasting theory}

NumPy shares PyTorch's shape-oriented semantics but with its own operator
vocabulary and broadcasting rules. A NumPy theory pack should create sites for:
\begin{itemize}[nosep]
  \item \py{np.reshape}, \py{np.resize}: cardinality viability and output shape.
  \item \py{np.dot}, \py{np.matmul}, \py{@}: inner-dimension alignment.
  \item \py{np.concatenate}, \py{np.stack}, \py{np.split}: dimension matching
        and output shape synthesis.
  \item \py{np.sum}, \py{np.mean}, \py{np.argmax}: reduction axis validity
        and non-emptiness.
  \item Broadcasting across all element-wise operations.
  \item \py{np.linalg.inv}, \py{np.linalg.solve}: square-matrix and
        full-rank requirements.
\end{itemize}

The backward pass should synthesize array shape requirements, and the forward
pass should express output shapes as symbolic functions of input shapes.
Because NumPy and PyTorch share many shape rules, the implementation should
factor common shape reasoning into a shared tensor-shape kernel that both
theory packs use.

\subsection{Pandas DataFrame and Series theory}

Pandas introduces a different kind of structural semantics: column schemas,
index alignment, and dtype-dependent operations.

\paragraph{Column-presence theory.}
A DataFrame has a column schema: a set of column names with associated dtypes.
Accessing a column that does not exist produces a \py{KeyError}. The theory
pack should:
\begin{itemize}[nosep]
  \item Create column-presence sites at every column access (\py{df["col"]},
        \py{df.col}, \py{df[["col1", "col2"]]}).
  \item Pull back column-presence requirements to the input boundary.
  \item Propagate column schemas through transformations: \py{df.rename},
        \py{df.drop}, \py{df.merge}, \py{df.assign}.
  \item Synthesize output column schemas for operations that add or remove
        columns.
\end{itemize}

\paragraph{Index alignment theory.}
Pandas operations align on indices, which can produce \py{NaN} values when
indices do not match. The theory should track whether operations require
index alignment and whether misaligned indices produce silent data loss or
explicit errors.

\paragraph{GroupBy and aggregation theory.}
\py{df.groupby("col").agg(func)} requires that \py{"col"} exists and that
\py{func} is applicable to the grouped columns' dtypes. The theory should
create sites for these requirements and propagate the resulting output schema
(grouped columns become the index; aggregated columns have modified dtypes).

\paragraph{Dtype-sensitive operations.}
Many Pandas operations are dtype-sensitive: \py{df.sum()} on string columns
produces concatenation rather than arithmetic; \py{df.mean()} on non-numeric
columns raises an error. The theory should track column dtypes and create
error sites for dtype-incompatible operations.

\subsection{Standard library: pathlib, os, and file I/O}

File I/O introduces resource semantics and path-validity requirements:
\begin{itemize}[nosep]
  \item \py{open(path, mode)}: creates a file handle with open/closed state.
  \item \py{f.read()}, \py{f.write()}: require the file to be open and in the
        correct mode (read vs. write).
  \item \py{f.close()}: transitions the resource to closed state; subsequent
        reads/writes should produce error sites.
  \item \py{pathlib.Path}: methods like \py{.read_text()}, \py{.exists()},
        \py{.mkdir()} have implicit filesystem preconditions.
\end{itemize}

The theory should model file handles as owned resources with state transitions
and create error sites for use-after-close patterns. Context managers
(\py{with open(...) as f}) should be recognized as safe boundaries.

\subsection{Standard library: collections}

The \py{collections} module provides specialized containers:
\begin{itemize}[nosep]
  \item \py{defaultdict}: eliminates \py{KeyError} for missing keys by providing
        a factory. The theory should recognize that \py{defaultdict} access
        sites are always viable but may produce default values.
  \item \py{OrderedDict}: maintains insertion order. Iteration sites should
        carry ordering semantics.
  \item \py{deque}: supports efficient append/pop from both ends. Length
        reasoning applies as with lists, but with additional methods
        (\py{appendleft}, \py{popleft}, \py{rotate}).
  \item \py{Counter}: maps elements to counts. The theory should track that
        values are non-negative integers and that \py{most_common()} returns
        sorted results.
  \item \py{namedtuple}: provides structured access by name. The theory
        should treat named tuples as immutable records with known fields.
\end{itemize}

\subsection{Standard library: itertools and functools}

The \py{itertools} module creates lazy iterators with specific semantic
properties:
\begin{itemize}[nosep]
  \item \py{itertools.chain(*iterables)}: concatenation; output length is
        sum of input lengths.
  \item \py{itertools.product(*iterables)}: Cartesian product; output length
        is the product of input lengths.
  \item \py{itertools.islice(it, n)}: truncation; output length is
        $\min(n, \mathrm{len}(it))$.
  \item \py{itertools.groupby(it, key)}: requires sorted input for correct
        grouping; the theory should create a sortedness requirement site.
  \item \py{itertools.accumulate(it, func)}: output has same length as input;
        func must be a binary operation compatible with element types.
\end{itemize}

The \py{functools} module provides higher-order utilities:
\begin{itemize}[nosep]
  \item \py{functools.reduce(func, it, init)}: requires the iterable to be
        nonempty if \py{init} is not provided; \py{func} must accept two
        arguments of compatible types.
  \item \py{functools.partial(f, *args)}: modifies the callable's signature;
        the theory should compute the resulting parameter list.
  \item \py{functools.lru_cache}: preserves the function's semantic section
        (caching does not change the logical contract).
\end{itemize}

\subsection{SQLAlchemy and database ORM theory}

Database ORMs introduce schema-dependent operations:
\begin{itemize}[nosep]
  \item Model classes define column schemas. Accessing a non-existent column
        should produce an error site.
  \item Query construction (\py{session.query(Model).filter(...)}) should
        verify that filter conditions reference valid columns.
  \item Result materialization (\py{.all()}, \py{.first()}, \py{.one()})
        has different nullability profiles: \py{.one()} raises if zero or
        multiple results exist; \py{.first()} may return \py{None}.
\end{itemize}

The theory should model ORM columns, query construction, and result types
so that the backward pass synthesizes valid query requirements and the
forward pass describes result types and nullability.

\subsection{Requests and HTTP client theory}

The \py{requests} library introduces network-boundary semantics:
\begin{itemize}[nosep]
  \item \py{requests.get(url)}: returns a \py{Response} object whose
        \py{.status_code} may indicate failure.
  \item \py{response.json()}: may raise \py{ValueError} if the response is
        not valid JSON; the theory should create an error site.
  \item \py{response.raise_for_status()}: raises an \py{HTTPError} for
        non-2xx status codes; this acts as a guard that refines the
        subsequent section to assume success.
\end{itemize}

\subsection{Theory-pack extensibility and composition}

A critical architectural property is that theory packs compose. When a function
uses both PyTorch and NumPy, or Pandas and SQLAlchemy, the site cover should
contain sites from all relevant packs, and overlaps should be created at values
that flow between libraries. For example, \py{torch.from_numpy(arr)} creates
an overlap between a NumPy shape site and a PyTorch shape site; the shapes must
agree. The descent engine should handle cross-library overlaps the same way it
handles within-library overlaps.

Adding a new theory pack should be additive: it creates new site families,
new viability predicates, and new transport maps, but it should not invalidate
existing sites or require re-analysis of code that does not use the new library.
This monotonicity property is a direct consequence of the sheaf-descent
architecture: a new pack extends the cover rather than replacing it.

\section{Deep dive: exception handling, control flow, and the Optional/None theory}

Python uses exceptions pervasively---not only for error conditions but also for
control flow in iterators, context managers, and pattern matching. A
sheaf-descent system must model exceptions as first-class semantic events rather
than treating them as an afterthought that escapes the type system.

\subsection{Exceptional control flow as site structure}

Every \py{try}/\py{except}/\py{finally} block defines a family of sites:
\begin{itemize}[nosep]
  \item The try-body site, where execution may raise exceptions.
  \item One except-handler site per handler clause, refined by the caught
        exception type.
  \item The else-body site, which sees the post-try state under the assumption
        that no exception was raised.
  \item The finally-body site, which executes regardless of outcome.
  \item The post-block site, which merges the normal and exceptional flows.
\end{itemize}

The key insight is that except-handler sites act as guard sites: inside
\py{except KeyError}, the system knows that a \py{KeyError} was raised, which
implies that some upstream operation failed in a specific way. This information
can refine the local section. For example, if a \py{KeyError} was caught after
\py{d[k]}, the handler body can infer that $k \notin \mathrm{keys}(d)$.

\subsection{Exception type hierarchies}

Python's exception hierarchy matters for handler matching. A handler for
\py{except Exception} catches almost everything; a handler for \py{except
ValueError} catches only \py{ValueError} and its subclasses. The system should
model the exception type hierarchy so that handler sites carry the correct type
refinement.

\paragraph{Bare except and broad handlers.}
A bare \py{except:} or \py{except Exception:} handler is a catch-all. The
system should recognize these as providing weak refinement---the handler body
knows an exception occurred but not which one. For safety analysis, the system
should warn when a broad handler masks a specific error that the user might
want to handle differently.

\subsection{Exception-raising operations as error sites}

Every operation that can raise an exception should create an error site. The
system already does this for \py{IndexError}, \py{KeyError}, etc. But the
exception theory should go further:
\begin{itemize}[nosep]
  \item User-defined \py{raise} statements create explicit error sites with
        the raised exception type.
  \item \py{assert condition} creates a conditional error site that raises
        \py{AssertionError} when the condition is false.
  \item Iterator exhaustion (\py{StopIteration}) creates error sites at
        \py{next()} calls without defaults.
  \item \py{NotImplementedError} in abstract methods creates error sites at
        dispatch sites when the receiver is abstract.
\end{itemize}

\subsection{The Optional/None theory}

\py{None} is the single most common source of \py{AttributeError} and
\py{TypeError} in Python. The Optional theory should track nullability as a
refinement coordinate on every site where a value might be \py{None}.

\paragraph{Sources of None.}
\begin{itemize}[nosep]
  \item Functions that may return \py{None}: \py{dict.get()},
        \py{list.pop()} with a default, \py{re.match()}, any function with
        an explicit \py{return None} or missing return path.
  \item Default parameter values: \py{def f(x=None)}.
  \item Optional type annotations: \py{x: Optional[int]}.
  \item Conditional initialization: a variable assigned only in one branch.
\end{itemize}

\paragraph{None guards and narrowing.}
The harvest layer should recognize None-checking patterns:
\begin{itemize}[nosep]
  \item \py{if x is not None}: narrows the true branch to non-None.
  \item \py{if x is None}: narrows the true branch to None.
  \item \py{if x}: for types where None is falsy, narrows to non-None in the
        true branch (with caveats for other falsy values like 0 or empty
        strings).
  \item \py{assert x is not None}: narrows the post-assertion site.
  \item \py{x = y if y is not None else default}: the assignment site carries
        non-None when the condition is true.
\end{itemize}

\paragraph{None-related error sites.}
An \py{AttributeError} site should be created at every attribute access or
method call on a value that might be \py{None}. The viability predicate
requires the value to be non-None. The backward pass should synthesize either a
non-None input requirement or require that an upstream guard narrows away the
None case.

\paragraph{Optional chaining and safe navigation.}
Python does not have a safe navigation operator like \py{?.} in other
languages, so Optional handling requires explicit guards. The system should
recognize common patterns:
\begin{lstlisting}[style=deppy,language=Python]
result = x.method() if x is not None else default
\end{lstlisting}
In this pattern, the method-call site is protected by the guard, and the
output site carries the union of the method's return type and the default's
type.

\subsection{Control-flow refinement across branches}

Beyond exceptions and None, the general control-flow story should handle:
\begin{itemize}[nosep]
  \item \py{isinstance(x, T)} guards: narrow the type of \py{x} in the true
        branch to \py{T}.
  \item \py{hasattr(obj, "field")} guards: establish field presence.
  \item Truthiness guards: \py{if xs} on a list narrows to nonempty.
  \item Comparison guards: \py{if n > 0} narrows to positive.
  \item Early returns: \py{if not valid: return None} narrows the
        continuation to the valid case.
\end{itemize}

Each of these patterns should be recognized by the harvest layer and converted
into site-level refinements that participate in the sheaf-descent computation.
The key principle is that control flow is evidence: every branch taken is a
local observation that narrows the admissible sections.

\subsection{Match statements (Python 3.10+)}

Python's \py{match}/\py{case} syntax provides structured pattern matching:
\begin{lstlisting}[style=deppy,language=Python]
match command:
    case {"action": "move", "x": x, "y": y}:
        handle_move(x, y)
    case {"action": "quit"}:
        handle_quit()
    case _:
        handle_unknown(command)
\end{lstlisting}

Each case arm should create a site refined by the pattern: the first arm
knows that \py{command} is a dict with keys \py{"action"}, \py{"x"}, \py{"y"}
and that \py{command["action"] == "move"}. The system should use these
pattern-derived refinements as local section constraints.

\subsection{Generators, yield, and async/await}

Generator functions (\py{yield}) and async functions (\py{async def}/\py{await})
introduce suspended-execution semantics:
\begin{itemize}[nosep]
  \item A generator function's ``return type'' is actually its yield type.
        The output section should describe the type of yielded values.
  \item \py{send()} and \py{throw()} introduce additional interaction sites
        between the generator and its consumer.
  \item \py{async} functions return coroutines; the actual return type is
        available only after \py{await}. The system should model the
        ``unwrapped'' return type at \py{await} sites.
  \item \py{async for} and \py{async with} combine iteration/context-manager
        protocols with async semantics.
\end{itemize}

These are important for real Python code but represent a stretch goal for the
implementation. The system should handle basic generator yield types and async
return types, with more complex interaction patterns generating residuals.

\section{Deep dive: numeric safety---overflow, NaN, division by zero, and floating-point contracts}

Numerical computation is the primary use case for both PyTorch and NumPy, and
numerical errors---NaN propagation, division by zero, overflow, underflow---are
among the most insidious bugs in scientific and machine-learning code. This
section describes how the sheaf-descent framework handles numeric safety as a
first-class theory.

\subsection{The numeric-safety site family}

Every arithmetic operation creates a numeric-safety site. The site's viability
predicate depends on the operation:
\begin{itemize}[nosep]
  \item \textbf{Division} (\py{x / y}, \py{x // y}, \py{x \% y}): viability
        requires $y \neq 0$. For floating-point division, the result may be
        $\pm\infty$ rather than raising \py{ZeroDivisionError}, but this is
        typically still a bug.
  \item \textbf{Square root} (\py{math.sqrt(x)}, \py{np.sqrt(x)},
        \py{torch.sqrt(x)}): viability requires $x \geq 0$ for real-valued
        results.
  \item \textbf{Logarithm} (\py{math.log(x)}, \py{np.log(x)},
        \py{torch.log(x)}): viability requires $x > 0$.
  \item \textbf{Power} (\py{x ** y}): when $x = 0$ and $y < 0$, produces
        \py{ZeroDivisionError} or $\infty$.
  \item \textbf{Integer overflow}: Python's arbitrary-precision integers
        cannot overflow, but NumPy and PyTorch fixed-width integers can.
        The theory should track whether operations on fixed-width types
        might exceed their range.
\end{itemize}

\subsection{NaN propagation theory}

NaN (Not a Number) is a floating-point value that propagates through
arithmetic: any operation involving NaN produces NaN. This makes NaN
a ``contagion'' that can silently corrupt entire computations.

\paragraph{NaN sources.}
Operations that can produce NaN include:
\begin{itemize}[nosep]
  \item \py{0.0 / 0.0}, \py{float('inf') - float('inf')},
        \py{float('inf') * 0.0}.
  \item \py{math.sqrt(-1.0)} in NumPy/PyTorch (returns NaN rather than
        raising).
  \item \py{math.log(0.0)} or \py{math.log(-1.0)}.
  \item User input or data loading from corrupted sources.
\end{itemize}

\paragraph{NaN-freedom as a refinement.}
The system should track a NaN-freedom refinement coordinate: a value is
\emph{NaN-free} if it is guaranteed to contain no NaN values. Operations
that preserve NaN-freedom (e.g., addition of two NaN-free values) should
propagate this property. Operations that might introduce NaN should create
error sites.

\paragraph{NaN checks as guards.}
Calls like \py{torch.isnan(x).any()}, \py{np.isnan(x)}, or
\py{math.isnan(x)} should be recognized as guard sites. After a NaN check
that confirms no NaN is present, the section should carry the NaN-free
refinement.

\subsection{Infinity tracking}

Similar to NaN, $\pm\infty$ can propagate through computations and cause
problems. The system should track finiteness as a refinement coordinate:
\begin{itemize}[nosep]
  \item \py{torch.isinf(x)} and \py{np.isinf(x)} are guard operations.
  \item Operations like \py{torch.clamp(x, min=-M, max=M)} establish
        finiteness and boundedness.
  \item \py{torch.softmax} outputs are always finite and in $[0, 1]$ when
        the input is finite, which the theory should capture.
\end{itemize}

\subsection{Division-by-zero analysis}

Division by zero is the most common numeric error. The theory should handle
several patterns:
\begin{itemize}[nosep]
  \item Direct division: \py{a / b} creates an error site with viability
        $b \neq 0$.
  \item Denominator computation: \py{a / (b + eps)} is a common safe
        pattern; the theory should recognize that adding a positive
        epsilon to the denominator prevents division by zero when $b \geq 0$.
  \item Reduction denominators: \py{x.sum() / x.size(0)} is safe when the
        tensor is nonempty, connecting numeric safety to shape theory.
  \item Normalization: \py{x / x.norm()} requires \py{x.norm() > 0}, which
        is equivalent to $x \neq \mathbf{0}$.
\end{itemize}

\paragraph{Backward synthesis for division safety.}
The backward pass should synthesize input requirements that prevent division
by zero. For \py{a / b}, the requirement is $b \neq 0$. For
\py{a / f(x)}, the requirement is $f(x) \neq 0$, which should be pulled back
through $f$ to a requirement on $x$.

\subsection{Overflow and underflow in fixed-width arithmetic}

While Python integers are arbitrary-precision, NumPy and PyTorch use
fixed-width types (\py{int32}, \py{float32}, etc.). The theory should track:
\begin{itemize}[nosep]
  \item \textbf{Integer overflow}: multiplication or addition of large values
        may wrap around. The theory should create overflow sites when
        operations on bounded types might exceed range.
  \item \textbf{Floating-point underflow}: very small values may be rounded
        to zero, which can interact with division-by-zero analysis.
  \item \textbf{Precision loss}: large floating-point values lose precision;
        \py{x + 1.0 == x} can be true for sufficiently large \py{x}.
  \item \textbf{Dtype promotion}: mixed-type arithmetic follows promotion
        rules (e.g., \py{int32 + float32 -> float64}). The theory should
        track dtype through computations and create sites where unexpected
        promotion occurs.
\end{itemize}

\subsection{Numeric-safety contracts for machine learning}

Machine learning code has domain-specific numeric safety patterns:
\begin{itemize}[nosep]
  \item \textbf{Loss functions}: cross-entropy loss requires probabilities in
        $(0, 1]$; the logarithm of zero produces $-\infty$. The theory should
        verify that softmax or sigmoid has been applied before log-loss.
  \item \textbf{Gradient computation}: gradients can be NaN or $\pm\infty$ when
        the loss function has singularities. Gradient clipping
        (\py{torch.nn.utils.clip_grad_norm_}) establishes boundedness.
  \item \textbf{Learning rate}: a learning rate of zero causes no parameter
        updates (not numerically unsafe but semantically problematic);
        a negative learning rate inverts the gradient direction.
  \item \textbf{Batch normalization}: requires batch size $> 1$ to compute
        variance; division by batch variance requires the batch to not be
        identical (variance $> 0$).
  \item \textbf{Softmax temperature}: \py{softmax(x / T)} requires $T > 0$;
        as $T \to 0$, the softmax approaches argmax and may cause numeric
        instability.
\end{itemize}

\subsection{Interaction with the solver boundary}

Numeric safety predicates are typically in decidable fragments:
\begin{itemize}[nosep]
  \item Linear arithmetic over integers: overflow bounds, array lengths.
  \item Nonlinear real arithmetic: square root domains, logarithm domains.
  \item Floating-point theory: Z3's FPA theory can reason about specific
        floating-point behaviors, though it is expensive.
\end{itemize}

The system should use the cheapest sufficient theory: linear integer arithmetic
for most overflow and division-by-zero checks, nonlinear real arithmetic for
domain constraints, and floating-point theory only when precise bit-level
reasoning is required. When the numeric predicate exceeds the solver boundary,
it should be surfaced as a residual obligation with an explanation.

\section{Deep dive: comprehensions, iteration semantics, and string safety}

Python's comprehension syntax, iterator protocol, and string operations are
ubiquitous in real code. Each introduces semantic constraints that the
sheaf-descent framework should capture.

\subsection{List, dict, set, and generator comprehensions}

Comprehensions are syntactic sugar for iteration with optional filtering.
Each comprehension creates a family of sites:

\paragraph{Iterable sites.}
The source expression (\py{for x in source}) creates an iterable site. The
viability predicate requires \py{source} to implement the iterator protocol
(\py{__iter__} and \py{__next__}). The element type of the iteration variable
$x$ is determined by the source's yield type.

\paragraph{Filter sites.}
If the comprehension has an \py{if} clause (\py{[f(x) for x in xs if pred(x)]}),
the filter condition creates a guard site. Inside the body, $x$ is refined by
the predicate.

\paragraph{Output length reasoning.}
The output of \py{[f(x) for x in xs]} has the same length as \py{xs} when there
is no filter. With a filter, the output length is bounded above by the input
length: $0 \leq |\mathrm{output}| \leq |\mathrm{input}|$. The system should
track these bounds.

\paragraph{Nested comprehensions.}
\py{[f(x, y) for x in xs for y in ys]} iterates over the Cartesian product.
The output length is $|\mathrm{xs}| \times |\mathrm{ys}|$. The system should
compose length constraints multiplicatively.

\paragraph{Dict comprehensions.}
\py{\{k: v for k, v in items\}} creates a dict whose keys are determined by
the comprehension. If the key expression can produce duplicates, later values
overwrite earlier ones. The system should create a key-uniqueness site when
the output dict's key set matters.

\subsection{For-loop semantics and loop-carried invariants}

For-loops are the general iteration mechanism. The system should handle:

\paragraph{Loop bounds.}
\py{for i in range(n)} creates a site where $i$ ranges over $[0, n)$.
Accesses like \py{xs[i]} inside the loop are safe when $n \leq \mathrm{len}(xs)$.
The system should synthesize this relationship as a backward requirement.

\paragraph{Loop-carried accumulation.}
A common pattern is accumulating a result:
\begin{lstlisting}[style=deppy,language=Python]
total = 0
for x in xs:
    total += x
\end{lstlisting}
The loop body creates an accumulation site. The output refinement for
\py{total} after the loop should express that it equals the sum of elements
in \py{xs}. In general, loop-carried invariants are undecidable, but common
patterns (summation, product, min, max, counting) should have dedicated
invariant templates.

\paragraph{Loop-carried list construction.}
\begin{lstlisting}[style=deppy,language=Python]
results = []
for x in xs:
    results.append(f(x))
\end{lstlisting}
The system should recognize this as equivalent to \py{[f(x) for x in xs]}
and apply the same length reasoning.

\paragraph{While-loop invariants.}
While-loops require explicit or synthesized invariants. The system should
attempt to synthesize invariants from the loop condition, body mutations,
and initialization. When synthesis fails, the invariant becomes a residual
obligation. Common patterns:
\begin{itemize}[nosep]
  \item Countdown: \py{while n > 0: n -= 1}---the invariant is $n \geq 0$.
  \item Search: \py{while not found: ...}---the post-loop state has
        \py{found = True}.
  \item Convergence: \py{while error > tol: ...}---the post-loop state has
        \py{error <= tol}.
\end{itemize}

\subsection{Iterator protocol and lazy evaluation}

Python's iterator protocol (\py{__iter__}, \py{__next__}) introduces
state-dependent behavior:
\begin{itemize}[nosep]
  \item An iterator is consumed: after iterating once, calling \py{next()}
        again produces the next element or raises \py{StopIteration}.
  \item Iterators are single-pass unless they implement \py{__iter__} by
        returning a fresh iterator.
  \item \py{list(iterator)} materializes the iterator; calling \py{list()}
        on an already-consumed iterator produces an empty list.
\end{itemize}

The system should model iterator state as a resource: creating an iterator
produces a fresh resource, consuming it depletes the resource, and using a
depleted iterator produces unexpected empty results.

\subsection{String operations and encoding safety}

Strings in Python are sequences of Unicode codepoints. Key safety properties:

\paragraph{String indexing and slicing.}
\py{s[i]} requires $0 \leq i < \mathrm{len}(s)$ (or $-\mathrm{len}(s) \leq i < 0$
for negative indices). The same bounds reasoning as list indexing applies.

\paragraph{String formatting.}
\py{f"\{x:.2f\}"} requires \py{x} to be a number. \py{"\{0\} \{1\}".format(a, b)}
requires at least two positional arguments. \py{"\{name\}".format(**d)} requires
the key \py{"name"} to be present in \py{d}. The system should create error sites
for format-string mismatches.

\paragraph{Encoding and decoding.}
\py{s.encode("utf-8")} is always safe for valid Python strings, but
\py{b.decode("utf-8")} can raise \py{UnicodeDecodeError} when the bytes are
not valid UTF-8. The system should create an error site for decoding operations
and allow backward synthesis of encoding requirements.

\paragraph{String methods with semantic guarantees.}
\begin{itemize}[nosep]
  \item \py{s.split(sep)} returns a list with at least one element.
  \item \py{s.strip()} returns a string no longer than \py{s}.
  \item \py{s.startswith(prefix)} and \py{s.endswith(suffix)} are boolean
        guards that refine the string's content.
  \item \py{s.replace(old, new)} preserves length when
        $\mathrm{len(old)} = \mathrm{len(new)}$.
  \item \py{int(s)} and \py{float(s)} require \py{s} to represent a valid
        number; the system should create error sites for parsing operations.
\end{itemize}

\paragraph{Regular expressions.}
\py{re.match(pattern, s)} returns a match object or \py{None}. A successful
match provides group access, where \py{m.group(i)} requires $i$ to be a valid
group index. The theory should model match results as optional-with-groups and
create error sites for invalid group access.

\subsection{Zip, enumerate, and iteration combinators}

\paragraph{zip.}
\py{zip(xs, ys)} truncates to the shorter input. The output length is
$\min(|\mathrm{xs}|, |\mathrm{ys}|)$. If the user expects equal-length
inputs, they should use \py{zip(xs, ys, strict=True)} (Python 3.10+), which
raises \py{ValueError} on length mismatch. The system should create a
length-equality site when \py{strict=True} is used.

\paragraph{enumerate.}
\py{enumerate(xs, start=0)} yields pairs $(i, x_i)$ where $i$ starts at
\py{start}. The system should track that the index $i$ satisfies
$\py{start} \leq i < \py{start} + |\mathrm{xs}|$.

\paragraph{map and filter.}
\py{map(f, xs)} applies \py{f} to each element; the output has the same
length as the input. \py{filter(pred, xs)} selects elements; the output
length is bounded by the input length.

\section{Deep dive: module-level state, imports, closures, and the class/object theory}

Real Python programs are organized into modules, use imports extensively, define
classes with mutable state, and exploit closures and nested functions. The
sheaf-descent framework must handle these organizational structures to scale
beyond single-function analysis.

\subsection{Module-level state and import semantics}

Python modules are objects with attributes. Module-level statements execute
sequentially when the module is first imported. Key considerations:

\paragraph{Module-level configuration.}
Many libraries use module-level state for configuration:
\begin{lstlisting}[style=deppy,language=Python]
import torch
torch.set_default_dtype(torch.float64)
device = torch.device("cuda" if torch.cuda.is_available() else "cpu")
\end{lstlisting}
The system should model module-level assignments as establishing global state
that subsequent functions can depend on. The device and dtype settings should
flow into the sections of functions that use the default device or dtype.

\paragraph{Import graph and dependency ordering.}
Circular imports and conditional imports create challenges:
\begin{itemize}[nosep]
  \item Forward references (\py{from __future__ import annotations}) defer
        type evaluation.
  \item Conditional imports (\py{try: import numpy; HAS_NP = True except:
        HAS_NP = False}) create configuration-dependent code paths.
  \item Lazy imports (inside function bodies) defer module loading.
\end{itemize}
The system should treat the import graph as a dependency structure for
theory-pack activation: importing \py{torch} activates the PyTorch theory
pack; importing \py{pandas} activates the Pandas pack.

\paragraph{Module attributes and monkey-patching.}
Python allows setting arbitrary attributes on modules:
\py{module.x = 42}. The system should model module-level attributes as a
record type with the keys established by module-level assignments and
\py{__all__} declarations.

\subsection{Class definitions and instance semantics}

Classes are the primary organization unit for stateful Python code. The system
should handle several aspects:

\paragraph{Class body as a namespace.}
The class body defines attributes, methods, class variables, and nested
classes. The system should extract:
\begin{itemize}[nosep]
  \item Instance attributes assigned in \py{__init__}.
  \item Class variables assigned at class scope.
  \item Methods with their signatures.
  \item Properties (\py{@property}) with their getter/setter contracts.
  \item Special methods (\py{__len__}, \py{__getitem__}, \py{__iter__}, etc.)
        that participate in protocol dispatch.
\end{itemize}

\paragraph{Constructor analysis.}
\py{__init__} establishes the initial state of an instance. The system should
analyze \py{__init__} to determine:
\begin{itemize}[nosep]
  \item Which instance attributes are always set (and their types).
  \item Which attributes are conditionally set.
  \item Parameter requirements for construction.
\end{itemize}
This produces an initial-state section for the class that describes the
guaranteed state after construction.

\paragraph{Method semantics.}
Each method operates on \py{self}, which carries the instance's current state
section. The method's local analysis should use the instance state as context:
\begin{lstlisting}[style=deppy,language=Python]
class Buffer:
    def __init__(self, capacity: int):
        self.capacity = capacity
        self.items: list = []

    def add(self, item):
        if len(self.items) >= self.capacity:
            raise OverflowError("full")
        self.items.append(item)
\end{lstlisting}
The \py{add} method has a precondition ($\mathrm{len(self.items)} <
\mathrm{self.capacity}$), and a postcondition (the item is added and the
length increases by one). The system should synthesize both.

\paragraph{Inheritance.}
Subclass methods may override parent methods. The system should verify that
overrides satisfy the Liskov substitution principle: the subclass method's
precondition should be no stronger than the parent's, and its postcondition
should be no weaker.

\paragraph{Multiple inheritance and MRO.}
Python's method resolution order (MRO) determines which method is called in
diamond-inheritance situations. The system should follow the MRO to determine
which method definition is active at a given call site.

\subsection{Closures and nested functions}

Closures capture variables from enclosing scopes:
\begin{lstlisting}[style=deppy,language=Python]
def make_multiplier(factor):
    def multiply(x):
        return x * factor
    return multiply
\end{lstlisting}
The inner function \py{multiply} captures \py{factor} from the enclosing
scope. The system should:
\begin{itemize}[nosep]
  \item Track which variables are captured by closures.
  \item Include captured-variable refinements in the closure's section.
  \item Propagate the enclosing scope's state into the closure's analysis.
\end{itemize}

\paragraph{Mutable closures.}
When a closure captures a mutable variable (via \py{nonlocal}), the variable
may be modified by both the enclosing scope and the closure. The system should
create shared-state sites for such variables and track mutations through both
scopes.

\paragraph{Closures as callbacks.}
A common pattern is passing closures as callbacks to higher-order functions:
\begin{lstlisting}[style=deppy,language=Python]
def train(model, data, on_step=None):
    for batch in data:
        loss = model(batch)
        if on_step is not None:
            on_step(loss)
\end{lstlisting}
The system should propagate the type requirements on \py{on_step} (accepts a
single argument of the same type as \py{loss}) into the callback's section.

\subsection{Dataclasses and attrs}

Dataclasses (\py{@dataclass}) and attrs (\py{@attrs}) generate boilerplate
methods automatically. The system should:
\begin{itemize}[nosep]
  \item Recognize dataclass field declarations and extract the field schema.
  \item Model the auto-generated \py{__init__}, \py{__eq__}, \py{__repr__},
        and \py{__hash__} methods.
  \item Handle \py{field(default=...)} and \py{field(default_factory=...)}
        for optional fields.
  \item Handle frozen dataclasses (\py{frozen=True}) as immutable records.
  \item Handle validators in attrs (\py{@x.validator}) as preconditions on
        field values.
\end{itemize}

\subsection{Metaclasses and dynamic class creation}

Metaclasses (\py{type(...)} or custom metaclasses) create classes dynamically.
The system should handle common patterns:
\begin{itemize}[nosep]
  \item \py{type("Name", (Base,), {"method": lambda self: ...})} creates a
        class with known attributes.
  \item \py{__init_subclass__} hooks modify subclass behavior at definition
        time.
  \item Abstract base classes (\py{abc.ABC}) with \py{@abstractmethod}
        require subclasses to implement certain methods; instantiating an
        abstract class should produce an error site.
\end{itemize}

For highly dynamic metaclass usage, the system should generate residuals rather
than attempting full static analysis. The key principle is graceful degradation:
the system should analyze what it can and surface what it cannot.

\subsection{Global state and singleton patterns}

Global mutable state is common in Python (module-level dicts, registries,
caches). The system should:
\begin{itemize}[nosep]
  \item Track which functions read or write global state.
  \item Create shared-state sites for globals that are accessed from multiple
        functions.
  \item Recognize singleton patterns (\py{_instance = None} with lazy
        initialization) and model the initialization state transition.
  \item Warn when global mutation creates implicit dependencies between
        otherwise independent functions.
\end{itemize}

\section{Deep dive: union types, type narrowing, pattern matching, and the complete type algebra}

Python's type system (as codified by PEP 484 and successors) includes union
types, literal types, type guards, and increasingly powerful narrowing
mechanisms. The sheaf-descent framework should subsume and extend these
features.

\subsection{Union types as disjunctive sections}

A value of type \py{int | str} (or \py{Union[int, str]}) may be either an
integer or a string. In the sheaf-descent framework, a union type corresponds
to a disjunctive section: the value's stalk at a given site includes both
alternatives until a guard narrows it.

\paragraph{Narrowing via isinstance.}
\begin{lstlisting}[style=deppy,language=Python]
def process(x: int | str) -> str:
    if isinstance(x, int):
        return str(x * 2)    # x: int in this branch
    else:
        return x.upper()     # x: str in this branch
\end{lstlisting}
The \py{isinstance} check creates a guard site that splits the section: the
true branch restricts to \py{int}, the false branch to \py{str}. The system
should automatically perform this narrowing for all union-typed variables.

\paragraph{Narrowing via type guards (PEP 647).}
\begin{lstlisting}[style=deppy,language=Python]
from typing import TypeGuard

def is_str_list(lst: list) -> TypeGuard[list[str]]:
    return all(isinstance(x, str) for x in lst)
\end{lstlisting}
When \py{is_str_list(xs)} returns \py{True}, the variable \py{xs} should be
narrowed to \py{list[str]} in the true branch. The system should recognize
\py{TypeGuard} annotations and use them as guard predicates.

\paragraph{Discriminated unions.}
A common pattern is using a discriminant field to select among alternatives:
\begin{lstlisting}[style=deppy,language=Python]
@dataclass
class Circle:
    kind: Literal["circle"] = "circle"
    radius: float = 0.0

@dataclass
class Rectangle:
    kind: Literal["rect"] = "rect"
    width: float = 0.0
    height: float = 0.0

Shape = Circle | Rectangle
\end{lstlisting}
Checking \py{shape.kind == "circle"} should narrow \py{shape} to \py{Circle}.
The system should recognize literal-typed discriminant fields and use them
for narrowing.

\subsection{Literal types and constant propagation}

\py{Literal["read", "write"]} restricts a value to specific constants. The
system should track literal types through assignments and function calls:
\begin{itemize}[nosep]
  \item \py{mode: Literal["r", "w"]} restricts the parameter to two values.
  \item Passing a literal string \py{"r"} satisfies the constraint.
  \item Conditional checks on literal-typed values narrow the section.
  \item \py{Literal[True]} and \py{Literal[False]} enable boolean narrowing.
\end{itemize}

\subsection{TypeVar, Generic, and parametric polymorphism}

Generic types use \py{TypeVar} to express parametric relationships:
\begin{lstlisting}[style=deppy,language=Python]
T = TypeVar("T")

def first(xs: list[T]) -> T:
    return xs[0]
\end{lstlisting}
The system should instantiate \py{T} at each call site based on the actual
argument type and propagate the instantiated return type. Bounded type
variables (\py{TypeVar("T", bound=Comparable)}) should additionally check
that the bound is satisfied.

\paragraph{ParamSpec and Callable types.}
\py{ParamSpec} (PEP 612) captures the parameter signature of a callable:
\begin{lstlisting}[style=deppy,language=Python]
P = ParamSpec("P")
R = TypeVar("R")

def decorator(f: Callable[P, R]) -> Callable[P, R]:
    def wrapper(*args: P.args, **kwargs: P.kwargs) -> R:
        return f(*args, **kwargs)
    return wrapper
\end{lstlisting}
The system should propagate \py{P} through decorator application so that the
wrapped function retains its original parameter types.

\subsection{Overloaded functions}

\py{@overload} (PEP 484) specifies different return types for different
argument patterns:
\begin{lstlisting}[style=deppy,language=Python]
@overload
def get(key: str, default: None = None) -> int | None: ...
@overload
def get(key: str, default: int) -> int: ...
\end{lstlisting}
The system should select the matching overload at each call site based on the
argument types and use the corresponding return type.

\subsection{NewType, TypeAlias, and nominal distinctions}

\py{NewType} creates a distinct type at the type-checking level:
\begin{lstlisting}[style=deppy,language=Python]
UserId = NewType("UserId", int)
\end{lstlisting}
A \py{UserId} is an \py{int} at runtime but should be treated as a distinct
type by the analysis. This prevents accidentally using a raw integer where a
\py{UserId} is expected. The system should respect these nominal distinctions.

\subsection{Final, ClassVar, and read-only refinements}

\begin{itemize}[nosep]
  \item \py{Final[int]} marks a variable as immutable after assignment. The
        system should create an error site for any subsequent assignment.
  \item \py{ClassVar[int]} marks an attribute as belonging to the class rather
        than instances. The system should track class-level vs. instance-level
        state separately.
  \item \py{ReadOnly} (PEP 705 for TypedDict) marks specific keys as read-only.
\end{itemize}

\subsection{TypedDict and structural record types}

\py{TypedDict} defines dicts with known keys and value types:
\begin{lstlisting}[style=deppy,language=Python]
class Point(TypedDict):
    x: float
    y: float
\end{lstlisting}
The system should model TypedDicts as structural record types: accessing a
declared key has the declared type; accessing an undeclared key produces an
error site. \py{total=False} makes all keys optional; \py{Required} and
\py{NotRequired} (PEP 655) fine-tune which keys are required.

\subsection{Self type and recursive types}

\py{Self} (PEP 673) refers to the current class in method return types:
\begin{lstlisting}[style=deppy,language=Python]
class Builder:
    def set_name(self, name: str) -> Self:
        self.name = name
        return self
\end{lstlisting}
The system should resolve \py{Self} to the actual receiver type at each call
site, enabling fluent-API type tracking through method chains.

Recursive types (e.g., \py{Tree = int | list[Tree]}) require the system to
handle fixed-point type definitions. The sheaf-descent framework handles this
naturally: the recursion creates a recursive site structure, and sections at
each depth are consistent with the unfolding.

\subsection{Intersection types and protocol conjunction}

While Python's type system does not have explicit intersection types, the
system should internally represent ``implements both Protocol A and Protocol
B'' as a conjunction. This arises when a variable is narrowed by multiple
\py{isinstance} checks or when a function parameter requires multiple
protocol implementations.

\subsection{Type algebra summary}

The complete type algebra supported by the system should include:
\begin{itemize}[nosep]
  \item Base types: \py{int}, \py{float}, \py{str}, \py{bool}, \py{bytes},
        \py{None}.
  \item Container types: \py{list[T]}, \py{dict[K, V]}, \py{set[T]},
        \py{tuple[T, ...]}, \py{tuple[T1, T2]}.
  \item Union types: \py{T1 | T2}.
  \item Optional types: \py{T | None}.
  \item Literal types: \py{Literal["a", "b"]}, \py{Literal[1, 2]}.
  \item Generic types: \py{Generic[T]}, bounded type variables.
  \item Callable types: \py{Callable[[T1, T2], R]}, \py{ParamSpec}.
  \item Protocol types: structural subtyping.
  \item TypedDict types: structural record types.
  \item NewType: nominal distinctions.
  \item Refinement types: $\{x : T \mid \phi(x)\}$ where $\phi$ is in a
        decidable fragment.
  \item Dependent types: $\Pi(x : A). B(x)$ where $B$ may depend on the
        value of $x$.
  \item Section types: the full sheaf-descent section structure.
\end{itemize}

This hierarchy subsumes Python's standard type annotations while extending them
with refinement and dependent types that capture the richer semantic information
the system infers.

\section{Deep dive: end-to-end worked example of sheaf-descent analysis}

This final deep-dive section walks through a complete sheaf-descent analysis of
a realistic Python program that uses multiple libraries, control flow, error
handling, and dependent-type-level reasoning. The goal is to show how all the
pieces fit together.

\subsection{The target program}

Consider a function that loads data, trains a model, and returns metrics:
\begin{lstlisting}[style=deppy,language=Python]
import torch
import torch.nn as nn
import numpy as np
from typing import Optional

def train_classifier(
    features: torch.Tensor,
    labels: torch.Tensor,
    hidden_dim: int = 128,
    epochs: int = 10,
    lr: float = 0.01,
    device: Optional[str] = None,
) -> dict:
    if device is None:
        device = "cuda" if torch.cuda.is_available() else "cpu"

    assert features.dim() == 2, "features must be 2D"
    assert labels.dim() == 1, "labels must be 1D"
    assert features.shape[0] == labels.shape[0], "batch mismatch"
    assert hidden_dim > 0, "hidden_dim must be positive"
    assert epochs > 0, "epochs must be positive"
    assert lr > 0, "lr must be positive"

    n_samples, input_dim = features.shape
    n_classes = int(labels.max().item()) + 1

    model = nn.Sequential(
        nn.Linear(input_dim, hidden_dim),
        nn.ReLU(),
        nn.Linear(hidden_dim, n_classes),
    )
    model = model.to(device)
    features = features.to(device)
    labels = labels.to(device)

    optimizer = torch.optim.SGD(model.parameters(), lr=lr)
    criterion = nn.CrossEntropyLoss()

    losses = []
    for epoch in range(epochs):
        logits = model(features)
        loss = criterion(logits, labels)
        optimizer.zero_grad()
        loss.backward()
        optimizer.step()
        losses.append(loss.item())

    with torch.no_grad():
        final_logits = model(features)
        preds = final_logits.argmax(dim=1)
        accuracy = (preds == labels).float().mean().item()

    return {
        "accuracy": accuracy,
        "final_loss": losses[-1],
        "losses": losses,
        "n_samples": n_samples,
        "n_classes": n_classes,
    }
\end{lstlisting}

\subsection{Site cover construction}

The system constructs the following site families:

\paragraph{Guard sites from assertions.}
Each \py{assert} statement creates a guard site. After the assertion block, the
section carries:
\begin{itemize}[nosep]
  \item $\mathrm{features.dim()} = 2$
  \item $\mathrm{labels.dim()} = 1$
  \item $\mathrm{features.shape}[0] = \mathrm{labels.shape}[0]$
  \item $\mathrm{hidden\_dim} > 0$
  \item $\mathrm{epochs} > 0$
  \item $\mathrm{lr} > 0$
\end{itemize}

\paragraph{None-narrowing site.}
The \py{if device is None} check creates a narrowing site. In the true branch,
\py{device} is assigned a string value; in the false branch, the input
\py{device} is used. After the merge, \py{device} is known to be a non-None
string.

\paragraph{Shape sites.}
\py{features.shape} creates a shape site with value
$(\mathrm{n\_samples}, \mathrm{input\_dim})$.
\py{labels.shape} creates a shape site with value $(\mathrm{labels.shape}[0],)$.
The destructuring \py{n_samples, input_dim = features.shape} binds symbolic
names to shape components.

\paragraph{PyTorch theory sites.}
\begin{itemize}[nosep]
  \item \py{nn.Linear(input_dim, hidden_dim)}: requires $\mathrm{input\_dim} > 0$
        and $\mathrm{hidden\_dim} > 0$; output shape is
        $(\mathrm{batch}, \mathrm{hidden\_dim})$.
  \item \py{nn.Linear(hidden_dim, n_classes)}: requires $\mathrm{hidden\_dim} > 0$
        and $\mathrm{n\_classes} > 0$; output shape is
        $(\mathrm{batch}, \mathrm{n\_classes})$.
  \item \py{model(features)}: the Sequential applies layers in order; the
        system composes shape transformations to verify shape compatibility.
  \item \py{criterion(logits, labels)}: CrossEntropyLoss requires
        $\mathrm{logits.shape}[0] = \mathrm{labels.shape}[0]$ and
        $\mathrm{logits.shape}[1] = \mathrm{n\_classes}$.
  \item \py{final_logits.argmax(dim=1)}: output shape is
        $(\mathrm{n\_samples},)$; values in $[0, \mathrm{n\_classes})$.
\end{itemize}

\paragraph{Numeric safety sites.}
\begin{itemize}[nosep]
  \item \py{labels.max().item() + 1}: requires the labels tensor to be
        nonempty (otherwise \py{max()} raises \py{RuntimeError}).
  \item \py{losses[-1]}: requires the losses list to be nonempty, which is
        guaranteed when $\mathrm{epochs} > 0$.
  \item \py{lr > 0}: prevents degenerate optimization.
\end{itemize}

\subsection{Backward synthesis: input requirements}

The backward pass pulls back all viability predicates to the function boundary,
synthesizing the weakest precondition. The result:
\begin{lstlisting}[style=deppy,language=Python]
@deppy.requires(
    features.dim() == 2,
    labels.dim() == 1,
    features.shape[0] == labels.shape[0],
    features.shape[0] > 0,          # nonempty for labels.max()
    labels.max() >= 0,              # n_classes >= 1
    hidden_dim > 0,
    epochs > 0,
    lr > 0,
    device is None or isinstance(device, str),
)
\end{lstlisting}

\subsection{Forward synthesis: output guarantees}

The forward pass computes the strongest postcondition. The synthesized output
section:
\begin{lstlisting}[style=deppy,language=Python]
@deppy.ensures(
    lambda result: isinstance(result, dict),
    lambda result: set(result.keys()) == {
        "accuracy", "final_loss", "losses", "n_samples", "n_classes"
    },
    lambda result: 0.0 <= result["accuracy"] <= 1.0,
    lambda result: isinstance(result["final_loss"], float),
    lambda result: len(result["losses"]) == epochs,
    lambda result: result["n_samples"] == features.shape[0],
    lambda result: result["n_classes"] == int(labels.max().item()) + 1,
)
\end{lstlisting}

\subsection{Runtime-error freedom certificate}

Given the synthesized preconditions, the system can certify that the function
is free of the following runtime errors:
\begin{itemize}[nosep]
  \item No \py{AssertionError}: all assertions are implied by the preconditions.
  \item No \py{RuntimeError} from empty tensor reduction: $\mathrm{features.shape}[0] > 0$
        guarantees nonempty tensors.
  \item No \py{IndexError} from \py{losses[-1]}: $\mathrm{epochs} > 0$
        guarantees the list is nonempty after the loop.
  \item No shape mismatch in \py{nn.Linear} or \py{CrossEntropyLoss}: the
        shape sites are consistent with the assertions.
\end{itemize}

\subsection{Residual obligations}

Some properties cannot be verified statically:
\begin{itemize}[nosep]
  \item Whether the device string is valid (e.g., \py{"cuda"} when CUDA is
        not available). This depends on runtime hardware.
  \item Whether the model converges (the loss decreases over epochs). This
        is a dynamic property that depends on data and hyperparameters.
  \item Whether the numeric values remain finite throughout training. The
        system can verify that no obvious division-by-zero or log-of-zero
        occurs, but cannot rule out all numeric instability.
\end{itemize}

These are surfaced as explicit residuals with explanations, allowing the user
to decide whether to add runtime checks, provide proofs, or accept the risk.

\subsection{Interprocedural context}

If \py{train_classifier} is called from another function, the caller inherits
the preconditions as obligations:
\begin{lstlisting}[style=deppy,language=Python]
def run_experiment(data_path: str) -> dict:
    data = np.load(data_path)
    features = torch.from_numpy(data["X"])
    labels = torch.from_numpy(data["y"])
    return train_classifier(features, labels, epochs=20)
\end{lstlisting}

The system pulls back the preconditions to \py{run_experiment}'s boundary:
the loaded data must have keys \py{"X"} and \py{"y"}; \py{data["X"]} must be
2D; \py{data["y"]} must be 1D; their first dimensions must match; the data
must be nonempty. These become requirements on the \py{.npz} file, surfaced
as either annotations or diagnostics.

This end-to-end example demonstrates the full sheaf-descent pipeline: from
source code through site construction, bidirectional synthesis, error-freedom
certification, residual surfacing, and interprocedural composition.

%% ======================================================================
%% PART IV: FULL DEPENDENT TYPING AND PROOF MODE
%% ======================================================================

\part{Full Dependent Typing and Proof Mode}

\chapter{The Proof-Mode Philosophy: From Inference to Verification}

The preceding parts of this monograph describe a system that \emph{infers}
rich dependent types with minimal annotation. The system harvests guards,
propagates sections, synthesizes preconditions and postconditions, and
certifies runtime-error freedom---all without requiring the programmer to
write a single type annotation beyond what Python already supports.

This part describes the \emph{proof mode}: a direct, seamless extension of
the inference engine that allows programmers to \emph{selectively} escalate
from inference to verification. The key design principle is:

\begin{quote}
\textbf{The proof mode is not a separate system.} It is the same
sheaf-descent engine with the same site covers, the same sections, and the
same solver backend. The only difference is that the programmer provides
additional annotations---preconditions, postconditions, invariants, lemmas,
or full proofs---at chosen sites, and the engine uses these annotations as
\emph{axioms} that enrich the local sections. Everything else continues to
be inferred.
\end{quote}

This design produces a unique value proposition: the programmer gets useful
static checking for \emph{all} of their code (via inference), and they can
invest annotation effort at any granularity, from a single assertion to a
full Agda-style proof, and the investment pays off immediately by enriching
the inferred sections of surrounding code.

\section{The annotation economy: cherry-picking depth}

In traditional proof assistants (Agda, Coq, Lean, F*), every function must
be fully annotated and every proof must be complete. This creates an
all-or-nothing barrier: either you commit to proving everything, or you
prove nothing. The result is that proof assistants are used only for
specialized, high-assurance code.

Our system inverts this economics. The programmer can operate at any point on
a continuous spectrum:

\paragraph{Level 0: Zero annotation.}
The programmer writes plain Python. The system infers whatever it can:
shapes, bounds, None-safety, basic type properties. This is the default mode
described in Parts I--III.

\paragraph{Level 1: Lightweight contracts.}
The programmer adds \py{@requires} and \py{@ensures} decorators to selected
functions. These are checked by the solver and used to enrich the sections of
callers and callees. No proofs are needed; the solver discharges obligations
automatically when they fall within decidable fragments.

\paragraph{Level 2: Invariant annotations.}
The programmer annotates loop invariants, class invariants, or module-level
invariants. These are checked inductively (base case + preservation) and used
to strengthen the sections of code that depends on the invariant.

\paragraph{Level 3: Lemma statements.}
The programmer states lemmas---properties they believe to be true---and the
system attempts to prove them automatically. If the solver succeeds, the
lemma becomes an axiom available to surrounding code. If the solver fails,
the lemma becomes a proof obligation.

\paragraph{Level 4: Proof sketches.}
The programmer provides a proof sketch: a sequence of intermediate assertions
that guide the solver. The system fills in the gaps between assertions
automatically.

\paragraph{Level 5: Full proofs.}
The programmer writes complete, step-by-step proofs in an Agda-like syntax
embedded in Python decorators or comments. Every step is verified.

The crucial property is that all five levels coexist in the same codebase.
A function at Level 0 can call a function at Level 5, and vice versa. The
system mediates between levels through the section mechanism: a fully-proved
function exports a maximally-rich section, which enriches the inferred
sections of its Level-0 callers.

\section{Why this is not a separate system}

Consider the alternative: building a separate proof language (like F* or
Dafny) and then connecting it to Python via an FFI or a translation layer.
This approach has fundamental problems:

\begin{enumerate}[nosep]
  \item \textbf{Semantic gap.} The proof language has its own semantics,
        which must be related to Python's semantics by a trusted translation.
        Any discrepancy between the two is a soundness hole.
  \item \textbf{Duplication.} The programmer must maintain two representations
        of the same code: the Python implementation and the proof-language
        specification. Changes to one must be reflected in the other.
  \item \textbf{Cliff.} There is no smooth transition from ``Python with
        types'' to ``proven Python.'' The programmer must cross a cliff from
        one tool to another.
  \item \textbf{Inference loss.} The proof language does not benefit from the
        inference engine's automatic section synthesis. The programmer must
        re-specify everything the engine already knows.
\end{enumerate}

Our approach avoids all four problems. The proof mode operates on the same
AST, the same SSA representation, the same site covers, and the same solver
backend as the inference mode. Annotations are additional constraints in the
same constraint language. Proofs are additional evidence in the same section
structure. There is no translation, no duplication, and no cliff.

\section{The annotation surface: Python-native proof syntax}

Proofs and specifications are written in Python syntax using decorators,
type annotations, and a small library of proof combinators. The design
principle is that every proof construct should be valid Python that can
be executed (for testing) or ignored (for production).

\subsection{Preconditions and postconditions}

\begin{lstlisting}[style=deppy,language=Python]
from deppy import requires, ensures, returns

@requires(lambda n: n >= 0)
@ensures(lambda n, result: result >= 1)
@ensures(lambda n, result: n <= 20 or returns.unchecked())
def factorial(n: int) -> int:
    if n == 0:
        return 1
    return n * factorial(n - 1)
\end{lstlisting}

The \py{@requires} decorator specifies a precondition; \py{@ensures}
specifies a postcondition. These are lambda expressions over the function's
parameters and (for \py{@ensures}) the return value. The system checks:
\begin{itemize}[nosep]
  \item The precondition is satisfiable (not vacuously false).
  \item The postcondition holds on all execution paths that satisfy the
        precondition.
  \item Callers satisfy the precondition at every call site.
\end{itemize}

The \py{returns.unchecked()} escape hatch allows the programmer to leave
some postcondition aspects unverified (here, the overflow behavior for
$n > 20$).

\subsection{Dependent function types}

Python's type annotations can express dependent types when combined with
\py{deppy.Refine}:
\begin{lstlisting}[style=deppy,language=Python]
from deppy import Refine

def safe_divide(
    a: float,
    b: Refine[float, lambda x: x != 0],
) -> float:
    return a / b
\end{lstlisting}

\py{Refine[T, pred]} is a refinement type: the type $\{x : T \mid
\mathrm{pred}(x)\}$. At call sites, the system checks that the actual
argument satisfies the refinement predicate. If the predicate cannot be
verified statically, the system can optionally insert a runtime check.

\subsection{Quantified properties}

For properties that quantify over collections, the system provides
\py{forall} and \py{exists}:
\begin{lstlisting}[style=deppy,language=Python]
from deppy import ensures, forall

@ensures(lambda xs, result:
    forall(range(len(result) - 1),
           lambda i: result[i] <= result[i + 1]))
def sort(xs: list[int]) -> list[int]:
    return sorted(xs)
\end{lstlisting}

The \py{forall} combinator expresses a universally-quantified property. The
solver can often discharge such properties for built-in operations (like
\py{sorted}) using library-specific axioms. For user-defined loops, the
programmer may need to provide an inductive argument.

\subsection{Proof blocks}

For complex properties that require multi-step reasoning, the system provides
\py{proof} blocks:
\begin{lstlisting}[style=deppy,language=Python]
from deppy import proof, have, hence, qed

@ensures(lambda xs, result: len(result) == len(xs))
@ensures(lambda xs, result:
    forall(range(len(xs)), lambda i: result[i] == xs[i] ** 2))
def square_all(xs: list[int]) -> list[int]:
    result = []
    for i, x in enumerate(xs):
        result.append(x ** 2)
        with proof("loop_invariant"):
            have(len(result) == i + 1)
            have(forall(range(i + 1),
                        lambda j: result[j] == xs[j] ** 2))
    with proof("postcondition"):
        have(len(result) == len(xs))
        hence(forall(range(len(xs)),
                     lambda j: result[j] == xs[j] ** 2))
        qed()
    return result
\end{lstlisting}

The \py{with proof(...)} block is a context manager that introduces a proof
scope. Inside the scope, \py{have(...)} asserts an intermediate fact,
\py{hence(...)} derives a consequence, and \py{qed()} closes the proof. The
system checks that each \py{have} follows from previous facts and the program
state, and that the final \py{hence} establishes the desired property.

At runtime, proof blocks are no-ops (they execute but do nothing). At
analysis time, they provide the solver with a proof skeleton that guides
verification.

\section{The trust gradient: how annotations flow through inference}

The most novel aspect of proof mode is how annotations interact with
inference. When a programmer annotates one function with a rich contract,
the benefit propagates to all code that interacts with that function:

\paragraph{Forward propagation.}
A function with a strong postcondition exports a rich section. Any caller
that receives the function's result inherits the postcondition as a known
fact, which may allow the caller's own postcondition to be stronger---even
if the caller has no annotations at all.

\paragraph{Backward propagation.}
A function with a strong precondition creates obligations at call sites. If
a caller cannot satisfy the precondition from its own inferred section, the
system propagates the requirement further back, potentially all the way to
the module boundary, where it becomes an input requirement.

\paragraph{Lateral propagation.}
A lemma proved about a data structure enriches the sections of all code that
uses that data structure. For example, proving that a custom list maintains
a sortedness invariant allows binary-search code elsewhere in the module to
be verified without additional annotation.

The result is that a small amount of annotation effort can have an outsized
impact on the quality of inference across the entire codebase. This is the
\emph{annotation leverage} property.


\chapter{Dependent Types in Python: The Full Theory}

This chapter develops the full dependent type theory that underpins proof
mode. The theory is designed to be (a) expressive enough to state and verify
real correctness properties, (b) compatible with Python's dynamic semantics,
and (c) decidable in common cases while allowing controlled escape to
interactive proof for hard cases.

\section{The type universe}

The type universe $\mathcal{U}$ is stratified into levels:

\begin{align*}
  \mathcal{U}_0 &= \{\py{int}, \py{float}, \py{str}, \py{bool},
    \py{bytes}, \py{None}\} \\
  \mathcal{U}_1 &= \mathcal{U}_0 \cup \{\py{list}[A], \py{dict}[A, B],
    \py{set}[A], \py{tuple}[A_1, \ldots, A_n] \mid A, B, A_i \in
    \mathcal{U}_0\} \\
  \mathcal{U}_2 &= \mathcal{U}_1 \cup \{A \to B, \Pi(x:A).B(x),
    \Sigma(x:A).B(x) \mid A \in \mathcal{U}_1, B \in \mathcal{U}_1\} \\
  \mathcal{U}_3 &= \mathcal{U}_2 \cup \{\{x : A \mid \phi(x)\} \mid
    A \in \mathcal{U}_2, \phi \in \mathcal{L}_{\mathrm{dec}}\} \\
  \mathcal{U}_\omega &= \bigcup_{n} \mathcal{U}_n
\end{align*}

where $\mathcal{L}_{\mathrm{dec}}$ is the decidable predicate language
(quantifier-free linear arithmetic, array theory, and uninterpreted
functions---the QF\_LIRA + arrays fragment of SMT-LIB).

\paragraph{Refinement types.}
$\{x : A \mid \phi(x)\}$ is the type of values of type $A$ that satisfy
predicate $\phi$. This is the workhorse of the system: most dependent types
are expressed as refinements.

\paragraph{Dependent function types.}
$\Pi(x:A).B(x)$ is the type of functions that take an argument $x$ of type
$A$ and return a value of type $B(x)$, where the return type may depend on
the argument value. In Python:
\begin{lstlisting}[style=deppy,language=Python]
def lookup(xs: list[int], i: Refine[int, lambda i: 0 <= i < len(xs)]
           ) -> int:
    return xs[i]
\end{lstlisting}
The return type is simply \py{int}, but the system infers the stronger
section: the result equals \py{xs[i]}, which depends on \py{i}.

\paragraph{Dependent pair types.}
$\Sigma(x:A).B(x)$ is the type of pairs $(a, b)$ where $a : A$ and
$b : B(a)$. This arises naturally when a function returns a value together
with a proof of some property:
\begin{lstlisting}[style=deppy,language=Python]
def find_min(xs: Refine[list[int], lambda l: len(l) > 0]
             ) -> tuple[int, int]:
    idx = 0
    for i in range(1, len(xs)):
        if xs[i] < xs[idx]:
            idx = i
    return idx, xs[idx]
    # inferred section: result[0] in range(len(xs))
    #                   result[1] == xs[result[0]]
    #                   forall j in range(len(xs)): result[1] <= xs[j]
\end{lstlisting}

\section{Propositions as types, proofs as programs}

Following the Curry-Howard correspondence, the system treats propositions as
types and proofs as programs. This means that proving a property is the same
as constructing a value of the corresponding type.

\paragraph{The proposition type.}
A proposition $P$ is a type in $\mathcal{U}_\omega$ that is either inhabited
(the proposition is true) or empty (the proposition is false). The system
uses the notation \py{Prop[P]} for proposition types.

\paragraph{Proof terms.}
A proof of $P$ is a value of type \py{Prop[P]}. The system provides several
ways to construct proof terms:
\begin{itemize}[nosep]
  \item \textbf{Solver discharge.} If $P$ is in the decidable fragment, the
        solver can construct the proof automatically.
  \item \textbf{Refl.} If $P$ is an equality $a = b$ and $a$ and $b$ are
        definitionally equal, the proof is \py{refl}.
  \item \textbf{Assumption.} If $P$ is in the current context (a precondition,
        a guard, or a previously-proved fact), the proof is the assumption
        itself.
  \item \textbf{Application.} If $f : P \to Q$ and $p : P$, then $f(p) : Q$.
  \item \textbf{Induction.} For properties of natural numbers or lists, the
        system supports inductive proofs.
\end{itemize}

\paragraph{The proof is erased at runtime.}
Proof terms exist only at analysis time. At runtime, they are erased to
\py{None}. This means that proofs have zero runtime cost.

\section{The decidable core and the proof frontier}

Not every proposition can be decided automatically. The system maintains a
clear boundary between what the solver can handle and what requires human
guidance:

\paragraph{The decidable core.}
The solver can automatically discharge propositions in the following
fragments:
\begin{itemize}[nosep]
  \item Quantifier-free linear integer arithmetic (QF\_LIA).
  \item Quantifier-free linear real arithmetic (QF\_LRA).
  \item Quantifier-free arrays with extensionality (QF\_AX).
  \item Quantifier-free bitvectors (QF\_BV).
  \item Quantifier-free uninterpreted functions (QF\_UF).
  \item Boolean combinations of the above.
  \item Bounded quantification: $\forall i \in [0, n).\ P(i)$ where $n$ is
        a known constant or a symbolically-bounded expression.
\end{itemize}

\paragraph{The proof frontier.}
When a proposition falls outside the decidable core, the system cannot
discharge it automatically. Instead, it reports the proposition as a
\emph{proof obligation}---a residual that the programmer must address. The
programmer can:
\begin{enumerate}[nosep]
  \item Provide a proof (Levels 4--5).
  \item Provide intermediate lemmas that reduce the obligation to decidable
        sub-goals (Level 3).
  \item Accept the obligation as a trusted axiom (\py{@assume}).
  \item Downgrade to a runtime check (\py{@check}).
\end{enumerate}

This design ensures that the system never silently gives up. Every
unverified property is either proved, assumed, or checked at runtime.

\section{Inductive types and recursive proofs}

Many interesting properties require induction. The system supports inductive
reasoning for:

\paragraph{Natural number induction.}
For a property $P(n)$ where $n : \py{int}$ and $n \geq 0$:
\begin{lstlisting}[style=deppy,language=Python]
from deppy import induction, base_case, inductive_step

@ensures(lambda n, result: result == n * (n + 1) // 2)
def triangle(n: Refine[int, lambda n: n >= 0]) -> int:
    total = 0
    for i in range(n + 1):
        total += i
    with induction(n):
        with base_case(n == 0):
            have(total == 0 == 0 * 1 // 2)
        with inductive_step(n > 0):
            assume(triangle(n - 1) == (n - 1) * n // 2)
            have(total == triangle(n - 1) + n)
            hence(total == n * (n + 1) // 2)
    return total
\end{lstlisting}

\paragraph{Structural induction on lists.}
For a property $P(\mathrm{xs})$ where $\mathrm{xs} : \py{list}[A]$:
\begin{itemize}[nosep]
  \item Base case: $P([])$.
  \item Inductive step: for all $x : A$ and $\mathrm{xs'} : \py{list}[A]$,
        if $P(\mathrm{xs'})$ then $P([x] + \mathrm{xs'})$.
\end{itemize}

\paragraph{Well-founded induction.}
For recursive functions that do not follow simple structural patterns, the
programmer can specify a well-founded measure:
\begin{lstlisting}[style=deppy,language=Python]
from deppy import decreases

@decreases(lambda n: n)  # measure function
@ensures(lambda n, result: result >= 1)
def factorial(n: Refine[int, lambda n: n >= 0]) -> int:
    if n == 0:
        return 1
    return n * factorial(n - 1)
\end{lstlisting}

The \py{@decreases} decorator specifies a measure that strictly decreases at
each recursive call. The system verifies that the measure is non-negative and
decreasing, which ensures termination.


\chapter{Cherry-Picked Annotation: The Selective Verification Workflow}

The power of proof mode lies not in forcing the programmer to prove
everything, but in allowing them to prove \emph{exactly what matters} while
the inference engine handles the rest. This chapter describes the practical
workflow for selective verification.

\section{Identifying verification targets}

Not all code deserves the same level of assurance. The programmer should
invest proof effort where:
\begin{itemize}[nosep]
  \item \textbf{Correctness is critical}: financial calculations, security
        checks, safety-critical control.
  \item \textbf{Bugs are expensive}: core algorithms, data structure
        invariants, protocol implementations.
  \item \textbf{Testing is insufficient}: corner cases, concurrency, numeric
        precision.
  \item \textbf{The property is non-obvious}: algorithmic invariants, complex
        shape transformations, protocol compliance.
\end{itemize}

The system helps identify targets by reporting the \emph{inferred section
quality} at each site. Sites where the inferred section is weak (few known
properties) are candidates for annotation investment.

\section{The annotation workflow}

\subsection{Step 1: Run inference on unannotated code}

The programmer runs the system on their existing Python code with zero
annotations. The system produces:
\begin{itemize}[nosep]
  \item Inferred preconditions and postconditions for each function.
  \item Detected potential runtime errors with their viability predicates.
  \item A list of residual obligations that could not be discharged.
  \item A section-quality map showing where inference is strong and where it
        is weak.
\end{itemize}

\subsection{Step 2: Triage residuals}

The programmer reviews the residual obligations and decides for each one:
\begin{itemize}[nosep]
  \item \textbf{Accept.} The residual is not important; mark it as accepted.
  \item \textbf{Check.} The residual should be checked at runtime; the system
        inserts a runtime assertion.
  \item \textbf{Prove.} The residual is important and should be proved; the
        programmer adds annotations.
  \item \textbf{Strengthen.} The inferred section is too weak; the programmer
        adds a contract to strengthen it.
\end{itemize}

\subsection{Step 3: Add targeted annotations}

The programmer adds annotations \emph{only} at the sites identified in
Step 2. The system re-runs inference with the new annotations, which may
discharge some residuals and strengthen some sections.

\subsection{Step 4: Iterate}

The programmer repeats Steps 2--3 until the remaining residuals are all
accepted or checked. There is no requirement to reach zero residuals; the
programmer chooses their own stopping point.

\section{Example: selective verification of a binary search}

Consider a binary search implementation. At Level 0 (no annotations), the
system infers basic properties: the function takes a sorted list and a
target, accesses list elements within bounds, and returns an integer. But it
cannot verify that the returned index is correct.

\begin{lstlisting}[style=deppy,language=Python]
def binary_search(xs: list[int], target: int) -> int:
    lo, hi = 0, len(xs)
    while lo < hi:
        mid = (lo + hi) // 2
        if xs[mid] < target:
            lo = mid + 1
        else:
            hi = mid
    return lo
\end{lstlisting}

\paragraph{Level 0 inference.}
The system infers:
\begin{itemize}[nosep]
  \item Input: \py{xs} is a list of int; \py{target} is int.
  \item The loop accesses \py{xs[mid]} where
        $0 \leq \py{mid} < \py{len(xs)}$---inferred safe when
        $\py{len(xs)} > 0$ and $0 \leq \py{lo} \leq \py{hi} \leq
        \py{len(xs)}$.
  \item Output: \py{lo} is an int in $[0, \py{len(xs)}]$.
  \item Residual: cannot verify that the returned index is the correct
        insertion point.
\end{itemize}

\paragraph{Level 1: add contracts.}
The programmer adds preconditions and postconditions:
\begin{lstlisting}[style=deppy,language=Python]
from deppy import requires, ensures, forall, Refine

@requires(lambda xs: forall(range(len(xs) - 1),
                            lambda i: xs[i] <= xs[i + 1]))
@ensures(lambda xs, target, result:
    0 <= result <= len(xs))
@ensures(lambda xs, target, result:
    forall(range(result), lambda i: xs[i] < target))
@ensures(lambda xs, target, result:
    forall(range(result, len(xs)), lambda i: xs[i] >= target))
def binary_search(xs: list[int], target: int) -> int:
    ...
\end{lstlisting}

\paragraph{Level 2: add the loop invariant.}
The postconditions require proving a loop invariant. The programmer annotates
the loop:
\begin{lstlisting}[style=deppy,language=Python]
    while lo < hi:
        with invariant():
            have(0 <= lo <= hi <= len(xs))
            have(forall(range(lo), lambda i: xs[i] < target))
            have(forall(range(hi, len(xs)),
                        lambda i: xs[i] >= target))
        mid = (lo + hi) // 2
        if xs[mid] < target:
            lo = mid + 1
        else:
            hi = mid
\end{lstlisting}

The system verifies that:
\begin{enumerate}[nosep]
  \item The invariant holds on loop entry (from the initialization and
        the sorted-input precondition).
  \item The invariant is preserved by each iteration.
  \item The invariant plus the loop-exit condition ($\py{lo} \geq \py{hi}$)
        implies the postconditions.
\end{enumerate}

All three checks reduce to decidable linear arithmetic over bounded
quantifiers, so the solver handles them automatically. The programmer needed
to provide only the loop invariant; everything else was inferred or checked
by the solver.

\section{Cost-benefit analysis of annotation depth}

The following table summarizes the annotation cost and verification benefit
at each level, using binary search as the running example:

\begin{center}
\begin{tabular}{lllp{5cm}}
\toprule
\textbf{Level} & \textbf{Annotation cost} & \textbf{Verification benefit} &
  \textbf{What is checked} \\
\midrule
0 & 0 lines & Bounds safety & Array accesses in bounds, integer ranges \\
1 & 6 lines & Functional contract & Pre/postconditions stated; callers checked \\
2 & 5 lines & Full correctness & Loop invariant verified; postcondition proved \\
3--5 & N/A & N/A & Not needed for this example \\
\bottomrule
\end{tabular}
\end{center}

The key observation is that 11 lines of annotation (Levels 1+2) suffice to
prove full functional correctness of binary search, while the inference
engine handles all the mechanical work: bounds checking, variable tracking,
invariant checking against the solver.

\chapter{The Proof Term Language: Agda-Style Reasoning in Python}

This chapter defines the complete proof term language available in proof mode.
The language is designed to be embeddable in Python (all constructs are valid
Python expressions or statements) while supporting the full power of
constructive dependent type theory.

\section{Proof terms as Python expressions}

Every proof term is a Python expression that, at analysis time, carries a
type (the proposition being proved) and, at runtime, evaluates to \py{None}
(proofs are erased). The proof term language is defined by the following
grammar:

\begin{enumerate}[nosep]
  \item \textbf{Reflexivity.} \py{refl()} proves $a = a$ for any expression
        $a$. The system checks that both sides of the equality are
        definitionally equal after normalization.
  \item \textbf{Symmetry.} \py{symm(p)} takes a proof $p : a = b$ and
        produces a proof of $b = a$.
  \item \textbf{Transitivity.} \py{trans(p, q)} takes proofs $p : a = b$ and
        $q : b = c$ and produces a proof of $a = c$.
  \item \textbf{Congruence.} \py{cong(f, p)} takes a function $f$ and a
        proof $p : a = b$ and produces a proof of $f(a) = f(b)$.
  \item \textbf{Substitution.} \py{subst(p, q)} takes a proof $p : a = b$
        and a proof $q : P(a)$ and produces a proof of $P(b)$.
  \item \textbf{Modus ponens.} \py{mp(p, q)} takes a proof $p : A \Rightarrow B$
        and a proof $q : A$ and produces a proof of $B$.
  \item \textbf{Introduction.} \py{intro(lambda x: ...)} introduces a
        universally-quantified variable or a function argument.
  \item \textbf{Elimination.} \py{elim(p, a)} applies a universally-quantified
        proof $p : \forall x.\ P(x)$ to a specific value $a$, producing
        $P(a)$.
  \item \textbf{And-introduction.} \py{and_intro(p, q)} takes proofs $p : A$
        and $q : B$ and produces a proof of $A \wedge B$.
  \item \textbf{And-elimination.} \py{fst(p)} and \py{snd(p)} project from
        a conjunction proof.
  \item \textbf{Or-introduction.} \py{inl(p)} and \py{inr(q)} inject into
        a disjunction.
  \item \textbf{Or-elimination.} \py{cases(p, f, g)} takes a proof
        $p : A \vee B$, a function $f : A \to C$, and a function
        $g : B \to C$, and produces a proof of $C$.
  \item \textbf{Absurdity.} \py{absurd(p)} takes a proof of $\bot$ (False)
        and produces a proof of anything.
  \item \textbf{Solver.} \py{auto()} asks the SMT solver to discharge the
        current goal. If the solver succeeds, the proof is complete.
  \item \textbf{Assumption.} \py{by_assumption(name)} references a named
        assumption in the current context.
\end{enumerate}

\section{Tactic-style proofs}

In addition to direct proof terms, the system supports a tactic-style
interface where the programmer manipulates a proof goal:

\begin{lstlisting}[style=deppy,language=Python]
from deppy.proof import Proof

with Proof("sum_zero") as p:
    p.goal("forall(range(n), lambda i: 0) == 0 * n")
    p.induction("n")
    with p.case("n == 0"):
        p.simp()  # simplifies to 0 == 0
    with p.case("n > 0"):
        p.rewrite("forall(range(n), lambda i: 0)",
                   "forall(range(n-1), lambda i: 0) + 0")
        p.apply_ih()  # applies induction hypothesis
        p.simp()      # simplifies arithmetic
\end{lstlisting}

The tactic interface provides:
\begin{itemize}[nosep]
  \item \py{p.goal(expr)}: sets the proof goal.
  \item \py{p.induction(var)}: initiates induction on a variable.
  \item \py{p.case(cond)}: introduces a case split.
  \item \py{p.simp()}: applies arithmetic and boolean simplification.
  \item \py{p.rewrite(lhs, rhs)}: rewrites an expression using a known
        equality.
  \item \py{p.apply_ih()}: applies the induction hypothesis.
  \item \py{p.apply(lemma)}: applies a previously-proved lemma.
  \item \py{p.auto()}: asks the solver to close the goal.
  \item \py{p.unfold(name)}: unfolds a definition.
  \item \py{p.split()}: splits a conjunction goal into two sub-goals.
  \item \py{p.left()} / \py{p.right()}: selects the left/right side of a
        disjunction goal.
  \item \py{p.intro(name)}: introduces a universally-quantified variable.
  \item \py{p.exact(term)}: provides an exact proof term.
  \item \py{p.contradiction()}: derives False from contradictory hypotheses.
\end{itemize}

\section{Proof by calculation}

For equational reasoning, the system supports calculation proofs:
\begin{lstlisting}[style=deppy,language=Python]
from deppy.proof import calc

@ensures(lambda n, result: result == n * (n + 1) * (2 * n + 1) // 6)
def sum_squares(n: Refine[int, lambda n: n >= 0]) -> int:
    total = 0
    for i in range(n + 1):
        total += i * i
    with calc("sum_squares_formula"):
        step("sum(i**2 for i in range(n+1))")
        eq("sum(i**2 for i in range(n)) + n**2",
           by="definition of range")
        eq("(n-1)*n*(2*(n-1)+1)//6 + n**2",
           by="induction hypothesis")
        eq("(n-1)*n*(2*n-1)//6 + n**2",
           by="arithmetic")
        eq("((n-1)*n*(2*n-1) + 6*n**2) // 6",
           by="common denominator")
        eq("(2*n**3 - 3*n**2 + n + 6*n**2) // 6",
           by="expansion")
        eq("(2*n**3 + 3*n**2 + n) // 6",
           by="simplification")
        eq("n * (n+1) * (2*n+1) // 6",
           by="factoring")
    return total
\end{lstlisting}

Each \py{eq(..., by=...)} step asserts that the previous expression equals
the new one, with a justification. The solver verifies each step
independently, so the programmer only needs to provide enough intermediate
steps for each gap to be decidable.

\section{Lemma libraries and reuse}

Proved lemmas are first-class objects that can be stored, imported, and
applied:
\begin{lstlisting}[style=deppy,language=Python]
from deppy import lemma, Prop

@lemma
def add_comm(a: int, b: int) -> Prop["a + b == b + a"]:
    return auto()

@lemma
def add_assoc(a: int, b: int, c: int
              ) -> Prop["(a + b) + c == a + (b + c)"]:
    return auto()

@lemma
def mul_dist(a: int, b: int, c: int
             ) -> Prop["a * (b + c) == a * b + a * c"]:
    return auto()
\end{lstlisting}

These lemmas are proved once and available for use in any proof via
\py{p.apply(add_comm)} or \py{mp(add_comm(x, y), ...)}.

\paragraph{Library lemma bundles.}
The system ships with pre-proved lemma bundles for common domains:
\begin{itemize}[nosep]
  \item \py{deppy.lemmas.arithmetic}: commutativity, associativity,
        distributivity, division properties.
  \item \py{deppy.lemmas.lists}: length properties, append, concatenation,
        reversal, sorting.
  \item \py{deppy.lemmas.sets}: membership, union, intersection, subset.
  \item \py{deppy.lemmas.order}: transitivity, antisymmetry, total order.
  \item \py{deppy.lemmas.tensors}: shape algebra, broadcasting rules,
        dimension arithmetic.
\end{itemize}

\section{Proof automation levels}

The system provides multiple levels of proof automation:

\paragraph{Level A: Full automation (\py{auto()}).}
The solver attempts to discharge the goal using its decidable theories. This
handles most arithmetic, boolean, and array-theory goals.

\paragraph{Level B: Simplification (\py{simp()}).}
The simplifier applies rewriting rules (commutativity, associativity,
identity elements, beta reduction) to normalize the goal. This often reduces
complex goals to forms that \py{auto()} can handle.

\paragraph{Level C: Heuristic search (\py{search()}).}
The system searches for a proof by trying combinations of available lemmas,
hypotheses, and proof strategies. This is more expensive but can find proofs
that \py{auto()} misses.

\paragraph{Level D: Interactive (\py{sorry()}).}
When all automation fails, \py{sorry()} marks the goal as unproved but
accepted. This allows the programmer to continue developing surrounding
proofs while deferring a difficult sub-goal. The system tracks all
\py{sorry()} uses and reports them as obligations.

\section{Proof-carrying code: embedding proofs in functions}

Proofs can be embedded directly in function bodies, interleaved with
computation:
\begin{lstlisting}[style=deppy,language=Python]
from deppy import requires, ensures, have, hence, Refine

@requires(lambda xs: len(xs) > 0)
@ensures(lambda xs, result: result in xs)
@ensures(lambda xs, result:
    forall(range(len(xs)), lambda i: result <= xs[i]))
def find_minimum(xs: list[int]) -> int:
    minimum = xs[0]
    have("minimum == xs[0]")

    for i in range(1, len(xs)):
        with invariant():
            have("minimum in xs[:i]")
            have("forall(range(i), lambda j: minimum <= xs[j])")
        if xs[i] < minimum:
            minimum = xs[i]
            have("minimum == xs[i]")
            have("minimum in xs[:i+1]")
        else:
            have("minimum <= xs[i]")
            have("minimum in xs[:i+1]")

    hence("minimum in xs")
    hence("forall(range(len(xs)), lambda j: minimum <= xs[j])")
    return minimum
\end{lstlisting}

The \py{have} and \py{hence} statements are proof checkpoints. Each one is
verified by the solver in the context of the current program state. Between
checkpoints, the system fills in the logical steps automatically. The result
is that the programmer provides a human-readable proof outline, and the
system verifies that each step follows.

\section{The interaction between proof terms and inference}

When proof mode is active, every proved property immediately enriches the
sheaf-descent sections:
\begin{itemize}[nosep]
  \item A proved postcondition becomes a guaranteed output section.
  \item A proved loop invariant becomes a fixed-point section at the loop
        header.
  \item A proved lemma becomes an available axiom for all sites that can
        reference its variables.
  \item A proved class invariant becomes a guaranteed property at every
        method entry and exit.
\end{itemize}

This enrichment propagates through the interprocedural transport mechanism.
If function $f$ has a proved postcondition and function $g$ calls $f$, then
$g$'s local section at the call return site is enriched with $f$'s
postcondition---even if $g$ has no annotations at all. This is the
\emph{annotation leverage} effect.


\chapter{Class Invariants and Object-Oriented Proofs}

Object-oriented code organizes state into objects with methods that maintain
invariants. The proof mode supports class-level invariants that are checked
at construction and preserved by every method.

\section{Class invariants as persistent sections}

A class invariant is a property that holds for every instance of the class
at every method boundary. In the sheaf-descent framework, the invariant is
a persistent section attached to the class's site family.

\begin{lstlisting}[style=deppy,language=Python]
from deppy import invariant_of, requires, ensures

@invariant_of(lambda self:
    0 <= self.size <= self.capacity and
    len(self.items) == self.size)
class BoundedBuffer:
    def __init__(self, capacity: Refine[int, lambda c: c > 0]):
        self.capacity = capacity
        self.items: list = []
        self.size = 0

    @requires(lambda self: self.size < self.capacity)
    @ensures(lambda self: self.size == old(self.size) + 1)
    def push(self, item):
        self.items.append(item)
        self.size += 1

    @requires(lambda self: self.size > 0)
    @ensures(lambda self, result:
        self.size == old(self.size) - 1)
    def pop(self):
        self.size -= 1
        return self.items.pop()

    @ensures(lambda self, result: result == (self.size == 0))
    def is_empty(self) -> bool:
        return self.size == 0

    @ensures(lambda self, result: result == (self.size == self.capacity))
    def is_full(self) -> bool:
        return self.size == self.capacity
\end{lstlisting}

The system verifies:
\begin{enumerate}[nosep]
  \item \textbf{Establishment.} The invariant holds after \py{__init__}
        completes.
  \item \textbf{Preservation.} Every public method preserves the invariant:
        if the invariant holds on entry (plus the method's precondition),
        it holds on exit.
  \item \textbf{Usage.} Every method can assume the invariant on entry, which
        enriches the local section. For example, \py{pop} knows that
        $\py{self.size} > 0$ (from the precondition) and that
        $\py{len(self.items)} = \py{self.size}$ (from the invariant), so the
        \py{self.items.pop()} call is safe.
\end{enumerate}

\section{The \py{old()} function for state change reasoning}

Postconditions often need to relate the post-state to the pre-state. The
\py{old(expr)} function captures the value of an expression at method entry:

\begin{itemize}[nosep]
  \item \py{old(self.size)}: the size before the method executed.
  \item \py{old(self.items)}: a snapshot of the items list before mutation.
  \item \py{old(len(xs))}: the length of \py{xs} before the method.
\end{itemize}

The system implements \py{old} by creating a shadow copy of the referenced
expressions at method entry and substituting these copies in the
postcondition. At analysis time, the pre-state values are symbolic; at
runtime (for checked contracts), they are actual snapshots.

\section{Inheritance and invariant strengthening}

When a subclass extends a class with an invariant:
\begin{lstlisting}[style=deppy,language=Python]
@invariant_of(lambda self:
    all(x >= 0 for x in self.items))
class NonNegativeBuffer(BoundedBuffer):
    @requires(lambda self, item: item >= 0)
    def push(self, item):
        super().push(item)
\end{lstlisting}

The subclass invariant is the \emph{conjunction} of the parent invariant and
the subclass-specific invariant. The system verifies that:
\begin{itemize}[nosep]
  \item The subclass \py{__init__} establishes both invariants.
  \item Overridden methods preserve both invariants.
  \item The Liskov substitution principle holds: the subclass's preconditions
        are no stronger than the parent's (for inherited methods), and the
        postconditions are no weaker.
\end{itemize}

\section{Protocol proofs: verifying structural contracts}

Python protocols (from \py{typing.Protocol}) define structural interfaces.
Proof mode can verify that a class correctly implements a protocol:

\begin{lstlisting}[style=deppy,language=Python]
from typing import Protocol, runtime_checkable
from deppy import implements

@runtime_checkable
class Stack(Protocol):
    @requires(lambda self: True)
    @ensures(lambda self: self.size() == old(self.size()) + 1)
    def push(self, item) -> None: ...

    @requires(lambda self: self.size() > 0)
    @ensures(lambda self, result:
        self.size() == old(self.size()) - 1)
    def pop(self) -> object: ...

    @ensures(lambda self, result: result >= 0)
    def size(self) -> int: ...

@implements(Stack)
class ListStack:
    def __init__(self):
        self.items: list = []

    def push(self, item) -> None:
        self.items.append(item)

    def pop(self):
        return self.items.pop()

    def size(self) -> int:
        return len(self.items)
\end{lstlisting}

The \py{@implements(Stack)} decorator instructs the system to verify that
\py{ListStack} satisfies the \py{Stack} protocol's contracts. This includes
signature compatibility, precondition weakening (the implementation's
preconditions must be at least as weak as the protocol's), and postcondition
strengthening (the implementation's postconditions must be at least as
strong).


\chapter{Proving Properties of Data Structures}

Data structures are the primary targets for proof-mode investment, because
their invariants are critical for correctness and are often violated by subtle
bugs that testing misses.

\section{Sorted lists and binary search trees}

\subsection{Proving sortedness preservation}

\begin{lstlisting}[style=deppy,language=Python]
from deppy import requires, ensures, forall, Refine, lemma

Sorted = Refine[list[int], lambda xs:
    forall(range(len(xs) - 1), lambda i: xs[i] <= xs[i+1])]

@requires(lambda xs: True)
@ensures(lambda xs, result: isinstance(result, list))
@ensures(lambda xs, result: len(result) == len(xs))
@ensures(lambda xs, result:
    forall(range(len(result) - 1),
           lambda i: result[i] <= result[i+1]))
@ensures(lambda xs, result:
    forall(xs, lambda x: x in result))
def insertion_sort(xs: list[int]) -> Sorted:
    result = []
    for x in xs:
        with invariant():
            have("isinstance(result, Sorted)")
            have("len(result) == len(xs[:i])")
        result = insert_sorted(result, x)
    return result

@requires(lambda xs: isinstance(xs, Sorted))
@ensures(lambda xs, x, result: isinstance(result, Sorted))
@ensures(lambda xs, x, result: len(result) == len(xs) + 1)
@ensures(lambda xs, x, result: x in result)
def insert_sorted(xs: Sorted, x: int) -> Sorted:
    for i in range(len(xs)):
        if x <= xs[i]:
            return xs[:i] + [x] + xs[i:]
    return xs + [x]
\end{lstlisting}

The system verifies:
\begin{enumerate}[nosep]
  \item \py{insert_sorted} preserves sortedness: given a sorted list, the
        output is sorted. This is verified by case analysis on the insertion
        point.
  \item \py{insertion_sort} produces a sorted list: the loop invariant states
        that \py{result} is always sorted, and \py{insert_sorted} preserves
        this.
  \item The output is a permutation of the input: every element of the input
        appears in the output.
\end{enumerate}

\subsection{Binary search tree invariant}

\begin{lstlisting}[style=deppy,language=Python]
from deppy import invariant_of

@invariant_of(lambda self:
    (self.left is None or
     forall(self.left.all_values(),
            lambda v: v < self.value)) and
    (self.right is None or
     forall(self.right.all_values(),
            lambda v: v > self.value)) and
    (self.left is None or self.left.is_bst()) and
    (self.right is None or self.right.is_bst()))
class BST:
    def __init__(self, value: int):
        self.value = value
        self.left: BST | None = None
        self.right: BST | None = None

    def is_bst(self) -> bool:
        """Check the BST invariant (used in the invariant itself)."""
        ...

    def all_values(self) -> list[int]:
        """Return all values in the tree."""
        result = [self.value]
        if self.left is not None:
            result = self.left.all_values() + result
        if self.right is not None:
            result = result + self.right.all_values()
        return result

    @ensures(lambda self, key, result:
        key in self.all_values() or result is None)
    @ensures(lambda self, key, result:
        result is not None and result.value == key
        or result is None and key not in self.all_values())
    def search(self, key: int) -> "BST | None":
        if key == self.value:
            return self
        elif key < self.value and self.left is not None:
            return self.left.search(key)
        elif key > self.value and self.right is not None:
            return self.right.search(key)
        return None
\end{lstlisting}

The BST invariant is recursive: each subtree must also be a BST. The system
handles this through structural induction on the tree, verifying that each
method preserves the invariant at every level of the recursion.

\section{Hash maps and amortized complexity}

While the system focuses on functional correctness rather than complexity,
some correctness properties depend on amortized arguments. For example, a
hash map with linear probing is correct only if the load factor remains
below 1. The proof mode can express this:

\begin{lstlisting}[style=deppy,language=Python]
@invariant_of(lambda self:
    self.count < self.capacity and
    self.capacity > 0 and
    self.count >= 0)
class HashMap:
    def __init__(self, initial_capacity: Refine[int, lambda c: c > 0]):
        self.capacity = initial_capacity
        self.count = 0
        self.table = [None] * initial_capacity

    @requires(lambda self: self.count < self.capacity - 1)
    def put(self, key, value):
        # linear probing always terminates when count < capacity
        idx = hash(key) % self.capacity
        while self.table[idx] is not None:
            with invariant():
                have("exists an empty slot in self.table")
            if self.table[idx][0] == key:
                self.table[idx] = (key, value)
                return
            idx = (idx + 1) % self.capacity
        self.table[idx] = (key, value)
        self.count += 1
\end{lstlisting}

The invariant $\py{self.count} < \py{self.capacity}$ ensures that the linear
probing loop always terminates (there is always at least one empty slot).

\chapter{Proving Properties of Algorithms}

This chapter demonstrates proof mode on a range of algorithms, showing how
the annotation effort scales with the complexity of the correctness property
and how the inference engine reduces the burden.

\section{Sorting algorithms: from specification to proof}

\subsection{The sorting specification}

A sorting function must satisfy three properties:
\begin{enumerate}[nosep]
  \item \textbf{Sorted output.} The output is in non-decreasing order.
  \item \textbf{Permutation.} The output contains exactly the same elements
        as the input (with the same multiplicities).
  \item \textbf{Length preservation.} The output has the same length as the
        input.
\end{enumerate}

In our notation:
\begin{lstlisting}[style=deppy,language=Python]
from deppy import ensures, forall, Refine
from collections import Counter

@ensures(lambda xs, result: len(result) == len(xs))
@ensures(lambda xs, result:
    forall(range(len(result) - 1),
           lambda i: result[i] <= result[i+1]))
@ensures(lambda xs, result: Counter(result) == Counter(xs))
def sort_spec(xs: list[int]) -> list[int]:
    """Specification for a correct sorting function."""
    ...
\end{lstlisting}

\subsection{Merge sort with full proof}

\begin{lstlisting}[style=deppy,language=Python]
from deppy import requires, ensures, forall, have, hence
from deppy import decreases, lemma, auto

@ensures(lambda xs, result: len(result) == len(xs))
@ensures(lambda xs, result:
    forall(range(len(result) - 1),
           lambda i: result[i] <= result[i+1]))
@ensures(lambda xs, result: Counter(result) == Counter(xs))
@decreases(lambda xs: len(xs))
def merge_sort(xs: list[int]) -> list[int]:
    if len(xs) <= 1:
        have("xs is trivially sorted")
        return xs[:]

    mid = len(xs) // 2
    have("0 < mid < len(xs)")

    left = merge_sort(xs[:mid])
    have("len(left) == mid")
    have("forall(range(len(left)-1), "
         "lambda i: left[i] <= left[i+1])")
    have("Counter(left) == Counter(xs[:mid])")

    right = merge_sort(xs[mid:])
    have("len(right) == len(xs) - mid")
    have("forall(range(len(right)-1), "
         "lambda i: right[i] <= right[i+1])")
    have("Counter(right) == Counter(xs[mid:])")

    result = merge(left, right)
    hence("len(result) == len(xs)")
    hence("Counter(result) == Counter(xs)")
    return result
\end{lstlisting}

The \py{@decreases(lambda xs: len(xs))} annotation ensures termination: at
each recursive call, the input length strictly decreases. The proof
checkpoints (\py{have}) document the key facts at each stage; the solver
verifies each one from the postconditions of the recursive calls and the
\py{merge} function.

\subsection{The merge function}

\begin{lstlisting}[style=deppy,language=Python]
@requires(lambda a, b:
    forall(range(len(a)-1), lambda i: a[i] <= a[i+1]))
@requires(lambda a, b:
    forall(range(len(b)-1), lambda i: b[i] <= b[i+1]))
@ensures(lambda a, b, result: len(result) == len(a) + len(b))
@ensures(lambda a, b, result:
    forall(range(len(result)-1),
           lambda i: result[i] <= result[i+1]))
@ensures(lambda a, b, result:
    Counter(result) == Counter(a) + Counter(b))
def merge(a: list[int], b: list[int]) -> list[int]:
    result = []
    i, j = 0, 0
    while i < len(a) and j < len(b):
        with invariant():
            have("len(result) == i + j")
            have("forall(range(len(result)-1), "
                 "lambda k: result[k] <= result[k+1])")
            have("(len(result) == 0 or "
                 "(i < len(a) and result[-1] <= a[i]) or "
                 "(j < len(b) and result[-1] <= b[j]))")
            have("Counter(result) == "
                 "Counter(a[:i]) + Counter(b[:j])")
        if a[i] <= b[j]:
            result.append(a[i])
            i += 1
        else:
            result.append(b[j])
            j += 1
    result.extend(a[i:])
    result.extend(b[j:])
    return result
\end{lstlisting}

The loop invariant captures four properties:
\begin{enumerate}[nosep]
  \item The result length tracks the progress ($i + j$).
  \item The result is always sorted.
  \item The last element of the result is less than or equal to the next
        unconsumed element from either input.
  \item The result is a merge of the consumed portions of both inputs.
\end{enumerate}

The system verifies each invariant clause inductively: it holds initially
(empty result, $i = j = 0$), and each iteration preserves it. This is
within the decidable fragment (bounded quantification, linear arithmetic,
array reasoning).

\section{Graph algorithms}

\subsection{Breadth-first search: reachability proof}

\begin{lstlisting}[style=deppy,language=Python]
from deppy import requires, ensures, forall, have, invariant

Graph = dict[int, list[int]]

@requires(lambda graph, start: start in graph)
@ensures(lambda graph, start, result:
    start in result)
@ensures(lambda graph, start, result:
    forall(result,
           lambda v: is_reachable(graph, start, v)))
@ensures(lambda graph, start, result:
    forall(graph.keys(),
           lambda v: is_reachable(graph, start, v)
                     or v not in result))
def bfs(graph: Graph, start: int) -> set[int]:
    visited = {start}
    queue = [start]

    while queue:
        with invariant():
            have("start in visited")
            have("forall(visited, "
                 "lambda v: is_reachable(graph, start, v))")
            have("forall(queue, lambda v: v in visited)")
            have("forall(visited, lambda v: "
                 "forall(graph.get(v, []), lambda w: "
                 "w in visited or w in queue or "
                 "v in queue))")

        node = queue.pop(0)
        for neighbor in graph.get(node, []):
            if neighbor not in visited:
                visited.add(neighbor)
                queue.append(neighbor)
                have("is_reachable(graph, start, neighbor)")

    hence("queue is empty, so all neighbors of visited "
          "nodes are in visited")
    hence("visited is the complete reachable set")
    return visited
\end{lstlisting}

The loop invariant captures the BFS frontier property: every visited node is
reachable from start, every queued node is visited, and every neighbor of a
visited node is either visited or queued. When the queue is empty, the last
property implies that the visited set is closed under the neighbor relation,
hence it contains all reachable nodes.

\subsection{Topological sort: acyclicity and ordering}

\begin{lstlisting}[style=deppy,language=Python]
from deppy import requires, ensures, forall, Refine

DAG = Refine[Graph, lambda g: is_acyclic(g)]

@requires(lambda graph: is_acyclic(graph))
@ensures(lambda graph, result: len(result) == len(graph))
@ensures(lambda graph, result:
    set(result) == set(graph.keys()))
@ensures(lambda graph, result:
    forall(range(len(result)),
           lambda i: forall(graph.get(result[i], []),
                            lambda dep: result.index(dep) < i)))
def topological_sort(graph: DAG) -> list[int]:
    in_degree = {v: 0 for v in graph}
    for v in graph:
        for w in graph[v]:
            in_degree[w] = in_degree.get(w, 0) + 1

    queue = [v for v in graph if in_degree[v] == 0]
    result = []

    while queue:
        with invariant():
            have("len(result) + len(queue) + "
                 "sum(1 for v in graph if in_degree[v] > 0 "
                 "and v not in result) == len(graph)")
            have("forall(range(len(result)), "
                 "lambda i: forall(graph.get(result[i],[]),"
                 "lambda d: d in result[:i]))")

        v = queue.pop(0)
        result.append(v)
        for w in graph[v]:
            in_degree[w] -= 1
            if in_degree[w] == 0:
                queue.append(w)

    have("len(result) == len(graph)")
    return result
\end{lstlisting}

The postcondition asserts topological ordering: for every node $v$ in the
result, all of $v$'s dependencies appear earlier. The acyclicity precondition
is essential for Kahn's algorithm to process all nodes.

\section{Numeric algorithms}

\subsection{Newton's method: convergence and precision}

For numeric algorithms, the system can verify safety properties (no division
by zero, no overflow) while leaving convergence properties as proof
obligations:

\begin{lstlisting}[style=deppy,language=Python]
from deppy import requires, ensures, Refine, have

@requires(lambda x: x > 0)
@requires(lambda tol: tol > 0)
@requires(lambda max_iter: max_iter > 0)
@ensures(lambda x, tol, max_iter, result:
    result > 0)
@ensures(lambda x, tol, max_iter, result:
    abs(result * result - x) <= tol
    or "max iterations reached")
def sqrt_newton(
    x: Refine[float, lambda v: v > 0],
    tol: Refine[float, lambda v: v > 0] = 1e-10,
    max_iter: Refine[int, lambda v: v > 0] = 100,
) -> float:
    guess = x / 2.0 if x > 1.0 else 1.0
    have("guess > 0")

    for _ in range(max_iter):
        with invariant():
            have("guess > 0")
        new_guess = (guess + x / guess) / 2.0
        have("guess > 0 implies x / guess is finite")
        have("new_guess > 0")
        if abs(new_guess - guess) < tol:
            return new_guess
        guess = new_guess

    return guess
\end{lstlisting}

The system verifies:
\begin{itemize}[nosep]
  \item No division by zero: \py{guess > 0} is maintained by the invariant,
        so \py{x / guess} is safe.
  \item The output is positive: follows from the invariant.
  \item Each iteration produces a finite result: follows from \py{guess > 0}
        and \py{x > 0}.
\end{itemize}

The convergence property ($|r^2 - x| \leq \epsilon$) is stated as a
postcondition but qualified with ``or max iterations reached.'' The system
cannot prove convergence automatically (it depends on real analysis), but
it can verify all the supporting numeric safety properties.

\subsection{Matrix multiplication: shape correctness with value proof}

\begin{lstlisting}[style=deppy,language=Python]
from deppy import requires, ensures, forall

@requires(lambda A, B:
    len(A) > 0 and len(A[0]) > 0 and
    len(B) > 0 and len(B[0]) > 0 and
    len(A[0]) == len(B))
@ensures(lambda A, B, result:
    len(result) == len(A) and
    len(result[0]) == len(B[0]))
@ensures(lambda A, B, result:
    forall(range(len(A)), lambda i:
        forall(range(len(B[0])), lambda j:
            result[i][j] == sum(
                A[i][k] * B[k][j]
                for k in range(len(A[0]))))))
def matmul(
    A: list[list[float]],
    B: list[list[float]],
) -> list[list[float]]:
    m, n, p = len(A), len(A[0]), len(B[0])
    result = [[0.0] * p for _ in range(m)]

    for i in range(m):
        for j in range(p):
            total = 0.0
            for k in range(n):
                with invariant():
                    have("total == sum(A[i][l] * B[l][j] "
                         "for l in range(k))")
                total += A[i][k] * B[k][j]
            result[i][j] = total
            have("result[i][j] == sum(A[i][k] * B[k][j] "
                 "for k in range(n))")

    return result
\end{lstlisting}

The innermost loop invariant tracks the partial sum. The solver verifies
that adding the $k$-th term extends the partial sum correctly. The outer
postcondition asserts that each element of the result matrix equals the
dot product of the corresponding row and column.

\section{String and parsing algorithms}

\subsection{Balanced parentheses checker}

\begin{lstlisting}[style=deppy,language=Python]
from deppy import ensures, have, invariant

@ensures(lambda s, result:
    result == is_balanced_spec(s))
def is_balanced(s: str) -> bool:
    depth = 0
    for ch in s:
        with invariant():
            have("depth >= 0")
            have("depth == count_open(s[:i]) - count_close(s[:i])")
        if ch == '(':
            depth += 1
        elif ch == ')':
            if depth == 0:
                return False
            depth -= 1
    return depth == 0
\end{lstlisting}

The invariant tracks that \py{depth} equals the number of unmatched open
parentheses seen so far. The system verifies that returning \py{False} when
\py{depth == 0} and \py{ch == ')'} is correct (an unmatched close paren),
and that returning \py{depth == 0} at the end correctly identifies balanced
strings.

\section{Annotation cost summary}

\begin{center}
\begin{tabular}{llrrp{4cm}}
\toprule
\textbf{Algorithm} & \textbf{Lines of code} & \textbf{Annotation lines} &
  \textbf{Ratio} & \textbf{What is proved} \\
\midrule
Binary search & 8 & 11 & 1.4x & Correct insertion point \\
Merge sort & 12 & 18 & 1.5x & Sorted, permutation, length \\
Merge & 14 & 12 & 0.9x & Sorted merge, permutation \\
BFS & 12 & 15 & 1.3x & Complete reachability \\
Topological sort & 14 & 10 & 0.7x & Topological ordering \\
Newton's method & 10 & 8 & 0.8x & Numeric safety, positivity \\
Matrix multiply & 10 & 12 & 1.2x & Value correctness \\
Balanced parens & 10 & 6 & 0.6x & Specification equivalence \\
\bottomrule
\end{tabular}
\end{center}

The annotation-to-code ratio ranges from 0.6x to 1.5x, with an average of
about 1.0x. This is dramatically less than traditional proof assistants, where
the proof-to-code ratio is typically 5--20x. The savings come from the solver
handling most arithmetic and array reasoning automatically, and the inference
engine providing the mechanical bookkeeping (variable tracking, control flow,
shape propagation).

\chapter{The Proof Engine: How Proofs Flow Through the Sheaf Machinery}

This chapter describes the internal architecture of the proof engine---how
proof obligations are generated, propagated, simplified, and discharged
within the sheaf-descent framework. The key insight is that proofs are
\emph{not} a separate layer bolted onto the inference engine; they are
additional constraints in the same section-optimization problem.

\section{Proof obligations as section constraints}

Recall that the inference engine computes sections (type assignments) at each
site by solving an optimization problem: maximize the information content of
sections subject to the constraint that sections agree on overlaps. Proof mode
adds a new kind of constraint: \emph{proof obligations}.

A proof obligation is a proposition $P$ that must hold at a specific site $s$.
In the section-optimization framework, this corresponds to the constraint
\[
  \sigma(s) \models P
\]
where $\sigma(s)$ is the section assigned to site $s$. The obligation may come
from:
\begin{itemize}[nosep]
  \item A postcondition: $P$ must hold at the function exit site.
  \item An invariant: $P$ must hold at the loop header site on every iteration.
  \item A \py{have} statement: $P$ must hold at the statement's site.
  \item A caller obligation: $P$ must hold at a call site to satisfy the
        callee's precondition.
\end{itemize}

\section{The obligation lifecycle}

Each proof obligation goes through a lifecycle:

\paragraph{Generation.}
When the system encounters a \py{@ensures}, \py{@requires}, \py{with
invariant()}, or \py{have()} construct, it generates an obligation. The
obligation records:
\begin{itemize}[nosep]
  \item The proposition $P$.
  \item The site $s$ where $P$ must hold.
  \item The source location in the Python code.
  \item The available context: all sections known at $s$.
\end{itemize}

\paragraph{Simplification.}
Before sending the obligation to the solver, the system simplifies it:
\begin{enumerate}[nosep]
  \item \textbf{Context reduction.} If $P$ is already implied by the
        current section at $s$ (from inference), the obligation is trivially
        discharged.
  \item \textbf{Substitution.} Symbolic values are substituted: if the
        section at $s$ knows that $x = 42$, then $P(x)$ becomes $P(42)$.
  \item \textbf{Normalization.} Arithmetic expressions are normalized
        (constant folding, associativity, commutativity).
  \item \textbf{Decomposition.} Conjunctive obligations ($P \wedge Q$) are
        split into separate obligations for $P$ and $Q$.
\end{enumerate}

\paragraph{Discharge.}
The simplified obligation is sent to the SMT solver. If the solver returns
SAT (satisfiable with a model showing the obligation holds) or UNSAT
(the negation is unsatisfiable, meaning the obligation always holds), the
obligation is discharged. In practice, the solver checks
$\Gamma \models P$ where $\Gamma$ is the conjunction of all known facts at
site $s$.

\paragraph{Escalation.}
If the solver returns UNKNOWN or times out, the obligation is escalated:
\begin{enumerate}[nosep]
  \item First, the system attempts more expensive proof strategies
        (lemma search, case splitting, bounded unrolling).
  \item If all automation fails, the obligation becomes a \emph{residual}
        reported to the programmer.
  \item The programmer can then provide a proof, a lemma, or mark the
        obligation as \py{sorry()} or \py{@assume}.
\end{enumerate}

\section{Obligation propagation through the site cover}

Proof obligations propagate through the site cover via the same transport
maps that propagate inferred sections:

\paragraph{Backward propagation of postconditions.}
A postcondition $P$ at the function exit site generates obligations at each
\py{return} site. At each return, the system checks whether the current
section implies $P$. If not, the system traces backward through the control
flow to find which earlier site must be strengthened.

\paragraph{Forward propagation of preconditions.}
A precondition $P$ at a callee's entry site generates an obligation at the
call site. The caller must ensure that $P$ holds given the caller's section
at the call site. If the caller's section is too weak, the obligation is
propagated further backward to the caller's entry, becoming a precondition of
the caller.

\paragraph{Loop invariant verification.}
A loop invariant $I$ generates three obligations:
\begin{enumerate}[nosep]
  \item \textbf{Initialization}: $I$ holds at the loop entry (before the
        first iteration), given the section at the loop's predecessor site.
  \item \textbf{Preservation}: If $I$ holds at the loop header and the loop
        body executes, then $I$ holds at the end of the loop body (which is
        the loop header's predecessor in the next iteration).
  \item \textbf{Usage}: After the loop exits, $I \wedge \neg C$ holds, where
        $C$ is the loop condition. This conjunction enriches the section at
        the post-loop site.
\end{enumerate}

\paragraph{Cross-function propagation.}
When a proved postcondition of function $f$ enriches the section at a call
site in function $g$, this enrichment may allow obligations in $g$ to be
discharged that previously could not be. The system re-checks pending
obligations whenever a section is strengthened.

\section{The proof context stack}

At each program point, the proof engine maintains a \emph{proof context}: an
ordered list of known facts (hypotheses) that can be used to discharge
obligations. The context is organized as a stack of scopes:

\begin{enumerate}[nosep]
  \item \textbf{Global scope}: axioms, proved lemmas, library facts.
  \item \textbf{Module scope}: module-level assertions, import-derived facts.
  \item \textbf{Function scope}: preconditions, parameter refinements.
  \item \textbf{Block scope}: guard refinements, pattern-match bindings.
  \item \textbf{Loop scope}: loop invariant (on each iteration).
  \item \textbf{Proof scope}: facts established by \py{have} statements within
        a \py{with proof()} block.
\end{enumerate}

When checking an obligation at site $s$, the system collects all hypotheses
from all enclosing scopes and passes them to the solver as the context $\Gamma$.

\section{Interaction with the inference engine}

The proof engine and the inference engine share the same data structures and
work in interleaved passes:

\paragraph{Pass 1: Inference.}
The inference engine computes initial sections for all sites, ignoring proof
annotations. This establishes the baseline: what the system can determine
without any programmer help.

\paragraph{Pass 2: Annotation integration.}
The system integrates programmer annotations (contracts, invariants, lemmas)
as additional section constraints. This may strengthen some sections and
create new obligations.

\paragraph{Pass 3: Obligation discharge.}
The proof engine attempts to discharge all obligations, using the sections
from Pass 2 as context. Successfully discharged obligations confirm the
annotations; failed obligations become residuals.

\paragraph{Pass 4: Section enrichment.}
Successfully verified annotations enrich the sections. Proved postconditions
become guaranteed output sections. Proved invariants become fixed-point
sections. The enriched sections propagate through the transport maps.

\paragraph{Pass 5: Re-inference.}
The inference engine re-runs with the enriched sections. This may infer
additional properties that were not available in Pass 1 (because the
programmer's annotations provided information the inference engine could not
derive on its own).

This five-pass architecture means that a single annotation can have cascading
effects: the annotation is verified (Pass 3), enriches a section (Pass 4),
which enables new inference (Pass 5), which may discharge obligations in other
functions that depend on the newly-inferred properties.

\section{The solver interface}

The proof engine communicates with the SMT solver through a typed interface
that translates between the section representation and SMT-LIB:

\paragraph{Section to SMT-LIB translation.}
\begin{itemize}[nosep]
  \item Integer refinements become QF\_LIA constraints.
  \item Floating-point refinements become QF\_LRA or QF\_FPA constraints.
  \item Array/list refinements become array-theory constraints.
  \item Boolean refinements become propositional logic.
  \item Quantified properties become bounded-quantifier expansions or
        array-theory axioms.
  \item Uninterpreted functions model user-defined pure functions.
\end{itemize}

\paragraph{Solver result interpretation.}
\begin{itemize}[nosep]
  \item \textbf{UNSAT}: The negation of the obligation is unsatisfiable under
        the context, meaning the obligation always holds. Discharge it.
  \item \textbf{SAT with counterexample}: The obligation can be violated.
        The counterexample provides a concrete input that triggers the
        violation. Report it as a verification failure with the
        counterexample.
  \item \textbf{UNKNOWN}: The solver could not determine the result within
        the timeout. Escalate to more expensive strategies or report as a
        residual.
\end{itemize}

\paragraph{Incremental solving.}
The system uses Z3's incremental solving features (\py{push()}/\py{pop()}) to
share context across related obligations. This is critical for performance:
when checking a loop invariant's preservation across many paths through the
loop body, the common context (the invariant itself, the function's
precondition, the class invariant) is shared, and only the path-specific
facts are pushed and popped.

\section{Proof caching and incrementality}

The proof engine caches discharged obligations so that re-running the analysis
after a small code change does not re-prove everything:

\begin{itemize}[nosep]
  \item Each obligation is keyed by its proposition and context hash.
  \item If the code at a site has not changed and the context sections are
        unchanged, the cached discharge result is reused.
  \item If the code changes but the obligation's SMT-LIB encoding is
        identical, the cached result is still valid.
  \item Proof cache invalidation is conservative: any change to a function
        invalidates all obligations in that function, but obligations in
        other functions are preserved.
\end{itemize}

This means that in a large codebase, re-verification after a local change is
fast: only the changed function and its direct callers are re-checked.

\section{Error reporting and diagnostics}

When an obligation cannot be discharged, the system provides rich diagnostics:

\begin{lstlisting}[language={}]
PROOF OBLIGATION FAILED at line 42:
  Obligation: result[i] <= result[i+1]
  Context:
    - result is a list of int (inferred)
    - len(result) == i + 1 (from loop invariant)
    - result[i] == xs[idx] (from assignment)
  Missing: no relationship between xs[idx] and result[i-1]
  Suggestion: add a loop invariant relating result elements
              to the insertion position
  Counterexample: xs = [3, 1, 2], i = 1
    result = [3, 1], result[0] = 3 > result[1] = 1
\end{lstlisting}

The diagnostics include:
\begin{itemize}[nosep]
  \item The obligation that failed and its source location.
  \item The context: what the system knows at the failure point.
  \item The gap: what information is missing to discharge the obligation.
  \item A suggestion: what annotation might help.
  \item A counterexample (if the solver found one): a concrete input that
        demonstrates the failure.
\end{itemize}


\chapter{Proving Properties of PyTorch Programs}

Machine learning programs present unique verification challenges: tensor
shapes, numeric precision, gradient correctness, and model architecture
properties. This chapter demonstrates how proof mode handles PyTorch-specific
verification.

\section{Shape-correct neural network architectures}

\subsection{Proving shape compatibility of a multi-layer model}

\begin{lstlisting}[style=deppy,language=Python]
from deppy import requires, ensures, have
import torch
import torch.nn as nn

@requires(lambda in_dim, hidden_dims, out_dim:
    in_dim > 0 and out_dim > 0 and
    len(hidden_dims) > 0 and
    forall(hidden_dims, lambda d: d > 0))
@ensures(lambda in_dim, hidden_dims, out_dim, result:
    result.input_dim == in_dim and
    result.output_dim == out_dim)
def build_mlp(
    in_dim: int, hidden_dims: list[int], out_dim: int
) -> nn.Sequential:
    layers = []
    prev_dim = in_dim

    for i, h_dim in enumerate(hidden_dims):
        with invariant():
            have("prev_dim > 0")
            have("len(layers) == 2 * i")
        layers.append(nn.Linear(prev_dim, h_dim))
        layers.append(nn.ReLU())
        prev_dim = h_dim
        have("prev_dim == hidden_dims[i]")
        have("prev_dim > 0")

    layers.append(nn.Linear(prev_dim, out_dim))
    have("the final linear layer maps from "
         "hidden_dims[-1] to out_dim")

    model = nn.Sequential(*layers)
    return model
\end{lstlisting}

The system verifies that each \py{nn.Linear} layer's input dimension matches
the previous layer's output dimension. The loop invariant tracks \py{prev_dim}
through the layer construction. The final layer's input dimension equals the
last hidden dimension, which is positive by the precondition.

\subsection{Attention mechanism shape proof}

\begin{lstlisting}[style=deppy,language=Python]
from deppy import requires, ensures, have

@requires(lambda query, key, value:
    query.dim() == 3 and key.dim() == 3 and
    value.dim() == 3)
@requires(lambda query, key, value:
    query.shape[0] == key.shape[0] == value.shape[0])
@requires(lambda query, key, value:
    query.shape[2] == key.shape[2])
@requires(lambda query, key, value:
    key.shape[1] == value.shape[1])
@ensures(lambda query, key, value, result:
    result.shape == (query.shape[0], query.shape[1],
                     value.shape[2]))
def scaled_dot_product_attention(
    query: torch.Tensor,
    key: torch.Tensor,
    value: torch.Tensor,
) -> torch.Tensor:
    d_k = query.shape[2]
    have("d_k > 0")  # from query.dim() == 3

    scores = torch.matmul(query, key.transpose(-2, -1))
    have("scores.shape == (query.shape[0], query.shape[1],"
         " key.shape[1])")

    scores = scores / (d_k ** 0.5)
    have("scores.shape unchanged, d_k > 0 prevents div-by-zero")

    weights = torch.softmax(scores, dim=-1)
    have("weights.shape == scores.shape")
    have("forall weights values: 0 <= w <= 1")
    have("weights sum to 1 along last dimension")

    output = torch.matmul(weights, value)
    have("output.shape == (query.shape[0], query.shape[1],"
         " value.shape[2])")

    return output
\end{lstlisting}

The proof traces shapes through each operation:
\begin{itemize}[nosep]
  \item \py{key.transpose(-2, -1)} swaps the last two dimensions: shape
        becomes $(B, D_k, S_k)$.
  \item \py{torch.matmul(query, key.T)}: matrix multiply on last two dims:
        $(B, S_q, D_k) \times (B, D_k, S_k) = (B, S_q, S_k)$.
  \item Division by $\sqrt{d_k}$: scalar operation, shape preserved.
  \item \py{torch.softmax}: element-wise, shape preserved.
  \item \py{torch.matmul(weights, value)}: $(B, S_q, S_k) \times
        (B, S_k, D_v) = (B, S_q, D_v)$.
\end{itemize}

Each step is within the decidable shape theory (linear arithmetic on
dimension sizes).

\section{Gradient correctness}

\subsection{Proving that gradients are well-defined}

The system can verify that gradient computation does not encounter
singularities:

\begin{lstlisting}[style=deppy,language=Python]
from deppy import requires, ensures, have

@requires(lambda x: x.requires_grad)
@requires(lambda x: x.dim() == 2)
@ensures(lambda x, result:
    result.shape == x.shape)
@ensures(lambda x, result:
    not torch.isnan(result).any())
def safe_softmax_cross_entropy(
    x: torch.Tensor,
    targets: torch.Tensor,
) -> torch.Tensor:
    # log-sum-exp trick for numerical stability
    max_x = x.max(dim=1, keepdim=True).values
    have("max_x.shape == (x.shape[0], 1)")
    have("no NaN: max of finite values is finite")

    shifted = x - max_x
    have("shifted.max(dim=1) == 0 for each row")
    have("shifted values are <= 0")

    exp_shifted = torch.exp(shifted)
    have("exp values are in (0, 1] since shifted <= 0")
    have("no overflow: exp(0) = 1 is the maximum")

    sum_exp = exp_shifted.sum(dim=1, keepdim=True)
    have("sum_exp > 0: sum of positive values")
    have("sum_exp >= 1: at least one exp(0) = 1 term")

    log_softmax = shifted - torch.log(sum_exp)
    have("log(sum_exp) is finite: sum_exp >= 1 > 0")
    have("no NaN in log_softmax")

    loss = -log_softmax.gather(1, targets.unsqueeze(1))
    have("loss is finite: log_softmax is finite")

    return loss.mean()
\end{lstlisting}

The proof traces numeric safety through the log-sum-exp computation:
\begin{enumerate}[nosep]
  \item Subtracting the max ensures all shifted values are $\leq 0$.
  \item $\exp(\text{shifted}) \in (0, 1]$: no overflow.
  \item The sum is $\geq 1$ (at least one term equals $\exp(0) = 1$): no
        log-of-zero.
  \item The result is finite throughout: no NaN propagation.
\end{enumerate}

This is a real-world pattern (the log-sum-exp trick) that is numerically
stable but whose stability is non-obvious. The proof mode makes the reasoning
explicit and machine-checked.

\section{Training loop correctness}

\subsection{Proving that a training loop maintains model invariants}

\begin{lstlisting}[style=deppy,language=Python]
from deppy import requires, ensures, have, invariant, old
import torch
import torch.nn as nn

@requires(lambda model, train_loader, epochs:
    epochs > 0 and len(train_loader) > 0)
@ensures(lambda model, train_loader, epochs, result:
    len(result) == epochs)
@ensures(lambda model, train_loader, epochs, result:
    forall(result, lambda loss: isinstance(loss, float)))
def train(
    model: nn.Module,
    train_loader,
    epochs: int,
    lr: float = 0.01,
) -> list[float]:
    optimizer = torch.optim.SGD(model.parameters(), lr=lr)
    criterion = nn.MSELoss()
    epoch_losses = []

    for epoch in range(epochs):
        with invariant():
            have("len(epoch_losses) == epoch")
            have("forall(epoch_losses, "
                 "lambda l: isinstance(l, float))")

        batch_losses = []
        for batch_x, batch_y in train_loader:
            pred = model(batch_x)
            have("pred.shape[0] == batch_x.shape[0]")

            loss = criterion(pred, batch_y)
            have("loss is a scalar tensor")
            have("loss.requires_grad")

            optimizer.zero_grad()
            loss.backward()
            have("all model.parameters() have .grad set")

            optimizer.step()
            batch_losses.append(loss.item())

        epoch_loss = sum(batch_losses) / len(batch_losses)
        have("len(batch_losses) > 0")  # train_loader nonempty
        have("isinstance(epoch_loss, float)")
        epoch_losses.append(epoch_loss)

    have("len(epoch_losses) == epochs")
    return epoch_losses
\end{lstlisting}

The proof verifies:
\begin{itemize}[nosep]
  \item The epoch losses list grows by exactly one per epoch.
  \item Each epoch loss is a float (not a tensor, not None).
  \item The division by \py{len(batch_losses)} is safe (the loader is
        nonempty).
  \item The model's output shape is compatible with the loss function.
\end{itemize}

\section{Model architecture validation}

\subsection{Proving that a ResNet skip connection is shape-compatible}

\begin{lstlisting}[style=deppy,language=Python]
from deppy import requires, ensures, invariant_of, have

@invariant_of(lambda self:
    self.conv1.in_channels == self.in_channels and
    self.conv2.out_channels == self.out_channels and
    (self.shortcut is None and
     self.in_channels == self.out_channels) or
    (self.shortcut is not None and
     self.shortcut.in_channels == self.in_channels and
     self.shortcut.out_channels == self.out_channels))
class ResidualBlock(nn.Module):
    def __init__(self, in_channels: int, out_channels: int):
        super().__init__()
        self.in_channels = in_channels
        self.out_channels = out_channels
        self.conv1 = nn.Conv2d(in_channels, out_channels, 3,
                               padding=1)
        self.conv2 = nn.Conv2d(out_channels, out_channels, 3,
                               padding=1)
        self.shortcut = None
        if in_channels != out_channels:
            self.shortcut = nn.Conv2d(in_channels, out_channels,
                                      1)

    @requires(lambda self, x:
        x.dim() == 4 and x.shape[1] == self.in_channels)
    @ensures(lambda self, x, result:
        result.shape == (x.shape[0], self.out_channels,
                         x.shape[2], x.shape[3]))
    def forward(self, x: torch.Tensor) -> torch.Tensor:
        identity = x

        out = self.conv1(x)
        have("out.shape[1] == self.out_channels")
        out = torch.relu(out)
        out = self.conv2(out)
        have("out.shape[1] == self.out_channels")

        if self.shortcut is not None:
            identity = self.shortcut(x)
            have("identity.shape[1] == self.out_channels")
        else:
            have("self.in_channels == self.out_channels")
            have("identity.shape[1] == self.out_channels")

        # The addition is shape-safe
        have("out.shape == identity.shape")
        out = out + identity
        return torch.relu(out)
\end{lstlisting}

The class invariant ensures that the shortcut connection exists exactly when
the input and output channel counts differ. The forward method proof shows
that the skip connection's output shape always matches the main path's output
shape, so the addition is safe.

\chapter{Proof Mode for Real-World Python: Effects, I/O, and the Trusted Computing Base}

Real Python programs are not pure functions. They read files, make network
requests, mutate global state, and interact with databases. This chapter
describes how proof mode handles effectful computation while maintaining
soundness.

\section{The effect discipline}

The system tracks effects using a lightweight effect system embedded in the
section structure. Each site carries an effect annotation that classifies the
operations at that site:

\begin{itemize}[nosep]
  \item \textbf{Pure}: no side effects; the result depends only on the inputs.
  \item \textbf{Read}: reads from external state (files, databases, network)
        but does not modify it.
  \item \textbf{Write}: modifies external state.
  \item \textbf{Alloc}: allocates a new resource (file handle, socket,
        database connection).
  \item \textbf{Free}: releases a resource.
  \item \textbf{Raise}: may raise an exception.
  \item \textbf{Diverge}: may not terminate.
\end{itemize}

Effects compose: a function that calls both a Read function and a Write
function has both Read and Write effects. The system infers effects
automatically; the programmer does not need to annotate them. However, the
programmer can constrain effects:

\begin{lstlisting}[style=deppy,language=Python]
from deppy import pure

@pure
def compute_hash(data: bytes) -> int:
    """Must be a pure function (no side effects)."""
    h = 0
    for b in data:
        h = (h * 31 + b) % (2**64)
    return h
\end{lstlisting}

The \py{@pure} decorator asserts that the function has no effects. The system
verifies this by checking that no operation in the function's body has a
Write, Alloc, Free, or Raise effect.

\section{Resource tracking and linear types}

Resources (file handles, database connections, locks) must be acquired and
released exactly once. The proof mode supports linear-type-style tracking:

\begin{lstlisting}[style=deppy,language=Python]
from deppy import requires, ensures, owned, consumed

@ensures(lambda path, result: owned(result))
def open_file(path: str) -> "FileHandle":
    return open(path, 'r')

@requires(lambda handle: owned(handle))
@ensures(lambda handle, result: consumed(handle))
def close_file(handle: "FileHandle") -> None:
    handle.close()

@requires(lambda handle: owned(handle))
@ensures(lambda handle, result: owned(handle))
def read_line(handle: "FileHandle") -> str:
    return handle.readline()
\end{lstlisting}

The \py{owned(handle)} predicate asserts that the resource is currently owned
(open, valid, not yet released). The \py{consumed(handle)} predicate asserts
that the resource has been released. The system tracks ownership through the
section structure:

\begin{itemize}[nosep]
  \item Opening a resource establishes ownership.
  \item Using a resource requires ownership.
  \item Closing a resource consumes ownership.
  \item Using a consumed resource creates an error site (use-after-close).
  \item Failing to close a resource before the owning scope exits creates a
        resource-leak warning.
\end{itemize}

\paragraph{Context managers as ownership boundaries.}
Python's \py{with} statement provides automatic resource management:
\begin{lstlisting}[style=deppy,language=Python]
with open("data.txt") as f:
    # f is owned here
    data = f.read()
# f is consumed here (automatically closed)
\end{lstlisting}

The system recognizes \py{with} as an ownership boundary: the resource is
owned inside the block and consumed at the block exit. This means that
well-structured code using context managers gets resource safety for free.

\section{State-dependent types for mutable objects}

Mutable objects have state that changes over time. The proof mode tracks
object state through \emph{state-dependent types}: the type of an object
includes its current state.

\begin{lstlisting}[style=deppy,language=Python]
from deppy import state, transition, requires, ensures

class StateMachine:
    @state("idle")
    def __init__(self):
        self.data = None

    @requires(lambda self: self.in_state("idle"))
    @transition("idle", "running")
    def start(self, data):
        self.data = data

    @requires(lambda self: self.in_state("running"))
    @transition("running", "done")
    def stop(self) -> object:
        result = self.data
        self.data = None
        return result

    @requires(lambda self: self.in_state("done"))
    @transition("done", "idle")
    def reset(self):
        pass
\end{lstlisting}

The \py{@state} and \py{@transition} decorators define a state machine.
The system verifies that:
\begin{itemize}[nosep]
  \item Every method is called only in its required state.
  \item State transitions are valid (only declared transitions are allowed).
  \item The state machine is used correctly at every call site.
\end{itemize}

This pattern is useful for protocol implementations, connection managers,
and any object whose methods must be called in a specific order.

\section{Proving properties across module boundaries}

In large codebases, functions are spread across modules. The proof mode
supports cross-module verification through \emph{module contracts}:

\begin{lstlisting}[style=deppy,language=Python]
# module: data_loader.py
from deppy import module_ensures

@module_ensures("""
    All functions in this module that return DataFrames
    guarantee that the DataFrame has the columns specified
    in the function's @ensures annotation.
    All functions handle missing files by raising FileNotFoundError.
""")

@ensures(lambda path, result:
    set(result.columns) >= {"id", "name", "value"})
@ensures(lambda path, result: len(result) > 0)
def load_data(path: str) -> "pd.DataFrame":
    ...
\end{lstlisting}

The \py{@module_ensures} decorator is a documentation-level annotation that
describes the module's guarantees. Individual function contracts are checked
mechanically; the module contract provides context for human readers.

When function \py{g} in module B calls function \py{f} in module A, the
system uses \py{f}'s contract as an interface. If \py{f} is fully verified,
the contract is trusted. If \py{f} is only inferred (Level 0), the contract
is the inferred section, which may be weaker. The system reports the trust
level at each cross-module boundary.

\section{The trusted computing base}

The \emph{trusted computing base} (TCB) is the set of components whose
correctness the system assumes without proof:

\begin{enumerate}[nosep]
  \item \textbf{The Python interpreter.} The system assumes that Python
        executes code according to the language specification.
  \item \textbf{The SMT solver (Z3).} The system assumes that Z3 returns
        correct results.
  \item \textbf{Library axioms.} The system assumes that library theory
        packs correctly model their libraries (e.g., that \py{len(xs)}
        returns a non-negative integer, that \py{torch.matmul} follows
        matrix multiplication shape rules).
  \item \textbf{The proof checker.} The system assumes that the proof
        checker correctly verifies proof terms.
  \item \textbf{Assumed axioms.} Properties marked with \py{@assume} are
        trusted without proof.
\end{enumerate}

The TCB is explicitly documented and minimized:
\begin{itemize}[nosep]
  \item Library axioms are generated from specifications, not from
        implementation details, and can be tested against the actual library.
  \item The proof checker is small (a few hundred lines) and can be
        independently audited.
  \item \py{@assume} usage is tracked and reported, so the programmer knows
        exactly what is trusted.
\end{itemize}

\section{Soundness and the escape hatches}

The system provides several escape hatches that compromise soundness in
exchange for practicality. Each escape hatch is tracked:

\paragraph{\py{@assume(P)}.}
Asserts that $P$ is true without proof. The system treats $P$ as an axiom.
If $P$ is false, the system may derive incorrect results.

\paragraph{\py{sorry()}.}
Marks a proof obligation as deferred. The system proceeds as if the
obligation were proved. All \py{sorry()} uses are reported in the final
verification summary.

\paragraph{\py{@unchecked}.}
Marks a function as unchecked. The system does not analyze the function's
body. The function's contract (if any) is trusted without verification.

\paragraph{\py{cast(T, x)}.}
Asserts that \py{x} has type \py{T} without verification. This is Python's
standard \py{typing.cast}, but the system tracks its use.

The verification summary reports all escape hatch uses:
\begin{lstlisting}[language={}]
VERIFICATION SUMMARY:
  Functions analyzed: 47
  Obligations discharged: 312 / 318
  Obligations deferred (sorry): 4
  Assumed axioms: 2
  Unchecked functions: 1
  Type casts: 3
  Trust level: 96.2% verified
\end{lstlisting}

This gives the programmer a clear picture of how much of their code is
fully verified and where the gaps are.


\chapter{Advanced Proof Techniques: Abstraction, Refinement, and Parametricity}

This chapter covers advanced proof techniques that enable reasoning about
complex software systems while keeping the annotation burden manageable.

\section{Data abstraction and representation invariants}

A key technique for managing proof complexity is \emph{data abstraction}:
proving properties about an abstract interface and then showing that the
concrete implementation satisfies the interface.

\begin{lstlisting}[style=deppy,language=Python]
from deppy import abstract, refines, rep_invariant

@abstract
class AbstractSet:
    """Abstract specification of a set."""
    @ensures(lambda self, result: result >= 0)
    def size(self) -> int: ...

    @ensures(lambda self, x, result:
        isinstance(result, bool))
    def contains(self, x: int) -> bool: ...

    @ensures(lambda self, x:
        self.contains(x) and
        self.size() == old(self.size()) +
        (0 if old(self.contains(x)) else 1))
    def add(self, x: int) -> None: ...

    @requires(lambda self, x: self.contains(x))
    @ensures(lambda self, x:
        not self.contains(x) and
        self.size() == old(self.size()) - 1)
    def remove(self, x: int) -> None: ...

@refines(AbstractSet)
@rep_invariant(lambda self:
    len(self.items) == len(set(self.items)) and
    all(isinstance(x, int) for x in self.items))
class SortedListSet:
    """Concrete implementation using a sorted list."""
    def __init__(self):
        self.items: list[int] = []

    def size(self) -> int:
        return len(self.items)

    def contains(self, x: int) -> bool:
        # binary search
        lo, hi = 0, len(self.items)
        while lo < hi:
            mid = (lo + hi) // 2
            if self.items[mid] == x:
                return True
            elif self.items[mid] < x:
                lo = mid + 1
            else:
                hi = mid
        return False

    def add(self, x: int) -> None:
        if not self.contains(x):
            # insert in sorted position
            pos = 0
            while pos < len(self.items) and self.items[pos] < x:
                pos += 1
            self.items.insert(pos, x)

    def remove(self, x: int) -> None:
        # binary search for position
        lo, hi = 0, len(self.items)
        while lo < hi:
            mid = (lo + hi) // 2
            if self.items[mid] == x:
                self.items.pop(mid)
                return
            elif self.items[mid] < x:
                lo = mid + 1
            else:
                hi = mid
\end{lstlisting}

The \py{@refines(AbstractSet)} decorator instructs the system to verify that
\py{SortedListSet} satisfies every contract in \py{AbstractSet}. The
\py{@rep_invariant} specifies the concrete representation invariant (no
duplicates, all elements are integers). The system verifies:

\begin{enumerate}[nosep]
  \item The representation invariant is established by \py{__init__}.
  \item Each method preserves the representation invariant.
  \item Each method satisfies the corresponding abstract contract.
  \item The abstraction function (mapping concrete state to abstract state)
        is well-defined.
\end{enumerate}

\section{Simulation and refinement proofs}

When replacing one implementation with another, the programmer can prove
that the new implementation refines (is at least as correct as) the old one:

\begin{lstlisting}[style=deppy,language=Python]
from deppy import simulates

@simulates(SortedListSet, simulation_relation=lambda old, new:
    set(old.items) == new.table.keys())
class HashSetImpl:
    def __init__(self):
        self.table: dict[int, bool] = {}

    def size(self) -> int:
        return len(self.table)

    def contains(self, x: int) -> bool:
        return x in self.table

    def add(self, x: int) -> None:
        self.table[x] = True

    def remove(self, x: int) -> None:
        del self.table[x]
\end{lstlisting}

The \py{@simulates} decorator asserts that \py{HashSetImpl} simulates
\py{SortedListSet} under the given simulation relation. The system verifies
that for every method, if the simulation relation holds before the call, it
holds after the call, and the observable behavior (return values, exceptions)
is the same.

\section{Parametricity and generic proofs}

Generic functions parameterized by type variables enjoy \emph{parametricity}:
their behavior is uniform across all type instantiations. This gives free
theorems that the system can exploit:

\begin{lstlisting}[style=deppy,language=Python]
from deppy import ensures, forall, TypeVar

T = TypeVar("T")

@ensures(lambda xs, result:
    len(result) == len(xs))
@ensures(lambda xs, result:
    forall(range(len(xs)),
           lambda i: result[i] == xs[len(xs) - 1 - i]))
def reverse(xs: list[T]) -> list[T]:
    return xs[::-1]
\end{lstlisting}

Because \py{reverse} is parametric in \py{T}, the system can derive:
\begin{itemize}[nosep]
  \item \py{reverse(reverse(xs)) == xs} (involution).
  \item \py{len(reverse(xs)) == len(xs)} (length preservation).
  \item For any function \py{f}, \py{map(f, reverse(xs)) == reverse(map(f, xs))}
        (naturality).
\end{itemize}

These ``free theorems'' follow from parametricity and do not require separate
proofs. The system derives them automatically for parametric functions.

\section{Proof by refinement: stepwise development}

The proof mode supports stepwise development, where the programmer starts
with a high-level specification and refines it into an implementation through
a series of verified steps:

\paragraph{Step 1: Specification.}
\begin{lstlisting}[style=deppy,language=Python]
@abstract
@ensures(lambda xs, result: is_sorted(result))
@ensures(lambda xs, result: is_permutation(result, xs))
def sort_spec(xs: list[int]) -> list[int]: ...
\end{lstlisting}

\paragraph{Step 2: Algorithm choice.}
\begin{lstlisting}[style=deppy,language=Python]
@refines(sort_spec)
def merge_sort(xs: list[int]) -> list[int]:
    if len(xs) <= 1:
        return xs[:]
    mid = len(xs) // 2
    return merge(merge_sort(xs[:mid]), merge_sort(xs[mid:]))
\end{lstlisting}

\paragraph{Step 3: Implementation details.}
\begin{lstlisting}[style=deppy,language=Python]
@refines(merge_sort)
def merge_sort_inplace(xs: list[int]) -> list[int]:
    """In-place merge sort with auxiliary buffer."""
    ...
\end{lstlisting}

Each refinement step is verified: the refined version satisfies the same
contract as the original, plus any additional properties the refined version
establishes.

\section{Proof composition and modular verification}

The system supports modular verification: each module is verified
independently, and cross-module verification uses contracts as interfaces.
This means that:

\begin{itemize}[nosep]
  \item Changing a function's implementation (without changing its contract)
        does not require re-verifying callers.
  \item Adding a new function does not affect existing verification results.
  \item Removing a function invalidates only code that calls it.
\end{itemize}

\subsection{Assume-guarantee reasoning}

For mutually-recursive modules, the system supports assume-guarantee
reasoning:

\begin{lstlisting}[style=deppy,language=Python]
# module A assumes module B's contract
from deppy import assume_module

@assume_module("B", contract={
    "process": {
        "requires": lambda x: x > 0,
        "ensures": lambda x, r: r >= 0,
    }
})
def compute(x: int) -> int:
    from B import process
    return process(x) + 1
\end{lstlisting}

The \py{@assume_module} decorator declares what module A assumes about
module B. The system verifies A under this assumption. Separately, the
system verifies that B actually satisfies the assumed contract. If both
verifications succeed, the combined system is correct.

\section{Proof of absence of common vulnerability patterns}

The proof mode can verify the absence of common security vulnerabilities:

\paragraph{SQL injection.}
\begin{lstlisting}[style=deppy,language=Python]
from deppy import requires, Refine, sanitized

@requires(lambda query: sanitized(query, "sql"))
def execute_query(query: str) -> list:
    cursor.execute(query)
    return cursor.fetchall()
\end{lstlisting}

The \py{sanitized(query, "sql")} predicate asserts that the query has been
sanitized against SQL injection. The system tracks sanitization status
through the program and creates error sites where unsanitized input reaches
\py{execute_query}.

\paragraph{Path traversal.}
\begin{lstlisting}[style=deppy,language=Python]
from deppy import requires

@requires(lambda path: ".." not in path.parts)
@requires(lambda path: path.is_relative_to(BASE_DIR))
def read_user_file(path: "Path") -> str:
    return path.read_text()
\end{lstlisting}

The precondition ensures that the path does not escape the base directory.
The system propagates this requirement to all callers.

\paragraph{Integer overflow in serialization.}
\begin{lstlisting}[style=deppy,language=Python]
from deppy import requires, Refine

@requires(lambda size: 0 <= size <= MAX_BUFFER_SIZE)
def allocate_buffer(
    size: Refine[int, lambda s: 0 <= s <= MAX_BUFFER_SIZE]
) -> bytes:
    return bytes(size)
\end{lstlisting}

The bounded-size requirement prevents denial-of-service via oversized
allocations.


\chapter{The Annotation Economy: Maximizing Proof Value with Minimal Effort}

This chapter addresses the central question of proof mode: how should a
programmer invest annotation effort to maximize the verification benefit?

\section{The information-theoretic view}

From an information-theoretic perspective, an annotation provides information
that resolves uncertainty in the section structure. The ``value'' of an
annotation is the number of obligations it helps discharge, weighted by the
importance of those obligations.

\paragraph{High-value annotations.}
\begin{itemize}[nosep]
  \item \textbf{Loop invariants}: A single loop invariant can discharge
        dozens of obligations (one per iteration, one for the postcondition,
        several for callers).
  \item \textbf{Library contracts}: A contract on a widely-used utility
        function propagates to every caller.
  \item \textbf{Class invariants}: A class invariant enriches every method's
        section.
  \item \textbf{Key lemmas}: A lemma about a core data structure
        enables verification of all code that uses the structure.
\end{itemize}

\paragraph{Low-value annotations.}
\begin{itemize}[nosep]
  \item \textbf{Annotations the inference engine already knows}: If the
        system already inferred that \py{len(xs) > 0}, annotating it
        provides zero additional information.
  \item \textbf{Annotations in dead code}: Verifying unreachable paths
        provides no runtime safety benefit.
  \item \textbf{Redundant postconditions}: If the postcondition follows
        trivially from the body, the annotation is documentation rather
        than verification.
\end{itemize}

\section{The annotation advisor}

The system includes an \emph{annotation advisor} that suggests where
annotation effort should be invested:

\begin{lstlisting}[language={}]
ANNOTATION ADVISOR REPORT:
  High-impact opportunities:
    1. Loop invariant at sort.py:45 (would discharge 8 obligations)
    2. Contract for utils.validate() (called 23 times)
    3. Class invariant for Database (5 methods, 12 obligations)

  Already well-covered:
    - math_utils.py: 94% of obligations auto-discharged
    - config.py: pure functions, fully inferred

  Residual clusters:
    - 6 residuals in parser.py relate to grammar validity
      (consider a grammar-correctness lemma)
    - 4 residuals in network.py relate to connection state
      (consider a state machine annotation)
\end{lstlisting}

The advisor analyzes:
\begin{enumerate}[nosep]
  \item \textbf{Obligation density}: functions with many undischarged
        obligations are candidates for annotation.
  \item \textbf{Call frequency}: functions called from many sites amplify the
        benefit of a contract.
  \item \textbf{Residual clustering}: when multiple residuals relate to the
        same underlying property, a single lemma can discharge all of them.
  \item \textbf{Inference gaps}: sites where the inferred section is
        significantly weaker than what the code actually guarantees.
\end{enumerate}

\section{Progressive verification: a worked example}

Consider a small project with 500 lines of Python across 5 modules. We walk
through the progressive verification workflow:

\paragraph{Phase 1: Zero-annotation baseline.}
The system analyzes all 500 lines and produces:
\begin{itemize}[nosep]
  \item 180 obligations auto-discharged (arithmetic, bounds, types).
  \item 42 residuals reported.
  \item Section quality: 81\% of sites have useful inferred sections.
\end{itemize}

\paragraph{Phase 2: Add contracts to the core utility module (8 functions).}
\begin{itemize}[nosep]
  \item 8 \py{@requires}/\py{@ensures} annotations added (24 lines).
  \item 15 additional obligations discharged (callers of the utilities).
  \item Residuals reduced to 27.
\end{itemize}

\paragraph{Phase 3: Add loop invariants to the main algorithm (2 loops).}
\begin{itemize}[nosep]
  \item 2 loop invariants added (6 lines).
  \item 9 additional obligations discharged.
  \item Residuals reduced to 18.
\end{itemize}

\paragraph{Phase 4: Add a class invariant to the central data structure.}
\begin{itemize}[nosep]
  \item 1 class invariant added (3 lines).
  \item 7 additional obligations discharged.
  \item Residuals reduced to 11.
\end{itemize}

\paragraph{Phase 5: Add 2 key lemmas.}
\begin{itemize}[nosep]
  \item 2 lemmas stated and proved (12 lines).
  \item 6 additional obligations discharged.
  \item Residuals reduced to 5.
\end{itemize}

\paragraph{Phase 6: Accept remaining residuals.}
The 5 remaining residuals involve properties about external data (user
input, file contents) that cannot be verified statically. The programmer
marks them as runtime-checked.

\paragraph{Summary.}
\begin{center}
\begin{tabular}{lrrrr}
\toprule
\textbf{Phase} & \textbf{Annotations} & \textbf{Discharged} &
  \textbf{Residuals} & \textbf{Trust} \\
\midrule
1: Baseline & 0 & 180 & 42 & 81\% \\
2: Utility contracts & 24 & 195 & 27 & 88\% \\
3: Loop invariants & 30 & 204 & 18 & 92\% \\
4: Class invariant & 33 & 211 & 11 & 95\% \\
5: Key lemmas & 45 & 217 & 5 & 98\% \\
6: Runtime checks & 45 & 222 & 0 & 100\% \\
\bottomrule
\end{tabular}
\end{center}

Total annotation effort: 45 lines for 500 lines of code (9\% overhead).
The system went from 81\% auto-verified to 100\% (verified + runtime-checked)
with targeted, incremental investment.

\section{When not to prove}

Not every property is worth proving. The proof mode should be used judiciously:

\begin{itemize}[nosep]
  \item \textbf{Don't prove UI layout properties.} UI correctness is visual
        and subjective; formal verification adds no value.
  \item \textbf{Don't prove performance properties.} While the system can
        verify termination, it cannot verify time complexity. Use benchmarks
        instead.
  \item \textbf{Don't prove external behavior.} The system cannot verify that
        a web API returns correct data or that a database query is
        semantically correct. Use integration tests.
  \item \textbf{Don't over-specify.} Over-specified contracts are brittle:
        they break when the implementation changes, even if the change is
        correct. Specify the essential properties and leave the rest to
        inference.
\end{itemize}

The golden rule: \emph{prove what matters, infer what doesn't, and test
what can't be proved.}

%% ======================================================================
%% PART V: ADVANCED SHEAF THEORY, OPTIMIZATION, AND RELEVANCE
%% ======================================================================

\part{Advanced Sheaf Theory, Optimization, and the Relevance Problem}

\chapter{The Full Sheaf-Theoretic Framework for Type Inference}

Parts I--IV developed the sheaf-descent idea informally: types are sections,
local type information glues into global types, and the optimization objective
maximizes information. This chapter makes the mathematics fully precise,
introducing the cohomological machinery that powers automatic deduction and
the obstruction theory that identifies when and why type inference fails.

\section{Sites, covers, and the Grothendieck topology}

\begin{definition}[Program site category]
Let $\mathbf{S}_P$ be the category whose objects are \emph{program sites}
(syntactic locations in a program $P$ where a value is produced, consumed, or
tested) and whose morphisms are \emph{data-flow edges} $s \to t$ indicating
that information at site $s$ is available at site $t$. Composition of
morphisms corresponds to transitive data flow.
\end{definition}

\begin{definition}[Cover]
A \emph{cover} of a program site $s$ is a family $\{u_i \to s\}_{i \in I}$
of data-flow edges such that the type of the value at $s$ is completely
determined by the types at the $u_i$. Formally, the cover satisfies the
\emph{local determination property}: any two sections that agree on all
$u_i$ must agree on $s$.
\end{definition}

\begin{definition}[Grothendieck topology on $\mathbf{S}_P$]
The collection of covers defines a Grothendieck topology $J$ on
$\mathbf{S}_P$. The triple $(\mathbf{S}_P, J)$ is a \emph{site} in the
sense of topos theory. The Grothendieck topology must satisfy:
\begin{enumerate}[nosep]
  \item \textbf{Identity cover.} $\{\mathrm{id}_s : s \to s\}$ is a cover
        of $s$.
  \item \textbf{Stability.} If $\{u_i \to s\}$ is a cover and $t \to s$ is
        a morphism, then $\{u_i \times_s t \to t\}$ is a cover of $t$.
  \item \textbf{Transitivity.} If $\{u_i \to s\}$ is a cover and for each
        $i$, $\{v_{ij} \to u_i\}$ is a cover, then $\{v_{ij} \to s\}$ is a
        cover.
\end{enumerate}
\end{definition}

These axioms ensure that type information composes correctly: if we can
determine a type from its parts, and each part from its sub-parts, then we
can determine the type from the sub-parts.

\section{The type presheaf and the sheaf condition}

\begin{definition}[Type presheaf]
A \emph{type presheaf} on $(\mathbf{S}_P, J)$ is a contravariant functor
\[
  \mathcal{T} : \mathbf{S}_P^{\mathrm{op}} \to \mathbf{Lat}
\]
where $\mathbf{Lat}$ is the category of complete lattices and monotone maps.
For each site $s$, $\mathcal{T}(s)$ is the lattice of possible type
assignments at $s$. For each morphism $f : s \to t$, the \emph{restriction
map} $\mathcal{T}(f) : \mathcal{T}(t) \to \mathcal{T}(s)$ restricts the
type information at $t$ to what is observable at $s$.
\end{definition}

\begin{definition}[Sheaf condition]
The presheaf $\mathcal{T}$ is a \emph{sheaf} if for every cover
$\{u_i \to s\}$ and every family of sections $\sigma_i \in \mathcal{T}(u_i)$
that agree on overlaps (i.e., $\mathcal{T}(p_1)(\sigma_i) =
\mathcal{T}(p_2)(\sigma_j)$ for all overlap morphisms), there exists a
unique section $\sigma \in \mathcal{T}(s)$ that restricts to each
$\sigma_i$.
\end{definition}

\paragraph{Interpretation.}
The sheaf condition says: if local type assignments are consistent on
overlaps, they glue into a unique global type. This is exactly what type
inference should do---combine local evidence into a global type assignment.

When the sheaf condition fails---when local sections do not glue---this
indicates a type error. The \emph{obstruction} to gluing identifies the
nature of the error.

\section{Cohomology of the type sheaf}

The sheaf cohomology groups $H^n(\mathbf{S}_P, \mathcal{T})$ measure the
obstructions to extending local type information to global type assignments.

\begin{definition}[\v{C}ech cohomology]
Given a cover $\mathcal{U} = \{u_i \to s\}$, define the \v{C}ech complex:
\begin{align*}
  C^0(\mathcal{U}, \mathcal{T}) &= \prod_i \mathcal{T}(u_i) \\
  C^1(\mathcal{U}, \mathcal{T}) &= \prod_{i < j} \mathcal{T}(u_i \times_s u_j) \\
  C^2(\mathcal{U}, \mathcal{T}) &= \prod_{i < j < k}
    \mathcal{T}(u_i \times_s u_j \times_s u_k) \\
  &\vdots
\end{align*}
with coboundary maps $\delta^n : C^n \to C^{n+1}$ defined by the alternating
sum of restriction maps.
\end{definition}

\paragraph{$H^0$: Global sections = valid type assignments.}
$H^0(\mathcal{U}, \mathcal{T}) = \ker(\delta^0)$ is the group of global
sections---families of local types that agree on all overlaps. This is
precisely the set of valid type assignments for the program.

\paragraph{$H^1$: Obstructions = type errors.}
$H^1(\mathcal{U}, \mathcal{T}) = \ker(\delta^1) / \mathrm{im}(\delta^0)$
measures the failure of the sheaf condition. A nonzero class in $H^1$
represents a \emph{type error}: local sections that are pairwise consistent
but cannot be glued into a global section. This happens when there is a
``twist''---the local sections form a non-trivial cocycle.

\paragraph{$H^2$: Higher obstructions = structural type incompatibilities.}
$H^2$ measures obstructions to extending partial gluings. In practice, $H^2$
detects situations where three or more type constraints interact in a way
that no pairwise analysis would catch---a higher-order type conflict.

\begin{theorem}[Type inference as cohomology vanishing]
A program $P$ is well-typed if and only if $H^1(\mathbf{S}_P, \mathcal{T}) = 0$.
That is, every collection of pairwise-consistent local type assignments glues
into a global type assignment.
\end{theorem}

\begin{proof}[Proof sketch]
The forward direction: if $P$ is well-typed, then a global type assignment
exists, which restricts to consistent local assignments, so every cocycle is
a coboundary. The reverse direction: if $H^1 = 0$, then every choice of
pairwise-consistent local types glues, so in particular the types inferred by
the local analysis glue into a global assignment.
\end{proof}

\section{Descent data and effective descent}

The gluing process---constructing a global section from local sections---is
formalized by \emph{descent data}.

\begin{definition}[Descent datum]
A \emph{descent datum} for a cover $\{u_i \to s\}$ is a pair
$(\{\sigma_i\}, \{\phi_{ij}\})$ where:
\begin{itemize}[nosep]
  \item $\sigma_i \in \mathcal{T}(u_i)$ is a local section at each cover
        element.
  \item $\phi_{ij} : \mathcal{T}(p_1)(\sigma_i) \xrightarrow{\sim}
        \mathcal{T}(p_2)(\sigma_j)$ is an isomorphism on overlaps.
  \item The cocycle condition holds: $\phi_{jk} \circ \phi_{ij} = \phi_{ik}$
        on triple overlaps.
\end{itemize}
\end{definition}

\begin{definition}[Effective descent]
Descent is \emph{effective} if every descent datum arises from a unique
global section. Equivalently, the natural map
\[
  \mathcal{T}(s) \to \mathrm{Desc}(\mathcal{U}, \mathcal{T})
\]
is an isomorphism, where $\mathrm{Desc}$ is the category of descent data.
\end{definition}

\paragraph{Type inference as effective descent.}
The type inference algorithm constructs descent data (local types with
consistency witnesses on overlaps) and then ``descends'' to a global type.
The effectiveness of descent guarantees that the inferred global type is
unique---there is no ambiguity in the type inference result.

When descent is not effective (because the sheaf condition fails at some
cover), the system reports the failure as a type error, with the cocycle
class in $H^1$ identifying the precise nature of the inconsistency.

\section{Stalks and the local-to-global principle}

\begin{definition}[Stalk]
The \emph{stalk} of $\mathcal{T}$ at a site $s$ is
\[
  \mathcal{T}_s = \varinjlim_{s \in U} \mathcal{T}(U)
\]
the colimit over all covers containing $s$. Informally, the stalk is the
``best possible type'' at $s$, combining all available information.
\end{definition}

\begin{theorem}[Local-to-global for types]
A morphism of sheaves $f : \mathcal{T} \to \mathcal{T}'$ is an isomorphism
if and only if $f_s : \mathcal{T}_s \to \mathcal{T}'_s$ is an isomorphism
at every stalk. That is, type equivalence can be checked locally.
\end{theorem}

This theorem justifies the system's architecture: type checking can be
decomposed into local checks at each site, and if all local checks pass,
the global type assignment is valid. This is the formal foundation for the
compositional verification strategy.

\section{Sheaf morphisms and type transformations}

When a function transforms data (e.g., \py{torch.reshape}), it induces a
morphism between type sheaves. The system must track how these morphisms
interact with sections.

\begin{definition}[Pushforward]
Given a program morphism $f : P \to Q$ (e.g., a function call), the
\emph{pushforward} $f_* \mathcal{T}$ is the sheaf on $\mathbf{S}_Q$ defined
by $(f_* \mathcal{T})(s) = \mathcal{T}(f^{-1}(s))$. This models how the
caller's type information is available at the callee's sites.
\end{definition}

\begin{definition}[Pullback]
The \emph{pullback} $f^* \mathcal{T}'$ is the sheaf on $\mathbf{S}_P$
defined by $(f^* \mathcal{T}')(s) = \mathcal{T}'(f(s))$. This models how the
callee's type information flows back to the caller.
\end{definition}

The adjunction $f^* \dashv f_*$ captures the bidirectional flow of type
information across function boundaries: forward (pushforward) and backward
(pullback) propagation.


\chapter{The Optimization Problem: Maximizing Type Information}

With the sheaf-theoretic framework in place, we can formalize the central
optimization problem: among all valid type assignments (global sections),
find the one that maximizes information content.

\section{The information lattice}

\begin{definition}[Information ordering]
Given two type assignments $\sigma, \tau$ at a site $s$, we say $\sigma$ is
\emph{more informative} than $\tau$ (written $\sigma \sqsupseteq \tau$) if
$\sigma$ refines $\tau$---every fact implied by $\sigma$ is also implied by
$\tau$, and $\sigma$ may imply additional facts.
\end{definition}

For refinement types, the ordering is:
\[
  \{x : A \mid \phi(x) \wedge \psi(x)\} \sqsupseteq \{x : A \mid \phi(x)\}
\]
A more refined type (stronger predicate) is more informative.

\begin{definition}[Information content]
The \emph{information content} of a type assignment $\sigma$ at site $s$ is
a measure $\mathcal{I}(\sigma, s) \in [0, \infty)$ that quantifies how much
the type constrains the possible values. Formally:
\[
  \mathcal{I}(\sigma, s) = -\log \frac{|\{v : v \models \sigma(s)\}|}
    {|\{v : v \models \top\}|}
\]
where $\top$ is the unconstrained type. A more restrictive type has higher
information content.
\end{definition}

In practice, we use an approximation: the number of atomic predicates in the
refinement, weighted by their discriminative power.

\section{The global optimization objective}

\begin{definition}[Maximum-information type inference]
Given a program $P$ with site category $\mathbf{S}_P$ and type sheaf
$\mathcal{T}$, the \emph{maximum-information type inference problem} is:
\[
  \sigma^* = \arg\max_{\sigma \in H^0(\mathbf{S}_P, \mathcal{T})}
    \sum_{s \in \mathbf{S}_P} w(s) \cdot \mathcal{I}(\sigma, s)
\]
where $w(s) \geq 0$ is a weight that reflects the importance of site $s$
(e.g., higher weight for sites involved in potential runtime errors).
\end{definition}

\paragraph{Why maximize?}
A more informative type assignment detects more bugs, provides more useful
IDE feedback, and enables more downstream verification. The trivial type
assignment ($\top$ at every site) is always valid but useless. The system
should find the most informative valid assignment.

\paragraph{Why weighted?}
Not all sites are equally important. Sites where a runtime error could occur
(division, indexing, attribute access) should be weighted more heavily, so the
optimizer prioritizes inferring precise types at those sites.

\section{Reformulation as a constraint satisfaction problem}

The optimization problem can be reformulated as a weighted MaxSAT or
MaxSMT problem:

\paragraph{Hard constraints (must be satisfied).}
\begin{itemize}[nosep]
  \item \textbf{Overlap consistency.} For every overlap $u_i \times_s u_j$,
        the restrictions of $\sigma_i$ and $\sigma_j$ must agree.
  \item \textbf{Semantic soundness.} For every operation site $s$ with
        semantics $\llbracket s \rrbracket$, the output type must be
        consistent with applying $\llbracket s \rrbracket$ to the input
        types.
  \item \textbf{Data-flow consistency.} For every data-flow edge $s \to t$,
        $\sigma(t)$ must be at least as general as $\sigma(s)$ restricted
        to $t$.
\end{itemize}

\paragraph{Soft constraints (maximize satisfaction).}
\begin{itemize}[nosep]
  \item For each site $s$ and each candidate predicate $\phi$ that could
        refine the type at $s$, add a soft clause $\sigma(s) \models \phi$
        with weight $w(s) \cdot \mathcal{I}(\phi)$.
  \item The MaxSMT solver maximizes the total weight of satisfied soft
        clauses.
\end{itemize}

\begin{theorem}[Soundness of MaxSMT formulation]
The MaxSMT solution $\sigma^*$ is a valid type assignment (all hard
constraints are satisfied) that maximizes the weighted information content
among all valid assignments.
\end{theorem}

\section{The lattice structure of solutions}

\begin{theorem}[Lattice of valid type assignments]
The set of valid type assignments $H^0(\mathbf{S}_P, \mathcal{T})$ forms a
complete lattice under the information ordering. The least element is
$\top$ (no information). The greatest element is the \emph{principal type}---
the most informative valid assignment.
\end{theorem}

\begin{proof}[Proof sketch]
The meet of two valid assignments restricts to the common information
(intersection of refinement predicates). The join takes the stronger of the
two at each site and verifies consistency; if both are valid and consistent,
their join is valid. The greatest element exists by Zorn's lemma applied to
the chain of increasingly-informative valid assignments.
\end{proof}

\paragraph{The principal type.}
The principal type is the unique most-informative valid type assignment. If
it exists and is computable, it is the ideal output of the inference engine.
In general, the principal type may not be computable (the predicate language
is too expressive), but within decidable fragments, the MaxSMT solver
converges to it.

\section{Approximation algorithms}

The exact MaxSMT problem may be intractable for large programs. The system
uses several approximation strategies:

\paragraph{Fixpoint iteration.}
Start with the least informative assignment ($\top$ at every site). At each
iteration, for each site $s$, strengthen the type at $s$ by adding predicates
that are implied by the current types at $s$'s predecessors. Repeat until
convergence (fixpoint). This is equivalent to abstract interpretation and
converges monotonically to a valid type assignment.

\paragraph{Widening.}
When fixpoint iteration is slow (large loops, deep recursion), apply widening
operators that sacrifice some precision for faster convergence. The widening
point is at loop headers: after a fixed number of iterations, generalize the
loop-header type to an over-approximation.

\paragraph{Stratified solving.}
Decompose the program into strongly-connected components (SCCs) of the
data-flow graph. Solve each SCC independently (from leaves to roots). Within
each SCC, use fixpoint iteration. This reduces the problem size and exploits
the acyclic structure of most programs.

\paragraph{Demand-driven inference.}
Instead of inferring types at all sites, infer types only at sites that are
``demanded''---sites involved in potential errors, sites with user
annotations, or sites explicitly queried. This reduces work for large programs
where most sites are uninteresting.

\section{Duality: the backward optimization problem}

The forward problem maximizes output information given inputs. The
\emph{backward} or \emph{dual} problem minimizes input requirements while
ensuring output properties:

\[
  \rho^* = \arg\min_{\rho \in \mathcal{T}(\text{input})}
    \mathcal{I}(\rho) \quad \text{subject to} \quad
    \forall s \in \text{error-sites}: \mathrm{viable}(\sigma^*_s, \rho)
\]

This finds the \emph{weakest precondition}---the least restrictive input
requirement that ensures all error sites are viable (no runtime errors).

\begin{theorem}[Duality]
The forward maximum-information problem and the backward
weakest-precondition problem are Lagrangian duals. The optimal Lagrange
multipliers correspond to the ``importance weights'' at each error site.
\end{theorem}

This duality is the formal justification for the bidirectional synthesis
described in Part I: the forward and backward passes are two views of the
same optimization problem.


\chapter{The Relevance Problem: Which Propositions Matter for Bug Detection?}

A dependent type system can decorate every term with an enormous number of
propositions: the integer 42 satisfies $x > 0$, $x < 100$, $x$ is even,
$x \equiv 0 \pmod{6}$, and infinitely many others. Most of these
propositions are irrelevant to detecting bugs in the surrounding code. This
chapter formalizes the \emph{relevance problem}: determining which
propositions are worth tracking.

\section{The relevance criterion}

\begin{definition}[Bug-relevant proposition]
A proposition $\phi$ at site $s$ is \emph{bug-relevant} if there exists a
reachable error site $e$ such that knowing $\phi$ at $s$ changes the
viability analysis at $e$. Formally, $\phi$ is bug-relevant if:
\[
  \exists e \in \mathrm{ErrorSites}(P) : \quad
  \mathrm{viable}(e \mid \sigma \cup \{\phi\}) \neq
  \mathrm{viable}(e \mid \sigma)
\]
where $\sigma$ is the current section without $\phi$.
\end{definition}

Intuitively, a proposition is bug-relevant if adding it to the type
information either (a) allows the system to prove that an error cannot occur
(the error site becomes viable), or (b) allows the system to prove that an
error must occur (the error site becomes non-viable and should be reported).

\paragraph{The relevance filter.}
The system should track only bug-relevant propositions. This dramatically
reduces the number of predicates in the MaxSMT problem and focuses the
analysis on what matters.

\section{Computing relevance: the backward slice}

Bug-relevance can be computed by a backward analysis from error sites:

\begin{definition}[Error-site backward slice]
The \emph{backward slice} from an error site $e$ is the set of sites
$s$ such that the type at $s$ can influence the viability predicate at $e$.
Formally:
\[
  \mathrm{Slice}(e) = \{s \in \mathbf{S}_P \mid
    \exists \text{data-flow path } s \leadsto e\}
\]
\end{definition}

A proposition $\phi$ at site $s$ can be bug-relevant only if $s \in
\mathrm{Slice}(e)$ for some error site $e$. This is a necessary but not
sufficient condition.

\begin{definition}[Predicate relevance]
A predicate $\phi$ over the value at site $s$ is \emph{relevant to error
site $e$} if:
\begin{enumerate}[nosep]
  \item $s \in \mathrm{Slice}(e)$.
  \item The viability predicate at $e$, when pulled back to $s$ through the
        data-flow path, involves the variable constrained by $\phi$.
  \item The decision procedure for the viability predicate at $e$ can use
        $\phi$ (i.e., $\phi$ is in a compatible theory).
\end{enumerate}
\end{definition}

\section{The predicate abstraction framework}

Rather than tracking arbitrary propositions, the system operates over a
finite \emph{predicate abstraction}: a finite set of predicates that are
sufficient to capture the relevant type information.

\begin{definition}[Predicate domain]
A \emph{predicate domain} $\Pi$ for a program $P$ is a finite set of
predicates $\{\phi_1, \ldots, \phi_n\}$ over the program's variables. The
\emph{abstract type} at site $s$ is a Boolean assignment to the predicates
in $\Pi$ that are defined at $s$.
\end{definition}

\begin{definition}[Sufficient predicate domain]
$\Pi$ is \emph{sufficient for bug detection} if for every error site $e$ in
$P$, the viability predicate at $e$ can be expressed as a Boolean
combination of predicates in $\Pi$ (possibly pulled back through the
data flow).
\end{definition}

\begin{theorem}[Existence of finite sufficient predicate domains]
If the viability predicates at all error sites are in a decidable theory
$\mathcal{L}$ with finite quantifier elimination, then a finite sufficient
predicate domain exists.
\end{theorem}

\begin{proof}[Proof sketch]
Each viability predicate involves a finite number of variables and constants.
Quantifier elimination for $\mathcal{L}$ produces a finite set of atomic
predicates. The union of these sets across all error sites is the sufficient
domain.
\end{proof}

\section{Automatic predicate discovery}

The system discovers predicates automatically through several mechanisms:

\paragraph{Error-site-driven discovery.}
For each error site $e$ with viability predicate $V_e$, the system:
\begin{enumerate}[nosep]
  \item Identifies the free variables of $V_e$.
  \item Traces each variable back through the data flow to find its origins.
  \item At each origin site, creates predicates corresponding to the atoms
        of $V_e$ specialized to the origin's context.
\end{enumerate}

For example, if the viability predicate at an indexing site is $0 \leq i <
\mathrm{len}(xs)$, the system traces $i$ and $\mathrm{len}(xs)$ back
through the data flow and creates predicates $i \geq 0$,
$i < \mathrm{len}(xs)$, and $\mathrm{len}(xs) > 0$ at the appropriate
upstream sites.

\paragraph{Guard-derived predicates.}
Every conditional branch (\py{if}, \py{while}, \py{assert}) introduces a
predicate. The system adds the guard's predicate (and its negation) to the
predicate domain.

\paragraph{Annotation-derived predicates.}
Every user annotation (\py{@requires}, \py{@ensures}, \py{@invariant_of})
introduces predicates. The system adds the annotation's predicates to the
domain.

\paragraph{Theory-derived predicates.}
Library theories (PyTorch shapes, Pandas columns) introduce domain-specific
predicates. For example, the PyTorch shape theory adds dimension-equality
and positivity predicates for all tensor-valued sites.

\section{Predicate domain refinement}

The initial predicate domain may be too coarse (missing predicates that are
needed) or too fine (containing predicates that are irrelevant). The system
refines the domain through a CEGAR-like (counterexample-guided abstraction
refinement) loop:

\paragraph{Step 1: Solve with current domain.}
Run the MaxSMT inference with the current predicate domain $\Pi$.

\paragraph{Step 2: Check residuals.}
For each residual obligation (an obligation that could not be discharged),
check whether adding new predicates to $\Pi$ would help.

\paragraph{Step 3: Refine.}
If a residual's viability predicate involves a sub-expression not represented
in $\Pi$, add the sub-expression as a new predicate. Re-run from Step 1.

\paragraph{Step 4: Coarsen.}
After convergence, remove predicates from $\Pi$ that were never used in
any discharged obligation. This keeps the domain minimal.

\begin{theorem}[CEGAR termination]
If the predicate theory has finite quantifier elimination, the CEGAR loop
terminates after finitely many refinement steps.
\end{theorem}

\section{Relevance scoring}

Not all bug-relevant predicates are equally important. The system assigns a
\emph{relevance score} to each predicate based on:

\begin{definition}[Relevance score]
The relevance score of predicate $\phi$ at site $s$ is:
\[
  R(\phi, s) = \sum_{e \in \mathrm{ErrorSites}(P)} w(e) \cdot
    \mathbb{1}[\phi \text{ contributes to discharging viability at } e]
\]
where $w(e)$ is the error site's importance weight.
\end{definition}

Predicates with high relevance scores should be tracked with high priority.
Predicates with zero relevance are irrelevant and can be dropped.

\paragraph{Importance weighting of error sites.}
The weights $w(e)$ can be set by:
\begin{itemize}[nosep]
  \item \textbf{Severity}: runtime crashes (e.g., division by zero) are
        weighted higher than warnings.
  \item \textbf{Reachability}: error sites on hot paths are weighted higher.
  \item \textbf{User priority}: the programmer can explicitly flag certain
        error sites as high-priority.
  \item \textbf{Historical data}: error sites that have caused bugs before
        are weighted higher.
\end{itemize}

\chapter{User-Guided Relevance and the Interactive Query Interface}

The automatic relevance analysis from the previous chapter determines which
propositions matter for bug detection. But users have domain knowledge that
the system lacks: they know which properties are important to their
application, which invariants they suspect might be violated, and which
components are most critical. This chapter describes the interactive
relevance interface.

\section{The \py{@track} annotation: user-declared relevance}

The simplest form of user guidance is declaring that a specific proposition
should be tracked, even if the automatic analysis does not identify it as
bug-relevant:

\begin{lstlisting}[style=deppy,language=Python]
from deppy import track

@track(lambda self: self.balance >= 0, name="non_negative_balance")
class BankAccount:
    def __init__(self, initial: float):
        self.balance = initial

    def withdraw(self, amount: float):
        self.balance -= amount

    def deposit(self, amount: float):
        self.balance += amount
\end{lstlisting}

The \py{@track} decorator instructs the system to:
\begin{enumerate}[nosep]
  \item Add the predicate $\py{self.balance} \geq 0$ to the predicate domain.
  \item Track this predicate through all methods of \py{BankAccount}.
  \item Report any method that might violate the predicate.
  \item Report any call sequence that might lead to a violation.
\end{enumerate}

The key difference from \py{@invariant_of} is that \py{@track} does not
assert that the property holds---it merely instructs the system to track
whether it holds. The system reports both confirmations (``this method
preserves non-negative balance when called with amount $\leq$ balance'') and
violations (``\py{withdraw} can make balance negative when
amount $>$ balance'').

\section{Property queries: asking the system what it knows}

The system provides a query interface where the user can ask about specific
properties:

\begin{lstlisting}[style=deppy,language=Python]
from deppy import query

# Does this function always return a sorted list?
query(merge_sort, "is_sorted(result)")
# -> VERIFIED: result is always sorted (given sorted inputs to merge)

# Can this function raise IndexError?
query(binary_search, "raises(IndexError)")
# -> SAFE: no IndexError possible (given len(xs) > 0)

# What is the shape of the output tensor?
query(attention, "result.shape", given={"query.shape": (B, S, D)})
# -> result.shape == (B, S, D_v) where D_v = value.shape[2]

# What preconditions does this function need?
query(train_classifier, "preconditions")
# -> features.dim() == 2, labels.dim() == 1,
#    features.shape[0] == labels.shape[0], ...
\end{lstlisting}

The query interface works by:
\begin{enumerate}[nosep]
  \item Parsing the query into a proposition or a query type (property check,
        error reachability, shape inference, precondition synthesis).
  \item Adding the query's predicates to the predicate domain if not already
        present.
  \item Running the inference/verification engine with the augmented domain.
  \item Reporting the result: verified, violated (with counterexample), or
        unknown (with explanation).
\end{enumerate}

\section{The exploration mode: discovering unknown properties}

Beyond checking specific properties, the system can \emph{discover} properties
that the user might not have thought to ask about:

\begin{lstlisting}[style=deppy,language=Python]
from deppy import explore

# What interesting properties does this function have?
explore(matrix_multiply)
# -> Discovered properties:
#    1. result.shape == (A.shape[0], B.shape[1])  [shape]
#    2. result[i][j] == sum(A[i][k]*B[k][j] for k) [value]
#    3. A.shape[1] == B.shape[0] is required      [precond]
#    4. No division by zero possible                [safety]
#    5. result is None-free if inputs are None-free [null]

# What could go wrong in this function?
explore(parse_config, mode="risks")
# -> Potential risks:
#    1. KeyError at line 12 if config missing "host"
#    2. ValueError at line 15 if port is not numeric
#    3. FileNotFoundError at line 8 if path invalid
#    4. TypeError at line 20 if timeout is not int/float
\end{lstlisting}

The exploration mode works by:
\begin{enumerate}[nosep]
  \item Generating candidate predicates from the function's operations
        (shapes, bounds, nullability, error conditions).
  \item Running the inference engine to determine which predicates hold.
  \item Ranking the discovered properties by relevance score.
  \item Presenting the top-$k$ properties to the user.
\end{enumerate}

\section{Relevance propagation across module boundaries}

When a function $f$ in module $A$ calls a function $g$ in module $B$, the
relevance of predicates in $g$ depends on what $f$ needs:

\begin{definition}[Cross-module relevance]
A predicate $\phi$ in module $B$ is \emph{relevant to module $A$} if:
\begin{enumerate}[nosep]
  \item There exists a call from $A$ to a function $g$ in $B$.
  \item $\phi$ appears in $g$'s postcondition or return-type section.
  \item $\phi$ is bug-relevant to some error site in $A$ (when pulled back
        through the call).
\end{enumerate}
\end{definition}

\paragraph{Demand-driven cross-module analysis.}
The system propagates relevance demands across module boundaries:
\begin{enumerate}[nosep]
  \item Module $A$ determines that it needs a specific property of $g$'s
        return value (e.g., the output tensor's shape).
  \item This demand is communicated to module $B$'s analysis.
  \item Module $B$ checks whether $g$ satisfies the demanded property.
  \item If yes, the property is exported as part of $g$'s section.
  \item If no, the failure is reported at the call site in $A$.
\end{enumerate}

This demand-driven approach avoids analyzing all possible properties of every
function in every module. Instead, only properties that are actually needed
by callers are analyzed.

\section{The relevance graph}

The system maintains a \emph{relevance graph} that tracks the flow of
relevance through the program:

\begin{definition}[Relevance graph]
The relevance graph $G_R = (V, E)$ has:
\begin{itemize}[nosep]
  \item Vertices $V$: predicate instances $(\phi, s)$ where $\phi$ is a
        predicate and $s$ is a site.
  \item Edges $E$: a directed edge $(\phi, s) \to (\psi, t)$ indicates
        that knowing $\phi$ at $s$ contributes to determining $\psi$ at $t$.
\end{itemize}
\end{definition}

The relevance graph supports several analyses:
\begin{itemize}[nosep]
  \item \textbf{Root cause analysis.} When a verification fails at site $t$,
        trace backward through the relevance graph to find which upstream
        sites contribute to the failure.
  \item \textbf{Impact analysis.} When a user adds an annotation at site $s$,
        trace forward through the relevance graph to see which downstream
        obligations are affected.
  \item \textbf{Annotation suggestion.} Find sites with high out-degree in
        the relevance graph (sites whose properties affect many downstream
        obligations). These are high-leverage annotation targets.
\end{itemize}

\section{Relevance under uncertainty: probabilistic relevance}

When the system cannot determine whether a predicate holds (the solver
returns UNKNOWN), it can assign a \emph{probabilistic relevance} based on
heuristics:

\begin{definition}[Probabilistic relevance]
The probabilistic relevance of predicate $\phi$ at site $s$ is:
\[
  R_{\mathrm{prob}}(\phi, s) = P(\phi \text{ is true at } s) \cdot
    R(\phi, s)
\]
where $P(\phi \text{ is true})$ is estimated from:
\begin{itemize}[nosep]
  \item Static analysis heuristics (e.g., most list accesses are in-bounds).
  \item Dynamic profiling (if available).
  \item Prior distributions based on common patterns.
\end{itemize}
\end{definition}

This allows the system to prioritize predicates that are both important
(high relevance) and likely to hold (high probability), focusing verification
effort where it is most likely to succeed.


\chapter{Formal Verification Integration: Connecting to External Provers}

While the Z3-based solver handles decidable fragments, some properties
require the full power of an interactive proof assistant. This chapter
describes how the system integrates with external formal verification tools.

\section{The verification backend interface}

The system defines an abstract \emph{verification backend} interface that
can be implemented by different provers:

\begin{lstlisting}[style=deppy,language=Python]
from abc import ABC, abstractmethod

class VerificationBackend(ABC):
    @abstractmethod
    def check(self, obligation: Obligation,
              context: Context) -> VerificationResult:
        """Check whether the obligation holds in the context."""
        ...

    @abstractmethod
    def supported_theories(self) -> set[str]:
        """Return the set of supported theories."""
        ...

class Z3Backend(VerificationBackend):
    """Z3 SMT solver for decidable fragments."""
    ...

class Lean4Backend(VerificationBackend):
    """Lean 4 proof assistant for interactive proofs."""
    ...

class CoqBackend(VerificationBackend):
    """Coq proof assistant for type-theoretic proofs."""
    ...

class DafnyBackend(VerificationBackend):
    """Dafny program verifier for Hoare-logic proofs."""
    ...
\end{lstlisting}

\section{Obligation routing}

The system routes each obligation to the most appropriate backend:

\begin{enumerate}[nosep]
  \item \textbf{Z3 first.} Every obligation is first sent to Z3. If Z3
        discharges it (UNSAT), done.
  \item \textbf{Theory matching.} If Z3 returns UNKNOWN, the system examines
        the obligation's theory. If it involves induction, it is routed to
        Lean4 or Coq. If it involves program logic, it is routed to Dafny.
  \item \textbf{User preference.} The user can specify a preferred backend
        for specific obligations or obligation types.
  \item \textbf{Fallback.} If all automated backends fail, the obligation
        becomes a residual for human review.
\end{enumerate}

\section{Translation to Lean 4}

For obligations that require interactive proof, the system generates Lean 4
proof obligations:

\begin{lstlisting}[style=deppy,language=Python]
# Python source:
@ensures(lambda xs, result:
    forall(range(len(xs)),
           lambda i: result[i] == xs[len(xs) - 1 - i]))
def reverse(xs: list[int]) -> list[int]:
    return xs[::-1]
\end{lstlisting}

The system generates:
\begin{lstlisting}[language={}]
-- Lean 4 translation:
theorem reverse_correct (xs : List Int) :
    forall i, i < xs.length ->
    (xs.reverse.get i) = (xs.get (xs.length - 1 - i)) := by
  intro i hi
  simp [List.get_reverse]
\end{lstlisting}

The translation preserves the logical structure while adapting to Lean's
syntax and type system. The system can:
\begin{itemize}[nosep]
  \item Generate Lean 4 theorem statements from Python postconditions.
  \item Translate Python data structures to Lean equivalents (lists, dicts,
        optional types).
  \item Import proved Lean theorems back as trusted axioms in the Python
        analysis.
\end{itemize}

\section{Proof portability and the semantic bridge}

The semantic bridge between Python and the proof assistant must be
trustworthy:

\begin{definition}[Semantic bridge]
A \emph{semantic bridge} from Python to a proof assistant $L$ is a pair
of translation functions:
\begin{align*}
  \mathrm{encode} &: \text{Python obligation} \to L\text{-theorem statement} \\
  \mathrm{decode} &: L\text{-proof certificate} \to \text{Python assumption}
\end{align*}
such that if $L \vdash P$ (the proof assistant proves $P$), then the
corresponding Python proposition holds for all values satisfying the
encoding assumptions.
\end{definition}

The bridge is trusted: if the translation is incorrect, the system may accept
invalid proofs. To mitigate this risk:
\begin{itemize}[nosep]
  \item The translation is kept simple and auditable.
  \item Each translation rule is documented with a justification.
  \item The system can generate test cases that exercise both the Python code
        and the Lean theorem, checking consistency dynamically.
\end{itemize}

\section{Hybrid verification workflows}

The integration enables hybrid workflows where some properties are verified
by Z3, some by Lean, and some by runtime checks:

\begin{lstlisting}[style=deppy,language=Python]
from deppy import ensures, lean_proof, runtime_check

@ensures(lambda xs, result: is_sorted(result),
         verified_by="z3")         # shape/bounds: Z3
@ensures(lambda xs, result: is_permutation(result, xs),
         verified_by="lean4")      # permutation: Lean4
@ensures(lambda xs, result: is_stable(result, xs),
         verified_by="runtime")    # stability: runtime check
def sort(xs: list[int]) -> list[int]:
    ...
\end{lstlisting}

The \py{verified_by} parameter specifies which backend should handle each
postcondition. This allows the programmer to match the verification tool to
the difficulty of the property.

\chapter{Information-Theoretic Optimization of Type Inference}

The previous chapters formalized type inference as an optimization problem
over sections of a sheaf. This chapter develops the information-theoretic
foundations in full detail, showing how entropy minimization, mutual
information, and rate-distortion theory provide principled answers to the
questions: How much type information should the system infer? Which
predicates should it track? What is the optimal trade-off between precision
and computational cost?

\section{Types as probability distributions}

Rather than treating types as deterministic sets of values, we adopt a
probabilistic viewpoint: a type specifies a \emph{distribution} over
possible values.

\begin{definition}[Type-as-distribution]
A type assignment $\sigma$ at site $s$ induces a distribution $P_\sigma$ over
the possible runtime values at $s$. The distribution is:
\[
  P_\sigma(v) = \begin{cases}
    1/|\{v : v \models \sigma(s)\}| & \text{if } v \models \sigma(s) \\
    0 & \text{otherwise}
  \end{cases}
\]
under the maximum-entropy assumption (uniform distribution over the admitted
values).
\end{definition}

\begin{definition}[Type entropy]
The \emph{entropy} of a type assignment $\sigma$ at site $s$ is:
\[
  H(\sigma, s) = -\sum_{v} P_\sigma(v) \log P_\sigma(v)
    = \log |\{v : v \models \sigma(s)\}|
\]
A more informative type has lower entropy (fewer admitted values).
\end{definition}

\paragraph{Type inference as entropy minimization.}
The maximum-information objective from the previous chapter is equivalent to
entropy minimization:
\[
  \sigma^* = \arg\min_{\sigma \in H^0(\mathbf{S}_P, \mathcal{T})}
    \sum_{s \in \mathbf{S}_P} w(s) \cdot H(\sigma, s)
\]

This reframing connects type inference to the extensive literature on
information theory, enabling the use of information-theoretic tools.

\section{Mutual information between sites}

The type at one site constrains the types at other sites through data flow.
The \emph{mutual information} between sites quantifies this coupling:

\begin{definition}[Inter-site mutual information]
The mutual information between sites $s$ and $t$ is:
\[
  I(s; t) = H(\sigma, s) + H(\sigma, t) - H(\sigma, s \cup t)
\]
where $H(\sigma, s \cup t)$ is the joint entropy of the types at $s$ and $t$.
\end{definition}

High mutual information means that knowing the type at $s$ strongly
constrains the type at $t$ (or vice versa). The system uses mutual
information to:
\begin{itemize}[nosep]
  \item Identify which sites are most tightly coupled (and should be analyzed
        together).
  \item Detect redundant predicates: if $I(\phi_1; \phi_2) \approx
        H(\phi_1)$, then $\phi_2$ is nearly determined by $\phi_1$ and may
        not need separate tracking.
  \item Guide the stratified solving order: analyze tightly-coupled SCCs
        together and loosely-coupled components independently.
\end{itemize}

\section{Rate-distortion theory for type precision}

Rate-distortion theory provides the optimal trade-off between the ``rate''
(amount of type information tracked) and the ``distortion'' (loss of
precision).

\begin{definition}[Type distortion]
The \emph{distortion} of a type assignment $\sigma$ relative to the
principal type $\sigma^*$ is:
\[
  D(\sigma) = \sum_{s \in \mathbf{S}_P} w(s) \cdot
    d(\sigma(s), \sigma^*(s))
\]
where $d(\cdot, \cdot)$ is a distance metric on types (e.g., the number of
predicates in $\sigma^*$ that are not in $\sigma$).
\end{definition}

\begin{definition}[Type rate]
The \emph{rate} of a type assignment $\sigma$ is the total amount of
information stored:
\[
  R(\sigma) = \sum_{s \in \mathbf{S}_P} |\{\phi \in \Pi :
    \sigma(s) \models \phi\}|
\]
(the total number of tracked predicate instances across all sites).
\end{definition}

\begin{definition}[Rate-distortion function]
The \emph{rate-distortion function} $R(D)$ is:
\[
  R(D) = \min_{\sigma : D(\sigma) \leq D} R(\sigma)
\]
the minimum rate (number of tracked predicates) needed to achieve distortion
at most $D$.
\end{definition}

\begin{theorem}[Optimal type precision allocation]
The rate-distortion-optimal type assignment allocates more predicates to
sites with higher importance weight $w(s)$ and more complex viability
predicates. Sites with simple viability predicates (e.g., $x > 0$) need
fewer tracked predicates; sites with complex viability predicates (e.g.,
involving multiple variables and array indices) need more.
\end{theorem}

This theorem justifies the system's demand-driven approach: allocate
analysis effort proportionally to the complexity and importance of each site.

\section{The predicate selection problem}

Given a budget of $k$ predicates (to keep the analysis tractable), which
predicates should the system track?

\begin{definition}[Predicate selection problem]
Given a set of candidate predicates $\Phi = \{\phi_1, \ldots, \phi_m\}$, a
budget $k \leq m$, and a set of error sites with viability predicates, find
the subset $\Pi \subseteq \Phi$ with $|\Pi| \leq k$ that maximizes the
number of dischargeable obligations.
\end{definition}

\begin{theorem}[NP-hardness of predicate selection]
The predicate selection problem is NP-hard by reduction from maximum
coverage.
\end{theorem}

\begin{proof}[Proof sketch]
Reduce from the maximum coverage problem: given a universe $U$, a collection
of sets $S_1, \ldots, S_m$, and a budget $k$, find $k$ sets whose union is
maximized. Map each set $S_i$ to a predicate $\phi_i$, each element $u \in U$
to an obligation, and say that $\phi_i$ ``helps discharge'' obligation $u$
if $u \in S_i$.
\end{proof}

\paragraph{Greedy approximation.}
Despite NP-hardness, the greedy algorithm achieves a $(1 - 1/e)$
approximation: iteratively select the predicate that discharges the most
additional obligations. This is the standard result for submodular
maximization.

\begin{theorem}[Submodularity of predicate selection]
The number of dischargeable obligations is a monotone submodular function of
the predicate set. Therefore, the greedy algorithm achieves a
$(1 - 1/e) \approx 0.632$ approximation ratio.
\end{theorem}

\begin{proof}[Proof sketch]
Monotonicity: adding a predicate can only help discharge more obligations
(it never removes useful information). Submodularity: the marginal benefit
of adding a predicate decreases as more predicates are already tracked
(diminishing returns). This is because obligations that can be discharged by
multiple predicates contribute less marginal benefit for each additional
predicate.
\end{proof}

\section{Pareto-optimal type assignments}

When multiple objectives compete (precision at different sites, or precision
vs. computation time), the system computes Pareto-optimal type assignments:

\begin{definition}[Pareto optimality]
A type assignment $\sigma$ is \emph{Pareto-optimal} if there is no other
valid type assignment $\sigma'$ that is strictly more informative at some
site and no less informative at any site. The set of Pareto-optimal
assignments is the \emph{Pareto frontier}.
\end{definition}

In practice, the Pareto frontier arises when:
\begin{itemize}[nosep]
  \item Increasing precision at one site requires decreasing precision at
        another (due to solver timeout budgets).
  \item Different error sites require different predicates, and the total
        predicate budget is limited.
  \item The user cares about different properties at different confidence
        levels.
\end{itemize}

The system presents the Pareto frontier to the user, who can choose the
desired trade-off:

\begin{lstlisting}[language={}]
PARETO FRONTIER:
  Option A: 95% of error sites verified, 12s analysis time
  Option B: 98% of error sites verified, 45s analysis time
  Option C: 100% of error sites verified, 180s analysis time
  Current: Option B (default)
\end{lstlisting}

\section{Online optimization: incremental type inference}

As the programmer edits code, the type assignments must be updated. Full
re-analysis is expensive; incremental analysis updates only the affected
sites.

\begin{definition}[Change set]
A code edit produces a \emph{change set} $\Delta \subseteq \mathbf{S}_P$:
the set of sites whose code has changed. The \emph{impact set}
$\mathrm{Impact}(\Delta)$ is the transitive closure of $\Delta$ under
data-flow edges.
\end{definition}

\begin{theorem}[Incremental soundness]
If the type assignment $\sigma$ was valid before the edit, and the edit
affects only sites in $\Delta$, then re-analyzing only
$\mathrm{Impact}(\Delta)$ produces a valid updated assignment. Sites
outside $\mathrm{Impact}(\Delta)$ retain their previous type assignments.
\end{theorem}

The incremental analysis:
\begin{enumerate}[nosep]
  \item Identifies the change set $\Delta$.
  \item Computes the impact set $\mathrm{Impact}(\Delta)$.
  \item Invalidates type assignments in $\mathrm{Impact}(\Delta)$.
  \item Re-runs fixpoint iteration restricted to $\mathrm{Impact}(\Delta)$.
  \item Re-checks obligations in $\mathrm{Impact}(\Delta)$.
\end{enumerate}

For typical edits (modifying a single function), the impact set is small,
and incremental analysis completes in milliseconds.

\section{The Lagrangian relaxation of the combined problem}

The combined problem---maximize type information subject to sheaf
consistency, budget constraints, and solver-decidability constraints---can be
approached via Lagrangian relaxation:

\begin{align*}
  \mathcal{L}(\sigma, \lambda, \mu) &=
    \sum_s w(s) \cdot \mathcal{I}(\sigma, s)
    - \sum_{(i,j)} \lambda_{ij} \cdot
      \mathrm{inconsistency}(\sigma_i, \sigma_j) \\
    &\quad - \mu \cdot (|\Pi(\sigma)| - k)
\end{align*}

where $\lambda_{ij}$ are Lagrange multipliers penalizing overlap
inconsistency and $\mu$ is the multiplier for the predicate budget.

The dual problem $\min_{\lambda, \mu \geq 0} \max_\sigma \mathcal{L}$ can
be solved by subgradient methods, alternating between:
\begin{enumerate}[nosep]
  \item Fixing the multipliers and solving for the best type assignment
        (the inner maximization).
  \item Updating the multipliers based on constraint violations
        (the outer minimization).
\end{enumerate}

This provides a principled approach to balancing precision, consistency, and
computational cost.


\chapter{Obstruction Theory for Type Errors}

When type inference fails, the sheaf-theoretic framework provides precise
diagnostic information through \emph{obstruction theory}: the mathematical
theory of why a global section cannot be constructed.

\section{Obstructions as cohomology classes}

Recall from the cohomology chapter that $H^1(\mathbf{S}_P, \mathcal{T})$
measures the obstructions to gluing local sections. Each nonzero class in
$H^1$ corresponds to a type error. This section makes the correspondence
precise.

\begin{definition}[Type-error cocycle]
A \emph{type-error cocycle} is an element $\alpha \in Z^1(\mathcal{U},
\mathcal{T})$ (a 1-cocycle that is not a 1-coboundary). Concretely,
$\alpha$ assigns to each overlap $u_i \cap u_j$ a ``type mismatch''---a
discrepancy between the local types at $u_i$ and $u_j$ that cannot be
resolved by adjusting the local types.
\end{definition}

\paragraph{Interpretation.}
The cocycle $\alpha_{ij}$ at the overlap $u_i \cap u_j$ is the ``gap''
between what site $u_i$ expects and what site $u_j$ provides. If $\alpha$ is
a coboundary ($\alpha = \delta \beta$ for some 0-cochain $\beta$), then the
gap can be closed by adjusting the local types by $\beta$. If $\alpha$ is
not a coboundary, the gap is structural and cannot be closed.

\begin{example}[Shape mismatch as a cocycle]
Consider two functions: $f$ outputs a tensor of shape $(n, 3)$, and $g$
expects a tensor of shape $(n, 4)$. The overlap is the call site
$g(f(x))$. The cocycle assigns the ``shape gap'' $(n, 3) \neq (n, 4)$ to
this overlap. This gap cannot be resolved by changing the types at other
sites (no adjustment to $f$'s output or $g$'s input can make $3 = 4$), so
it is a non-trivial cohomology class.
\end{example}

\begin{example}[Resolvable inconsistency as a coboundary]
Consider $f$ outputs an \py{int | None} and $g$ expects an \py{int}. The
overlap has a ``None gap.'' But if there is a guard (\py{if x is not None})
between $f$ and $g$, the gap is resolvable: the guard narrows the type,
so the cocycle is a coboundary (the 0-cochain $\beta$ is the guard's
narrowing effect).
\end{example}

\section{Minimal obstruction representatives}

The cohomology class $[\alpha] \in H^1$ may have many representatives.
The system should find the \emph{minimal} representative: the one that
identifies the most localized cause of the error.

\begin{definition}[Minimal representative]
A representative $\alpha$ of a cohomology class $[\alpha]$ is
\emph{minimal} if the support of $\alpha$ (the set of overlaps where
$\alpha_{ij} \neq 0$) has minimum cardinality among all representatives.
\end{definition}

\begin{theorem}[Minimal error localization]
The minimal representative of a type-error cohomology class identifies the
smallest set of code locations that jointly cause the type error.
\end{theorem}

This is the formal version of ``error localization'': the system should blame
the smallest set of locations, not the entire program. The minimal
representative provides this.

\section{Error classification by obstruction type}

Different types of type errors correspond to different cohomology structures:

\paragraph{$H^1$ errors: pairwise inconsistencies.}
These are the most common errors: two sites disagree on the type of a shared
value. Examples: shape mismatch, type mismatch, arity mismatch.

\paragraph{$H^2$ errors: three-way inconsistencies.}
These are rarer: three sites are pairwise consistent but jointly
inconsistent. Example: function $f$ returns a value that satisfies both
predicates $\phi$ and $\psi$, but the conjunction $\phi \wedge \psi$ is
unsatisfiable (the predicates contradict when combined).

\paragraph{Torsion errors: periodic inconsistencies.}
In loop-carried data flow, the type assignment at the loop header must be
consistent with itself after one iteration. If the loop body changes the
type in a way that cannot be undone (e.g., narrowing a type that the
loop condition requires to be wide), the inconsistency is ``periodic'' and
shows up as a torsion element in the cohomology.

\section{Obstruction-guided repair}

When a type error is detected, the obstruction class guides the repair:

\begin{definition}[Repair as coboundary selection]
A \emph{repair} for a type error $[\alpha] \in H^1$ is a 0-cochain
$\beta \in C^0$ such that $\alpha' = \alpha - \delta\beta$ has smaller
support than $\alpha$. Ideally, $\alpha' = 0$ (the error is fully repaired).
The 0-cochain $\beta$ specifies how to adjust the local types to resolve the
inconsistency.
\end{definition}

The system computes repair suggestions by solving:
\[
  \beta^* = \arg\min_\beta \|\alpha - \delta\beta\|
\]
where $\|\cdot\|$ measures the ``size'' of the remaining inconsistency (e.g.,
the number of overlaps with nonzero type gap).

\paragraph{Repair ranking.}
Multiple repairs may exist. The system ranks them by:
\begin{itemize}[nosep]
  \item \textbf{Locality}: repairs that change fewer sites are preferred.
  \item \textbf{Minimality}: repairs that make smaller type changes are
        preferred.
  \item \textbf{Semantic plausibility}: repairs that correspond to common
        coding patterns (adding a guard, adding a cast, changing a function
        signature) are preferred.
\end{itemize}

\begin{lstlisting}[language={}]
TYPE ERROR at line 42: shape mismatch
  f() returns shape (n, 3) but g() expects shape (n, 4)

  Repair suggestions (ranked):
    1. Add a projection: result = f(x)[:, :4]  [if n >= 4]
    2. Change f() to return shape (n, 4)
    3. Change g() to accept shape (n, 3)
    4. Add a reshape: result = f(x).reshape(n, 4)
       [requires n*3 divisible by 4]
\end{lstlisting}

\section{Spectral sequences for complex type errors}

For complex programs with deeply nested function calls and multiple
interacting type constraints, the cohomology may need to be computed using
\emph{spectral sequences}.

\begin{definition}[Spectral sequence for type inference]
The \emph{Leray spectral sequence} for a composition of covers
$\mathcal{U} \to \mathcal{V} \to \mathbf{S}_P$ has $E_2$ page:
\[
  E_2^{p,q} = H^p(\mathcal{V}, H^q(\mathcal{U}/\mathcal{V}, \mathcal{T}))
\]
converging to $H^{p+q}(\mathbf{S}_P, \mathcal{T})$.
\end{definition}

The practical interpretation: first compute local cohomology within each
function ($H^q$ over the function's internal cover), then compute global
cohomology across functions ($H^p$ over the inter-function cover). This
decomposes the global type-error analysis into manageable local analyses,
analogous to the stratified solving strategy.

\paragraph{Differential diagnosis.}
The differentials $d_r : E_r^{p,q} \to E_r^{p+r, q-r+1}$ in the spectral
sequence correspond to type errors that ``cross abstraction boundaries'':
errors that cannot be detected within a single function but arise from the
interaction between functions. These are the hardest errors to diagnose, and
the spectral sequence provides a systematic way to enumerate them.

\chapter{Topos-Theoretic Semantics: The Internal Logic of the Type Topos}

The sheaf-theoretic framework developed so far uses sheaves as a technical
device for gluing local type information. This chapter goes deeper: the
category of sheaves on the program site forms a \emph{topos}, and the
internal logic of this topos provides a powerful language for reasoning about
types, propositions, and their relevance to bug detection.

\section{The type topos}

\begin{definition}[Type topos]
The \emph{type topos} $\mathbf{Sh}(\mathbf{S}_P, J)$ is the category of
sheaves of types on the site $(\mathbf{S}_P, J)$. Objects are type sheaves;
morphisms are natural transformations (type-preserving maps between type
sheaves).
\end{definition}

As a topos, $\mathbf{Sh}(\mathbf{S}_P, J)$ has:
\begin{itemize}[nosep]
  \item A terminal object $\mathbf{1}$: the trivial type sheaf (a single
        type at every site).
  \item Products $\mathcal{T} \times \mathcal{T}'$: pair types.
  \item Exponentials $\mathcal{T}'^{\mathcal{T}}$: function types.
  \item A subobject classifier $\Omega$: the ``truth value'' sheaf.
  \item An internal logic: a higher-order logic that can express propositions
        about types.
\end{itemize}

\section{The subobject classifier and propositions}

\begin{definition}[Subobject classifier]
The subobject classifier $\Omega$ in the type topos is the sheaf defined by:
\[
  \Omega(s) = \{\text{sieves on } s\}
\]
where a \emph{sieve} on $s$ is a downward-closed set of morphisms into $s$.
A global section of $\Omega$ is a ``truth value'' that may vary across the
program.
\end{definition}

\paragraph{Interpretation.}
In a standard set-theoretic context, $\Omega = \{\mathrm{true},
\mathrm{false}\}$. In the type topos, truth values are more nuanced: a
proposition may be true at some sites and false (or unknown) at others. The
subobject classifier captures this \emph{local truth}.

\begin{example}[Local truth for bounds checking]
Consider the proposition ``$i < \mathrm{len}(xs)$'' in a program with
a guard \py{if i < len(xs)}. This proposition is:
\begin{itemize}[nosep]
  \item True at sites inside the true branch.
  \item False at sites inside the false branch.
  \item Unknown at sites before the guard.
\end{itemize}
The truth value is a section of $\Omega$, not a single Boolean.
\end{example}

\section{Propositions as subobjects of the type sheaf}

In the topos, a proposition about a type is a \emph{subobject}: a
monomorphism $\mathcal{P} \hookrightarrow \mathcal{T}$ that picks out the
values satisfying the proposition.

\begin{definition}[Refinement as subobject]
The refinement type $\{x : A \mid \phi(x)\}$ corresponds to the subobject
$\phi^{-1}(\mathrm{true}) \hookrightarrow A$ in the type topos. The
``characteristic map'' $\chi_\phi : A \to \Omega$ sends each value to its
local truth value under $\phi$.
\end{definition}

This formalization captures the key insight: a refinement type is not just a
set constraint; it is a \emph{site-dependent} constraint whose truth may vary
across the program.

\section{The internal logic for type reasoning}

The internal logic of the type topos is an intuitionistic higher-order logic
with the following connectives:

\begin{itemize}[nosep]
  \item \textbf{Conjunction} $\phi \wedge \psi$: a value satisfies both $\phi$
        and $\psi$ at a site iff it satisfies both propositions at that site.
  \item \textbf{Disjunction} $\phi \vee \psi$: a value satisfies $\phi$ or
        $\psi$ at a site. Note: this is constructive---disjunction does not
        imply that we know which disjunct holds.
  \item \textbf{Implication} $\phi \Rightarrow \psi$: at every site where
        $\phi$ holds, $\psi$ also holds. This is the \emph{local implication}.
  \item \textbf{Universal quantification} $\forall x. \phi(x)$: for all
        values $x$ at the site, $\phi(x)$ holds.
  \item \textbf{Existential quantification} $\exists x. \phi(x)$: there
        exists a value $x$ at the site satisfying $\phi(x)$.
  \item \textbf{Negation} $\neg \phi = (\phi \Rightarrow \bot)$: $\phi$
        implies absurdity. Note: $\neg\neg\phi$ does not imply $\phi$
        (the logic is intuitionistic).
\end{itemize}

\paragraph{Why intuitionistic?}
The excluded middle ($\phi \vee \neg\phi$) does not hold in general in the
type topos, because a proposition may be neither provable nor refutable at a
given site. This matches the reality of type inference: the system may not be
able to determine whether a property holds or not (the solver may time out).
Intuitionistic logic is the natural logic for this setting.

\paragraph{Recovering classical reasoning.}
For decidable predicates (those in the solver's decidable fragment), the
excluded middle does hold: the solver either proves $\phi$ or proves
$\neg\phi$. The system tracks which predicates are decidable and allows
classical reasoning for them while maintaining intuitionistic reasoning for
undecidable predicates.

\section{The forcing semantics: what is ``true at a site''?}

\begin{definition}[Forcing]
A site $s$ \emph{forces} a proposition $\phi$ (written $s \Vdash \phi$) if
$\phi$ holds in the section assigned to $s$. The forcing relation satisfies:
\begin{align*}
  s \Vdash \phi \wedge \psi &\iff s \Vdash \phi \text{ and } s \Vdash \psi \\
  s \Vdash \phi \vee \psi &\iff \text{there is a cover } \{u_i \to s\}
    \text{ such that each } u_i \Vdash \phi \text{ or } u_i \Vdash \psi \\
  s \Vdash \phi \Rightarrow \psi &\iff \forall (t \to s) : t \Vdash \phi
    \implies t \Vdash \psi \\
  s \Vdash \forall x. \phi(x) &\iff \forall (t \to s), \forall a \in
    \mathcal{T}(t) : t \Vdash \phi(a) \\
  s \Vdash \exists x. \phi(x) &\iff \text{there is a cover } \{u_i \to s\}
    \text{ and } a_i \in \mathcal{T}(u_i) \text{ with } u_i \Vdash \phi(a_i)
\end{align*}
\end{definition}

\paragraph{Interpretation for type inference.}
\begin{itemize}[nosep]
  \item $s \Vdash (x > 0)$ means: at program point $s$, the system has
        established that $x > 0$.
  \item $s \Vdash (\mathrm{len}(xs) = n)$ means: at program point $s$, the
        system knows the list's length.
  \item The forcing relation propagates through data flow: if $s \Vdash \phi$
        and $s \to t$ is a data-flow edge that preserves $\phi$, then
        $t \Vdash \phi$.
\end{itemize}

\section{Relevance in the internal logic}

The internal logic provides a precise language for expressing relevance:

\begin{definition}[Bug-relevance in the internal logic]
A proposition $\phi$ at site $s$ is \emph{bug-relevant} if:
\[
  s \Vdash \left(\phi \Rightarrow \bigwedge_{e \in E(s)}
    \mathrm{viable}(e)\right) \wedge \left(\neg\phi \Rightarrow
    \bigvee_{e \in E(s)} \neg\mathrm{viable}(e)\right)
\]
where $E(s)$ is the set of error sites reachable from $s$. That is: $\phi$
being true implies all reachable errors are viable (safe), and $\phi$ being
false implies some error is non-viable (a bug).
\end{definition}

This definition captures the intuition precisely: a bug-relevant proposition
is one that \emph{separates} safe behavior from buggy behavior.

\section{The Lawvere-Tierney topology for decidable propositions}

\begin{definition}[Lawvere-Tierney topology]
A \emph{Lawvere-Tierney topology} on the type topos is a morphism
$j : \Omega \to \Omega$ satisfying:
\begin{enumerate}[nosep]
  \item $j \circ \mathrm{true} = \mathrm{true}$ (truth is preserved).
  \item $j \circ j = j$ (idempotent).
  \item $j \circ \wedge = \wedge \circ (j \times j)$ (commutes with
        conjunction).
\end{enumerate}
\end{definition}

The decidability boundary defines a Lawvere-Tierney topology: $j(\phi) =
\mathrm{true}$ if $\phi$ is decidable (the solver can determine its truth),
and $j(\phi) = \phi$ otherwise. The \emph{$j$-sheaves} are the types whose
propositions are all decidable---the ``decidable fragment'' of the type
system.

\paragraph{The decidable subtoposes.}
The subtoposes $\mathbf{Sh}_j(\mathbf{S}_P) \subseteq
\mathbf{Sh}(\mathbf{S}_P)$ correspond to different decidable fragments:
\begin{itemize}[nosep]
  \item $j_{\mathrm{LIA}}$: linear integer arithmetic.
  \item $j_{\mathrm{LRA}}$: linear real arithmetic.
  \item $j_{\mathrm{BV}}$: bitvector arithmetic.
  \item $j_{\mathrm{ARR}}$: array theory.
\end{itemize}

The system operates primarily within $\mathbf{Sh}_{j_{\mathrm{LIA}}}$ (where
most practical type reasoning lives) and escapes to the full topos when
necessary.

\section{Modalities for annotation levels}

The annotation levels (0--5 from Part IV) correspond to modalities in the
topos:

\begin{itemize}[nosep]
  \item \textbf{Level 0} ($\Box_{\mathrm{infer}}$): only propositions that
        the inference engine can derive.
  \item \textbf{Level 1} ($\Box_{\mathrm{contract}}$): propositions stated by
        contracts and verified by the solver.
  \item \textbf{Level 2} ($\Box_{\mathrm{invariant}}$): propositions
        maintained by invariants and verified inductively.
  \item \textbf{Level 3--5} ($\Box_{\mathrm{proof}}$): propositions with
        human-provided proofs.
\end{itemize}

Each modality defines a subtopos of ``known truths'' at that level. Higher
levels include all lower levels. The full type topos includes all levels,
plus propositions that are true but unknown at any level.

This topos-theoretic perspective reveals the mathematical structure of the
annotation economy: the annotation levels form a chain of subtoposes, and the
``value'' of an annotation is its contribution to extending the subtopos from
one level to the next.


\chapter{Automated Deduction Strategies for Richer Dependent Types}

This chapter presents advanced deduction strategies that go beyond standard
SMT solving to infer richer dependent types automatically, without user
annotations.

\section{Abductive inference: hypothesizing missing predicates}

Standard forward inference derives conclusions from known facts. \emph{Abductive
inference} works in reverse: given a desired conclusion (a viability predicate
that must hold), hypothesize the missing premises.

\begin{definition}[Abductive inference for types]
Given a viability predicate $V$ at an error site $e$ and the current section
$\sigma$ at $e$ (which does not imply $V$), the \emph{abductive inference
problem} is:
\[
  \text{Find the weakest } \phi \text{ such that }
    \sigma \cup \{\phi\} \models V
\]
\end{definition}

\paragraph{Algorithm.}
\begin{enumerate}[nosep]
  \item Negate $V$ and conjoin with $\sigma$: form $\sigma \wedge \neg V$.
  \item Use Z3's \py{unsat_core()} to find the minimal set of facts in
        $\sigma$ that participate in the (un)satisfiability check.
  \item If $\sigma \wedge \neg V$ is satisfiable, the satisfying model
        identifies the ``gap'': the conditions under which $V$ fails.
  \item Negate the gap conditions to get $\phi$: the missing predicate.
  \item Trace $\phi$ backward through data flow to find where $\phi$
        should be established.
\end{enumerate}

\begin{example}
Error site: \py{xs[i]} with viability $0 \leq i < \mathrm{len}(xs)$.
Current section: $i = n - 1$ (from a loop), $\mathrm{len}(xs) = m$.
The gap: $n - 1 < m$ is not implied. The abduced predicate: $n \leq m$.
Trace backward: $n$ comes from a loop bound, $m$ comes from the input.
The system synthesizes the precondition $n \leq \mathrm{len}(xs)$.
\end{example}

\section{Craig interpolation for interface synthesis}

When two modules interact, the system needs to find the ``right'' interface:
a set of predicates that is strong enough for the caller's needs and weak
enough for the callee to satisfy.

\begin{definition}[Craig interpolant]
Given formulas $A$ and $B$ such that $A \models B$, a \emph{Craig
interpolant} is a formula $I$ such that:
\begin{enumerate}[nosep]
  \item $A \models I$ (the implementation implies the interface).
  \item $I \models B$ (the interface implies the caller's needs).
  \item The variables of $I$ are shared between $A$ and $B$ (the interface
        uses only shared vocabulary).
\end{enumerate}
\end{definition}

\paragraph{Application to type inference.}
Let $A$ be the callee's postcondition (what the function actually
guarantees) and $B$ be the caller's requirement (what the caller needs to
know about the return value). The interpolant $I$ is the optimal function
signature: it exports exactly the information the caller needs, no more and
no less.

Z3 supports Craig interpolation for linear arithmetic and some other
theories. The system uses this to automatically synthesize function
signatures that are both sufficient for callers and provable from the
implementation.

\section{Quantifier instantiation for bounded domains}

Many dependent type predicates involve bounded quantification:
$\forall i \in [0, n). \phi(i)$. The solver handles these by instantiation:

\paragraph{Finite instantiation.}
When $n$ is a known constant (or small symbolic bound), the quantifier is
expanded into a finite conjunction:
$\phi(0) \wedge \phi(1) \wedge \cdots \wedge \phi(n-1)$.

\paragraph{Inductive instantiation.}
When $n$ is symbolic but the predicate has an inductive structure, the
system uses inductive reasoning:
\begin{itemize}[nosep]
  \item Base case: $\phi(0)$.
  \item Inductive step: $\phi(k) \Rightarrow \phi(k+1)$ for $0 \leq k < n$.
  \item Conclusion: $\forall i \in [0, n). \phi(i)$.
\end{itemize}

The system attempts inductive instantiation automatically for quantified
postconditions of loops and recursive functions.

\paragraph{E-matching.}
For quantifiers over complex expressions, the system uses E-matching: it
identifies ground instances of the quantified variable that appear in the
current context and instantiates the quantifier at those instances. This is
Z3's standard approach for quantifier handling.

\section{Abstract interpretation with predicate abstraction}

The system uses abstract interpretation to infer loop invariants and
function summaries automatically:

\begin{definition}[Predicate abstract domain]
Given a predicate domain $\Pi = \{\phi_1, \ldots, \phi_k\}$, the
\emph{predicate abstract domain} is the Boolean lattice $2^\Pi$ (subsets of
$\Pi$). An abstract value $S \subseteq \Pi$ represents the conjunction
$\bigwedge_{\phi \in S} \phi$.
\end{definition}

The abstract interpretation computes, for each program point, the set of
predicates that hold:
\begin{enumerate}[nosep]
  \item Initialize all program points with $\Pi$ (all predicates hold).
  \item For each program point $s$, compute the abstract transfer function:
        which predicates are preserved by the operations at $s$?
  \item Iterate until fixpoint.
  \item The fixpoint gives the strongest invariant expressible in $\Pi$.
\end{enumerate}

This is the standard approach from model checking (SLAM, BLAST), adapted to
the sheaf-descent framework. The key difference is that the predicate domain
$\Pi$ is not fixed a priori but is dynamically refined by the CEGAR loop
described earlier.

\section{Symbolic execution for path-sensitive inference}

Path-sensitive analysis tracks different type information along different
execution paths:

\begin{lstlisting}[style=deppy,language=Python]
def process(x: int | str) -> int:
    if isinstance(x, int):
        return x + 1       # path 1: x is int
    else:
        return len(x)      # path 2: x is str
\end{lstlisting}

The path-insensitive analysis would merge the two paths and lose the type
narrowing. The path-sensitive analysis maintains separate sections:
\begin{itemize}[nosep]
  \item Path 1: $x : \py{int}$, result $= x + 1 : \py{int}$.
  \item Path 2: $x : \py{str}$, result $= \mathrm{len}(x) : \py{int}$.
\end{itemize}

The system merges paths at join points but preserves path conditions that
affect type-critical properties.

\paragraph{Bounded symbolic execution.}
For loops, the system unrolls a bounded number of iterations and analyzes
each unrolled path separately. After the bound, it widens to an
over-approximation. The bound is configurable (default: 3 iterations for
inference, unlimited for proof checking).

\section{Liquid types: type inference as predicate abduction}

The system incorporates ideas from \emph{Liquid Types} (Rondon, Kawaguchi,
Jhala): automatically inferring refinement types using predicate abduction
and Horn clause solving.

\begin{definition}[Liquid type inference]
Given a program with types annotated as ``holes'' $\kappa_i$ (unknown
refinements), the system:
\begin{enumerate}[nosep]
  \item Generates subtyping constraints between types with holes.
  \item Converts the constraints to Horn clauses over the predicate domain.
  \item Solves the Horn clauses using a fixpoint solver.
  \item The solution assigns predicates to each hole, giving the strongest
        valid refinement.
\end{enumerate}
\end{definition}

The integration with the sheaf framework adds a key capability: the predicate
domain is not fixed but dynamically refined based on the relevance analysis.
This means the system discovers predicates that are both expressible and
useful, rather than requiring a predefined set.

\section{Machine-learning-guided predicate suggestion}

For predicates that fall outside the decidable fragment, the system can use
machine learning to suggest likely predicates:

\begin{itemize}[nosep]
  \item \textbf{Pattern matching}: recognize common code patterns (e.g.,
        ``loop over range(len(xs))'') and suggest standard predicates
        (e.g., ``loop variable is a valid index'').
  \item \textbf{Historical data}: learn from previously verified programs
        which predicates were useful in similar contexts.
  \item \textbf{Natural language}: use the function's docstring or comments
        to suggest predicates (e.g., ``returns a sorted list'' suggests
        $\forall i < \mathrm{len}(r)-1. r[i] \leq r[i+1]$).
\end{itemize}

Machine-learning suggestions are never trusted; they are treated as
candidates that must be verified by the solver. The value is in reducing the
search space for predicate discovery.

\chapter{Enriched Category Theory for Type Propagation and Kan Extensions}

The interprocedural transport of type information---how types flow across
function calls, module boundaries, and library interfaces---can be formalized
using enriched category theory and Kan extensions. This chapter develops the
mathematics and shows how it enables optimal type propagation.

\section{Enriched categories of types}

Standard categories have sets of morphisms. \emph{Enriched categories} have
morphism objects in a monoidal category $\mathcal{V}$, allowing richer
structure.

\begin{definition}[$\mathcal{V}$-enriched type category]
Let $\mathcal{V} = (\mathbf{Lat}, \wedge, \top)$ be the monoidal category
of complete lattices with meet as the tensor product. A
$\mathcal{V}$-enriched category $\mathbf{T}$ has:
\begin{itemize}[nosep]
  \item Objects: program sites $s, t, \ldots$
  \item Hom-objects: $\mathbf{T}(s, t) \in \mathbf{Lat}$, the lattice of
        \emph{type transfer functions} from $s$ to $t$.
  \item Composition: a morphism
        $\circ : \mathbf{T}(t, u) \otimes \mathbf{T}(s, t) \to
        \mathbf{T}(s, u)$ in $\mathcal{V}$ (composition of transfer
        functions).
\end{itemize}
\end{definition}

\paragraph{Interpretation.}
A type transfer function $f \in \mathbf{T}(s, t)$ is a lattice morphism that
maps a type at $s$ to a type at $t$. For example, the operation
\py{len(xs)} induces a transfer function that maps the type of \py{xs}
(a list with refinements) to the type of the result (an integer with the
refinement $\geq 0$).

The enrichment in $\mathbf{Lat}$ means that the set of transfer functions
is itself a lattice: there is a most informative transfer function (the one
that preserves the most type information) and a least informative one (the
one that loses all information).

\section{Weighted colimits for type synthesis}

Type inference at a site $s$ with multiple inputs $u_1, \ldots, u_n$ is a
\emph{weighted colimit} in the enriched category:

\begin{definition}[Type synthesis as weighted colimit]
The inferred type at site $s$ with inputs from sites $u_1, \ldots, u_n$ is:
\[
  \sigma(s) = \mathrm{colim}_{W} F
\]
where:
\begin{itemize}[nosep]
  \item $F : \{u_1, \ldots, u_n\} \to \mathbf{T}$ maps each input site to
        its type.
  \item $W : \{u_1, \ldots, u_n\}^{\mathrm{op}} \to \mathcal{V}$ is the
        \emph{weight} functor that assigns importance to each input.
  \item The colimit synthesizes the optimal type at $s$ by combining the
        input types according to the weights and transfer functions.
\end{itemize}
\end{definition}

\paragraph{Concrete construction.}
For a binary operation $s = f(u_1, u_2)$, the weighted colimit is:
\[
  \sigma(s) = \bigsqcup \{f_*(\sigma(u_1), \sigma(u_2)) \mid
    f_* \text{ is the transfer function of } f\}
\]
the join (strongest common type) of all types that can be derived from the
inputs through $f$'s semantics.

\section{Kan extensions for interprocedural transport}

The key challenge of interprocedural analysis is: given a function $g$ that
calls function $f$, how should $f$'s type information be transported into
$g$'s context?

\begin{definition}[Left Kan extension for forward transport]
Given the inclusion functor $i : \mathbf{S}_f \hookrightarrow \mathbf{S}_g$
(embedding $f$'s sites into $g$'s context at the call site) and the type
sheaf $\mathcal{T}_f$ on $\mathbf{S}_f$, the \emph{left Kan extension}
\[
  \mathrm{Lan}_i \mathcal{T}_f : \mathbf{S}_g \to \mathbf{Lat}
\]
is the optimal way to extend $f$'s type information to all of $g$'s sites.
At each site $s \in \mathbf{S}_g$:
\[
  (\mathrm{Lan}_i \mathcal{T}_f)(s) = \mathrm{colim}_{(i \downarrow s)}
    \mathcal{T}_f
\]
the colimit over all sites in $f$ that can reach $s$.
\end{definition}

\begin{definition}[Right Kan extension for backward transport]
Dually, the \emph{right Kan extension}
$\mathrm{Ran}_i \mathcal{T}_g$ transports $g$'s type requirements backward
into $f$'s context:
\[
  (\mathrm{Ran}_i \mathcal{T}_g)(s) = \lim_{(s \downarrow i)}
    \mathcal{T}_g
\]
the limit over all sites in $g$ that can be reached from $s$.
\end{definition}

\paragraph{Interpretation.}
\begin{itemize}[nosep]
  \item The left Kan extension is \emph{forward propagation}: it pushes
        $f$'s output type into $g$'s sites as strongly as possible.
  \item The right Kan extension is \emph{backward propagation}: it pulls
        $g$'s requirements back into $f$'s sites as weakly as possible.
\end{itemize}

\begin{theorem}[Optimality of Kan extensions]
The left Kan extension gives the strongest possible forward propagation, and
the right Kan extension gives the weakest possible backward propagation,
among all extensions that respect the data flow and the type transfer
functions.
\end{theorem}

\begin{proof}[Proof sketch]
By the universal property of Kan extensions: $\mathrm{Lan}_i$ is left adjoint
to restriction $i^*$, so it is the \emph{best} (strongest) extension. Dually,
$\mathrm{Ran}_i$ is right adjoint to $i^*$, so it is the best (weakest)
backward extension.
\end{proof}

\section{The pointwise formula and computational realization}

In practice, Kan extensions are computed pointwise:

\begin{theorem}[Pointwise Kan extension]
If $\mathcal{V}$ is complete and cocomplete (which $\mathbf{Lat}$ is), then
the Kan extensions exist and are computed pointwise:
\begin{align*}
  (\mathrm{Lan}_i \mathcal{T}_f)(s) &= \int^{t \in \mathbf{S}_f}
    \mathbf{S}_g(i(t), s) \otimes \mathcal{T}_f(t) \\
  (\mathrm{Ran}_i \mathcal{T}_g)(s) &= \int_{t \in \mathbf{S}_f}
    [\mathbf{S}_g(s, i(t)), \mathcal{T}_g(i(t))]
\end{align*}
where $\int^t$ denotes the coend and $\int_t$ the end.
\end{theorem}

\paragraph{Algorithmic realization.}
The pointwise formula translates directly to an algorithm:
\begin{enumerate}[nosep]
  \item For each site $s$ in the caller, find all sites $t$ in the callee
        that can reach $s$ (the ``comma category'' $i \downarrow s$).
  \item For each such $t$, apply the transfer function along the path
        $i(t) \to s$ to transport $\mathcal{T}_f(t)$.
  \item Take the join (for forward propagation) or meet (for backward
        propagation) of all transported types.
\end{enumerate}

This is exactly what the interprocedural transport already does in the
implementation, but the Kan extension framework proves it is \emph{optimal}.

\section{Profunctors for type interfaces}

The interface between two modules can be modeled as a \emph{profunctor}:

\begin{definition}[Type profunctor]
A \emph{type profunctor} $H : \mathbf{S}_A \nrightarrow \mathbf{S}_B$ is a
$\mathcal{V}$-enriched functor
$H : \mathbf{S}_A^{\mathrm{op}} \times \mathbf{S}_B \to \mathcal{V}$.
For each pair $(a, b)$ of sites in modules $A$ and $B$, $H(a, b)$ is the
lattice of ``type relationships'' between $a$ and $b$.
\end{definition}

\paragraph{Composition of interfaces.}
If module $A$ calls module $B$ which calls module $C$, the composed interface
$A \to C$ is the profunctor composition:
\[
  (H_2 \circ H_1)(a, c) = \int^{b \in \mathbf{S}_B}
    H_1(a, b) \otimes H_2(b, c)
\]
This captures transitive type propagation across chains of function calls.

\paragraph{Identity profunctor.}
The identity profunctor $\mathrm{Id}(a, b) = \mathbf{S}(a, b)$ represents
``no interface overhead'': direct data flow with no loss of type information.

\section{The fibered category of type families}

As programs evolve, the type sheaf changes. The collection of all type
sheaves (one per program version) forms a \emph{fibered category}:

\begin{definition}[Type fibration]
The \emph{type fibration} $\pi : \mathbf{T} \to \mathbf{Prog}$ has:
\begin{itemize}[nosep]
  \item Base category $\mathbf{Prog}$: programs and program transformations
        (edits, refactorings).
  \item Fiber $\mathbf{T}_P$ over each program $P$: the type topos
        $\mathbf{Sh}(\mathbf{S}_P, J_P)$.
  \item Transition functors: for each program transformation $f : P \to Q$,
        a functor $f^* : \mathbf{T}_Q \to \mathbf{T}_P$ that ``translates''
        types between program versions.
\end{itemize}
\end{definition}

\paragraph{Incremental analysis as base change.}
When the programmer edits the program ($f : P \to Q$), the transition functor
$f^*$ identifies which types can be preserved (unchanged sites) and which
must be recomputed (changed or affected sites). This is the formal
foundation for incremental type inference.


\chapter{The Relevance Optimization Problem: Formal Treatment}

This chapter formalizes the relevance optimization problem with full
mathematical precision and develops algorithms that determine, for each
proposition the system could attach to a term, whether that proposition is
worth tracking for the purpose of bug detection.

\section{The space of decorations}

\begin{definition}[Decoration]
A \emph{decoration} of a term $t$ at site $s$ is a proposition $\phi$ such
that $s \Vdash \phi(t)$ (the system can establish that $t$ satisfies $\phi$
at site $s$). The \emph{decoration space} $\mathcal{D}(t, s)$ is the set of
all propositions that the system can prove about $t$ at $s$.
\end{definition}

The decoration space is typically enormous: for an integer variable, it
includes $x > 0$, $x \leq 100$, $x$ is even, $x \equiv 3 \pmod{7}$, and
infinitely many others. The relevance problem is to select the
\emph{useful} subset.

\begin{definition}[Useful decoration]
A decoration $\phi \in \mathcal{D}(t, s)$ is \emph{useful} if removing it
from the type at $s$ would cause at least one obligation to become
undischargeable. Formally:
\[
  \phi \in \mathcal{D}_{\mathrm{useful}}(t, s) \iff
  \exists e \in \mathrm{Obligations}(P) :
  \sigma \setminus \{\phi\} \not\models V_e \wedge
  \sigma \models V_e
\]
\end{definition}

\section{Minimal useful decoration sets}

\begin{definition}[Minimal useful set]
A set $\Phi \subseteq \mathcal{D}(t, s)$ is a \emph{minimal useful
decoration set} if:
\begin{enumerate}[nosep]
  \item Every useful decoration is in $\Phi$ or is implied by $\Phi$.
  \item No proper subset of $\Phi$ has property 1.
\end{enumerate}
\end{definition}

\begin{theorem}[Existence and computation]
The minimal useful decoration set exists and can be computed as the set of
unsatisfiable core predicates over all obligations:
\[
  \Phi^* = \bigcup_{e \in \mathrm{Obligations}} \mathrm{unsat\_core}(e, \sigma)
\]
where $\mathrm{unsat\_core}(e, \sigma)$ is the minimal subset of $\sigma$
that implies the viability predicate $V_e$.
\end{theorem}

\begin{proof}[Proof sketch]
Each obligation $e$ depends on a subset of the section's predicates. The
unsat core identifies this subset. The union across all obligations gives the
predicates that are needed by at least one obligation. Any predicate not in
this union is not needed by any obligation and can be dropped.
\end{proof}

\section{The relevance lattice}

\begin{definition}[Relevance lattice]
The \emph{relevance lattice} $\mathcal{R}$ at a site $s$ is the lattice of
subsets of the decoration space, ordered by inclusion, with the additional
structure:
\begin{itemize}[nosep]
  \item Bottom: $\emptyset$ (no decorations tracked).
  \item Top: $\mathcal{D}(t, s)$ (all decorations tracked).
  \item Meet: intersection (common decorations).
  \item Join: union (combined decorations).
\end{itemize}
\end{definition}

The minimal useful set $\Phi^*$ is a specific element of this lattice. The
system should compute $\Phi^*$ or a close approximation.

\section{Relevance propagation as a Galois connection}

The relationship between decorations at different sites is captured by a
Galois connection:

\begin{definition}[Relevance Galois connection]
For sites $s$ and $t$ connected by a data-flow edge $f : s \to t$, define:
\begin{align*}
  \alpha &: \mathcal{R}(s) \to \mathcal{R}(t), \quad
    \alpha(\Phi) = \{f_*(\phi) \mid \phi \in \Phi\} \\
  \gamma &: \mathcal{R}(t) \to \mathcal{R}(s), \quad
    \gamma(\Psi) = \{f^*(\psi) \mid \psi \in \Psi\}
\end{align*}
where $f_*$ pushes decorations forward and $f^*$ pulls decorations backward.
The pair $(\alpha, \gamma)$ forms a Galois connection:
$\alpha(\Phi) \subseteq \Psi \iff \Phi \subseteq \gamma(\Psi)$.
\end{definition}

\paragraph{Forward relevance propagation.}
If a decoration $\phi$ at $s$ is relevant, then the pushed-forward decoration
$f_*(\phi)$ at $t$ may also be relevant (it contributes to the same
downstream obligations through $t$).

\paragraph{Backward relevance propagation.}
If a decoration $\psi$ at $t$ is relevant (needed for some obligation at or
beyond $t$), then the pulled-back decoration $f^*(\psi)$ at $s$ is also
relevant (it is needed to establish $\psi$ at $t$).

\section{The relevance fixpoint}

Combining forward and backward propagation gives a fixpoint computation:

\begin{algorithm}[H]
\caption{Relevance fixpoint computation}
\begin{algorithmic}[1]
\State Initialize: $\mathcal{R}(s) \gets \emptyset$ for all sites $s$
\State Seed: For each error site $e$, add $V_e$'s atoms to $\mathcal{R}(e)$
\Repeat
  \ForAll{data-flow edges $f : s \to t$}
    \State $\mathcal{R}(s) \gets \mathcal{R}(s) \cup f^*(\mathcal{R}(t))$
      \Comment{backward}
    \State $\mathcal{R}(t) \gets \mathcal{R}(t) \cup f_*(\mathcal{R}(s))$
      \Comment{forward}
  \EndFor
\Until{convergence}
\State \Return $\mathcal{R}$
\end{algorithmic}
\end{algorithm}

\begin{theorem}[Convergence of relevance fixpoint]
If the decoration space at each site is finite (which it is when the
predicate domain is finite), the relevance fixpoint converges in at most
$|\Pi| \cdot |\mathbf{S}_P|$ iterations.
\end{theorem}

\section{User interaction with the relevance system}

The user can interact with the relevance system at several levels:

\paragraph{Querying relevance.}
\begin{lstlisting}[style=deppy,language=Python]
from deppy import why_relevant, why_irrelevant

# Why is this predicate tracked?
why_relevant(x, "x > 0", at_line=42)
# -> "x > 0" is relevant because:
#    - It is needed to prove xs[x-1] is in bounds (line 45)
#    - It is needed to prove sqrt(x) is real-valued (line 48)

# Why is this predicate NOT tracked?
why_irrelevant(x, "x % 2 == 0", at_line=42)
# -> "x % 2 == 0" is irrelevant because:
#    - No error site reachable from line 42 depends on
#      the parity of x
#    - Dropping this predicate does not affect any obligation
\end{lstlisting}

\paragraph{Forcing relevance.}
\begin{lstlisting}[style=deppy,language=Python]
from deppy import force_relevant

# Track this predicate even if automatic analysis
# considers it irrelevant:
force_relevant(matrix, "matrix.is_symmetric()",
               reason="symmetric matrices have real eigenvalues")
\end{lstlisting}

\paragraph{Suppressing relevance.}
\begin{lstlisting}[style=deppy,language=Python]
from deppy import suppress

# Don't bother tracking this (known imprecision):
suppress(x, "x < 2**63", reason="Python ints don't overflow")
\end{lstlisting}

\section{Relevance-weighted MaxSMT}

The relevance scores feed directly into the MaxSMT objective function:

\[
  \sigma^* = \arg\max_{\sigma \in H^0(\mathbf{S}_P, \mathcal{T})}
    \sum_{(\phi, s) \in \mathcal{R}} R(\phi, s) \cdot
    \mathbb{1}[\sigma(s) \models \phi]
\]

This weighted objective ensures that the solver spends its effort on
predicates that are both provable and useful for bug detection. Irrelevant
predicates have $R = 0$ and are ignored by the solver.

\begin{theorem}[Optimality of relevance-weighted inference]
The relevance-weighted MaxSMT solution discharges the maximum possible
number of obligations (weighted by error-site importance) among all type
assignments that are consistent and use at most $k$ predicates per site.
\end{theorem}

\chapter{Connections to Homotopy Type Theory and Higher-Dimensional Type Structure}

This chapter explores deeper connections between the sheaf-descent framework
and homotopy type theory (HoTT). While the practical system operates in a
decidable fragment, the HoTT perspective reveals structural insights that
improve automatic deduction and clarify the space of type equivalences.

\section{Types as spaces, programs as paths}

In HoTT, types are interpreted as spaces and terms as points in those spaces.
Equalities between terms are paths; equalities between equalities are
homotopies. This provides a geometric intuition for type equivalence.

\paragraph{Program paths.}
Two programs $f$ and $g$ of the same type $A \to B$ are \emph{path-connected}
if there exists a continuous transformation from $f$ to $g$ that preserves
typing. In our framework, this means that the type sheaves of $f$ and $g$ are
equivalent (their sections carry the same information).

\paragraph{Practical consequence.}
Two implementations of the same specification are path-connected. The system
can use this to transfer proofs between implementations: if we prove a
property of implementation $f$ and show that $g$ is path-connected to $f$
(via the refinement proofs from Part IV), the property transfers to $g$
automatically.

\section{The univalence principle for type equivalence}

\begin{definition}[Univalence for dependent types]
Two types $A$ and $B$ are \emph{equivalent} ($A \simeq B$) if there exist
functions $f : A \to B$ and $g : B \to A$ with $g \circ f \sim \mathrm{id}_A$
and $f \circ g \sim \mathrm{id}_B$. The univalence axiom states:
\[
  (A \simeq B) \simeq (A =_{\mathcal{U}} B)
\]
Equivalence of types \emph{is} equality of types.
\end{definition}

\paragraph{Application to type inference.}
Univalence means that the system should not distinguish between equivalent
types. If $A = \{x : \py{int} \mid x > 0\}$ and $B = \{x : \py{int} \mid
x \geq 1\}$, then $A \simeq B$ (they have the same elements), so the
system should treat them as the same type. This is captured by the
normalization step in the solver: equivalent refinement predicates are
simplified to a canonical form before comparison.

\paragraph{Automatic equivalence detection.}
The system should detect type equivalences automatically:
\begin{itemize}[nosep]
  \item $\{x : A \mid \phi(x) \wedge \psi(x)\} \simeq
        \{x : A \mid \psi(x) \wedge \phi(x)\}$
        (conjunction is commutative).
  \item $\{x : \py{int} \mid x > 0 \wedge x > 0\} \simeq
        \{x : \py{int} \mid x > 0\}$ (idempotence).
  \item $\{x : \py{int} \mid x > 0 \wedge x < 10\} \simeq
        \{x : \py{int} \mid 0 < x < 10\}$ (rewriting).
  \item $\py{list}[A]$ where $\mathrm{len} = 0 \simeq$ the unit type (an
        empty list carries no element information).
\end{itemize}

Normalizing types to canonical forms reduces the number of distinct types
the system must track, improving both efficiency and readability of the
inferred types.

\section{Higher inductive types for recursive data structures}

Higher inductive types (HITs) generalize ordinary inductive types by allowing
constructors that produce paths (equalities) in addition to points (elements).

\paragraph{Application: quotient types.}
A common programming pattern is working with values modulo an equivalence
relation:
\begin{lstlisting}[style=deppy,language=Python]
# Rational numbers: (p, q) with q != 0, modulo (p,q) ~ (p*k, q*k)
class Rational:
    def __init__(self, p: int, q: Refine[int, lambda q: q != 0]):
        g = gcd(abs(p), abs(q))
        self.p = p // g
        self.q = abs(q) // g
\end{lstlisting}

The type of \py{Rational} is a \emph{quotient type}: two pairs $(p_1, q_1)$
and $(p_2, q_2)$ represent the same rational if $p_1 q_2 = p_2 q_1$. As a
HIT, this is:
\begin{align*}
  \py{Rational} &= \mathrm{HIT} \\
  &\quad \mathrm{mk} : (p : \mathbb{Z}) \times (q : \mathbb{Z}_{> 0})
    \to \py{Rational} \\
  &\quad \mathrm{path} : \forall p_1, q_1, p_2, q_2.\
    p_1 q_2 = p_2 q_1 \to \mathrm{mk}(p_1, q_1) = \mathrm{mk}(p_2, q_2)
\end{align*}

The system should recognize quotient types and track the canonical
representative (after normalization by GCD), ensuring that the type
refinements are invariant under the equivalence relation.

\paragraph{Application: sets as quotients of lists.}
Python's \py{set} type is a quotient of \py{list} by permutation and
deduplication. The system should track set properties (membership, cardinality)
without being sensitive to the internal representation order.

\section{Transport along type paths}

When a type changes along a data-flow path (e.g., a list is reversed, a
tensor is transposed), the predicates must be \emph{transported} along the
type path:

\begin{definition}[Transport]
Given a path $p : A =_{\mathcal{U}} B$ and a predicate $\phi : A \to
\mathrm{Prop}$, the \emph{transport} of $\phi$ along $p$ is:
\[
  p_*(\phi) : B \to \mathrm{Prop}, \quad p_*(\phi)(b) = \phi(p^{-1}(b))
\]
\end{definition}

\begin{example}[Transport through list reversal]
If $\phi = (xs[0] = 42)$ and $p$ is the reversal operation
$\py{reverse} : \py{list} \to \py{list}$, then:
\[
  p_*(\phi) = (\mathrm{reverse}(xs)[0] = 42) = (xs[-1] = 42)
\]
The predicate ``first element is 42'' becomes ``last element is 42'' after
reversal.
\end{example}

Transport is what the system does when propagating refinement predicates
through operations: the predicate's variables are rewritten according to the
operation's semantics. The HoTT perspective formalizes this as path transport,
ensuring that the rewriting is coherent (composing two transports gives the
transport along the composed path).

\section{Identity types and equality reasoning}

In HoTT, the identity type $a =_A b$ is the type of proofs that $a$ and $b$
are equal. The system uses identity types for equality reasoning:

\begin{itemize}[nosep]
  \item \textbf{Reflexivity}: $\mathrm{refl} : a =_A a$. Every value is
        equal to itself.
  \item \textbf{Symmetry}: if $p : a =_A b$, then $p^{-1} : b =_A a$.
  \item \textbf{Transitivity}: if $p : a =_A b$ and $q : b =_A c$, then
        $p \cdot q : a =_A c$.
  \item \textbf{Congruence}: if $p : a =_A b$ and $f : A \to B$, then
        $\mathrm{ap}_f(p) : f(a) =_B f(b)$.
\end{itemize}

These are the proof terms from Part IV, now grounded in the HoTT framework.
The advantage of the HoTT perspective is that it provides a principled account
of when equality reasoning is valid across type transformations.

\section{The fundamental groupoid of a type}

\begin{definition}[Fundamental groupoid]
The \emph{fundamental groupoid} $\Pi_1(A)$ of a type $A$ has:
\begin{itemize}[nosep]
  \item Objects: elements of $A$.
  \item Morphisms: paths (equalities) between elements.
  \item Composition: path concatenation.
  \item Inverses: path reversal.
\end{itemize}
\end{definition}

For practical types, the fundamental groupoid captures the ``transformation
group'' of the type:
\begin{itemize}[nosep]
  \item For integers: the trivial groupoid (each integer is distinct).
  \item For tensors: the groupoid includes shape-preserving transformations
        (transpositions, reshapes that preserve element count).
  \item For sets: the symmetric group (any permutation is a valid path).
  \item For quotient types: the equivalence relation's groupoid.
\end{itemize}

The system uses the groupoid structure to determine which predicates are
invariant under type transformations and which need to be transported.

\section{Truncation levels and proof irrelevance}

HoTT introduces \emph{truncation levels} that classify types by the
complexity of their equality structure:

\begin{itemize}[nosep]
  \item \textbf{$(-2)$-truncated}: contractible types (a single value).
  \item \textbf{$(-1)$-truncated}: propositions (at most one proof).
        The truth value is what matters, not how it was proved.
  \item \textbf{$0$-truncated}: sets (equality is a proposition).
        This is the level of ordinary types like \py{int} and \py{str}.
  \item \textbf{$1$-truncated}: groupoids.
  \item \textbf{$n$-truncated}: higher groupoids.
\end{itemize}

\paragraph{Proof irrelevance.}
Propositions ((-1)-truncated types) have the property that all proofs are
equal. This is critical for the system's efficiency: when a refinement
predicate is proved, the system does not need to track \emph{how} it was
proved---only that it was proved. The proof can be erased at runtime and
discarded during analysis after the obligation is discharged.

\paragraph{Sets and decidable equality.}
Most Python types are 0-truncated (sets): equality between values is a
proposition (either they are equal or they are not). This justifies the
system's use of SMT equality reasoning for most type comparisons.


\chapter{Synthesis of Relevant Propositions: Algorithms and Complexity}

This chapter presents concrete algorithms for synthesizing the propositions
that should be tracked at each program site, combining the theoretical
framework from the preceding chapters with practical computational
strategies.

\section{The proposition synthesis pipeline}

The full pipeline for determining which propositions to track at each site:

\begin{enumerate}[nosep]
  \item \textbf{Error site enumeration.} Identify all sites where a runtime
        error could occur and extract the viability predicate at each.
  \item \textbf{Backward slicing.} For each error site $e$, compute the
        backward slice: all sites that could influence $e$'s viability.
  \item \textbf{Predicate extraction.} For each site $s$ in a backward slice,
        extract candidate predicates from the viability predicate $V_e$
        pulled back to $s$.
  \item \textbf{Relevance filtering.} Remove predicates that cannot affect
        any viability predicate (zero relevance score).
  \item \textbf{Predicate selection.} If the remaining predicates exceed the
        budget, select the top-$k$ by greedy submodular maximization.
  \item \textbf{Fixpoint inference.} Run the MaxSMT inference with the
        selected predicates to compute the optimal type assignment.
  \item \textbf{CEGAR refinement.} If residuals remain, add new predicates
        and repeat from step 3.
\end{enumerate}

\section{Complexity analysis}

\begin{theorem}[Complexity of the proposition synthesis pipeline]
Let $n$ be the number of program sites, $m$ the number of error sites, and
$k$ the predicate budget. The pipeline has the following complexity:
\begin{enumerate}[nosep]
  \item Error site enumeration: $O(n)$ (one pass over the AST).
  \item Backward slicing: $O(n \cdot m)$ (one backward traversal per error
        site, amortized via shared computation).
  \item Predicate extraction: $O(m \cdot p)$ where $p$ is the average number
        of atoms per viability predicate.
  \item Relevance filtering: $O(m \cdot p \cdot n)$ (check each predicate
        against each error site's backward slice).
  \item Predicate selection: $O(k \cdot m \cdot p)$ (greedy selection with
        $k$ iterations, each evaluating $O(m \cdot p)$ marginal gains).
  \item Fixpoint inference: $O(n \cdot k^2 \cdot T_{\mathrm{SMT}})$ where
        $T_{\mathrm{SMT}}$ is the average SMT query time.
  \item CEGAR refinement: $O(r)$ refinement rounds, each repeating steps
        3--6 with an expanded predicate set.
\end{enumerate}
Total: $O(r \cdot n \cdot k^2 \cdot T_{\mathrm{SMT}})$, dominated by the
fixpoint inference step.
\end{theorem}

\section{Optimizations for large programs}

\paragraph{Incremental slicing.}
After a code edit, only the changed function's backward slices need to be
recomputed. Other slices are cached.

\paragraph{Parallelization.}
Independent SCCs (strongly connected components of the call graph) can be
analyzed in parallel. Within each SCC, the fixpoint iterations are sequential,
but the SMT queries within each iteration can be parallelized across sites.

\paragraph{Tiered analysis.}
\begin{enumerate}[nosep]
  \item \textbf{Tier 1 (fast, imprecise)}: flow-insensitive analysis with
        a small predicate domain. Runs in seconds for 10K-line programs.
        Catches simple bugs (type mismatches, obvious null dereferences).
  \item \textbf{Tier 2 (moderate, precise)}: flow-sensitive analysis with
        the full predicate domain. Runs in minutes for 10K-line programs.
        Catches shape errors, bounds violations, numeric safety issues.
  \item \textbf{Tier 3 (slow, complete)}: path-sensitive analysis with
        CEGAR refinement. Runs in minutes to hours for 10K-line programs.
        Handles complex interactions and proves functional properties.
\end{enumerate}

The system defaults to Tier 2 and escalates to Tier 3 only for functions with
proof annotations or unresolved residuals.

\section{The relevance oracle: user-system interaction protocol}

When the automatic analysis cannot determine relevance, it consults the user
through a structured protocol:

\paragraph{Ambiguous relevance query.}
\begin{lstlisting}[language={}]
RELEVANCE QUERY at line 42:
  The predicate "x > 0" at line 42 MIGHT be relevant to
  the IndexError at line 58. The system cannot determine
  whether x's positivity affects the array access at line 58.

  Options:
    [1] Mark as relevant (track it, may find additional bugs)
    [2] Mark as irrelevant (drop it, faster analysis)
    [3] Show me the data flow from line 42 to line 58
    [4] Let me provide additional context
\end{lstlisting}

\paragraph{Insufficient context query.}
\begin{lstlisting}[language={}]
CONTEXT NEEDED at line 30:
  The function "load_data" returns a DataFrame, but the system
  cannot determine whether the DataFrame has the required
  columns {"name", "age", "score"} needed at line 45.

  Options:
    [1] The DataFrame always has these columns
        (add as @ensures)
    [2] The DataFrame might not have these columns
        (add runtime check)
    [3] Show me the implementation of load_data
    [4] I'll add a contract manually
\end{lstlisting}

The interaction protocol is designed to minimize user effort: the system
presents binary or multiple-choice questions rather than open-ended queries,
and each response immediately refines the analysis.

\section{Relevance for proof obligations}

When the system is in proof mode (Part IV), relevance extends to proof
obligations:

\begin{definition}[Proof-relevant proposition]
A proposition $\phi$ at site $s$ is \emph{proof-relevant} if it is used in
the proof of at least one obligation. Formally:
\[
  \phi \in \mathcal{D}_{\mathrm{proof-relevant}}(t, s) \iff
  \exists e \in \mathrm{Obligations}_{\mathrm{proved}}(P) :
  \phi \in \mathrm{proof\_dependencies}(e)
\]
\end{definition}

Proof-relevant propositions are a superset of bug-relevant propositions: a
proposition might not directly affect a viability predicate but might be needed
as a lemma in a proof chain that eventually establishes viability.

\begin{example}
The proposition ``the list is sorted'' is not directly needed for bounds
checking (bug relevance), but it is needed to prove that binary search
returns the correct result (proof relevance). If the user has asked for a
correctness proof, sortedness becomes relevant; if they have only asked for
safety checking, it is irrelevant.
\end{example}

The system distinguishes between bug-relevance (always computed) and
proof-relevance (computed only when proof mode is active) to avoid tracking
unnecessary propositions in the common case.

\end{document}
